% concatenated from 2026-7-28-2D-beam.tex

% !TeX encoding = ISO-8859-1
\documentclass[reqno]{amsart}


\parindent = 16 pt
\parskip = 0pt
\setlength{\textheight}{8.7in}\setlength{\textwidth}{6.5in}\setlength{\oddsidemargin}{0cm}\setlength{\evensidemargin}{0cm}\setlength{\topmargin}{0mm}
\renewcommand\baselinestretch{1.05}
\widowpenalty=10000


\usepackage{color}
\usepackage{mathrsfs}
\usepackage{amscd}
\usepackage{amsmath}
\usepackage{latexsym}
\usepackage{amsfonts}

\usepackage{amssymb}
\usepackage{amsthm}
\usepackage{graphicx}
\usepackage{hyperref}
\usepackage{makecell}
\usepackage{array,color}
\usepackage{booktabs}
\usepackage{multirow}
\usepackage{bm}
\usepackage{bbm}
\usepackage{verbatim}
\usepackage{booktabs}
\usepackage{tabularx}
\usepackage{array}
\usepackage{booktabs}
\usepackage{multirow}
\usepackage{makecell}
\usepackage{float}
\setlength{\intextsep}{6pt}
\setlength{\textfloatsep}{8pt}
\setlength{\floatsep}{6pt}
%\baselineskip 16pt
%\hfuzz=6pt
\usepackage{booktabs}
\usepackage{longtable}
\usepackage{pdflscape}
\usepackage{booktabs}
\usepackage{longtable}
\usepackage{ragged2e}
\usepackage{tabularx}
\newcolumntype{Y}{>{\RaggedRight\arraybackslash}X}

\usepackage{array}
\usepackage{multirow}

\newcolumntype{C}[1]{%
	>{\centering\arraybackslash}m{#1}}

\newtheorem{cl}{Claim}
\newtheorem{theorem}{Theorem}[section]
\newtheorem{proposition}[theorem]{Proposition}
\newtheorem{corollary}[theorem]{Corollary}
\newtheorem{lemma}[theorem]{Lemma}
\newtheorem{definition}[theorem]{Definition}
\newtheorem{example}[theorem]{Example}
\newtheorem{remark}[theorem]{Remark}
\newcommand{\ra}{\rightarrow}
\renewcommand{\theequation}{\thesection.\arabic{equation}}
\newcommand{\norm}[1]{{\|#1\|}}
\newcommand{\ccc}{{\mathcal C}}
\newcommand{\one}{1\hspace{-4.5pt}1}
\newcommand\mcS{\mathcal{S}}
\newcommand\mcB{\mathcal{B}}
\newcommand\D{\mathcal{D}}
\newcommand\C{\mathbb{C}}
\newcommand\N{\mathbb{N}}
\newcommand\R{\mathbb{R}}
\newcommand\G{\mathbb{G}}
\newcommand\T{\mathbb{T}}
\newcommand\Z{\mathbb{Z}}
\newcommand{\bracket}[1]{\langle #1 \rangle}
\newcommand{\cqd}{\hfill$\Box$}
\newcommand{\medoplus}
{\mathbin{\scalebox{1.2}{$\oplus$}}}
%\newcommand{\EVstrut}{\rule[-2.0ex]{0pt}{5.6ex}}

\newenvironment{pcompactmatrix}[1] 
  {\left(\addtolength{\arraycolsep}{#1}\begin{matrix}}
  {\end{matrix}\right)}

\def\sgn{\mathop{\mathrm{sgn}}\nolimits}
\def\Id{\mathop{\mathrm{Id}}\nolimits}
\def\<{{\langle}}
\def\>{{\rangle}}

\numberwithin{equation}{section}


%\iffalse
\title[Decay estimates for Beam equations with potentials in dimension two]
{Decay Estimates for the two-dimensional Beam Equation with Potentials}

\author{Shuangshuang Chen, Han Cheng, Zijun Wan and Xiaohua Yao\textsuperscript{\dag}}

\address{Shuangshuang Chen, Department of Mathematics, Central China Normal University, Wuhan, 430079, P.R. China}
\email{chenss@mails.ccnu.edu.cn}

\address {Han Cheng, Institute of Applied Physics and Computational Mathematics, Beijing, 100088, People's Republic of China}
\email{chmathh@163.com}

\address{Zijun Wan, Department of Mathematics, Central China Normal University, Wuhan, 430079, P.R. China}
\email{zijunwan@mails.ccnu.edu.cn}

\address{Xiaohua Yao, Department of Mathematics and Key Laboratory of Nonlinear Analysis and Applications(Ministry of Education), Central China Normal University, Wuhan, 430079, P.R. China}
\email{yaoxiaohua@ccnu.edu.cn}

%\thanks{\textsuperscript{\dag} Corresponding author}

\date{\today}
\keywords{Decay estimate; Beam equation; Fourth order Schr\"odinger operator; Zero resonance}



\begin{document}



\begin{abstract}
%\baselineskip=10pt
 
This paper establishes the  sharp time-decay estimates for the following
two-dimensional beam  (plate)  equation with a real-valued
decaying potential $V$:
\[
\partial_t^2u+(\Delta^2+V)u=0,
\qquad
u(0,x)=f(x),\quad
\partial_tu(0,x)=g(x).
\]
When zero is \textit{a regular point or a first-kind resonance} of $H=\Delta^2+V$, we  prove the sharp $L^1\to L^\infty$ estimates for the solution operators:
	\begin{align*}
		\left\|\cos(t\sqrt{H})P_{\mathrm{ac}}(H)\right\|_{L^1\to L^\infty} +
		\left\|\frac{\sin(t\sqrt{H})}{t\sqrt{H}}P_{\mathrm{ac}}(H)\right\|_{L^1\to L^\infty} \lesssim \frac{1}{|t|},
	\end{align*}
	and obtain an enhanced decay $(|t|\log|t|)^{-1}$ in logarithmically weighted spaces $L^1_\omega\to L^\infty_{-\omega}$ with $\omega(x)=\log(2+|x|)$. This reveals that the potential $V$ effectively shifts the intrinsic zero resonance of  $\Delta^2$, endowing the  perturbed  sine evolution with a logarithmic gain $(\log|t|)^{-1}$ unattainable in the free case.
	
	
	
	For \textit{second-kind resonances}, 
    the unweighted decay of  both propagators deteriorates to the sharp rate
$
|t|^{-1}(\log|t|)^2
$
exactly when  one nonzero resonance
state $\phi_1$ has the trace moment
\(\langle |x|^2V,\phi_1\rangle\neq 0\), and another nonzero
state $\phi_2$ makes all the second-order moments
\(\langle x_ix_jV,\phi_2\rangle=0\). This two channel
conditions are the only sources of the logarithmic loss. If either
 is absent, the unweighted decay returns to 
\(|t|^{-1}\) and  the weighted bound reaches $(|t|\log|t|)^{-1}$ for the cosine, perfectly matching the free $\Delta^2$ dynamics.

	For \textit{third-kind resonances or zero eigenvalue}, a \(d\)-wave
 gives the strongest threshold obstruction, reducing  both propagators to the logarithmic decay
\((\log|t|)^{-1}\). Several improved estimates are  obtained  without $d$-waves. In particular, for the pure  eigenvalue case (i.e., neither $d$- nor $p$-wave resonance),  both propagators recover the optimal unweighted  rate $|t|^{-1}.$ 

%they both achieve the best weighted bound $(|t|\log|t|)^{-1}$ if and only if the fourth-order tensor $T_{ijkl}(\phi)=\langle x_i x_j x_kx_l V,\phi\rangle$ is  isotropic for each zero eigenfunction $\phi.$
	
	These results provide a complete picture of the time-decay behavior for the two-dimensional plate equation with a potential, covering all possible threshold singularities at zero energy.
\end{abstract}  


	\maketitle
	\tableofcontents
	%\baselineskip=15pt
	%%%%%%%%%%%%%%%%%%%%%%%%%%%%%%%%%%%%%%%%
	\section{Introduction and main results}
	\subsection{Introduction}	In this paper, we investigate time decay estimates for the following
two-dimensional beam (or plate) equation with a real-valued potential
$V(x)$:
\begin{equation}\label{beam_equation}
\begin{cases}
\partial_t^2u+\Delta^2u+V(x)u=0, \qquad x\in\mathbb R^2,\\
u(0,x)=f(x), \qquad \partial_tu(0,x)=g(x).
\end{cases}
\end{equation}
The two-dimensional beam (or plate) equation is a fundamental model
originating from classical elasticity theory, dating back to
Kirchhoff and Love, and describes the transverse vibrations of thin
elastic plates, see, e.g., 
\cite{Kirchhoff1850,Love1927,TimoshenkoWoinowskyKrieger1959}. In the
unperturbed case, the dynamics are governed by the biharmonic
operator $\Delta^2$, which reflects the intrinsic bending stiffness
of the material. In realistic contexts, spatially decaying potential
perturbations $V(x)$ are typically used as a simplified model for
structural heterogeneities, localized defects, or material
impurities.


    
Let $V(x)$ satisfy $ |V(x)| \lesssim \langle x \rangle^{-\mu}$ for some $\mu > 0$. By the Kato-Rellich theorem, $H=\Delta^2+V$ is a self-adjoint operator on $L^2(\mathbb{R}^n)$ with domain $H^4(\mathbb{R}^n)$ (Sobolev space). The solution to the beam equation \eqref{beam_equation} can be explicitly expressed as
\[
	u(t, x) = \cos (t \sqrt{H }) f(x) + \frac{\sin (t \sqrt{H })}{\sqrt{H}} g(x).
\]
Since the negative eigenvalues of $H$ can lead to exponentially growing modes, we will restrict our attention to the evolution projected onto the absolutely continuous spectrum. That is, we focus on the continuous part of the solution:
\[
	u_c(t, x) = \cos (t \sqrt{H })P_{\mathrm{ac}}(H) f(x) + \frac{\sin (t \sqrt{H })}{\sqrt{H}} P_{\mathrm{ac}}(H) g(x),
\]
where $P_{\mathrm{ac}}(H)$ denotes the orthogonal projection onto the absolutely continuous spectrum space of $H$. Under the assumption that $H$ has no embedded positive eigenvalues (see Subsection \ref{eigenvalue} below for a sufficient condition), there exist several interesting works about the decay estimates of the continuous part $u_{c}(t, x)$.
 In dimension $n=3$, Chen et al. \cite{CLSY_ArXiv} proved that the $L^1 \to L^\infty$ decay rates for the solution operators are ${O}(|t|^{-3/2})$ when zero is a regular point or a first-kind resonance, while the decay rate changes when other types of zero-energy resonances occur. In dimension $n=1$, Chen-Wan-Yao \cite{CWY} established an $O(|t|^{-1/2})$ decay rate, which remains unchanged  in the presence of resonances. For higher dimensions ($n\geq5$), the decay rate $O(|t|^{-n/2})$ in the regular case can be deduced using wave operator methods developed by Erdo\u{g}an and Green \cite{EG21,Erdogan-Green23}.

It is worth noting that decay estimates in even dimensions, particularly $n=2$, are more challenging. This difficulty arises from logarithmic singularities and delicate threshold obstructions at zero energy. For the Schr\"odinger operator $-\Delta+V$ in two dimensions, Schlag \cite{Schlag-2005} first established $L^1\to L^\infty$ estimates when zero is regular. Later, Erdo\u{g}an and Green \cite{Erdogan-Green} systematically treated zero-energy resonances and obtained improved decay estimates in log-weighted spaces for the regular case \cite{EG-2013}.
For the beam equation with potentials, we need to consider the fourth-order Schr\"odinger operator $H=\Delta^2+V$ in $\mathbb{R}^2$. However, the bi-Laplacian $\Delta^2$ exhibits high-order degeneracy at zero energy. This degeneracy, combined with the inherent logarithmic singularities in two dimensions, makes the dispersive analysis of $\Delta^2+V$ a challenging task.

In the free case (i.e., $V=0$), this  simplifies to $\sqrt{H} = -\Delta$, and
it is known that the following optimal decay estimates  in $\mathbb{R}^2$ hold (see Proposition \ref{free_case}):
   \[
\|\cos(t\Delta)\|_{L^1\to L^\infty} + \Big\|\frac{\sin(t\Delta)}{t\Delta}\Big\|_{L^1\to L^\infty} \sim {|t|}^{-1}.  
    \]
More interestingly, in the weighted space setting $L^1_\omega (\R^2)\to L^{\infty}_{-\omega}(\R^2)$ with $\omega(x) = \log(2+|x|)$, the decay rate of cosine evolution operator can be improved: 
 \begin{align} \label{free_weighted2} \|\cos(t\Delta)\|_{L^1_\omega\to L^\infty_{-\omega}} \sim ({|t|\log |t|})^{-1},       \end{align} 
which is strictly unattainable for both $e^{it\Delta}$ and $\sin(t\Delta)/(t\Delta)$ in the  weighted space. This enhanced phenomenon \eqref{free_weighted2} exactly happens in these dimensions $n \equiv 2 \pmod{4}$ (see Remark \ref{free_cos_higher_dim}). 

For a decaying potential $V$, understanding how $\cos(t\sqrt{H})$ and $\sin(t\sqrt{H})/(t\sqrt{H})$ decay over time requires focusing on delicate behavior at low frequencies. This low-frequency dynamics becomes significant due to the creation of new zero resonances of the operator $H$, which occur  when  $V$ strongly couples with $\Delta^2$ in low-energy regimes.
As described in Definition \ref{definition1}, zero can be a regular point, one of three types of resonances, or an eigenvalue. These different threshold configurations lead to distinct dispersive behaviors.






When zero is either a regular point or a first-kind resonance of $H$,  we  not only establish  the following sharp $L^1 \to L^\infty$ estimates for the perturbed  operators (see Theorem \ref{main_theorem-1}):
 \begin{align*}
    \left\|\cos (t \sqrt{H}) P_{\mathrm{ac}}(H)\right\|_{L^1 \to L^{\infty}} +
    \bigg\|\frac{\sin (t \sqrt{H})}{t\sqrt{H}} P_{\mathrm{ac}}(H)\bigg\|_{L^1 \to L^{\infty}} \lesssim {|t|}^{-1},
    \end{align*}
but also obtain the following improved weighted decay estimates:    
 \begin{align}\label{the first and regular weigthed estimate}
    \left\|\cos (t \sqrt{H}) P_{\mathrm{ac}}(H)\right\|_{L_\omega^1 \to L_{-\omega}^{\infty}}+
	\bigg \|\frac{\sin (t \sqrt{H})}{t\sqrt{H}} P_{\mathrm{ac}}(H)\bigg \|_{L^1_\omega \rightarrow L^{\infty}_{-\omega}}\lesssim ({|t|\log |t|})^{-1}.
    \end{align}
It is worth noting that the free sine evolution $\sin(t\Delta)/(t\Delta)$ only satisfies a bound $\sim |t|^{-1}$ in logarithmically weighted spaces $L^1_\omega \to L^\infty_{-\omega}$. Hence this reveals that  the potential $V$  effectively modified the intrinsic zero resonance of $\Delta^2$, so that  the perturbed sine evolution $\sin(t\sqrt{H})/(t\sqrt{H})$ gains an additional logarithmic factor $(\log|t|)^{-1}$ in the regular or the first kind resonance case that is distinct from the free case $V=0$.




\begin{comment}
The second-kind resonance case exhibits a more subtle threshold
phenomenon.
Although the free bi-Laplacian $\Delta^2$ itself has a second-kind
resonance at zero, the perturbed operator may either retain the free
decay rate or develop a sharp logarithmic loss. Let
\[
\mathfrak M_2(\phi)
=
\bigl(
\langle x_ix_jV,\phi\rangle
\bigr)_{i,j=1}^2,
\qquad
\tau(\phi)
=
\operatorname{tr}\mathfrak M_2(\phi)
=
\langle |x|^2V,\phi\rangle .
\]
Theorem \ref{main_theorem-2} shows that the sharp logarithmic loss
$
|t|^{-1}(\log|t|)^2
$
occurs precisely when two distinct second-moment channels coexist:
the resonance space contains a mode with nonzero radial second
moment, \(\tau(\phi)\neq0\), and also a nonzero mode whose full
quadratic moment tensor vanishes, \(\mathfrak M_2(\phi)=0\).
 If either channel is absent,
the unweighted decay returns to 
\(|t|^{-1}\) and  the weighted bound achieves $(|t|\log|t|)^{-1}$ for the cosine, perfectly matching the the free bi-Laplacian dynamics.
\end{comment}


The second-kind resonance case exhibits a more subtle threshold
phenomenon. Although the free bi-Laplacian $\Delta^2$ itself has a
second-kind resonance at zero, the perturbed operator may either
retain the free decay rate or develop a sharp logarithmic loss. To
describe the criterion, define
\[
\mathfrak M_2(\phi)
=
\bigl(
\langle x_ix_jV,\phi\rangle
\bigr)_{i,j=1}^2,
\qquad
\tau(\phi)
=
\operatorname{tr}\mathfrak M_2(\phi)
=
\langle |x|^2V,\phi\rangle,
\]
where \(\phi\) belongs to the second-kind resonance space of \(H\).
Theorem \ref{main_theorem-2} shows that the unweighted
\(L^1\to L^\infty\) decay estimate has the sharp order
$
|t|^{-1}(\log|t|)^2
$
precisely when two distinct second-moment channels coexist: the
resonance space contains a state \(\phi_1\) with
$
\tau(\phi_1)\neq0,
$
and also a nonzero state \(\phi_2\) with
$
\mathfrak M_2(\phi_2)=0.
$
Thus the logarithmic loss is a mixed-channel effect.  If either
channel is absent, the unweighted decay returns to 
\(|t|^{-1}\) and  the weighted bound achieves $(|t|\log|t|)^{-1}$ for the cosine, perfectly matching the free $\Delta^2$ dynamics.



Finally, we consider third-kind resonances and zero eigenvalues.
Theorem \ref{main_theorem-31} shows that the presence of a
\(d\)-wave resonance gives the strongest threshold obstruction:
it only satisfies the logarithmic decay:
\begin{equation*}  
    \left\|\cos (t \sqrt{H}) P_{\mathrm{ac}}(H)\right\|_{L^1 \to L^{\infty}} +    
    \bigg\|\frac{\sin (t \sqrt{H})}{t\sqrt{H}} P_{\mathrm{ac}}(H)\bigg\|_{L^1 \to L^{\infty}} \sim ({\log |t|})^{-1}.    
    \end{equation*}
Moreover, we  obtain several improved  estimates in the absence of a $d$-wave. In particular, in the pure $L^2$ eigenvalue case (i.e., neither $d$-wave nor $p$-wave), we not only recover the unweighted $|t|^{-1}$ bound,  but also
obtain the enhanced  bound $(|t|\log|t|)^{-1}$ as in \eqref{the first and regular weigthed estimate} if the fourth-order
tensor
$
T_{ijkl}(\phi)
=
\langle x_ix_jx_kx_lV,\phi\rangle
$
is isotropic,  namely,  there exists a scalar \(\gamma_\phi\) such that
$
T_{ijkl}(\phi)
=
\gamma_\phi
\bigl(
\delta_{ij}\delta_{kl}
+
\delta_{ik}\delta_{jl}
+
\delta_{il}\delta_{jk}
\bigr).
$


These results provide a complete classification of
time-decay estimates for 2D plate equation with a
decaying potential. They cover all possible threshold singularities at zero and identify the precise moment
conditions responsible for logarithmic losses or logarithmic gains
in the long-time dynamics.


\subsection{Some notations}\label{Notations}
In this subsection, we collect some notations used throughout the paper.
	\begin{itemize}
		
		
		%\vskip0.2cm
\item For $a,b\in\mathbb{R}^+$, $a\lesssim b$ (resp. $a\gtrsim b$) means $a\le cb$ (resp. $a\ge cb$) for some $c>0$. In particular, if $b\lesssim a\lesssim b$, we write $a\sim b$. Moreover, for $a\in \mathbb{R}$,   $a\pm$ denote $a\pm \epsilon$   for any arbitrarily small  $\epsilon >0$.
		%\vskip0.2cm
        %\item  $\mathbb{N}_0=\{0,1,\cdots\}$,  $\mathbb{N}^+=\{1,2,\cdots\}$.
        \vskip0.2cm
		\item We write $f(\lambda) = {O}_k(g(\lambda))$ if
    \[
        \frac{d^\ell}{d\lambda^\ell} f(\lambda) = O\Bigl( \frac{d^\ell}{d\lambda^\ell} g(\lambda) \Bigr), \quad \ell = 0, 1, \cdots,k.
    \]
   \vskip 0.2cm
   \item The notation $\mathcal{O}(f(t))$ refers to an operator satisfying $\|\mathcal{O}(f(t))\|_{L^1_\omega \rightarrow L^{\infty}_{-\omega}} \lesssim |f(t)|.$
   \vskip 0.2cm
   \item  Let \(\mathbb{B}(L^2)\) denote the space of bounded linear operators on \(L^2\). For \(A\in\mathbb{B}(L^2)\), we write \(A^*\) for its adjoint.
        \vskip0.2cm
		
		\item Let $\chi_1 \in C_c^{\infty}(\mathbb{R})$ such that $\chi_1(\lambda)=1$ for $|\lambda|<\lambda_0\ll1$ and $\chi_1(\lambda)=0$ if $|\lambda|>2\lambda_0$, where $\lambda_0$ is some sufficiently small positive constant depending on low energy expansion of $\left( M^{\pm}(\lambda)\right)^{-1}$ in Theorem \ref{thm:M_inverse}. Define $\chi_2(\lambda):=1-\chi_1(\lambda)$ and $\widetilde{\chi}_j(\lambda):=\chi_j(\lambda^4)$  for  $j=1,2.$
\vskip0.2cm
\item
    $ \langle f, g \rangle := \int_{\mathbb{R}^2} f(x) \overline{g(x)} \, dx$ denotes the inner product  or the dual pair.
	For $x\in\mathbb{R}^n$, write $\langle x \rangle := 1 + | x |$.
		For $s \in \mathbb{R}$, we  define the weighted $L^2(\mathbb{R}^2) $ spaces:
		$$L^{2,s}(\mathbb{R}^2) := \{ f \in L_{\text {loc}}^2(\mathbb{R}^2) \mid  \langle \cdot \rangle^s f \in L^2(\mathbb{R}^2) \} .$$
\end{itemize}
	\subsection{Main results}\label{main result}
  In the following, we first give the precise definition of zero resonance types of $H=\Delta^2+V(x)$ on $\mathbb{R}^2$.
For $\sigma \in \mathbb{R}$, let $W_{-\sigma}\left(\mathbb{R}^2\right)$ denote the intersection space
$$
W_{-\sigma}\left(\mathbb{R}^2\right)=\bigcap_{s>\sigma} L^{2,-s}\left(\mathbb{R}^2\right) ,
$$
which is increasing in $\sigma\in\R$ and satisfies $L^{2,-\sigma}(\mathbb{R}^2) \subset W_{-\sigma}(\mathbb{R}^2)$. In particular, $W_0\left(\mathbb{R}^2\right) \supset L^2\left(\mathbb{R}^2\right)$.
\begin{definition}\label{definition1}
{\rm Let $H = \Delta^2 + V(x)$ with $|V(x)| \lesssim \langle x \rangle^{-\mu}$ for some $\mu > 0$. We say that
\begin{itemize}
\item[(i)] Zero is a {\it first kind resonance}  of $H$ if the equation $H\phi = 0$  has only zero solution  in $W_{-1}(\mathbb{R}^2)$, but has a nonzero solution  $\phi \in W_{-2}(\mathbb{R}^2)$ and there exists a polynomial $\sum_{|\alpha|=1} C_\alpha x^\alpha$ such that $\phi-\sum_{|\alpha|=1} C_\alpha x^\alpha\in W_{-1}(\mathbb{R}^2)$.  
\vskip0.1cm

\item[(ii)] Zero is a {\it second kind resonance} of $H$ if
the equation $H\phi = 0$ has a nonzero solution  in $W_{-1}(\mathbb{R}^2)$, but only has a zero solution  $\phi$ in $W_{0}(\mathbb{R}^2)$.
\vskip0.1cm

\item[(iii)] Zero is a {\it third kind resonance} of $H$ if the equation $H\phi = 0$ has a nonzero solution  in $W_0(\mathbb{R}^2)$, but there exists only a zero solution  $\phi$ in  $L^2(\mathbb{R}^2)$.

\vskip0.1cm
\item[(iv)] Zero is an \emph{eigenvalue}  of $H$ if there exists a nonzero $\phi \in L^2(\mathbb{R}^2)$ satisfying $H\phi = 0$.
\vskip0.1cm

\item[(v)] Zero is a {\it regular point} of $H$ if $H$ has neither zero eigenvalues nor zero resonances.
		\end{itemize}}
	\end{definition}
For convenience, we use $\mathbf{k} \in \{0,1,2,3,4\}$ to denote zero energy types: $\mathbf{k}=0$ (regular), $\mathbf{k}=1,2,3$ (resonances), $\mathbf{k}=4$ (eigenvalue).
Moreover, two equivalent characterizations of zero resonances for $H$ will be provided in Proposition~\ref{characterizations} and  Proposition~\ref{characterizations Q} in Section~\ref{subsec:Q}.
 

 We now state our main results, which are organized into the following four theorems  based on the nature of the zero-energy resonances. To begin, we define the logarithmic weight function $\omega(x) = \log(2+|x|)$ and the associated weighted spaces
\[
L^p_{\pm \omega} := \bigl\{ f \in L_{\mathrm{loc}}^p(\mathbb{R}^2) \bigm| \omega(x)^{\pm 1} f \in L^p(\mathbb{R}^2) \bigr\}, \qquad p = 1, \infty.
\]




\begin{theorem}\label{main_theorem-1}
Let $H = \Delta^2 + V$ and $|V(x)| \lesssim \langle x \rangle^{-11-}$. 
Assume that $H$ has no positive embedded eigenvalues and that zero is \textbf{a regular point or a first-kind resonance} of $H$.
Then 
\begin{equation}\label{regular_first}
    \left\|\cos (t \sqrt{H}) P_{\mathrm{ac}}(H)\right\|_{L^1 \to L^{\infty}} +
    \bigg\|\frac{\sin (t \sqrt{H})}{t\sqrt{H}} P_{\mathrm{ac}}(H)\bigg\|_{L^1 \to L^{\infty}} \lesssim |t|^{-1},
\end{equation}where $P_{\mathrm{ac}}(H)$ denotes the projection onto the absolutely continuous spectrum space of $H$. Moreover, in the weighted space setting $L^1_\omega \to L^{\infty}_{-\omega}$, we can obtain the improved decay estimate:
\begin{equation}\label{regular_first_weight}
    \left\|\cos (t \sqrt{H}) P_{\mathrm{ac}}(H)\right\|_{L^1_\omega \to L^{\infty}_{-\omega}} +
    \bigg\|\frac{\sin (t \sqrt{H})}{t\sqrt{H}} P_{\mathrm{ac}}(H)\bigg\|_{L^1_\omega \to L^{\infty}_{-\omega}} \lesssim \frac{1}{|t| \log|t|}, \quad |t|\ge2.
\end{equation}
\end{theorem}


\begin{remark}
{\rm
We provide several comments regarding the decay estimates \eqref{regular_first} and \eqref{regular_first_weight}:  
\begin{itemize}  
    \item  
    For the unweighted estimate \eqref{regular_first}, the $|t|^{-1}$ decay rate exactly matches the intrinsic dispersive rate of the free operator in $\mathbb{R}^2$ (see Proposition \ref{free_case}):   
    \[
\|\cos(t\Delta)\|_{L^1\to L^\infty} + \Big\|\frac{\sin(t\Delta)}{t\Delta}\Big\|_{L^1\to L^\infty} \sim |t|^{-1}.  
    \]
    This suggests that the estimate \eqref{regular_first} is expected to be optimal in the unweighted setting $L^1\to L^\infty$.
    \begin{comment}
    Consequently, we expect the bound \eqref{regular_first} to be sharp in $L^1\to L^\infty$. Heuristically, since a short-range potential $V(x)$ is highly localized in space, it is not expected to improve the decay rate in the unweighted setting.
  % it is not expected to alter the global  dispersion in the unweighted setting. 
  \end{comment}
    \item    
    For the weighted estimate \eqref{regular_first_weight}, it is noteworthy that the perturbed sine evolution $\sin (t \sqrt{H})/(t\sqrt{H})$ exhibits an extra logarithmic decay factor of $(\log|t|)^{-1}$ compared to the following free weighted estimates (see Proposition \ref{free_case}):  
       \[\Big\|\frac{\sin(t\Delta)}{t\Delta}\Big\|_{L^1_\omega\to L^\infty_{-\omega}} \sim |t|^{-1}.   \]
   \begin{comment}  
  This logarithmic gain reflects the transition of the threshold type
from a second-kind resonance of the free operator $\Delta^2$ to a regular point
or a first-kind resonance of the perturbed operator  $\Delta^2+V$.
\end{comment}
This reflects that the potential $V$ effectively shifts the intrinsic zero resonance of  $\Delta^2$, endowing the  perturbed  sine evolution with the logarithmic gain.
  
\item  %Moreover, we remark on the decay assumption of the potential
%$V$ appearing  in Theorem~\ref{main_theorem-1}( %including Theorems~\ref{main_theorem-2}, %\ref{main_theorem-31}, and
%\ref{main_theorem-32}).
 Moreover, we point out that the decay rate $\mu$ of the potential $V$   in Theorem~\ref{main_theorem-1}  is determined by the types of zero energy. The decay requirement guarantees the
	asymptotic expansion  in Theorem~\ref{thm:M_inverse}.	In contrast, the high-energy analysis requires only the weaker decay
	condition $\mu>4$ (see Theorem~\ref{main_theorem_high}). Hence the decay rate in Theorem~\ref{main_theorem-1}  may not be optimal, even in the regular case.
    \end{itemize}}  
\end{remark}
%\vskip0.3cm

Next to consider the case of \textbf{the second-kind zero resonance},   let us introduce the space of resonance solutions in $W_{-1}(\mathbb{R}^2)$:  
\[
\mathcal{E} = \big\{ \phi \in W_{-1}(\mathbb{R}^2) \ \big| \ H\phi = 0 \big\}. 
\] 
When zero is a second-kind resonance of $H$, Definition \ref{definition1} implies that $\mathcal{E} \neq \{0\}$ and  $\mathcal{E} \cap W_0(\mathbb{R}^2) = \{0\}$.     
To classify the dispersive behavior based on the second-order potential moments on $\mathcal{E},$
we define the
second-moment map
\begin{equation}\label{eq:second-moment-map}
	\mathfrak{M}_2:
	\mathcal{E}
	\longrightarrow
	\operatorname{Sym}_2(\mathbb{C}),
	\qquad
	\mathfrak{M}_2(\phi)
	=
	\left(
	\langle x_ix_jV,\phi\rangle
	\right)_{i,j=1}^2,
\end{equation}
and its trace
\begin{equation}\label{eq:radial-second-moment}
	\tau(\phi)
	:=
	\operatorname{tr}\mathfrak{M}_2(\phi)
	=
	\langle |x|^2V,\phi\rangle.
\end{equation}

The condition
$
\tau\not\equiv0
$ on $\mathcal{E}$
means that the resonance space contains a state with nonzero trace
second moment, whereas
$
\ker\mathfrak{M}_2\neq\{0\}
$
means that it contains a nonzero state whose  all second-order moments 
vanish.
 Theorem \ref{main_theorem-2} shows that
a logarithmic loss occurs only when these two distinct channels
coexist. 
\begin{comment}
Thus the second-kind resonance case exhibits a genuinely
non-monotone phenomenon with respect to moment cancellation: the
loss is not caused by either channel separately, but by their
interaction in the low-energy resolvent expansion.
\end{comment}
\begin{theorem}\label{main_theorem-2}
	Let $H=\Delta^2+V$ with
	$
	|V(x)|
	\lesssim
	\langle x\rangle^{-14-}.
	$
	Assume that $H$ has no positive embedded eigenvalues and that zero
	is a second-kind resonance of $H$. Let $\mathfrak{M}_2$ and $\tau$
	be defined by \eqref{eq:second-moment-map} and
	\eqref{eq:radial-second-moment}, respectively.
	
	\begin{itemize}
		\item[(i)]
		Assume that
	$\tau\not\equiv0$
			on $\mathcal{E}$ and $
			\ker\mathfrak{M}_2\neq\{0\}.$
		Then 
	\begin{align}\label{log-cos-sin}
			&
			\left\|
			\cos(t\sqrt H)P_{\mathrm{ac}}(H)
			\right\|_{L^1\to L^\infty}
			+
			\bigg\|
			\frac{\sin(t\sqrt H)}{t\sqrt H}
			P_{\mathrm{ac}}(H)
			\bigg\|_{L^1\to L^\infty}
			\lesssim
			\frac{\bigl(\log(2+|t|)\bigr)^2}{|t|},\end{align}
            and in the weighted setting $L^1_\omega\to L^\infty_{-\omega},$ for $|t|\geq 2,$
            \begin{align*}
		\left\|
			\cos(t\sqrt H)P_{\mathrm{ac}}(H)
			\right\|_{L^1_\omega\to L^\infty_{-\omega}}\lesssim \frac{\log|t|}{|t|}\ \ \text{and} \ \ \bigg\|
		\frac{\sin(t\sqrt H)}{t\sqrt H}
			P_{\mathrm{ac}}(H)
			\bigg\|_{L^1_\omega\to L^\infty_{-\omega}}
			\sim
			\frac{(\log|t|)^2}{|t|}. 
		\end{align*}
		
		\item[(ii)]
		Assume that
			$\tau\equiv0$
			on $\mathcal{E}$
			or
			$\ker\mathfrak{M}_2=\{0\}.$
	Then   
	the  estimates further improve to match the optimal bounds for the free case $\Delta^2$, namely,  
	\begin{align}  
		%\label{regular_first}  
		&\ \ \ \ \ \ \ \ \ \ \ \ \ \ \  \ \  \ \ \ \  \ \ \  \left\| \cos (t \sqrt{H})  P_{\mathrm{ac}}(H)\right\|_{L^1 \to L^{\infty}} +  
		\bigg\|\frac{\sin (t \sqrt{H})}{t\sqrt{H}} P_{\mathrm{ac}}(H)\bigg\|_{L^1 \to L^{\infty}} \lesssim |t|^{-1}, \nonumber \\
		\label{2-best}  
		& \ \ \ \ \left \|\cos (t \sqrt{H}) P_{\mathrm{ac}}(H)\right\|_{L^1_\omega\to L^\infty_{-\omega}} \lesssim\frac{1}{|t|\log |t|}\ \ \text{and}\ \ \bigg\|\frac{\sin (t \sqrt{H})}{t\sqrt{H}} P_{\mathrm{ac}}(H)\bigg\|_{L^1_\omega\to L^\infty_{-\omega}}\sim|t|^{-1}, \quad |t| \ge 2.    
	\end{align} 
	\end{itemize}
\end{theorem}
%\begin{comment}
\begin{remark}
	{\rm In Theorem \ref{main_theorem-2}, the logarithmic loss is instead a mixed-channel effect.  The loss occurs only when the resonance space contains both a 
	nonzero trace-moment state $\phi_1$ and a fully vanishing quadratic-moment nonzero state $\phi_2$.
	In particular, this  happens only when
	$
	\dim\mathcal{E}\geq2.
	$}
\end{remark}

%\end{comment}


  %\vskip0.2cm


Finally, it remains to analyze the scenario where zero is \textbf{a third-kind resonance or an eigenvalue} of $H$.
By Definition \ref{definition1}, this is equivalent to   
$$ \mathcal{E}_1 := \big\{ \phi \in W_{0}(\mathbb{R}^2) \ \big| \ H\phi = 0 \big\} \neq \{0\},$$   
which means there exists at least one nonzero solution $\phi\in W_{0}(\mathbb{R}^2)$ to $H\phi=0.$   
Below we split our results into Theorems \ref{main_theorem-31} and \ref{main_theorem-32} by the spatial decay of the zero-energy states of $\mathcal{E}_1$:  
\begin{itemize}  
    \item \textbf{Presence of a $d$-wave resonance:}  $\mathcal{E}_1$ is not fully contained in $L^2$, meaning there exists a nonzero solution $\phi \in W_{0}(\mathbb{R}^2) \setminus L^2(\mathbb{R}^2)$ to $H\phi = 0$.  
    \vskip0.1cm
    \item \textbf{Absence of a $d$-wave resonance:} Zero is an eigenvalue of $H$ but exhibits no $d$-wave  (i.e., every solution $\phi \in \mathcal{E}_1$ strictly belongs to $L^2(\mathbb{R}^2)$).  
\end{itemize}  
  
\begin{theorem}\label{main_theorem-31}    
Let $H = \Delta^2 + V$ with $|V(x)| \lesssim \langle x \rangle^{-18-}$.   
Assume that $H$ has no positive embedded eigenvalues and that zero is either a third-kind resonance or an eigenvalue accompanied by a d-wave resonance. Then 
    \[ 
        \left\|\cos (t \sqrt{H}) P_{\mathrm{ac}}(H)\right\|_{L^1 \to L^{\infty}} +    
        \bigg\|\frac{\sin (t \sqrt{H})}{t\sqrt{H}} P_{\mathrm{ac}}(H)\bigg\|_{L^1 \to L^{\infty}} \sim \frac{1}{\log |t|}, \quad |t| \gg 1.    
    \]   
\end{theorem}  



We next assume that zero is an eigenvalue but that the $d$-wave
resonance is absent. In this case
\begin{align}\label{E--2}
\mathcal{E}_1
=
\mathcal{E}_2,
\qquad
\mathcal{E}_2
:=
\left\{
\phi\in L^2(\mathbb{R}^2)
\mid
H\phi=0
\right\},
\end{align}
namely, all threshold states in $W_0(\mathbb{R}^2)$ are 
$L^2$-eigenfunctions.

In the absence of a $d$-wave resonance, a possible $p$-wave  is the remaining source of logarithmic loss in the unweighted estimates. If no $p$-wave  occurs, the decay  is determined by the fourth-order potential moments of the zero eigenfunctions. 
To describe this case, define
\begin{equation}\label{eq:E3}
	\mathcal{E}_3
	:=
	\left\{
	\phi\in\mathcal{E}_2
	\ \middle|\
	\begin{aligned}
		&
		\langle x_1^3x_2V,\phi\rangle
		=
		\langle x_1x_2^3V,\phi\rangle
		=
		0,
		\\
		&
		\langle x_1^4V,\phi\rangle
		=
		\langle x_2^4V,\phi\rangle
		=
		3\langle x_1^2x_2^2V,\phi\rangle
	\end{aligned}
	\right\}.
\end{equation}
Equivalently, $\mathcal{E}_3$ consists of those zero eigenfunctions
whose fourth-order potential moment tensor
$
T_{ijkl}(\phi)
:=
\langle x_ix_jx_kx_lV,\phi\rangle
$
is isotropic, i.e. for $\phi\in\mathcal{E}_3$ there
exists a scalar $\gamma_\phi\in\mathbb C$ such that
\[
T_{ijkl}(\phi)
=
\gamma_\phi
\left(
\delta_{ij}\delta_{kl}
+
\delta_{ik}\delta_{jl}
+
\delta_{il}\delta_{jk}
\right),
\qquad
 i,j,k,l=1,2.
\]
Indeed,  this tensorial condition is equivalent to
$
T_{1112}=T_{1222}=0,$ and
$
T_{1111}=T_{2222}=3T_{1122},
$
which is exactly the set of moment conditions in
\eqref{eq:E3}. Thus
$
\mathcal{E}_2=\mathcal{E}_3
$
means that every zero eigenfunction has isotropic fourth-order moments tensor form.


\begin{theorem}\label{main_theorem-32}


Let $H = \Delta^2 + V$ with $|V(x)| \lesssim \langle x \rangle^{-22-}$.   
	Assume that $H$ has no positive embedded eigenvalues and that zero is an eigenvalue of $H$ but exhibits no d-wave resonance $(\text{i.e.,}\  \mathcal{E}_1 = \mathcal{E}_2 )$.  
	\begin{enumerate}   
		\item If $H$ admits a p-wave resonance $($i.e., there exists a nonzero $\phi \in W_{-1}(\mathbb{R}^2) \setminus W_0(\mathbb{R}^2)$ such that $H\phi = 0 )$, then the $L^1 \to L^{\infty}$ estimate \eqref{log-cos-sin}  holds. 
		\vskip0.2cm
		\item If $H$ lacks a p-wave resonance, then  the $L^1 \to L^{\infty}$ estimate \eqref{regular_first} holds.  
		Moreover, we can  further distinguish the following two subcases:
		\begin{itemize}
			\item If $\mathcal{E}_2\neq \mathcal{E}_3$, then we obtain  the same weighted estimate \eqref{2-best}  as in the free case. 
			\vskip0.1cm
			\item If $\mathcal{E}_2=\mathcal{E}_3,$ then the  enhanced weighted estimate \eqref{regular_first_weight} holds as in the regular case.
		\end{itemize}  
	\end{enumerate}    
\end{theorem}  


\begin{comment}
\begin{remark}\label{rem:hierarchy_eigenvalue} {\rm     
We provide some further comments for Theorems \ref{main_theorem-31} and \ref{main_theorem-32}:    
\begin{itemize}    
\item  Theorem \ref{main_theorem-31} shows that a $d$-wave resonance causes the strongest threshold obstruction. It slows down  to the  logarithmic decay $(\log |t|)^{-1}$,  regardless of any eigenvalue contributions. When the \(d\)-wave  is absent, the next
possible obstruction is a \(p\)-wave resonance, which gives the
logarithmic loss  in the unweighted
estimates.
			
\item  If both \(d\)-wave and \(p\)-wave resonances are absent, threshold eigenvalues do not affect the unweighted decay, which remains optimal at $|t|^{-1}$.   In weighted spaces, if $\mathcal{E}_2 \neq \mathcal{E}_3$,
the lack of  isotropy in the fourth-order moments of $V\phi$ restricts the sine evolution to $|t|^{-1}$.
If \(\mathcal E_2=\mathcal E_3\), 
both propagators satisfy the enhanced weighted decay
\((|t|\log |t|)^{-1}\), matching the same weighted bounds as for regular or first-kind threshold.
\end{itemize}    
}  
\end{remark}  

\end{comment}
For the reader's convenience, we summarize the main dispersive estimates obtained above in Table~\ref{tab:main-results}.
\vspace{-5mm}
\newcommand{\Tcell}[1]{\makecell[c]{\rule[-1.9ex]{0pt}{5.5ex}#1}}
\begin{table}[H]
\centering
\caption{Summary of the long-time decay estimates}
\label{tab:main-results}
\vspace{1mm}
\renewcommand{\arraystretch}{1.35}
\setlength{\tabcolsep}{4pt}
\footnotesize
\begin{tabular}{|
>{\centering\arraybackslash}m{0.18\textwidth}|
>{\centering\arraybackslash}m{0.23\textwidth}|
>{\centering\arraybackslash}m{0.20\textwidth}|
>{\centering\arraybackslash}m{0.15\textwidth}|
>{\centering\arraybackslash}m{0.15\textwidth}|}
\hline
\multirow{2}{*}{\makecell[c]{\textbf{Zero-energy type}}}
&
\multirow{2}{*}{\makecell[c]{\textbf{Sub-condition}}}
&
\multirow{2}{*}{\makecell[c]{\textbf{Unweighted}
\(\boldsymbol{L^1\to L^\infty}\)\\$\boldsymbol{U_{\cos}(t)}$ \ \bf{\&} \ $\boldsymbol{U_{\sin}(t)}$ }}
&
\multicolumn{2}{c|}{\textbf{Weighted }
\(\boldsymbol{L^1_\omega\to L^\infty_{-\omega}}\)}
\\
\cline{4-5}
&
&
&
\(\boldsymbol{U_{\cos}(t)}\)
&
\(\boldsymbol{U_{\sin}(t)}\)
\\
\hline

\Tcell{Regular / 1st-kind \\ \((\mathbf k=0,1)\)}
&
\Tcell{None}
&
\Tcell{\( |t|^{-1}\)}
&
\Tcell{\((|t|\log|t|)^{-1}\)}
&
\Tcell{\( (|t|\log|t|)^{-1}\)}
\\
\hline

\multirow{2}{=}{\Tcell{2nd-kind Resonance\\ \((\mathbf k=2)\)}}
&
\Tcell{
\(\tau\not\equiv0\)\ \text{and}\ 
\(\ker\mathfrak M_2\neq\{0\}\)}
&
\Tcell{\( |t|^{-1}(\log|t|)^2\)}
&
\Tcell{\( |t|^{-1}(\log|t|)\)}
&
\Tcell{\(\sim |t|^{-1}(\log|t|)^2\)}
\\
\cline{2-5}

&
\Tcell{
\(\tau\equiv0\) \ or \ 
\(\ker\mathfrak M_2=\{0\}\)}
&
\Tcell{\( |t|^{-1}\)}
&
\Tcell{\( (|t|\log|t|)^{-1}\)}
&
\Tcell{\(\sim |t|^{-1}\)}
\\
\hline

\Tcell{3rd-kind or\\ Eigenvalue+\(d\)-wave\\
\((\mathbf k=3\) or \(4)\)}
&
\Tcell{\(d\)-wave present}
&
\Tcell{\(\sim(\log|t|)^{-1}\)}
&
\Tcell{--}
&
\Tcell{--}
\\
\hline

\multirow{3}{=}{\Tcell{Eigenvalue+no \(d\)-wave\\ \((\mathbf k=4)\)}}
&
\Tcell{\(p\)-wave present}
&
\Tcell{\( |t|^{-1}(\log|t|)^2\)}
&
\Tcell{--}
&
\Tcell{--}
\\[2mm]
\cline{2-5}

&
\Tcell{No \(p\)-wave and
\(\mathcal E_2\neq\mathcal E_3\)}
&
\Tcell{\( |t|^{-1}\)}
&
\Tcell{\( (|t|\log|t|)^{-1}\)}
&
\Tcell{\(\sim |t|^{-1}\)}
\\[2mm]
\cline{2-5}

&
\Tcell{No \(p\)-wave and
\(\mathcal E_2=\mathcal E_3\)}
&
\Tcell{\( |t|^{-1}\)}
&
\Tcell{\( (|t|\log|t|)^{-1}\)}
&
\Tcell{\((|t|\log|t|)^{-1}\)}
\\[2mm]
\hline
\end{tabular}
\vskip0.2cm
\parbox{0.95\textwidth}{\footnotesize
Here
\(U_{\cos}(t)=\cos(t\sqrt H)\) and
\(U_{\sin}(t)=\sin(t\sqrt H)/(t\sqrt H)\).
The notations
\(\mathfrak M_2\), \(\tau\), \(\mathcal E_2\), and \(\mathcal E_3\)
are defined in \eqref{eq:second-moment-map},
\eqref{eq:radial-second-moment}, \eqref{E--2}, and \eqref{eq:E3},
respectively.}
\end{table}











\subsection{Further background}
		
\subsubsection{Embedded positive eigenvalue}\label{eigenvalue}
According to Kato's theorem \cite{Kat59}, the Schr\"odinger operator $-\Delta+V$ has no positive eigenvalues if $V$ is bounded and $V=o(|x|^{-1})$ as $|x|\to\infty$ (see also \cite{FHHH82, IJ03, KT06, Simon69}). This criterion, however, fails for the fourth-order operator $H=\Delta^2+V$. For any $n\ge1$, one can construct $V\in C_0^\infty(\R^n)$ such that $H$ possesses positive eigenvalues (see Feng et al. \cite[Section 7.1]{FSWY20}). These results show that the absence of positive eigenvalues for fourth-order Schr\"odinger operators is more subtle and less stable than in the second-order case under potential perturbations.
Nevertheless, if $V\in C^1(\R^2)\cap L^\infty(\R^2)$ is repulsive, i.e., $(x\cdot\nabla)V\le0$, then $H$ has no eigenvalues (see \cite[Theorem 1.11]{FSWY20}). This criterion also applies to general higher-order elliptic operators $P(D)+V$ in any dimension $n\ge1$.

\subsubsection{Wave equations }

For the classical wave equation with a real-valued decaying potential,
\[
\left\{\begin{array}{l}
	\partial_t^2 u + (-\Delta + V) u = 0, \quad (t,x) \in \mathbb{R} \times \mathbb{R}^n, \\
	u(0,x) = f(x), \quad \partial_t u(0,x) = g(x),
\end{array}\right.
\]
decay estimates have been extensively studied. When $V=0$, the solution operators $\cos(t\sqrt{-\Delta})$ and $\sin(t\sqrt{-\Delta})/\sqrt{-\Delta}$ satisfy  $W^{k+1,1}\to L^\infty$ and $W^{k,1}\to L^\infty$ estimates  with $O(|t|^{-1/2})$  in dimension two for $k>\frac12$. The endpoint
\(k=\frac{1}{2}\) is more delicate and typically requires a refined
formulation in Hardy-, Besov-, or BMO-type spaces; see \cite{MSW,Green14}.   By contrast, in odd spatial dimensions \(n\ge3\), the strong Huygens principle
provides a more localized representation of the free wave
propagators; see \cite{Str05}.

For \(V\neq0\) and \(n\ge3\), early systematic \(L^p\) and
\(L^\infty\)-decay estimates were obtained by Beals--Strauss
\cite{B-S} and Beals \cite{Beals}; see also \cite{GV03,DP}.
In dimension two, comparatively fewer results are available for
\(W^{k,1}\to L^\infty\) or regularized \(L^1\to L^\infty\)
estimates. Moulin \cite{Mou} proved high-frequency dispersive
estimates for the wave and Schr\"odinger groups. Assuming that zero
is regular, Kopylova \cite{Kop} obtained the decay rate
\(t^{-1}(\log t)^{-2}\) in polynomially weighted energy spaces.
Green \cite{Green14} established regularized \(L^1\to L^\infty\)
bounds with decay \(|t|^{-1/2}\), proved the weighted improvement
\(|t|^{-1}(\log|t|)^{-2}\) in the regular case, and analyzed
zero-energy obstructions. Beceanu \cite{Bec16} subsequently established pointwise, weighted,
and Strichartz estimates for almost scaling-critical potentials,
assuming that zero is regular.


In higher dimensions, Cardoso--Vodev \cite{CV} obtained optimal
dispersive estimates for \(4\le n\le7\), while
Erdo\u{g}an--Goldberg--Green \cite{EGG14} studied the influence of
zero-energy obstructions on the four-dimensional Schr\"odinger and
wave evolutions. Related Strichartz, smoothing, and maximal estimates
for rough, inverse-square, and critically decaying potentials can be
found in \cite{Bec-Gold,GV95,BPSZ03,BPSZ04}.






\subsubsection{Higher-order Schr\"odinger operators.}
Recent progress has been made on time decay estimates for $e^{-itH}$ with $H=\Delta^2+V$. Feng--Soffer--Yao \cite{FSY} 
established Jensen--Kato  estimates  $O(|t|^{-n/4})$ for \(n\ge5\) and  the  long-time \(L^1\to L^\infty\)
estimate  $O(|t|^{-1/2})$  for $n=3$ when zero is a regular point. Subsequently, Erdo\u{g}an--Green--Toprak \cite{EGT19} and Green--Toprak \cite{GT19} showed the $L^1\to L^\infty$ estimates  $O(|t|^{-n/4})$  for $n=3$ and $n=4$, when zero is regular or a first-kind resonance, and observed rate changes for other zero-energy resonances.

For $n=1$, Soffer--Wu--Yao \cite{SWY22} proved $O(|t|^{-1/4})$ estimates regardless of zero resonance types. For $n=2$, Li--Soffer--Yao \cite{LSY21} obtained optimal $O(|t|^{-1/2})$ decay in the regular and first-kind resonance cases. More generally, for \(H=(-\Delta)^m+V\), \(m>1\), and
\(2m<n<4m\), Erdo\u{g}an--Goldberg--Green
\cite{EGG_ArXiv} obtained the optimal
\(O(|t|^{-n/(2m)})\) \(L^1\to L^\infty\) bound for
scaling-critical potentials under suitable spectral assumptions.




Substantial advances have also been made on $L^p$-boundedness of wave operators for $H=(-\Delta)^m+V$ with $m>1$: see \cite{GG21b, MWY22, MWY23, MWY_ArXiv23_2} for $n<2m$ with $m=2$; \cite{Yajima_2024JST} for $n=2m$ with $m=2$; \cite{EG21, Erdogan-Green23, EGG23, EGL} for $n>2m$ with $m\in\mathbb{N}^+$. Recently, Cheng et al. \cite{CSWY25} established sharp boundedness of wave operators for all zero-resonance singularities in odd dimensions $1\le n\le 4m-1$ with $m\ge2$.
\subsection{The ideas of the proof}
In this subsection, we outline the main ideas underlying the proofs of our main theorems.
The starting point is the Stone formulas:
\begin{align*}
	\cos (t \sqrt{H}) P_{\mathrm{ac}}(H) & =\frac{2}{\pi i} \int_0^{\infty} \lambda^3 \cos (t \lambda^2)\left[R_V^{+}(\lambda^4)-R_V^{-}(\lambda^4)\right]  d \lambda, \\
	\frac{\sin (t \sqrt{H})}{t\sqrt{H}} P_{\mathrm{ac}}(H) & =\frac{2}{\pi i t} \int_0^{\infty} \lambda \sin (t \lambda^2)\left[R_V^{+}(\lambda^4)-R_V^{-}(\lambda^4)\right] d \lambda.
\end{align*}
To make these identities precise, we recall that $R_0(z)=(\Delta^2-z)^{-1}$ for $z\in \mathbb{C}\setminus[0,\infty)$ and $R_V(z)=(H-z)^{-1}$ for $z\in \mathbb{C}\setminus\sigma(H)$ denote the resolvents of the free operator $\Delta^2$ and the perturbed operator $H=\Delta^2+V$ respectively.
Their boundary values (limiting resolvents) on $(0,\infty)$ are defined as
\begin{align*}
	R^\pm_0(\lambda)=\lim_{\epsilon \searrow 0}R_0(\lambda\pm i\epsilon ),\qquad  R_V^\pm(\lambda)=\lim_{\epsilon  \searrow 0}R_V(\lambda\pm i\epsilon ),\quad \lambda>0.
\end{align*}
The existence of $R^\pm_0(\lambda)$ as bounded operators from $L^{2,s}(\mathbb{R}^2)$ to $L^{2,-s}(\mathbb{R}^2)$ for any $s>1/2$ follows from the limiting absorption principle for the resolvent $(-\Delta-z)^{-1}$ (see, e.g., Agmon \cite{Agmon}) and the splitting identity
$$
R_0(z)=\frac{1}{2\sqrt z}\left[(-\Delta-\sqrt z)^{-1}-(-\Delta+\sqrt z)^{-1}\right],\quad z\in \mathbb{C}\setminus[0,\infty), \ \Im\sqrt z>0.
$$
The existence of $R^\pm_V(\lambda)$ under suitable potential decay conditions is established in the existing literature (see, e.g., \cite{FSY}).

To evaluate the Stone formulas, we decompose the integrals into a low-energy part ($0 < \lambda \ll 1$) and a high-energy part ($\lambda \gtrsim 1$). The high-energy part is straightforward to treat since the free resolvent (see \eqref{reso-big} and \eqref{reso_small}) possesses no singularities for $\lambda \gtrsim 1$. The primary analytical difficulty lies in the low-energy part, which requires a delicate analysis of the singularities of $R_V^\pm(\lambda^4)$ near $\lambda=0$.

Setting $v(x) = \sqrt{|V(x)|}$ and $U(x) = \operatorname{sgn}(V(x))$, we define the Birman-Schwinger operator $M^\pm(\lambda) = U + vR_0^\pm(\lambda^4)v$. The symmetric resolvent identity,
\begin{equation}\label{resolvent identity}
	R_V^\pm(\lambda^4) = R_0^\pm(\lambda^4) - R_0^\pm(\lambda^4) v \big(M^\pm(\lambda)\big)^{-1} v R_0^\pm(\lambda^4),
\end{equation}
reduces the problem to analyzing the asymptotic expansion of $(M^{\pm}(\lambda))^{-1}$ as $\lambda \to 0^+$.

If zero is a resonance of the $\mathbf{k}$-th kind ($0\leq\mathbf{k}\leq 4$), Theorem~\ref{thm:M_inverse} guarantees that $(M^\pm(\lambda))^{-1}$ admits the following expansion:
\begin{equation}\label{asymptotic expansion}
	\bigl(M^{\pm}(\lambda)\bigr)^{-1}= \sum_{0\leq\alpha,\beta\leq\mathbf{k}+2}
	\lambda^{2-k_\alpha-k_\beta} Q_\alpha \mathcal{M}_{\alpha,\beta}^{\pm}(\lambda) Q_\beta,
\end{equation}
where the singularity exponents $k_\alpha $ are given by \eqref{tab:k_alpha}. The operators $\mathcal{M}_{\alpha,\beta}^{\pm}(\lambda)\in\mathbb{B}(L^2)$ exhibit additional decay as $\lambda \to 0^+$ for certain indices $\alpha, \beta$. The orthogonal projections $Q_\alpha$ satisfy the important cancellation properties (see Remark~\ref{cancellationQ}), which are applied to  establish Lemma~\ref{lemma_projection} to analyze the expansion for  $Q_\alpha v R_0^\pm(\lambda^4)$. Lemma~\ref{lemma_projection} demonstrates how the projections extract compensatory positive powers of $\lambda$, effectively annihilating the zero-energy singularities of the free resolvent.

Substituting the identities \eqref{resolvent identity} and \eqref{asymptotic expansion} into the low-energy part of the Stone formulas, we reduce the analysis to bounding a finite number of kernels
\begin{equation*}
	\big(\mathcal{K}_{B}^+ - \mathcal{K}_{B}^-\big)(t,x,y) \quad \text{and} \quad \big(\mathcal{N}_{B}^+ - \mathcal{N}_{B}^-\big)(t,x,y) ,
\end{equation*}
where the component kernels are given by
\begin{align*}
	\mathcal{K}_{B}^\pm(t,x,y) &= \frac{2}{\pi i}\int_0^{\infty} \cos(t \lambda^2)\lambda \widetilde{\chi}_1(\lambda) \left[\lambda^{4-k_\alpha-k_\beta} R_0^\pm(\lambda^{4})v Q_\alpha B^\pm(\lambda)Q_\beta vR_0^\pm(\lambda^{4})\right](x,y) \, d \lambda, \\
	\mathcal{N}_{B}^\pm(t,x,y) &= \frac{2}{\pi it}\int_0^{\infty} \sin(t\lambda^2)\lambda^{-1}\widetilde{\chi}_1(\lambda)
	\left[\lambda^{4-k_\alpha-k_\beta} R_0^\pm(\lambda^{4})v Q_\alpha B^\pm(\lambda) Q_\beta vR_0^\pm(\lambda^{4})\right](x,y) \, d \lambda,
\end{align*}
with $B^\pm(\lambda) \in \{ \mathcal{M}_{\alpha, \beta}^{\pm}(\lambda) \mid 0 \le \alpha,\beta \le \mathbf{k}+2 \}$. By applying the expansion of $(Q_\alpha v R_0^\pm(\lambda^4))(x,y)$ established in Lemma~\ref{lemma_projection}, the intricate integrals $\mathcal{K}_{B}^\pm$ and $\mathcal{N}_{B}^\pm$ can be reduced to canonical oscillatory forms systematically studied in Lemma~\ref{oscillatory}.
 Consequently, the desired temporal decay estimates are ultimately derived by invoking the  bounds from Lemma~\ref{oscillatory}.

To illustrate our general strategy for controlling these oscillatory integrals, let us examine a representative term where $B^\pm(\lambda)=\mathcal{M}_{1,1}^{\pm}(\lambda)$.
Using the Taylor expansion for $F_\pm$ and the cancellation property $Q_1 v = 0$, we extract an additional positive power of $\lambda$ (see Lemma~\ref{lemma_projection}):
$$
\bigl(Q_1 v R_0^\pm(\lambda^4)\bigr)(x,y)
= -\lambda^{-1} Q_1\Bigl(|\cdot| v \int_0^1 F_\pm'(\lambda|\theta\cdot-y|)\cos\xi\,d\theta\Bigr)(x)
=: \lambda^{-1} \Omega_{1,\pm}(\lambda,x,y),
$$
where $R^\pm_0(\lambda^4)(x,y)=\lambda^{-2}F_\pm(\lambda|x-y|)$ (see \eqref{def:F_pm}). Hence, the term $Q_1 v R_0^\pm(\lambda^4)$ is less singular than the free resolvent. Furthermore, we derive the  bound
\[
\bigl\|\partial_\lambda^\ell\bigl(e^{\mp i\lambda|y|} \Omega_{1,\pm}(\lambda,\cdot,y)\bigr)\bigr\|_{L^2}
\lesssim \langle \lambda y\rangle^{-1/2} \lambda^{-\ell}, \qquad \ell=0,1,2.
\]
With these properties, the kernels can be recast into the canonical forms appearing in Lemma~\ref{oscillatory}:
\begin{align*}
	\mathcal{K}^\pm(t,x,y) &= \frac{2}{\pi i}\int_0^{\infty} \cos (t \lambda^2)e^{\pm i\lambda r} \lambda \mathcal{E}^\pm(\lambda, x,y) \, d\lambda, \\
	\mathcal{N}^\pm(t,x,y) &= \frac{2}{\pi i t}\int_0^{\infty} \sin(t\lambda^2) e^{\pm i\lambda r} \lambda^{-1} \mathcal{E}^\pm(\lambda, x,y) \, d\lambda,
\end{align*}
with  $r(x,y) = |x| + |y|$ and the effective amplitude factor $$\mathcal{E}^\pm(\lambda,x,y) = e^{\mp i\lambda r} \widetilde{\chi}_1(\lambda) \langle \mathcal{M}_{1,1}^{\pm}(\lambda) \Omega_{1,\pm}(\cdot, y), \Omega_{1,\mp}(\cdot, x) \rangle.$$ Applying Leibniz's rule along with the bound $\|\partial_\lambda^{\ell} \mathcal{M}_{1,1}^{\pm}(\lambda)\|_{\mathbb{B}(L^2)}\lesssim\lambda^{-\ell}$ yields
\[
\bigl|\partial_\lambda^{\ell} \mathcal{E}^\pm(\lambda,x,y) \bigr| \lesssim \bigl\langle \lambda (|x|+|y|)\bigr\rangle^{-1/2} \lambda^{-\ell} \widetilde{\chi}_1(\lambda/2),  \qquad  \ell=0,1,2.
\]
Then $\mathcal{E}^\pm$ satisfies condition \eqref{condit 2} of Lemma~\ref{oscillatory} (ii) with parameters $(\sigma, \nu)=(0,0)$ and $ h(x,y)=1$, yielding  the bound $|\mathcal{K}^\pm|+|\mathcal{N}^\pm |\lesssim \langle t \rangle^{-1}$ uniformly in $x, y$.


To establish the spatially weighted estimates for $|t| \ge 2$, we utilize the following expansion:
\begin{align*}
	R^\pm_0(\lambda^4)(x,y) &= \frac{a_0^\pm}{\lambda^2} + \bigl(a_1^\pm+b_0\log\lambda\bigr)|x-y|^2 + b_0|x-y|^2\log|x-y| + {O}_2\big(\lambda^2|x-y|^4\big) \nonumber \\
	&=: \frac{a_0^\pm}{\lambda^2} + E_1^\pm(\lambda, x, y).
\end{align*}
The cancellation property $Q_1 v = 0$ annihilates the highly singular leading term $\lambda^{-2}$, meaning $\bigl( Q_1 v R_0^{\pm}(\lambda^4) \bigr) = \bigl( Q_1 v E_1^{\pm} \bigr)$. This cancellation mitigates the low-energy singularity by a factor of $\lambda^{2-}$ relative to $R_0^\pm$ but introduces  spatial growth.  Provided that the potential $V$ decays sufficiently fast, we deduce the  bound:
\[
\|\partial_\lambda^\ell  \bigl(Q_1 v R_0^{\pm}(\lambda^4)\bigr)(\cdot,y)\|_{L^2}
\lesssim \|v \partial_\lambda^\ell E_1^\pm(\lambda,\cdot,y)\|_{L^2}
\lesssim \langle y\rangle^4\lambda^{-\ell-}, \quad \ell = 0, 1, 2.
\]
In this weighted scenario, we align the phase to $r(x,y)=0$ and define the amplitude factor as
$$
\mathcal{E}^\pm(\lambda,x,y)
= \widetilde{\chi}_1(\lambda) \lambda^{2}
\bigl\langle \mathcal{M}_{1,1}^\pm(\lambda) \bigl( Q_1 v E_1^{\pm} \bigr)(\cdot,y),
\bigl( Q_1 v E_1^{\mp} \bigr)(\cdot,x) \bigr\rangle.
$$
Consequently, we derive the weighted bound
\[
\bigl| \partial_\lambda^\ell \mathcal{E}^\pm(\lambda,x,y) \bigr|
\lesssim \lambda^{1/2-\ell}\widetilde{\chi}_1(\lambda/2) \langle x\rangle^{4} \langle y\rangle^{4}, \quad   \ell=0,1,2.
\]
An application of Lemma~\ref{oscillatory}(ii) with parameters $(\sigma, \nu)=(-1/2,0)$ and  $h(x,y)=\langle x\rangle^{4} \langle y\rangle^{4}$ yields the improved temporal decay $|\mathcal{K}^\pm|+|\mathcal{N}^\pm |\lesssim |t|^{-5/4}\langle x \rangle^4 \langle y \rangle^4$. Combining this with the uniform $O(|t|^{-1})$ bound, we  obtain the logarithmically weighted bound:
$$
\min\left(\frac{1}{|t|}, \frac{\langle x\rangle^4\langle y\rangle^4}{|t|^{5/4}}\right) = \frac{1}{|t|} \min\left(1, \frac{\langle x\rangle^4\langle y\rangle^4}{|t|^{1/4}}\right) \lesssim \frac{\omega(x)\omega(y)}{|t|\log|t|},
$$
which gives the desired result $\omega(x)\omega(y)(|t|\log|t|)^{-1}$.

For the remaining operators $B^\pm(\lambda) \in \{\mathcal{M}_{\alpha,\beta}^{\pm}(\lambda) \mid 0 \le \alpha,\beta \le \mathbf{k}+2\}$, the essential strategy remains unchanged. By exploiting the structural bounds of $Q_\alpha v R_0^{\pm}(\lambda^4)$ in conjunction with the operator norm bounds for $\mathcal{M}_{\alpha, \beta}^{\pm}(\lambda)$, the analysis is reduced to estimating canonical oscillatory integrals in Lemma~\ref{oscillatory}.

However, the analysis becomes more intricate and subtle when higher-order resonances occur ($\mathbf{k} \ge 2$).   To achieve the sharp   estimates for $\mathbf{k} \ge 2$,  one must conduct a  deeper investigation into the expansion structures of $\mathcal{M}_{\alpha,\beta}^{\pm}(\lambda)$ and  show that certain intricate  operators do not  vanish (see Propositions \ref{sine-sharp}, \ref{sine-sharp-2}, \ref{2-sine-sharp-2}, \ref{sharp-3-4} and Lemma \ref{lemma:non-zero_1}). Despite these  technical challenges, our unified framework  resolves all threshold singularities, thereby providing a complete picture of the time-decay behavior for the two-dimensional plate equation across all resonance scenarios.

\begin{comment}
\subsection{Organization of the paper}
The remainder of this paper is organized as follows.
Section~\ref{sec:reso-expan} studies the free and perturbed resolvent expansions near zero energy. This section also provides the equivalent characterizations of zero resonances and introduces two key lemmas (Lemmas~\ref{lemma_projection} and \ref{oscillatory}).

Sections~\ref{sec:free_case} and \ref{sec:high-energy} establish decay estimates for the free case and the high-energy part, respectively.

Sections~\ref{sec:regular_first} through \ref{sec:third_fourth} constitute the  low-energy analysis, systematically treating the regular case and the first (Section~\ref{sec:regular_first}), second (Section~\ref{sec:second}), and third kind resonance/eigenvalue cases (Section~\ref{sec:third_fourth}).

Finally, Sections~\ref{sec:proof_characterization} through \ref{section 8} are reserved for the deferred proofs concerning the resonance characterizations, the expansions of $(M^\pm(\lambda))^{-1}$, and the oscillatory integral estimates.
\end{comment}



\section{Resolvent expansions}\label{sec:reso-expan}
\subsection{Free resolvent  }\label{subsec:free_resolvent}
In this subsection, we derive the asymptotic expansions of the free resolvent $R_0^\pm(\lambda^4)$ for the bi-Laplace operator $\Delta^2$.

Recall the expression of the free resolvent for the Laplacian (see, e.g., \cite{Erdogan-Green}):
\begin{equation}\label{resolvent-R}
	R^\pm (-\Delta;\lambda^2)(x,y) = \pm\frac{i}{4}H_0^\pm(\lambda|x-y|),
\end{equation}
where $H_0^\pm (z)= J_0(z)\pm i Y_0(z) $ are Hankel functions of order zero. From the series expansions for Bessel functions
(see, e.g., \cite{AS64}), we have
\begin{gather*}
    J_0(z)=1-\frac14z^2+\frac1{64}z^4-\frac1{2304}z^6+ \frac{1}{147456}z^8+O(z^{10}),\\
    Y_0(z)=\frac{2}{\pi}\bigl(\log\frac{z}{2}+\gamma\bigr)J_0(z)
          +\frac{2}{\pi}\Bigl(\frac14z^2-\frac3{128}z^4+\frac{11}{13824}z^6-\frac{25}{1769472}z^8+O(z^{10})\Bigr),
\end{gather*}
where $\gamma $ is Euler's constant. In addition, for $|z| \gtrsim 1$,
\begin{equation}\label{id-H0big}
H^\pm_0(z)=e^{\pm iz}w_\pm(z),\qquad |w_\pm^{(\ell)}(z)|\lesssim(1+|z|)^{-1/2-\ell}, \
\ell= 0,1,2,\cdots.
\end{equation}

Note that
\begin{equation*}
    R^\pm_0(\lambda^4) = \frac{1}{2\lambda^2} \left( R^\pm(-\Delta; \lambda^2) - R(-\Delta; -\lambda^2) \right), \quad \lambda > 0.
\end{equation*}
Then by \eqref{resolvent-R} and \eqref{id-H0big}, it follows that
\begin{equation}\label{reso-big}
    R_0^\pm(\lambda^4)(x,y) = \frac{i}{8\lambda^2} \left[ \pm e^{\pm i\lambda |x - y|} w_\pm(\lambda|x - y|) - e^{-\lambda |x - y|} w_+(i\lambda|x - y|) \right],\ \ \  \lambda|x - y| \gtrsim 1.
\end{equation}
 For $\lambda|x-y| \ll 1 $,  the series expansions of Bessel functions give
   \begin{equation}\label{reso_small}
   \begin{split}
    R^\pm_0(\lambda^4)(x,y)=&\frac{a_0^\pm}{\lambda^2}G_0(x,y)+g_0^\pm(\lambda)G_2(x,y)+b_0G_2^0(x,y)+a_2^\pm\lambda^2G_4(x,y)\\
&+g_1^\pm(\lambda)\lambda^4G_6(x,y)+b_1\lambda^4G_6^0(x,y)+a_4^\pm\lambda^6G_8(x,y)
    +O_2\bigl(\lambda^{8-\varepsilon}|x-y|^{10-\varepsilon}\bigr),
    \end{split}
    \end{equation}
where
$ b_0 = \frac{1}{8\pi}, \ b_1 = \frac{1}{4608\pi},\
a_0^\pm = \pm \frac{i}{8}, \  
a_2^\pm = \pm \frac{i}{512}, \      
a_4^\pm=\pm \frac{i}{1179648},$ and
$$
g_0^\pm(\lambda)=b_0\log\lambda+a_1^\pm, \ g_1^\pm(\lambda)=b_1\log\lambda+a_3^\pm,\ \ a_1^\pm, a_3^\pm\in\mathbb{C}\setminus\mathbb{R}, $$
\begin{equation}\label{def:G_k}
   G_{k}(x,y)= |x-y|^k, \,
   G_k^0(x,y)= |x-y|^k\log|x-y|.
\end{equation}
Combining \eqref{reso-big} and \eqref{reso_small}, we obtain the following expansion of $R^\pm_0(\lambda^4)(x,y)$ for any $\lambda > 0$:
\begin{lemma}[{\cite[Lemma 2.2]{LSY21}}]\label{lem-reso}
For any $\lambda> 0 $, we have the following expansions:
\begin{equation}\label{free-R0pmlambda4}
   \begin{split}
	  R^\pm_0(\lambda^4)(x,y)=&\frac{a_0^\pm}{\lambda^2}G_0(x,y)+g_0^\pm(\lambda)G_2(x,y)+b_0G_2^0(x,y)+a_2^\pm\lambda^2G_4(x,y)\\
&+g_1^\pm(\lambda)\lambda^4G_6(x,y)+b_1\lambda^4G_6^0(x,y)+O_2\bigl(\lambda^{6}|x-y|^{8}\bigr).\nonumber
		\end{split}
	\end{equation}
\end{lemma}
For convenience in the subsequent analysis, we define
\begin{equation}\label{def:F_pm}
R^\pm_0(\lambda^4)(x,y):=\lambda^{-2}F_\pm(\lambda|x-y|),\quad
\widetilde{R}^\pm(\lambda|x-y|):=\lambda^{2}e^{\mp i\lambda|x-y|}R^\pm_0(\lambda^4)(x,y).
\end{equation}
Then $F_\pm(p)=e^{\pm ip}\widetilde{R}^\pm(p)$. From \eqref{reso-big} and \eqref{reso_small},    we have
\begin{equation}\label{free-widetilde_R}
\begin{aligned}
\widetilde{R}^\pm(p)=&\,e^{\mp ip}\Bigl[a_0^\pm+b_0p^2\log p+a_1^\pm p^2
                     +a_2^\pm p^4+b_1p^6\log p+a_3^\pm p^6+O_2\bigl(p^{8}\bigr)\Bigr]\chi_1(p)\\
                     &+\frac{i}{8}\Bigl(\pm w_\pm(p)-e^{\mp ip}e^{-p}w_+(ip)\Bigr)\chi_2(p).
\end{aligned}
\end{equation}

%{
%\color{red} estimate of $\widetilde{R}^\pm(p)$??
%}

In the sequel,  we denote by $G_k^0$ (resp. $G_k$) the integral operator with kernel $G_k^0(x,y)$ (resp. $G_k(x,y)$) defined in \eqref{def:G_k}.
In particular, $b_0 G_2^0$ is a fundamental solution of $\Delta^2$, i.e., $\Delta^2 (b_0 G_2^0) = \delta$ in the sense of distributions. Consequently, on the orthogonal complement of $\ker(\Delta^2)$, we may identify $b_0 G_2^0$ with $(\Delta^2)^{-1}$.


\subsection{Asymptotic expansion of $R_V^\pm(\lambda^4)$ near zero}\label{subsec:Q}  
This subsection is devoted to the asymptotic expansion of the resolvent $R_V^\pm(\lambda^4)$ as $\lambda\to 0^+$. We begin by introducing several key projection subspaces of $L^2(\mathbb{R}^2)$.  

\begin{definition}\label{defS}
	Let $x = (x_1, x_2) \in \mathbb{R}^2$, $v=\sqrt{|V|}$.  Define $P = \|V\|_{L^1(\mathbb{R}^2)}^{-1} \langle \cdot, v\rangle v$ and $T_0 : = U + b_0 v G_2^0 v $. Set  $S_{-1} = \Id$, $S_6 = 0$, and for $j=0,\dots,5$, let $S_j$ be the orthogonal projection onto the subspace $S_jL^2$ defined as follows:
	\begin{itemize}
			\item $S_0L^2 = \{ f \in L^2 \mid \langle  v, f \rangle = 0\}$.
		\item $S_1L^2 = \{ f \in S_0L^2 \mid \langle x_j v, f \rangle = 0,\ j = 1,2 \}$.
		\vskip0.1cm
		\item $S_2L^2 = \{ f \in S_1L^2 \mid S_1T_0f = 0 \}$.
		\vskip0.1cm
		\item $S_3L^2 = \{ f \in S_2L^2 \mid S_0T_0f = 0 \}$.
		\vskip0.1cm
		\item $S_4^0L^2 = \{ f \in S_3L^2 \mid \langle |x|^2 v, f \rangle = 0 \}$.
        \vskip0.1cm
		\item $S_4^1L^2 = \{ f \in S_3L^2 \mid \langle x_ix_j v, f \rangle = 0,\ i,j = 1,2\}$.
		\vskip0.1cm
		\item $S_4L^2 = \{ f \in S_3L^2 \mid \langle x_ix_j v, f \rangle = 0,\ i,j = 1,2,\ \text{and}\ PT_0f = 0 \}$.
		\vskip0.1cm
		\item $S_5L^2 = \{ f \in S_4L^2 \mid \langle x_ix_jx_k v, f \rangle = 0,\ i,j,k = 1,2 \}$.
	\end{itemize}
\end{definition}

Note that these projections $S_j$ are well-defined due to the decay condition \eqref{condition} on $V$.  Since  $vG_2^0 v$ is a Hilbert-Schmidt operator and
$$
 T_0  = U + b_0 v G_2^0 v
$$
is a compact perturbation of $U$,  Fredholm alternative then ensures that $S_2$ is a finite-rank projection. Moreover, observe that $P+S_0=\Id$, then it follows that $S_jT_0=T_0S_j=0$ for $j=4,5.$

From the definitions above, we have the chain of inclusions
\[
S_0 L^2 \supseteq S_1 L^2 \supseteq S_2 L^2 \supseteq S_3 L^2 \supseteq S_4^0 L^2  \supseteq S_4^1 L^2 \supseteq S_4 L^2 \supseteq S_5 L^2,
\]
and the projections $S_j$ (for $j=2,\cdots,5$) and $S_4^0, S_4^1$ are finite-rank operators.

We now introduce the following complementary family of orthogonal projections.
\begin{definition}\label{defQ}
Define
$$
	Q_j := S_{j-1} - S_j \quad \text{for } j = 0,1,\dots,6, \quad \text{and} \quad Q_4^0 = S_3 - S_4^0, \quad Q_4^1 = S_4^0 - S_4.
$$
and further decompose $Q_4^1$ into the sum of the following projectors:
$$
Q_{41}^1 := S_4^0 - S_4^1, \qquad Q_{42}^1 := S_4^1 - S_4. 
$$
\end{definition}
In particular, we observe that $Q_0=P.$
According to Definitions \ref{defS} and \ref{defQ}, we obtain the following cancellation properties:
\begin{remark}\label{cancellationQ}
	The orthogonal projections $Q_4^0$, $Q_4^1$, and $Q_\alpha$ (for $\alpha \in \mathbb{N}_0$ with $1 \le \alpha \le 6$) satisfy:
	\begin{align*}
		& Q_{1}v = 0 \quad (\text{i.e., } \langle v, Q_{1}f \rangle = 0,\ \forall f \in L^2), \\
		& Q_{\alpha}v = Q_{\alpha}(x_jv) =Q_{4}^{0}v = Q_{4}^{0}(x_jv) =  0 \quad (\alpha = 2,3,4), \\
		& Q_{4}^{1}v = Q_{4}^{1}(x_jv )= Q_{4}^{1}(|x|^2v )= Q_{41}^{1}v = Q_{41}^{1}(x_jv )= Q_{41}^{1}(|x|^2v )= 0,   \\
        & Q_{42}^{1}v = Q_{42}^{1}(x_jv )= Q_{42}^{1}(x_ix_jv )=Q_{5}v = Q_{5}(x_jv) = Q_{5}(x_ix_jv) = 0,\quad Q_5T_0=T_0Q_5=0,\\
		& Q_{6}v = Q_{6}(x_jv) = Q_{6}(x_ix_jv) = Q_{6}(x_ix_jx_kv) = 0,\ Q_6T_0=T_0Q_6=0.
	\end{align*}
\end{remark}

 Besides Definition \ref{definition1}, zero-energy resonances of $H$ can also be characterized by the orthogonal projections $S_j$; the proof is given in Section \ref{sec:proof_characterization}. Recall that $\mathbf{k} \in \{0,1,2,3,4\}$  denotes  the zero-energy spectral types: $\mathbf{k}=0$ for a regular point, $\mathbf{k} \in \{1,2,3\}$ for zero resonances, and $\mathbf{k}=4$ for a zero eigenvalue. 
Throughout this paper, the decay index $\mu$ of the potential $V$ is required to satisfy 
\begin{equation}\label{condition}
	\mu> \begin{cases}11, & \mathbf{k}=0,1, \\
14, &  \mathbf{k}=2,\\
18, & \mathbf{k}=3,4. \end{cases}
		\end{equation}
\begin{proposition}\label{characterizations}
Assume that $\left|V(x)\right| \lesssim \left \langle  x\right \rangle^{-\mu}$ for some $\mu>0$
satisfying the condition \eqref{condition}.
Then 
\begin{itemize}
\item[(i)] Zero  is a regular point of $H$  $( \text{i.e.,}\  \mathbf{k}=0 )$  if and only if $S_{2}=0$;
\item[(ii)] Zero is a $\mathbf{k}$-th kind resonance with $1\leq\mathbf{k}\le 4$ if and only if  $S_{\mathbf{k}+1} \neq 0$ and $S_{\mathbf{k}+2}=0$.
\end{itemize}
\end{proposition}


The relationship between the orthogonal projections $Q_j$ and $S_j$, as given in Definition \ref{defQ}, yields the following equivalent characterizations in terms of these $\{Q_j\}$:

\begin{proposition}\label{characterizations Q}
	Assume that $|V(x)| \lesssim \langle x \rangle^{-\mu}$ for some $\mu > 0$ satisfying the condition \eqref{condition}. Then
	\begin{itemize}
		\item[(i)] Zero is a regular point of $H$ $(i.e., \mathbf{k} = 0) $ if and only if
		\[
		Q_0 + Q_1 + Q_2 = \Id.
		\]
		\item[(ii)]  Zero is a $\mathbf{k}$-th kind resonance of $H$ with $1 \le \mathbf{k} \le 4$, if and only if
		\[
		\sum_{j=0}^{\mathbf{k}+1} Q_j \neq \Id \quad \text{and} \quad \sum_{j=0}^{\mathbf{k}+2} Q_j = \Id.
		\]
	\end{itemize}
\end{proposition}
Next we turn to study the asymptotic expansions of $R^\pm_V(\lambda^4)$ near $\lambda=0.$
Let $U(x)=\operatorname{sgn}(V(x))$ and $v(x)=|V(x)|^{1/2}$, then we have $ V=Uv^2$ and
the following symmetric resolvent identity:
 \begin{equation}\label{id-RV}
R^\pm_V(\lambda^4) = R_0^\pm(\lambda^4) -R^\pm_0(\lambda^4)v(M^\pm(\lambda))^{-1}vR^\pm_0(\lambda^4),
\end{equation}
where $M^{\pm}(\lambda)=U+ vR^\pm_0(\lambda^4)v$. Hence, we need to establish the expansions of
$(M^\pm(\lambda))^{-1}$.

Recall that  $T_0=U+b_0 v G_2^0 v$.
Using the expansions of $R^\pm_0(\lambda^4)(x,y)$ from Lemma \ref{lem-reso}, we obtain the following asymptotic expansions of $M^\pm(\lambda)$:
\begin{itemize}
\item If $\mu > 11$, then
\begin{equation}\label{Mpm-1}
\begin{split}
M^\pm(\lambda) = \frac{a_0^\pm}{\lambda^2}v G_{0} v + g_0^\pm(\lambda)v G_{2} v + T_0 + \Gamma_2^\pm (\lambda);
\end{split}
\end{equation}

\item If $\mu > 14$, then
\begin{equation}\label{Mpm-2}
\begin{split}
M^\pm(\lambda) = \frac{a_0^\pm}{\lambda^2}v G_{0} v + g_0^\pm(\lambda)v G_{2} v + T_0 + a_2^\pm\lambda^2 vG_4 v + \Gamma_{4-\epsilon}^\pm (\lambda);
\end{split}
\end{equation}

\item
If $\mu > 18$, then
\begin{equation}\label{Mpm-3}
\begin{split}
M^\pm(\lambda) = \frac{a_0^\pm}{\lambda^2}v G_{0} v + g_0^\pm(\lambda)v G_{2} v + T_0 + a_2^\pm\lambda^2 vG_4 v + g_1^\pm(\lambda)\lambda^4vG_6v + b_1\lambda^4 vG_6^0v  + \Gamma_6^\pm (\lambda).
\end{split}
\end{equation}
\end{itemize}
Here, $\Gamma_s^\pm (\lambda)$ denote $\lambda$-dependent Hilbert-Schmidt operators satisfying
\[
\big\| \Gamma_s^\pm (\lambda) \big\|_{\mathbb{B}(L^2)} +\lambda \big\| \partial_\lambda \Gamma_s^\pm (\lambda) \big\|_{\mathbb{B}(L^2)}+\lambda^2 \big\| \partial_\lambda^2 \Gamma_s^\pm (\lambda) \big\|_{\mathbb{B}(L^2)}
\lesssim \lambda^{s}.
\]

\begin{comment}
We now provide the asymptotic expansion of $(M^\pm(\lambda))^{-1}$ as $\lambda \to 0^+$; the proof can be seen in Section \ref{subsec:proof_inverse}. 

\begin{theorem}\label{thm:M_inverse}
	Assume that zero is a $\mathbf{k}$-th kind resonance with $0 \leq \mathbf{k} \leq 4$, and $|V(x)| \lesssim \left \langle  x\right \rangle^{-\mu}$ for some $\mu>0$ satisfying the condition \eqref{condition}.
	Then there exists $0 < \lambda_0 \ll 1$ such that for all $0 < \lambda \le \lambda_0$, the operator $\left(M^{\pm}(\lambda)\right)^{-1}$ on $L^2(\mathbb{R}^2)$ admits the asymptotic expansion:
	\begin{equation}\label{eq:M_inverse}  
		\left(M^{\pm}(\lambda)\right)^{-1} = \sum_{0 \leq \alpha, \beta \leq \mathbf{k}+2} \lambda^{2-k_\alpha-k_\beta} Q_\alpha \mathcal{M}_{\alpha, \beta}^{\pm}(\lambda) Q_\beta ,  
	\end{equation}  
	where the values of the parameter $k_\alpha$ are assigned as follows:  
	\begin{equation}\label{tab:k_alpha}  
		\begin{array}{c|ccccccc}  
			\alpha & 0 & 1 & 2 & 3 & 4 & 5 & 6 \\  
			\hline  
			k_\alpha & 0 & 1 & 1 & 1 & 2 & 3 & 3  
		\end{array}  
	\end{equation}  
	Moreover, the  operators $\mathcal{M}_{\alpha, \beta}^{\pm}(\lambda)$ possess the following properties:  
	
	\vskip0.5cm  
	{\rm (i)} \textbf{For all $0 \le \alpha, \beta \le 3$ and $\ell=0,1,2$},
\[
\Big\| \partial_\lambda^\ell \mathcal{M}_{\alpha,\beta}^\pm(\lambda) \Big\|_{\mathbb{B}(L^2)} \lesssim
\begin{cases}
	\lambda^{-\ell}, & \text{for } ( \alpha,\beta) \neq (3,3), \\
	|\log \lambda| \lambda^{-\ell}, &  \text{for } (\alpha,\beta) = (3,3).
\end{cases}
\]
	In particular, for  $(\alpha, \beta) = (0,0)$, 
	 under the additional condition $Q_4^0 = 0$,  the following refined asymptotic expansion holds:
		\begin{align}\label{D00-1-1}  
\mathcal{M}_{0,0}^\pm(\lambda) = \mathcal{D}_{0,0}^\pm + \lambda (-\log \lambda)^{3/2}\Gamma_{0,0}^{\pm}(\lambda),  
		\end{align}		  
		where $\big\|\partial_\lambda^{\ell} \Gamma_{0,0}^{\pm}(\lambda) \big\|_{\mathbb{B}(L^2)} \lesssim \lambda^{-\ell}$ for $\ell = 0,1,2$, and the leading term is given by  
		\begin{equation}\label{Doo}  
			\mathcal{D}_{0,0}^{\pm} = (a_0^\pm)^{-1} \|V\|_{L^1}^{-1}Q_0 + (a_0^\pm)^{-2}\mathcal{W} (\boldsymbol{d}^\pm)^{-1} \mathcal{W}^* := (a_0^\pm)^{-1} \|V\|_{L^1}^{-1}Q_0 + \widetilde{\mathcal{D}}_{0,0}^{\pm}.  
		\end{equation}  
		Here, $\mathcal{W} = \|V\|_{L^1}^{-1}Q_0 T_0 Q_4^1$ and  
		\begin{equation}\label{eq:inverse_d}  
			(\boldsymbol{d}^\pm)^{-1} := \bigl[a_2^\pm Q_4^1 v G_4 v Q_4^1 - (a_0^\pm)^{-1} \|V\|_{L^1}^{-1} Q_4^1 T_0  Q_0 T_0 Q_4^1\bigr]^{-1}.  
		\end{equation}  
		
	
	
	\vskip0.5cm
	{\rm (ii)} \textbf{For $\alpha \in \{1,2,3\}$}, we have  
	\[  
	\big\|\partial_\lambda^\ell \mathcal{M}_{4,\alpha}^\pm(\lambda)\big\|_{\mathbb{B}(L^2)} + \big\|\partial_\lambda^\ell \mathcal{M}_{\alpha,4}^\pm(\lambda)\big\|_{\mathbb{B}(L^2)} \lesssim \lambda^{1-\ell}|\log \lambda|^{4}, \quad \ell = 0,1,2.  
	\]  
	 Let $\Lambda^\pm(\lambda)$ denote a generic operator in $\mathbb{B}(L^2)$ (possibly varying at each occurrence) satisfying
	$
	\left\|\partial_\lambda^\ell \Lambda^\pm(\lambda)\right\|_{\mathbb{B}(L^2)} \lesssim \lambda^{-\ell},\ \ell = 0, 1, 2.$
	For  $(\alpha, \beta) \in \{(0,4), (4,0), (4,4)\}$, we have
	\vskip0.3cm
 \textbf{Case 1: $Q_4^0 = 0$.} We have  
	\[  
\mathcal{M}_{\alpha,\beta}^\pm(\lambda) = \mathcal{D}_{\alpha,\beta}^{\pm}+ \lambda (-\log \lambda)^{3/2} \Lambda^\pm(\lambda),  
	\]  
	where the operators $\mathcal{D}_{\alpha,\beta}^{\pm}$ are explicitly given by  
	\begin{equation}\label{eq:expre_D_00} 
		\mathcal{D}_{4,4}^{\pm} = (\boldsymbol{d}^\pm)^{-1}, \quad  
		\mathcal{D}_{0,4}^{\pm} = - (a_0^\pm)^{-1}  \mathcal{W} (\boldsymbol{d}^\pm)^{-1}, \quad  
		\mathcal{D}_{4,0}^{\pm} = - (a_0^\pm)^{-1} (\boldsymbol{d}^\pm)^{-1} \mathcal{W}^*,  
	\end{equation}  
	with $\mathcal{W}$ and  $\boldsymbol{d}^\pm$ defined in \eqref{eq:inverse_d}.  
	\vskip0.3cm
 \textbf{Case 2: $Q_4^0 \neq 0$.} This case further splits based on $Q_{42}^1$:  
	\begin{itemize}  
		\item \textit{Subcase $Q_{42}^1 = 0$:}  
			$$
		\begin{aligned}
		&\mathcal{M}_{4,0}^\pm(\lambda)=(\log \lambda)^{-1}  Q_4 \Lambda^\pm(\lambda) Q_0,  \quad
		\mathcal{M}_{0,4}^\pm(\lambda)=(\log \lambda)^{-1}  Q_0 \Lambda^\pm(\lambda) Q_4,  \\
		&\mathcal{M}_{4,4}^\pm(\lambda) = \sum_{h,l\in\{0,1\}}(\log \lambda)^{-(2-h-l)} Q_4^h \Lambda^\pm(\lambda) Q_4^l;
		\end{aligned}
		$$	
		\item \textit{Subcase $Q_{42}^1 \neq 0$:}  
		\begin{equation*}  
			\begin{aligned}  
				\mathcal{M}_{0,4}^{\pm}(\lambda) &= (\log \lambda)^{-1} Q_0 \Lambda^\pm(\lambda) Q_{4}^0 + Q_0 \Lambda^\pm(\lambda) Q_{4}^1, \\  
	\mathcal{M}_{4,0}^{\pm}(\lambda) &= (\log \lambda)^{-1} Q_{4}^0 \Lambda^\pm(\lambda) Q_0 + Q_{4}^1 \Lambda^\pm(\lambda) Q_0, \\  
	\mathcal{M}_{4,4}^{\pm}(\lambda) &= Q_4^0\mathcal{C}^\pm Q_4^0 +(\log \lambda) \bigl(Q_4^0 \mathcal{C}_1^\pm Q_{42}^1 + Q_{42}^1 \mathcal{C}_2^\pm Q_4^0 \bigr)+ (\log \lambda)^2 (a_2^\pm)^{-1} Q_{42}^1 \mathcal{S}^{-1} Q_{42}^1  \\
&  \quad + (\log \lambda)^{-1} Q_4^0\Lambda^\pm(\lambda) Q_4^0 + \bigl(Q_4^0 \Lambda^\pm(\lambda) Q_{4}^1 + Q_{4}^1 \Lambda^\pm(\lambda) Q_4^0\bigr)  +(\log \lambda) Q_{4}^1 \Lambda^\pm(\lambda) Q_{4}^1.  
			\end{aligned}  
		\end{equation*}  
		Here, setting $\mathcal{W}_1 = (b_0 Q_0 v G_2 v Q_4^0)^{-1} Q_0 T_0 Q_{42}^1,$ the operators $\mathcal{C}^\pm, \mathcal{C}_1^\pm, \mathcal{C}_2^\pm$ are given by  
	\begin{equation}\label{eq:mathcal_C}  
			\mathcal{C}^\pm := (a_2^\pm)^{-1}\mathcal{W}_1 \mathcal{S}^{-1} \mathcal{W}_1^*, \quad  
			\mathcal{C}_1^\pm := - (a_2^\pm)^{-1}\mathcal{W}_1 \mathcal{S}^{-1}, \quad  
			\mathcal{C}_2^\pm := - (a_2^\pm)^{-1}\mathcal{S}^{-1}\mathcal{W}_1^*,
		\end{equation}  
		and $\mathcal{S}$ is strictly positive on $Q_{42}^1 L^2$
		defined by  $$
		\mathcal{S} =\mathcal{W}_1^* \bigl[Q_4^0 v G_4 v Q_4^0 - Q_4^0 v G_4 v Q_{41}^1 (Q_{41}^1 v G_4 v Q_{41}^1)^{-1} Q_{41}^1 v G_4 v Q_4^0\bigr] \mathcal{W}_1,
		$$
      where  the inverse $(Q_{41}^1 v G_4 v Q_{41}^1)^{-1}$ is interpreted
		as the zero operator  
        if $Q_{41}^1L^2=\{0\}$.
	\end{itemize}  

	
	\vskip0.5cm
	{\rm (iii)} \textbf{For $\alpha \in \{5,6\}$ and $0 \le \beta \le 4$, we have for $\ell=0,1,2$}:  
	\[  
	\big\|\partial_\lambda^{\ell} \mathcal{M}_{\alpha,\beta}^\pm(\lambda)\big\|_{\mathbb{B}\left(L^2\right)}  
	+ \big\|\partial_\lambda^{\ell} \mathcal{M}_{\beta,\alpha}^\pm(\lambda)\big\|_{\mathbb{B}\left(L^2\right)}  
	\lesssim\lambda^{-\ell}.  
	\]  
	For the remaining principal pairs, we have  
	\[  
\mathcal{M}_{\alpha,\beta}^\pm(\lambda) = 
	\begin{cases}  
		(\log \lambda)^{-1} \mathcal{A}_{\alpha,\beta}(\lambda) + (\log \lambda)^{-2}\Gamma_{\alpha,\beta}^{\pm}(\lambda), & \text{for } (\alpha,\beta) \in \{(5,5), (5,6), (6,5)\}, \\[1ex]  
		\mathcal{A}_{6,6}(\lambda) + (\log \lambda)^{-2}\Gamma_{6,6}^{\pm}(\lambda), & \text{for } (\alpha,\beta) = (6,6).  
	\end{cases}  
	\]  
	Here, the principal terms $\mathcal{A}_{\alpha,\beta}(\lambda)$ are independent of the sign $\pm$, and the remainders take the form  
	\[  
	\Gamma_{\alpha,\beta}^{\pm}(\lambda) = \mathcal{L}_{\alpha,\beta}^{\pm} + (-\log \lambda)^{-1/2} \Lambda_{\alpha,\beta}^\pm(\lambda),  
	\]  
	with the leading term defined by  
	\begin{equation}\label{mathcal_L}  
		\begin{aligned}  
			\mathcal{L}_{5,5}^{\pm} &= -a_3^\pm b_1^{-2}(Q_5 v G_6 v Q_5)^{-1}, \ \  
			&&\mathcal{L}_{5,6}^{\pm} = -a_3^\pm \mathcal{D}_{5,5} (Q_5 v G_6^0 v Q_6) \mathcal{D}_{6,6}, \\  
			\mathcal{L}_{6,5}^{\pm} &= -a_3^\pm \mathcal{D}_{6,6} (Q_6 v G_6^0 v Q_5) \mathcal{D}_{5,5}, \ \  
			&&\mathcal{L}_{6,6}^{\pm} = a_3^\pm b_1 \mathcal{D}_{6,6} (Q_6 v G_6^0 v Q_5) \mathcal{D}_{5,5} (Q_5 v G_6^0 v Q_6) \mathcal{D}_{6,6},  
		\end{aligned}  
	\end{equation}  
	where $\mathcal{D}_{5,5} = -(b_1 Q_5 v G_6 v Q_5)^{-1}$ and $\mathcal{D}_{6,6} = (b_1 Q_6 v G_6^0 v Q_6)^{-1}$.   
	Moreover, all $\mathcal{A}_{\alpha,\beta}(\lambda)$, $\Gamma_{\alpha,\beta}^{\pm}(\lambda)$, and $\Lambda_{\alpha,\beta}^\pm(\lambda)$ belong to $\mathbb{B}(L^2)$ and satisfy
\[
\big\|\partial_\lambda^\ell \mathcal{A}_{\alpha,\beta}(\lambda)\big\|_{\mathbb{B}(L^2)} + \big\|\partial_\lambda^\ell \Gamma_{\alpha,\beta}^{\pm}(\lambda)\big\| _{\mathbb{B}(L^2)}+ \big\|\partial_\lambda^\ell \Lambda_{\alpha,\beta}^{\pm}(\lambda)\big\|_{\mathbb{B}(L^2)} \lesssim \lambda^{-\ell}, \quad \ell = 0,1,2.
\]
\end{theorem} 
\end{comment}


We now provide the asymptotic expansion of $(M^\pm(\lambda))^{-1}$ as $\lambda \to 0^+$; the proof can be seen in Section \ref{subsec:proof_inverse}. 

\begin{theorem}\label{thm:M_inverse}
	Assume that zero is a $\mathbf{k}$-th kind resonance with $0 \leq \mathbf{k} \leq 4$, and $|V(x)| \lesssim \left \langle  x\right \rangle^{-\mu}$ for some $\mu>0$ satisfying the condition \eqref{condition}.
	Then there exists $0 < \lambda_0 \ll 1$ such that for all $0 < \lambda \le \lambda_0$, the operator $\left(M^{\pm}(\lambda)\right)^{-1}$ on $L^2(\mathbb{R}^2)$ admits the asymptotic expansion:
	\begin{equation}\label{eq:M_inverse}  
		\left(M^{\pm}(\lambda)\right)^{-1} = \sum_{0 \leq \alpha, \beta \leq \mathbf{k}+2} \lambda^{2-k_\alpha-k_\beta} Q_\alpha \mathcal{M}_{\alpha, \beta}^{\pm}(\lambda) Q_\beta ,  
	\end{equation}  
	where the values of the parameter $k_\alpha$ are assigned as follows:  
	\begin{equation}\label{tab:k_alpha}  
		\begin{array}{c|ccccccc}  
			\alpha & 0 & 1 & 2 & 3 & 4 & 5 & 6 \\  
			\hline  
			k_\alpha & 0 & 1 & 1 & 1 & 2 & 3 & 3  
		\end{array}  
	\end{equation}  
	Moreover, the  operators $\mathcal{M}_{\alpha, \beta}^{\pm}(\lambda)$ possess the following properties:  
	
	\vskip0.3cm  
	\textbf{{\rm (I)}\  If zero is a regular point or a first-kind resonance of $H$} $(\text{i.e.,}\  \mathbf{k}=0,1 )$, then 
	\begin{align}\label{M0-1}
	\Big\| \partial_\lambda^\ell \mathcal{M}_{\alpha,\beta}^\pm(\lambda) \Big\|_{\mathbb{B}(L^2)} \lesssim
	\begin{cases}
		\lambda^{-\ell}, & \text{for } ( \alpha,\beta) \neq (3,3), \\
		|\log \lambda| \lambda^{-\ell}, &  \text{for } (\alpha,\beta) = (3,3),
	\end{cases}
	\end{align}
for  all $0 \le \alpha, \beta \le 3$ and $\ell=0,1,2$.	In particular, for  $(\alpha, \beta) = (0,0)$,
	\[
		\mathcal{M}_{0,0}^\pm(\lambda) = (a_0^\pm)^{-1} \|V\|_{L^1}^{-1}Q_0  + \lambda (-\log \lambda)^{3/2}\Lambda^\pm(\lambda).
	\]	  
Here and hereafter,  $\Lambda^\pm(\lambda)$ denotes a generic operator in $\mathbb{B}(L^2)$ (possibly varying at each occurrence) satisfying
$$
\left\|\partial_\lambda^\ell \Lambda^\pm(\lambda)\right\|_{\mathbb{B}(L^2)} \lesssim \lambda^{-\ell},\quad \ell = 0, 1, 2.$$
	\vskip0.1cm
	 \textbf{{\rm (II)}\ If  zero is a  second-kind resonance of $H$} $(\text{i.e.,}\  \mathbf{k}=2 )$,  then
	for all $0 \le \alpha, \beta \le 3$ and $\ell=0,1,2$, the estimates \eqref{M0-1} remain valid. Moreover,
for $\alpha \in \{1,2,3\}$,   
	\[  
	\big\|\partial_\lambda^\ell \mathcal{M}_{4,\alpha}^\pm(\lambda)\big\|_{\mathbb{B}(L^2)} + \big\|\partial_\lambda^\ell \mathcal{M}_{\alpha,4}^\pm(\lambda)\big\|_{\mathbb{B}(L^2)} \lesssim \lambda^{1-\ell}|\log \lambda|^{4}, \quad \ell = 0,1,2.  
	\]  
	It remains to specify the case
	$
	(\alpha,\beta)
	\in
	\{(0,4),(4,0),(4,4)\}
,$
	which are given as follows:
	\vskip0.3cm
	\textbf{Case 1: $Q_4^0 = 0$.}   For all $(\alpha,\beta)
	\in
	\{(0,4),(4,0),(4,4)\}$, we have
	\[
	\mathcal{M}_{\alpha,\beta}^\pm(\lambda) = Q_\alpha \Lambda^\pm(\lambda)Q_\beta.
	\]
	%\vskip0.3cm
    
	\textbf{Case 2: $Q_4^0 \neq 0$.}   The case for $(\alpha, \beta) \in \{ (0,4), (4,0), (4,4)\}$  further splits based on $Q_{42}^1$:  
	\begin{itemize}  
		\item \textit{ When $Q_{42}^1 = 0$:}  
		$$
		\begin{aligned}
			&\mathcal{M}_{4,0}^\pm(\lambda)=(\log \lambda)^{-1}  Q_4 \Lambda^\pm(\lambda) Q_0,  \quad
			\mathcal{M}_{0,4}^\pm(\lambda)=(\log \lambda)^{-1}  Q_0 \Lambda^\pm(\lambda) Q_4,  \\
			&\mathcal{M}_{4,4}^\pm(\lambda) = \sum_{h,l\in\{0,1\}}(\log \lambda)^{-(2-h-l)} Q_4^h \Lambda^\pm(\lambda) Q_4^l;
		\end{aligned}
		$$	
		\item \ \textit{When $Q_{42}^1 \neq 0$:}  
		\begin{equation*}  
			\begin{aligned}  
				\mathcal{M}_{0,4}^{\pm}(\lambda) &= (\log \lambda)^{-1} Q_0 \Lambda^\pm(\lambda) Q_{4}^0 + Q_0 \Lambda^\pm(\lambda) Q_{4}^1, \\  
				\mathcal{M}_{4,0}^{\pm}(\lambda) &= (\log \lambda)^{-1} Q_{4}^0 \Lambda^\pm(\lambda) Q_0 + Q_{4}^1 \Lambda^\pm(\lambda) Q_0, \\  
				\mathcal{M}_{4,4}^{\pm}(\lambda) &= Q_{4} \Lambda^\pm(\lambda) Q_{4} + (\log\lambda)\bigl(Q_{4} \Lambda^\pm(\lambda) Q_{4}^1+Q_{4}^1 \Lambda^\pm(\lambda) Q_{4}\bigr) + (\log\lambda)^2 Q_{4}^1 \Lambda^\pm(\lambda) Q_{4}^1.
			\end{aligned}  
		\end{equation*}  
	\end{itemize}  
	
	
	\vskip0.1cm
	 \textbf{{\rm (III)} If  zero  is a  third-kind resonance or an eigenvalue of $H$} $(\text{i.e.,}\  \mathbf{k}=3,4 )$,  then
	for all $0 \le \alpha, \beta \le 4$ and $\ell=0,1,2$, all the conclusions established in the second resonance case   remain valid. Moreover,
	for $\alpha \in \{5,6\}$ and $0 \le \beta \le 4$,
	\[  
	\big\|\partial_\lambda^{\ell} \mathcal{M}_{\alpha,\beta}^\pm(\lambda)\big\|_{\mathbb{B}\left(L^2\right)}  
	+ \big\|\partial_\lambda^{\ell} \mathcal{M}_{\beta,\alpha}^\pm(\lambda)\big\|_{\mathbb{B}\left(L^2\right)}  
	\lesssim\lambda^{-\ell},\quad \ell=0,1,2.
	\]  
	For the remaining principal pair $(\alpha,\beta) \in \{(5,5), (5,6), (6,5), (6,6)\},$  
	\[  
	\mathcal{M}_{\alpha,\beta}^\pm(\lambda) = 
	 \mathcal{A}_{\alpha,\beta}(\lambda) + (\log \lambda)^{-2}\Gamma_{\alpha,\beta}^{\pm}(\lambda).  
	\]  
	Here, the principal terms $\mathcal{A}_{\alpha,\beta}(\lambda)$ are independent of the sign $\pm$ and 
	all $\mathcal{A}_{\alpha,\beta}(\lambda), \Gamma_{\alpha,\beta}^{\pm}(\lambda)\in \mathbb{B}(L^2)$  satisfy
	\[
	\big\|\partial_\lambda^\ell \mathcal{A}_{\alpha,\beta}(\lambda)\big\|_{\mathbb{B}(L^2)} + \big\|\partial_\lambda^\ell \Gamma_{\alpha,\beta}^{\pm}(\lambda)\big\| _{\mathbb{B}(L^2)} \lesssim \lambda^{-\ell}, \quad \ell = 0,1,2.
	\]
\end{theorem}

To clarify the singular hierarchy in Theorem~\ref{thm:M_inverse},
Table~\ref{tab:Q-leading-order} records the zero-threshold types,
their \(S\)-projection characterizations, and the corresponding
worst leading orders of \((M^\pm(\lambda))^{-1}\) as
\(\lambda\to0^+\).

\vskip-0.3cm
\begin{table}[H]
	\centering
	\caption{Zero-threshold types and leading orders of
		\((M^\pm(\lambda))^{-1}\).}
	\label{tab:Q-leading-order}
	\vskip0.1cm
	\renewcommand{\arraystretch}{1.45}
	\setlength{\tabcolsep}{4pt}
	\footnotesize
	
	\newcommand{\Qcell}[1]{\makecell[c]{\rule[-1.8ex]{0pt}{6.0ex}#1}}
	
	\begin{tabular}{|
			>{\centering\arraybackslash}m{0.20\textwidth}|
			>{\centering\arraybackslash}m{0.34\textwidth}|
			>{\centering\arraybackslash}m{0.36\textwidth}|}
		\hline
		\textbf{Zero-threshold type}
		&
		\textbf{\(S\)-characterization}
		&
		\textbf{The leading order of \((M^\pm(\lambda))^{-1}\)}
		\\[1mm]
		\hline
		
		\Qcell{Regular point\\ \((\mathbf k=0)\)}
		&
		\Qcell{\(S_2=0\)}
		&
		\Qcell{\(O(1)\)}
		\\
		\hline
		
		\Qcell{1st-kind resonance\\ \((\mathbf k=1)\)}
		&
		\Qcell{
			\(\displaystyle S_2\neq 0\),
			\ \ \(\displaystyle S_3=0\)}
		&
		\Qcell{\(O(|\log\lambda|)\)}
		\\
		\hline
		
		\multirow{2}{=}{\Qcell{2nd-kind \\ \((\mathbf k=2)\)}}
		&
		\multirow{2}{=}{\Qcell{
				\(\displaystyle S_3\neq 0\),
				\ \ \(\displaystyle S_4=0\)}}
		&
		\Qcell{
			\(O\!\left(\lambda^{-2}|\log\lambda|^2\right)\)\ \
			if\ \  \(Q_4^0\neq0\) and \(Q_{42}^1\neq0\)}
		\\
		\cline{3-3}
		
		&
		&
		\Qcell{
		\ \   \ \   	\(O(\lambda^{-2})\)
		  \  \  \ \ \ \   	if \ \ \(Q_4^0=0\) or \(Q_{42}^1=0\)}
		\\
		\hline
		
		\Qcell{3rd-kind \\ \((\mathbf k=3)\)}
		&
		\Qcell{
			\(\displaystyle S_4\neq 0\),
			\ \ \(\displaystyle S_5=0\)}
		&
		\Qcell{\(O\!\left(\lambda^{-4}|\log\lambda|^{-1}\right)\)}
		\\
		\hline
		
		\Qcell{Zero eigenvalue\\ \((\mathbf k=4)\)}
		&
		\Qcell{
			\(\displaystyle S_5\neq0\)}
		&
		\Qcell{\(O(\lambda^{-4})\)}
		\\
		\hline
	\end{tabular}
	
	\vskip0.2cm
	\parbox{0.95\textwidth}{\footnotesize
The projections \(S_j\) are defined in Definition \ref{defS}. The
\(Q\)-projections are the successive orthogonal differences of 
\(S\):
\[
Q_j=S_{j-1}-S_j,\qquad 0\le j\le6.
\]
 In the second-kind case, 
\(Q_4\) is further decomposed into $Q_4^0+Q_4^1$ where
\[
Q_4^0=S_3-S_4^0,\quad
Q_4^1=S_4^0-S_4,\quad
Q_4^1=Q_{41}^1+Q_{42}^1,\quad
Q_{41}^1=S_4^0-S_4^1,\quad
Q_{42}^1=S_4^1-S_4.
\]
}
\end{table}







\begin{comment}
Although the singularity of \((M^\pm(\lambda))^{-1}\) becomes
increasingly severe as the threshold obstruction deepens, it is not the final singularity of the perturbed resolvent.
Indeed, after inserting \eqref{eq:M_inverse} into the symmetric
resolvent identity \eqref{id-RV}, each \(Q_\alpha\)-\(Q_\beta\) block
appears in the form
\[
R_0^\pm(\lambda^4)vQ_\alpha\,
\mathcal M_{\alpha,\beta}^\pm(\lambda)\,
Q_\beta vR_0^\pm(\lambda^4).
\]
Thus the singular inverse block is always coupled with the projected
free resolvent kernels \(Q_\alpha vR_0^\pm(\lambda^4)\).
\end{comment}


The cancellation properties of the projections \(Q_\alpha\) listed
in Remark \ref{cancellationQ}, eliminate the corresponding 
moment terms in the free  expansion of \(R_0^\pm(\lambda^4)\).
Consequently, the projected kernels \(Q_\alpha vR_0^\pm(\lambda^4)\) exhibit improved \(\lambda\)-decay than the unprojected free resolvent $R_0^\pm(\lambda^4)$. This improvement partially
compensates for the singularity of \((M^\pm(\lambda))^{-1}\). 
The following Lemma \ref{lemma_projection}  captures this compensation mechanism through structural decompositions and \(L^2\)-bounds for \(Q_\alpha vR_0^\pm(\lambda^4)\).

\subsection{Two fundamental lemmas}\label{subsec:osci_integral}
This subsection presents two key lemmas: the first analyzes the kernel of $Q_\alpha v R_0^\pm(\lambda^4)$, and the second establishes oscillatory integral estimates.

We begin with the definition of the operator $Q_\alpha v R_0^\pm(\lambda^4)$:
\[
\bigl(Q_\alpha v R_0^\pm(\lambda^4) f\bigr)(x) = Q_\alpha\!\Bigl( v(\cdot) \int_{\mathbb{R}^2} R_0^\pm(\lambda^4)(\cdot, y) f(y) \, dy \Bigr)(x).
\]
Its integral kernel is given explicitly by
\(
\bigl(Q_\alpha v R_0^\pm(\lambda^4)\bigr)(x, y) = Q_\alpha\!\left( v\, R_0^\pm(\lambda^4)(\cdot, y) \right)(x).
\)
The structural decompositions and $L^2$-estimates of this integral kernel are summarized in the following Lemma \ref{lemma_projection}.

\begin{lemma}\label{lemma_projection}
	Let $Q_\alpha$ $(\alpha\in\mathbb{N}_0, 0\le\alpha\le6)$ and $Q_4^0, Q_4^1, Q_{41}^1, Q_{42}^1$ be defined by \eqref{defQ}. Assume $\left|V(x)\right| \lesssim \left \langle  x\right \rangle^{-\mu}$ for some $\mu>0$ satisfying the condition \eqref{condition}. Then, for $0<\lambda\ll1$, the integral kernels admit the following decompositions:
	\begin{enumerate}
		\item[(i)] For $0\le\alpha\le6$,
		\[
		\bigl(Q_\alpha v R_0^\pm(\lambda^4)\bigr)(x,y)=\lambda^{-1-\delta_{\alpha0}}\,\Omega_{\alpha,\pm}(\lambda,x,y),
		\]
		where $\delta_{\alpha0}=1$ if $\alpha=0$ and $0$ otherwise.
		
		\item[(ii)] For $2\le\alpha\le6$, 
		\begin{align*}
		\bigl(Q_\alpha v R_0^\pm(\lambda^4)\bigr)(x,y)&=\mathcal{J}_{\alpha, \pm} (\lambda,x,y).
		\end{align*}
		\item[(iii)] For $2\le\alpha\le6$,
		setting $m_\alpha=0$ for $\alpha=2,3,4$, $m_5=1$, and $m_6=2$:
		$$
			\bigl(Q_\alpha v R_0^\pm(\lambda^4)\bigr)(x,y) =(b_0Q_\alpha v G_2^0)(x,y)+b_0(\log\lambda)\,Q_\alpha(v|\cdot|^2)(x)+ \lambda^{m_\alpha} (\mathcal{T}_{\alpha, \pm} + \mathcal{T}_\alpha)(\lambda, x, y),
        $$
		where the term $b_0 (\log\lambda) \, Q_\alpha(v|\cdot|^2)(x)$ vanishes
        for $\alpha=5,6$.
		In particular,
		\begin{align*}
\bigl(Q_4^0 v R_0^\pm(\lambda^4)\bigr)(x,y) &=(b_0Q_4^0 v G_2^0)(x,y)+b_0(\log\lambda)\,Q_4^0(v|\cdot|^2)(x)+(\mathcal{T}_{4,\pm}^0+\mathcal{T}_4^0)(\lambda,x,y),\\
\bigl(Q_4^1 v R_0^\pm(\lambda^4)\bigr)(x,y) &=(b_0Q_4^1 v G_2^0)(x,y)+(\mathcal{T}_{4,\pm}^1+\mathcal{T}_4^1)(\lambda,x,y),\\
\bigl(Q_{41}^1  v R_0^\pm(\lambda^4)\bigr)(x,y) &=(b_0Q_{41}^1  v G_2^0)(x,y)+(\mathcal{T}_{4,\pm}^{1,1}+\mathcal{T}_4^{1,1})(\lambda,x,y),\\
\bigl(Q_{42}^1  v R_0^\pm(\lambda^4)\bigr)(x,y) &=(b_0Q_{42}^1  v G_2^0)(x,y)+\lambda\bigl(\mathcal{T}_{4,\pm}^{1,2}+\mathcal{T}_4^{1,2})(\lambda,x,y).
		\end{align*}
	\end{enumerate}
	Furthermore, for $\ell=0,1,2$, these integral kernels satisfy the following bounds:
	\begin{align}\label{J}
		&\|\partial_\lambda^\ell\big(e^{\mp i\lambda|y|}\mathcal{J}_{\pm}(\lambda,\cdot,y)\big)\|_{L^2} \lesssim\lambda^{-\ell}|\log\lambda|\langle\lambda y\rangle^{-1/2},\  \text{for } \mathcal{J}_{\pm} \in \{\mathcal{J}_{\alpha,\pm}|_{2\le\alpha\le6}\},\\
&\bigl\|(b_0Q v G_2^0)(\cdot,y)\bigr\|_{L^2} \lesssim
		\begin{cases}
		\log(2+|y|), & \text{for } Q \in \{Q_\alpha|_{2 \le \alpha \le 4}, Q_4^0\}, \\[2pt]
		1, & \text{for } Q \in \{Q_4^1, Q_{41}^1, Q_{42}^1, Q_5, Q_6\},
		\end{cases} \label{lemma_projection_G} \\
&\|\partial_\lambda^\ell\mathcal{T}(\lambda,\cdot,y)\|_{L^2} \lesssim
		\begin{cases}
		\lambda^{-\ell}\log(2+|y|), & \text{for } \mathcal{T} \in \{\mathcal{T}_\alpha|_{2\le\alpha\le4}, \, \mathcal{T}_4^0\}, \\[2pt]
		\lambda^{-\ell}, & \text{for } \mathcal{T} \in \{\mathcal{T}_4^{1}, \mathcal{T}_4^{1,1}, \mathcal{T}_4^{1,2}, \mathcal{T}_5, \mathcal{T}_6\},
		\end{cases} \label{lemma_projection_T} \\   &\|\partial_\lambda^\ell\big(e^{\mp i\lambda|y|}K_{\pm}(\lambda,\cdot,y)\big)\|_{L^2} \lesssim\lambda^{-\ell}\langle\lambda y\rangle^{-1/2}, \label{lemma_projection_1}
	\end{align}
	where $K_{\pm}$ generically denotes any  term among $\Omega_{\alpha,\pm} \, (0\le\alpha\le6)$, $\mathcal{T}_{\alpha,\pm} \, (2\le\alpha\le6)$, $\mathcal{T}_{4,\pm}^j \, (j=0,1)$, or $\mathcal{T}_{4,\pm}^{1,j} \, (j=1,2)$.
\end{lemma}

The proof of Lemma \ref{lemma_projection} relies heavily on the cancellation properties of the projections $Q_{\alpha}$. We defer it to Section \ref{section 8}.
Next, we state Lemma \ref{oscillatory} which  serves as a fundamental tool for the oscillatory integrals encountered in the subsequent decay analysis. Its proof is also given in Section \ref{section 8}.

\begin{lemma}\label{oscillatory}    Let \(h(x,y)\) and \(r:=r(x,y)\) be real-valued functions on \(\mathbb{R}^2 \times \mathbb{R}^2\) with \(h(x,y)>0\). Define the oscillatory integrals    
	\[
	\begin{aligned}
		\mathcal{K}^\pm(t,x,y) &= \int_0^{\infty} \cos (t \lambda^2)e^{\pm i\lambda r} \lambda \mathcal{E}^\pm(\lambda, x,y) \, d\lambda, \\
		\mathcal{N}^\pm(t,x,y) &= t^{-1}\int_0^{\infty} \sin(t\lambda^2) e^{\pm i\lambda r} \lambda^{-1} \mathcal{E}^\pm(\lambda, x,y) \, d\lambda.
	\end{aligned}
    \]  
	\begin{itemize}    
		\item[(i)] Suppose that for $\ell = 0, 1, 2$, the amplitude factor $\mathcal{E}^\pm$ satisfies    
		\begin{align}\label{condit 1}    
		\left| \partial_\lambda^{\ell} \mathcal{E}^\pm(\lambda, x,y) \right| \lesssim h(x, y) \langle \lambda r \rangle^{-\frac{1}{2}} \lambda^{-\ell}.    
		\end{align}    
		Then    
		\[    
		\left| \mathcal{K}^\pm(t,x,y) \right| + \left| \mathcal{N}^\pm(t,x,y) \right| \lesssim h(x, y) |t|^{-1}.    
		\]    
		
		\item[(ii)] Suppose that for $\ell = 0, 1, 2$, $\mathcal{E}^\pm$ exhibits a  low-frequency behavior localized by the smooth cutoff $\widetilde{\chi}_1(\lambda/2)$ supported in $[0, 2(2\lambda_0)^{1/4}]$ and    
		\begin{align}\label{condit 2}    
		\left| \partial_\lambda^{\ell} \mathcal{E}^\pm(\lambda, x,y) \right| \lesssim h(x, y) \langle \lambda r \rangle^{-\frac{1}{2}} |\log \lambda|^{-\nu} \lambda^{-\sigma - \ell} \widetilde{\chi}_1 (\lambda/2).    
		\end{align}    
		Then the following estimates hold:    
		
		$\bullet$ For $0 < \sigma < 2$ with $\nu \in \mathbb{R}$, or $\sigma = 0$ with $\nu \le 0$:    
		\[    
		\left| \mathcal{K}^\pm(t,x,y) \right| + \left| \mathcal{N}^\pm(t,x,y) \right| \lesssim \frac{h(x, y)}{\langle t \rangle^{1-\frac{\sigma}{2}} (\log(2+|t|))^{\nu}}.    
		\]    
Moreover, if $r=0,$  this estimate  remains
   valid for the larger parameter range
   $
   -2<\sigma<2$ and 
   $\nu\in\mathbb R.
   $		
		
        $\bullet$ For $\sigma = 2$ and $\nu > 1$:    
		\[    
		\left| \mathcal{K}^\pm(t,x,y) \right| + \left| \mathcal{N}^\pm(t,x,y) \right| \lesssim  \frac{h(x, y)}{\left(\log(2+|t|)\right)^{\nu-1}}.    
		\]    
	\end{itemize} 
\end{lemma}    
  

\section{The decay estimates for the free case}\label{sec:free_case}
In this section, we establish the $L^1 \to L^{\infty}$ and weighted $L^1_\omega \to L^{\infty}_{-\omega}$ bounds for the free case $V=0$. The corresponding unperturbed operator is $\sqrt{H}=-\Delta$,  which corresponds to the evolution propagators $\cos(t\Delta)$ and $\sin(t\Delta)/(t\Delta)$. 
  By Stone's formula, their kernels can be expressed as
 \begin{align}\label{free-stone_cos}
    \cos (t\Delta)(x,y) & =\frac{2}{\pi i} \int_0^{\infty} \lambda^3 \cos (t \lambda^2)\left[R_0^{+}(\lambda^4)-R_0^{-}(\lambda^4)\right](x,y)  d \lambda,\\
    \label{free-stone_sin}
     \frac{\sin (t \Delta)}{t\Delta}(x,y) & =\frac{2}{\pi i t} \int_0^{\infty} \lambda \sin (t \lambda^2)\left[R_0^{+}(\lambda^4)-R_0^{-}(\lambda^4)\right](x,y) d \lambda.
     \end{align}
On the other side, the kernels of $\cos(t\Delta)$ and $\sin(t\Delta)/(t\Delta)$ can be written explicitly as
\begin{align}\label{eq:free_kernels}
K_t^{\cos,0}(x,y) = \frac{1}{4\pi t} \sin \frac{|x-y|^2}{4t},\quad
K_t^{\sin,0}(x,y) =\frac{1}{4\pi |t|}\int_{\frac{|x-y|^2}{4|t|}}^{\infty}\frac{\sin u}{u}\,du .
\end{align}
Indeed, the well-known  kernel of the free Schr\"{o}dinger group \(e^{it\Delta}\): $$e^{it\Delta}(x,y) = \frac{1}{4\pi i t} \exp\bigl(i\frac{|x-y|^2}{4t}\bigr),$$ and the identity $\cos(t\Delta) = (e^{it\Delta}+e^{-it\Delta})/2$ yield $K_t^{\cos,0}(x,y)$. The sine kernel $K_t^{\sin,0}(x,y)$ then follows  from the  relation $\frac{\sin(t\Delta)}{t\Delta} = \frac{1}{t} \int_0^t \cos(s\Delta)\,ds$ via the substitution $u = \frac{|x-y|^2}{4s}$.

Based on these representations, we obtain the following sharp decay estimates.

\begin{proposition}\label{free_case}
The free  operators satisfy the following sharp bounds:
\vskip 0.1cm
\begin{enumerate}
\item[(i)] In the unweighted space $L^1 \to L^{\infty}$:
\[
\|\cos (t \Delta) \|_{L^1 \to L^{\infty}} \sim |t|^{-1}, \quad
\Bigl\|\frac{\sin (t \Delta)}{t\Delta}\Bigr\|_{L^1 \to L^{\infty}} \sim |t|^{-1}.
\]
\vskip 0.2cm
\item[(ii)] In the weighted space $L^1_\omega \to L^\infty_{-\omega}$:
\[
\|\cos (t \Delta) \|_{L^1_\omega \to L^{\infty}_{-\omega}} \sim |t|^{-1}(\log (2+|t|))^{-1}, \quad
\Bigl\|\frac{\sin (t \Delta)}{t\Delta} \Bigr\|_{L^1_\omega \to L^{\infty}_{-\omega}} \sim |t|^{-1}.
\]
\end{enumerate}
\end{proposition}
\begin{comment}
{\color{blue}
For a positive function $\omega(\lambda)$ that is monotone increasing and convex, satisfying $\omega(1)\ge 1$. If $\frac{\lambda^2}{\omega(\lambda)}$ is also monotone increasing, we have
\[
\|\cos (t \Delta) \|_{L^1_{\omega(\langle x\rangle)} \to L^{\infty}_{-{\omega(\langle x\rangle)}}} \sim |t|^{-1}\omega^{-1}(\sqrt{t}).
\]
}
\end{comment}
\begin{proof}
	(i) The sharp uniform bounds $\|\cos(t\Delta)\|_{L^1\to L^\infty}\sim|t|^{-1}$ and $\|\sin(t\Delta)/(t\Delta) \|_{L^1 \to L^{\infty}} \sim |t|^{-1}$ follow trivially by taking the supremum over the explicit kernels in \eqref{eq:free_kernels}. However, to set the stage for subsequent perturbation arguments, we provide an alternative proof for the upper bound via oscillatory integrals.
	
	In view of \eqref{free-stone_cos} and \eqref{free-stone_sin}, substituting $R^\pm_0(\lambda^4)(x,y) = \lambda^{-2}e^{\pm i\lambda|x-y|}\widetilde{R}^\pm(\lambda|x-y|)$ (see \eqref{def:F_pm} and \eqref{free-widetilde_R}) reduces the problem to estimating
	\begin{align*}
	\mathcal{K}^{ \pm}(t, x, y) &:= \frac{2}{\pi i}\int_0^{\infty} \cos(t \lambda^2) \lambda e^{\pm i\lambda|x-y|}\widetilde{R}^\pm(\lambda|x-y|) \, d\lambda, \\
	\mathcal{N}^{ \pm}(t, x, y) &:= \frac{2}{\pi i t}\int_0^{\infty} \sin (t \lambda^2) \lambda^{-1} e^{\pm i\lambda|x-y|}\widetilde{R}^\pm(\lambda|x-y|) \, d\lambda.
	\end{align*}
	By virtue of \eqref{free-widetilde_R}, it is straightforward to verify that for $\ell=0,1,2$,
	\[
	\left| \partial_\lambda^{\ell} \widetilde{R}^\pm(\lambda |x-y|) \right| \lesssim \langle \lambda |x-y| \rangle^{-1/2} \lambda^{-\ell}.
	\]
	Applying Lemma \ref{oscillatory}(i) with $r(x,y)=|x-y|$ and $\mathcal{E}^\pm(\lambda, x,y)=\widetilde{R}^\pm(\lambda|x-y|)$, we deduce that both $|\mathcal{K}^{\pm}|$ and $|\mathcal{N}^{\pm}|$ are uniformly bounded by ${O}(|t|^{-1})$, concluding the proof of (i).
	
	(ii) We evaluate the weighted operator norm for $\cos(t\Delta)$ by taking the spatial supremum:
	\begin{equation*}
	\|\cos(t\Delta)\|_{L^1_\omega\to L^\infty_{-\omega}} = \sup_{x,y\in\mathbb{R}^2} \frac{|K_t^{\cos,0}(x,y)|}{\omega(x)\omega(y)}.
	\end{equation*}
	Setting $\rho = |x-y|$ and utilizing the basic weight inequality $\omega(x)\omega(y) \ge \omega(0)\omega(\rho/2)$, we have
	\begin{equation}\label{free-weight}
	\|\cos(t\Delta)\|_{L^1_\omega\to L^\infty_{-\omega}}
	\lesssim \sup_{\rho \ge 0} \frac{|\sin(\rho^2/(4t))|}{|t|\,\omega(\rho/2)}
	= \sup_{u\ge 0} \frac{|\sin u|}{|t|\log\bigl(2+\sqrt{|t| u}\bigr)}.
	\end{equation}
	For $u>1$, we use $|\sin u|\le1$ and $\log(2+\sqrt{|t|u}) \ge \log(2 + \sqrt{|t|}) \sim \log (2 + |t|)$. For $0 \le u \le 1$, using $|\sin u|\le u$ and the substitution $v=\sqrt{|t|u}$ yields
	\[
	\sup_{0\le u\le 1} \frac{|\sin u|}{\log(2+\sqrt{|t|u})}
	\le \sup_{0\le v\le \sqrt{|t|}} \frac{v^2/|t|}{\log(2+v)} = \frac{1}{\log(2+\sqrt{|t|})} \sim \frac{1}{\log (2 + |t|)}.
	\]
	Inserting these into \eqref{free-weight} establishes the upper bound $\|\cos(t\Delta)\|_{L^1_\omega\to L^\infty_{-\omega}} \lesssim |t|^{-1}\bigl(\log (2 + |t|)\bigr)^{-1}$. The matching lower bound follows by testing the kernel at $y=0$ and $|x| = \sqrt{2\pi|t|}$.
	
	For $\sin(t\Delta)/(t\Delta)$, the unweighted  bound from part (i) implies the weighted upper bound. Evaluating $K_t^{\sin,0}(0, 0) = 1/(8|t|)$ provides the corresponding lower bound, completing the proof.
\end{proof}

\begin{remark}\label{free_cos_higher_dim}{\rm 
For a general dimension $n$, the kernel of $\cos(t\Delta)$ admits the explicit form
\[
\cos(t\Delta)(x,y) = \frac{1}{(4\pi t)^{n/2}} \cos\!\left(\frac{|x-y|^2}{4t} - \frac{n\pi}{4}\right),
\]
which trivially yields the sharp decay estimate $\|\cos(t\Delta)\|_{L^1\to L^\infty} \lesssim |t|^{-n/2}$. 
In the weighted setting, the enhanced decay bound
\[
\|\cos(t\Delta)\|_{L^1_\omega \to L^\infty_{-\omega}} \sim |t|^{-n/2} \bigl(\log(2+|t|)\bigr)^{-1}
\]
holds if and only if $n \equiv 2 \pmod{4}$ (e.g., $n = 2, 6, 10, \dots$). For all other dimensions, the logarithmic gain is absent, yielding purely polynomial decay. This dimension-dependent behavior arises because the dimensional phase shift $n\pi/4$ determines whether the oscillatory kernel  vanishes along the diagonal $x=y$.
}\end{remark}
\begin{remark}\label{rmk:free_case}{\rm
(i) The same analysis as in Proposition \ref{free_case} for $\mathcal{K}^{\pm}(t,x,y)$ and $\mathcal{N}^{\pm}(t,x,y)$ yields
\[
\big\|\cos(t\Delta)\chi_j(\Delta)\big\|_{L^1\to L^{\infty}}
+ \Big\|\frac{\sin(t\Delta)}{t\Delta}\chi_j(\Delta)\Big\|_{L^1\to L^{\infty}} \lesssim |t|^{-1},
\qquad j=1,2.
\]
For the high-frequency part, the sine propagator actually satisfies
the stronger bound
\[
\Big\|\frac{\sin(t\Delta)}{t\Delta}\,\chi_2(\Delta)\Big\|_{L^1\to L^{\infty}} \lesssim |t|^{-2}.
\]
Indeed, when $\lambda\gtrsim1$, Remark \ref{remark:osci} (i) gives
\[
\Bigl|\int_0^{\infty} e^{-it\lambda^2}\,\lambda\,\widetilde{\chi}_2(\lambda)R_0^{\pm}(\lambda^4)(x,y)\,d\lambda\Bigr| \lesssim |t|^{-1},
\]
which after multiplication by the factor $t^{-1}$ present in $\bigl(\sin(t\Delta)/(t\Delta)\bigr)\chi_2(\Delta)$ produces the $|t|^{-2}$ decay. Consequently,
the same bound $|t|^{-2}$ also holds in the weighted space $L^1_\omega \rightarrow L^{\infty}_{-\omega}$.
\vskip0.2cm

(ii)
$\cos(t\Delta)\chi_j(\Delta)$ satisfies the same weighted decay as $\cos(t\Delta)$:
\[
\|\cos(t\Delta)\chi_j(\Delta)\|_{L^1_\omega\to L^{\infty}_{-\omega}} \lesssim (|t|\log(2+|t|))^{-1}, \qquad j=1,2.
\]
}
\end{remark}

Finally, we  establish the asymptotic expansion for the low-frequency free sine kernel,
which is essential for the subsequent log-weighted estimates in the regular and the first kind resonance cases.
\begin{proposition}\label{prop_free_sin}
For $|t| \ge 2,$ the low-frequency free sine kernel admits the asymptotic expansion:
$$
\begin{aligned}
\frac{\sin(t\Delta)}{t\Delta}\chi_1(\Delta)(x,y)
=\frac{1}{8 |t|}+O\left(\frac{\omega(x)\omega(y)}{|t|\log|t|}\right).
\end{aligned}
$$
\end{proposition}
\begin{proof}
	Recall the kernel representation of the full free sine propagator from \eqref{eq:free_kernels}:
	\begin{equation}\label{free-whole}
		\frac{\sin(t\Delta)}{t\Delta}(x,y)
		= \frac{1}{4\pi |t|}\int_{\frac{|x-y|^2}{4|t|}}^{\infty}\frac{\sin u}{u}\,du
		= \frac{1}{8|t|} - \frac{1}{4\pi |t|}\int^{\frac{|x-y|^2}{4|t|}}_{0}\frac{\sin u}{u}\,du,
	\end{equation}
	where we used the standard Dirichlet integral $\int_0^\infty \frac{\sin u}{u} du = \frac{\pi}{2}$.
	Using the uniform bound $|\frac{\sin u}{u}| \le 1$ and the inequality $|x-y|^2 \le \langle x \rangle^2 \langle y \rangle^2$, we derive 
	\begin{equation}\label{free-est}
		 \frac{1}{4\pi |t|}\Big| \int_{0}^{\frac{|x-y|^2}{4|t|}} \frac{\sin u}{u} \, du \Big| \le \frac{|x-y|^2}{16\pi |t|^2} \lesssim \frac{\langle x \rangle^2 \langle y \rangle^2}{|t|^2}.
	\end{equation}
To isolate the low-frequency part, we decompose the operator as $\chi_1(\Delta) = \mathrm{Id} - \chi_2(\Delta)$.	According to Remark~\ref{rmk:free_case}\,(i), the high-frequency part satisfies
	\begin{equation}\label{free-high}
		\Big| \frac{\sin(t\Delta)}{t\Delta}\chi_2(\Delta)(x,y) \Big| \lesssim |t|^{-2}.
	\end{equation}
	Combining \eqref{free-whole}, \eqref{free-est}, and \eqref{free-high}, we deduce that
	\[
	\frac{\sin(t\Delta)}{t\Delta}\chi_1(\Delta)(x,y) -\frac{1}{8 |t|} ={O}\!\left(\frac{\langle x\rangle^2\langle y\rangle^2}{|t|^2}\right).
	\]
    
	Finally, minimizing with the uniform bound $|t|^{-1}$ via using $\min(1,a/b) \lesssim \log a/\log b$ for $a,b\ge 2$, we obtain
	\[
	\min\Bigl(\frac{1}{|t|}, \frac{\langle x\rangle^2\langle y\rangle^2}{|t|^2}\Bigr)
	= \frac{1}{|t|} \min\Bigl(1, \frac{\langle x\rangle^2\langle y\rangle^2}{|t|}\Bigr)
	\lesssim \frac{\omega(x)\omega(y)}{|t|\log|t|}, \qquad |t| \ge 2,
	\]
	which yields the desired result.
\end{proof}
\section{High energy decay estimates}\label{sec:high-energy}
In this section, we focus on the high-energy part of $\cos(t\sqrt{H})$ and $\frac{\sin(t\sqrt{H})}{t\sqrt{H}}$, establishing the bounds stated in Theorem~\ref{main_theorem_high} below.
In view of the identities 
\[
	\cos (t \sqrt{H})=\frac{e^{i t \sqrt{H}}+e^{-i t \sqrt{H}}}{2}\ \ \text{and}\ \ \frac{\sin (t \sqrt{H})}{\sqrt{H}}=\frac{e^{i t \sqrt{H}}-e^{-i t \sqrt{H}}}{2i\sqrt{H}},
	\]
   it suffices to analyze the operator $H^{\frac{\gamma}{2}} e^{-it\sqrt{H}}$ for $\gamma \in \{-1, 0\}$.
By using Stone's formula,
\begin{equation}\label{stone_high_1}
  H^{\frac{\gamma}{2}}e^{-i t \sqrt{H}}P_{\mathrm{ac}}(H)\chi_2(H)
= \frac{2}{\pi i} \int_0^{\infty} e^{-it\lambda^2}\lambda^{3+2\gamma} \widetilde{\chi}_2(\lambda) \left[R_V^{+}(\lambda^4)-R_V^{-}(\lambda^4)\right] d \lambda,
\end{equation}
where $\widetilde{\chi}_2(\lambda)=\chi_2(\lambda^4)(\lambda>0)$ so that $\operatorname{supp}\widetilde{\chi}_2 \subset [\lambda_0^{1/4}, \infty)$.
\begin{theorem}\label{main_theorem_high}
Let $|V(x)| \lesssim \langle x \rangle^{-4-}$. Assume that $H=\Delta^2+V$ has no positive embedded eigenvalue and $P_{\mathrm{ac}}(H)$ denotes the projection onto the absolutely continuous spectrum of $H$. Then
$$
\begin{aligned}
\Bigl\|\cos (t \sqrt{H}) P_{\mathrm{ac}}(H)\chi_2(H)\Bigr\|_{L^1\to L^\infty}
+ \Bigl\|\frac{\sin(t\sqrt{H})}{t\sqrt{H}}P_{\mathrm{ac}}(H)\chi_2(H)\Bigr\|_{L^1\to L^\infty}
\lesssim |t|^{-1}.
\end{aligned}
$$
Moreover, for $|t|\ge 2$,
$$
\begin{aligned}
 \Bigl\|\cos (t \sqrt{H}) P_{\mathrm{ac}}(H)\chi_2(H)\Bigr\|_{L^1_\omega\to L^\infty_{-\omega}}
\lesssim ({|t|\log |t|})^{-1}, \quad
\Bigl\|\frac{\sin(t\sqrt{H})}{t\sqrt{H}}P_{\mathrm{ac}}(H)\chi_2(H)\Bigr\|_{L^1\to L^\infty}
\lesssim |t|^{-2}.
\end{aligned}
$$
\end{theorem}
In this section, we use the following resolvent identity:
$$
R_V^{\pm}(\lambda^4)=R_0^{\pm}(\lambda^4)-R_0^{ \pm}(\lambda^4)VR_0^{\pm}(\lambda^4)+R_0^{ \pm}(\lambda^4)VR_V^{\pm}(\lambda^4)VR_0^{ \pm}(\lambda^4).
$$
Substituting this expansion into \eqref{stone_high_1} and 
invoking the free estimates in Remark \ref{rmk:free_case},
we only need to analyze the following integral kernels for $\gamma \in \{-1,0\}$:
$$
\begin{aligned}
\nonumber
& L_{1}^{\pm}(\gamma,t,x,y):=\int_0^{\infty} e^{-i t \lambda^2} \lambda^{3+2\gamma} \widetilde{\chi}_2(\lambda)[R_0^{ \pm}(\lambda^4)VR_0^{ \pm}(\lambda^4)](x, y) d\lambda, \\
&L_{2}^{\pm}(\gamma, t,x,y):=\int_0^{\infty} e^{-i t \lambda^2} \lambda^{3+2\gamma} \widetilde{\chi}_2(\lambda)[R_0^{\pm}(\lambda^4)VR_V^{\pm}(\lambda^4)VR_0^{\pm}(\lambda^4)](x, y) d\lambda.
\end{aligned}
$$
More precisely, the high-energy part of $\cos(t\sqrt{H})$ is dictated by $L_{1}^{\pm}(0,t,x,y)$ and $L_{2}^{\pm}(0,t,x,y)$, while the high-energy part for ${\sin(t\sqrt{H})}/({t\sqrt{H}})$ corresponds to $ t^{-1} L_{1}^{\pm}(-1,t,x,y)$ and $ t^{-1} L_{2}^{\pm}(-1,t,x,y)$. The desired bounds for these kernels are established in Propositions \ref{high_1} and \ref{high_2}, respectively.


\begin{proposition}\label{high_1}
Assume that $|V(x)| \lesssim \left \langle  x\right \rangle^{-4-}$. Then for $\gamma\in \{-1,0\},$
$$
\sup _{x, y \in \mathbb{R}^2}\left|L_{1}^{\pm}(\gamma,t,x,y)\right| \lesssim |t|^{-1}, \ \ \ \
\left|L_{1}^{\pm}(\gamma,t,x,y)\right| \lesssim |t|^{-2} \langle x \rangle^2 \langle y \rangle^2.
$$
Moreover, $|L_1^\pm(-1,t,x,y)|\lesssim 1$  uniformly in $t,x,y.$
\end{proposition}
\begin{proof}
Recalling that $R^\pm_0(\lambda^4)(x,y) = \lambda^{-2}e^{\pm i\lambda|x-y|}\widetilde{R}^\pm(\lambda|x-y|)$ (see \eqref{def:F_pm} and \eqref{free-widetilde_R}), we rewrite
\[
L_{1}^{\pm}(\gamma,t,x,y) = \int_{\mathbb{R}^2} K_1^\pm(t,x,y,y_1)\, V(y_1)\, dy_1,
\]
where
\begin{equation}\label{eq:high_1}
K_1^\pm(t,x,y,y_1):=\int_0^\infty e^{-it\lambda^2}\lambda\,e^{\pm i\lambda r}\, \mathcal{E}^\pm(\lambda, x, y,y_1)\, d\lambda,
\end{equation}
with $r=r(x,y,y_1)=|x-y_1|+|y-y_1|$ and
\begin{equation}\label{eq:def_E}
\mathcal{E}^\pm(\lambda,x,y,y_1) = \lambda^{-2+2\gamma}\widetilde{\chi}_2(\lambda)\,
\widetilde{R}^\pm(\lambda|x-y_1|)\,\widetilde{R}^\pm(\lambda|y-y_1|).
\end{equation}
By virtue of \eqref{free-widetilde_R}, it can be checked that for $\ell=0,1,2,$
$$
\left| \partial_\lambda^{\ell}  \widetilde{R}^\pm(\lambda|x-y|)  \right| \lesssim \langle \lambda |x-y| \rangle^{-\frac{1}{2}} \lambda^{-\ell}.
$$
Then for $\ell=0,1,2$,
\[
\left|\partial_\lambda^{\ell} \mathcal{E}^\pm(\lambda, x, y,y_1)\right|
\lesssim \lambda^{-2+2\gamma-\ell}\langle \lambda |x-y_1| \rangle^{-\frac{1}{2}}
\langle \lambda |y-y_1| \rangle^{-\frac{1}{2}}
\lesssim \lambda^{-2+2\gamma-\ell}\langle \lambda r\rangle^{-\frac{1}{2}}.
\]
Recalling Remark ~\ref{remark:osci} (i), we conclude that
$
|K_1^\pm(t,x,y,y_1)| \lesssim |t|^{-1},
$
uniformly in $x,y,y_1$. Combining this with $\int |V(y_1)| dy_1 < \infty$, 
we deduce that for $\gamma=-1,0,$
$$\sup_{x,y\in\mathbb{R}^2} |L_1^\pm(\gamma,t,x,y)| \lesssim |t|^{-1}.$$
Moreover,  when $\gamma=-1$, it follows directly from \eqref{eq:def_E} that \( |\lambda \mathcal{E}^{\pm}(\lambda,x,y,y_1) | \lesssim \lambda^{-3} \). Then by \eqref{eq:high_1}, \( |K_1^\pm(t,x,y,y_1)| \lesssim 1 \) when $\gamma=-1$. Hence $|L_1^\pm(-1,t,x,y)|\lesssim 1$  uniformly  in $t,x,y.$

 For the pointwise estimate for $L_{1}^{\pm}(\gamma,t,x,y)$, note that for $\ell=0,1,2$,
\[
\Bigl|\partial_\lambda^\ell \bigl( e^{\pm i\lambda r} \mathcal{E}^\pm(\lambda,x,y,y_1)\bigr)\Bigr|
\lesssim\sum_{\ell_1+\ell_2+\ell_3 \leq\ell}\lambda^{-2+2\gamma-\ell_1}|x-y_1|^{\ell_2}|y-y_1|^{\ell_3}
\lesssim\lambda^{-2+2\gamma}\langle x\rangle^\ell\langle y\rangle^\ell\langle y_1\rangle^\ell.
\]
Applying integration by parts twice to the integral \eqref{eq:high_1}(the boundary terms vanish) yields
\[
|K_1^\pm(t,x,y,y_1)| \lesssim |t|^{-2}\langle x\rangle^2\langle y\rangle^2\langle y_1\rangle^2
\int_{\lambda_0}^\infty\lambda^{-3+2\gamma}d\lambda
\lesssim |t|^{-2}\langle x\rangle^2\langle y\rangle^2\langle y_1\rangle^2.
\]
Hence by $\int \langle y_1\rangle^2|V(y_1)| dy_1 < \infty$, it follows that $\left|L_{1}^{\pm}(\gamma,t,x,y)\right| \lesssim |t|^{-2} \langle x \rangle^2 \langle y \rangle^2.$
\end{proof}

We next consider the integral kernels associated with $[R_0^{\pm}(\lambda^4)VR_V^{\pm}(\lambda^4)VR_0^{\pm}(\lambda^4)](x, y).$
To deal with these integral kernels, we utilize the following estimates.
\begin{lemma}[{\cite[Theorem 2.23]{FSY}}]\label{RRR}
Assume that $|V(x)| \lesssim \langle x \rangle^{-\mu}$ for some $\mu > k+1$ with $k \in \mathbb{N}_0$, and that $H = \Delta^2 + V$ has no positive eigenvalues. Then, for any $\sigma>k+\frac{1}{2}, R_V^{ \pm}(\lambda) \in \mathbb{B}\left(L^{2,\sigma}(\mathbb{R}^2), L^{2,-\sigma}(\mathbb{R}^2)\right)$ are $C^k$-continuous for all $\lambda>0$. Furthermore,
$$
\left\|\partial_\lambda^{k} R_V^{ \pm}(\lambda)\right\|_{L^{2,\sigma}(\mathbb{R}^2) \rightarrow L^{2,-\sigma}(\mathbb{R}^2)}=O\left(|\lambda|^{-\frac{3(k+1)}{4}}\right),~\lambda \rightarrow+\infty.
$$
\end{lemma}


\begin{proposition}\label{high_2}
Assume that $|V(x)| \lesssim \langle x \rangle^{-4-}$. Then for $\gamma=-1,0$,
$$
 \sup _{x, y \in \mathbb{R}^2}\left|L_{2}^{\pm}(\gamma,t,x,y)\right| \lesssim \langle t\rangle^{-1}, \ \ \ \
\left|L_{2}^{\pm}(\gamma,t,x,y)\right| \lesssim |t|^{-2} \langle x \rangle^2 \langle y \rangle^2.
$$
\end{proposition}
\begin{proof}
Recalling that $R^\pm_0(\lambda^4)(x,y) = \lambda^{-2}e^{\pm i\lambda|x-y|}\widetilde{R}^\pm(\lambda|x-y|)$, we express the inner product as
\begin{align*}
\big\langle R_V^{\pm}(\lambda^4) V(R_0^{\pm}(\lambda^4)(*, y))(\cdot), VR_0^{\mp}(\lambda^4)(x, \cdot)\big\rangle
= \lambda^{-4} e^{\pm i \lambda(|x|+|y|)} E_{1}^{\pm}(\lambda,x,y),
\end{align*}
where
\[
E_{1}^{\pm}(\lambda,x,y) = \big\langle  R_V^{\pm} V\big(e^{\pm i \lambda(|*-y|-|y|)} \widetilde{R}^{\pm}(\lambda|*-y|)\big)(\cdot), V\big(e^{\mp i \lambda(|x-\cdot|-|x|)} \widetilde{R}^{\mp}(\lambda|x-\cdot|)\big)\big\rangle.
\]

Then we rewrite
\begin{equation}\label{eq:high_2}
  L_{2}^{\pm}(\gamma,t,x,y):=\int_0^{\infty} e^{-it \lambda^2}  e^{\pm i\lambda r} \lambda\mathcal{E}^\pm(\lambda,x,y)  d \lambda, \\
\end{equation}
with $r=r(x,y)=|x|+|y|$ and
$$
\mathcal{E}^{\pm}(\lambda,x,y)=\lambda^{-2+2\gamma}\widetilde{\chi}_2(\lambda) E_{1}^{\pm}(\lambda,x,y).
$$

Note that for $\ell=0,1,2$,
\begin{equation}\label{eq:inequality_1}
 \left|\partial_\lambda^{\ell} e^{\pm i \lambda(|x-\cdot|-|x|)}\right| \lesssim   \langle \cdot \rangle^\ell, \ \ \ 
 \left|\partial_\lambda^{\ell} \widetilde{R}^{\pm}(\lambda |x-\cdot|)\right|
\lesssim \langle \lambda x \rangle^{-\frac{1}{2}} \langle \cdot \rangle^{\frac{1}{2}}\lambda^{\frac{1}{2}-\ell},
\end{equation}
where the second bound utilizes the inequality  $\langle \lambda(x-\cdot) \rangle^{-1/2} \lesssim \langle \lambda x \rangle^{-1/2} \langle  \cdot \rangle^{1/2} \lambda^{1/2}$ for $\lambda \gtrsim 1$.
By Lemma \ref{RRR}, for $\sigma>\ell+\frac{1}{2}$ with $\ell=0,1,2$:
\begin{equation}\label{eq:inequality_2}
\left\|\partial_\lambda^{\ell} R_V^{ \pm}(\lambda^4)\right\|_{L^{2,\sigma} \rightarrow L^{2,-\sigma}} \lesssim \lambda^{-3}.
\end{equation}
Given $|V(x)| \lesssim \langle x \rangle^{-4-}$, by \eqref{eq:inequality_1} and \eqref{eq:inequality_2}, H\"older's inequality yields for $\lambda \gtrsim 1$ and $\ell=0,1,2$,
$$
\begin{aligned}
\left|\partial_\lambda^{\ell} E_{1}^{\pm}(\lambda,x,y)\right|
& \lesssim \frac{1}{\left\langle \lambda x\right\rangle^{\frac{1}{2}}\left\langle \lambda y\right\rangle^{\frac{1}{2}}}\sum_{\ell_1+\ell_2 \le \ell}
\lambda \left\|V(\cdot)\langle\cdot\rangle^{\ell_1+\sigma+\frac{1}{2}}\right\|_{L^2}^2 \cdot\left\|\partial_\lambda^{\ell_2} R_V^{\pm}(\lambda^{4} )\right\|_{L^{2,\sigma} \rightarrow L^{2,-\sigma}}  \\
 & \lesssim \lambda^{-2}  {{\left\langle \lambda x\right\rangle}^{-\frac{1}{2}}{\left\langle \lambda y\right\rangle}^{-\frac{1}{2}}}
 \lesssim \left\langle \lambda (|x|+|y|)\right\rangle^{-\frac{1}{2}}\lambda^{-2},
 \end{aligned}
 $$
with $\sigma > \ell_2 + \frac{1}{2}$.
Hence for $\gamma=-1, 0,$ $\lambda \gtrsim 1$ and $\ell = 0,1,2$,
$$
\left|\partial_\lambda^{\ell} \mathcal{E}^\pm(\lambda,x,y)\right|
=\left|\partial_\lambda^{\ell} \left(\lambda^{-2+2\gamma}\widetilde{\chi}_2(\lambda)E_{1}^{\pm}(\lambda,x,y) \right)\right|
 \lesssim \left\langle \lambda (|x|+|y|)\right\rangle^{-\frac{1}{2}}\lambda^{-\ell}.
$$
Invoking Remark~\ref{remark:osci}(i), we immediately deduce that
 for $\gamma=-1, 0$,
$$\sup_{x,y\in\mathbb{R}^2}|L_2^{\pm}(\gamma,t,x,y)|\lesssim|t|^{-1}.$$
On the other hand, the bound $|E_{1}^{\pm}(\lambda,x,y)| \lesssim \lambda^{-2}$ implies that $|\lambda \mathcal{E}^{\pm}(\lambda,x,y)| \lesssim \lambda^{-3+2\gamma}$. Hence by \eqref{eq:high_2} and $\lambda \gtrsim 1$, we derive that for $\gamma=-1,0,$  $|L_2^\pm(\gamma,t,x,y)|\lesssim 1$  uniformly in $t,x,y.$

 For another point-wise estimate for $L_{2}^{\pm}(\gamma,t,x,y)$,
it can be checked that for $\ell=0,1,2,$
$$
\begin{aligned}
\left|\partial_\lambda^{\ell}\bigl( e^{\pm i\lambda r} \mathcal{E}^\pm(\lambda, x, y) \bigr)\right|
 & \lesssim \sum_{\ell_1+\ell_2+\ell_3 \leq \ell}\lambda^{-2+2\gamma-\ell_1}(|x|+|y|)^{\ell_2} \left|\partial_\lambda^{\ell_3} E_{1}^{\pm}(\lambda,x,y)\right|\lesssim \lambda^{-4+2\gamma}\langle x \rangle^{\ell} \langle y \rangle^{\ell}.
\end{aligned}
$$
Integrating by parts twice in the integral \eqref{eq:high_2}  again gives $\left|L_{2}^{\pm}(\gamma,t,x,y)\right| \lesssim |t|^{-2} \langle x \rangle^2 \langle y \rangle^2$ for $\gamma=-1, 0$.
%\[
%|L_2^{\pm}(\gamma,t,x,y)| \lesssim %|t|^{-2}\langle x\rangle^{2}\langle y\rangle^{2}
%\int_{\lambda_0}^{\infty}\lambda^{-4+2\gamma}\,d%\lambda\lesssim |t|^{-2}\langle %x\rangle^{2}\langle y\rangle^{2}.\]
\end{proof}


\noindent\textbf{Proof of Theorem \ref{main_theorem_high}:} With Propositions \ref{high_1} and \ref{high_2} at hand, we
immediately obtain the pointwise estimate $|t|^{-2} \langle x \rangle^2 \langle y \rangle^2
$
and the uniform  bound $|t|^{-1}$ of $L_j^{\pm}(\gamma,t,x,y)$ for $j=1,2,$ and $\gamma=-1,0.$  In particular, we have
$ |t|^{-1}|L_j^{\pm}(-1,t,x,y)|\lesssim |t|^{-1}\langle t\rangle^{-1}$  for  $j=1,2.$ 
Moreover, by the inequality  $\min(1,a/b) \lesssim \log a/\log b$ (valid for $a,b\ge 2$),
it follows that for $|t|\ge 2$,
\[|L_j^{\pm}(\gamma,t,x,y)|\lesssim
\min\!\Bigl(\frac1{|t|},\frac{\langle x\rangle^{2}\langle y\rangle^{2}}{|t|^{2}}\Bigr)
\lesssim\frac{\omega(x)\omega(y)}{|t|\log|t|}.
\]
Combine these with Remark \ref{rmk:free_case}, and recall that \(\cos(t\sqrt{H})P_{\mathrm{ac}}(H)\chi_2(H)\) corresponds to \(L_{j}^{\pm}(0,t,x,y)\), while ${\sin(t\sqrt{H})}/({t\sqrt{H}})P_{\mathrm{ac}}(H)\chi_2(H)$  corresponds to \( t^{-1} L_{j}^{\pm}(-1,t,x,y)\), then we complete the proof of Theorem \ref{main_theorem_high}.


\section{Low energy estimates in the regular and first kind resonance cases}\label{sec:regular_first}
With  the high-energy estimates in Theorem~\ref{main_theorem_high} established, the proofs of our main results (Theorems~\ref{main_theorem-1}, \ref{main_theorem-2}, \ref{main_theorem-31}, and \ref{main_theorem-32}) are reduced to the low-energy part. 
We begin with a
general strategy, valid for all threshold resonances, and then
specialize in this section to the cases where zero is either a
regular point or a first-kind resonance.

Recall Stone's formulas for the low-energy spectral multipliers:
\begin{align}
	\cos(t\sqrt{H})P_{\mathrm{ac}}(H)\chi_1(H)
	&= \frac{2}{\pi i} \int_0^{\infty} \cos(t\lambda^2)\,\lambda^{3}\,\widetilde{\chi}_1(\lambda)
	\bigl[R_V^{+}(\lambda^4)-R_V^{-}(\lambda^4)\bigr] \, d\lambda,   \label{stone_low_1} \\[4pt]
	\frac{\sin(t\sqrt{H})}{t\sqrt{H}}P_{\mathrm{ac}}(H)\chi_1(H)
	&= \frac{2}{\pi i t} \int_0^{\infty} \sin(t\lambda^2)\,\lambda\,\widetilde{\chi}_1(\lambda)
	\bigl[R_V^{+}(\lambda^4)-R_V^{-}(\lambda^4)\bigr] \, d\lambda, \label{stone_low_2}
\end{align}
where $\widetilde{\chi}_1(\lambda) := \chi_1(\lambda^4)(\lambda>0)$ is supported in $[0,(2\lambda_0)^{1/4}]$.
Substituting the expansion \eqref{eq:M_inverse} into the identity \eqref{id-RV} yields:
\begin{equation*}
	R_V^{\pm}(\lambda^4) = R_0^\pm(\lambda^{4}) - R_0^\pm(\lambda^{4})v \bigg( \sum_{0 \leq \alpha, \beta \leq \mathbf{k}+2} \lambda^{2-k_\alpha-k_\beta} Q_\alpha \mathcal{M}_{\alpha, \beta}^{\pm}(\lambda) Q_\beta \bigg) vR_0^\pm(\lambda^{4}).
\end{equation*}
Inserting this representation back into \eqref{stone_low_1} and \eqref{stone_low_2}, and recalling the free low-energy  estimates from Remark~\ref{rmk:free_case}, 
we reduce the low-energy analysis to bounding a finite sum of perturbed integral kernels.  Specifically, it suffices to estimate the kernel differences
\begin{equation*}
	\big(\mathcal{K}_{B}^+ - \mathcal{K}_{B}^-\big)(t,x,y) \quad \text{and} \quad \big(\mathcal{N}_{B}^+ - \mathcal{N}_{B}^-\big)(t,x,y) ,
\end{equation*}
where the component kernels are explicitly given by
\begin{align}
	\label{K_1}
	& \mathcal{K}_{B}^\pm(t,x,y)=\frac{2}{\pi i}\int_0^{\infty} \cos(t \lambda^2)\lambda \widetilde{\chi}_1(\lambda) \lambda^{4-k_\alpha-k_\beta}\bigl\langle  B^\pm(\lambda)Q_\beta vR_0^\pm(\lambda^{4})(\cdot,y), Q_\alpha vR_0^\mp(\lambda^{4})(\cdot,x)\bigr\rangle \, d \lambda, \\
	\label{Lambda_1}
	& \mathcal{N}_{B}^\pm(t,x,y)=\frac{2}{\pi it}\int_0^{\infty} \sin(t\lambda^2)\lambda^{-1}\widetilde{\chi}_1(\lambda)\lambda^{4-k_\alpha-k_\beta}
	\bigl\langle  B^\pm(\lambda)Q_\beta vR_0^\pm(\lambda^{4})(\cdot,y), Q_\alpha vR_0^\mp(\lambda^{4})(\cdot,x)\bigr\rangle \, d \lambda,
\end{align}
with $B^\pm(\lambda) \in \{ \mathcal{M}_{\alpha, \beta}^{\pm}(\lambda) \mid \alpha,\beta\in\mathbb{Z},\ 0 \le \alpha,\beta \le \mathbf{k}+2 \}.$
\vskip0.3cm
We now specialize to \textbf{the regular and first-kind resonance case},
i.e. $\mathbf k\in\{0,1\}$. The goal of this section is to prove the following  Theorem~\ref{main_theorem_low_1}.
\begin{theorem}\label{main_theorem_low_1}
	Let $H = \Delta^2 + V$ and $|V(x)| \lesssim \langle x \rangle^{-11-}$.
	Assume that $H$ has no positive embedded eigenvalues and that zero is a regular point or a first-kind resonance of $H$. Then
	\begin{equation}\label{eq:thm_unweighted}
		\left\|\cos (t \sqrt{H}) P_{\mathrm{ac}}(H)\chi_1(H)\right\|_{L^1 \rightarrow L^{\infty}}+
		\bigg \|\frac{\sin (t \sqrt{H})}{t\sqrt{H}} P_{\mathrm{ac}}(H)\chi_1(H)\bigg \|_{L^1 \rightarrow L^{\infty}} \lesssim \langle t \rangle^{-1}.
	\end{equation}
	Moreover, in the weighted space setting $L^1_\omega \to L^{\infty}_{-\omega}$, the decay rate improves:
	\begin{equation}\label{eq:thm_weighted}
		\left\|\cos (t \sqrt{H}) P_{\mathrm{ac}}(H)\chi_1(H)\right\|_{L^1_\omega \rightarrow L^{\infty}_{-\omega}}
		+\bigg\|\frac{\sin (t \sqrt{H})}{t\sqrt{H}} P_{\mathrm{ac}}(H)\chi_1(H)\bigg\|_{L^1_\omega \rightarrow L^{\infty}_{-\omega}} \lesssim \frac{1}{ |t| \log |t|}, \quad |t| \ge 2.
	\end{equation}
\end{theorem}

The proof proceeds in two steps: we first establish the unweighted $L^1\to L^\infty$ estimates \eqref{eq:thm_unweighted}, and then  derive the weighted bounds \eqref{eq:thm_weighted}.
Based on the reductions made above, this requires bounding the kernels $(\mathcal{K}_{B}^+ - \mathcal{K}_{B}^-)$ and $(\mathcal{N}_{B}^+ - \mathcal{N}_{B}^-)$  for  $B^\pm(\lambda) = \mathcal{M}_{\alpha,\beta}^{\pm}(\lambda)$ with $0 \le \alpha,\beta \le 3$.


\subsection{The $L^1 \rightarrow L^{\infty}$ estimates}\label{subsec:regular_first_1}
Recall from Theorem~\ref{thm:M_inverse} (I) that for $0 \le \alpha, \beta \le 3$ and $\ell = 0,1,2$,
\[
\Big\| \partial_\lambda^\ell \mathcal{M}_{\alpha,\beta}^\pm(\lambda) \Big\|_{\mathbb{B}(L^2)} \lesssim
\begin{cases}
	\lambda^{-\ell}, & \text{for } ( \alpha,\beta) \neq (3,3), \\
	|\log \lambda| \lambda^{-\ell}, &  \text{for } (\alpha,\beta) = (3,3).
\end{cases}
\]

\begin{proposition}\label{prop_K}
	Let $B^\pm(\lambda) =\mathcal{M}_{\alpha,\beta}^\pm(\lambda)$ with $0 \le \alpha, \beta \le 3$. Then
	\[
	\sup_{x,y\in\mathbb{R}^2} \bigl| \mathcal{K}_{B}^\pm(t,x,y) \bigr| +\sup_{x,y\in\mathbb{R}^2} \bigl| \mathcal{N}_{B}^\pm(t,x,y) \bigr| \lesssim \langle t\rangle^{-1}.
	\]
\end{proposition}

\begin{proof}
	By Lemma~\ref{lemma_projection}, the integral kernel of
\(Q_jvR_0^\pm(\lambda^4)\) admits the representations
	\begin{align}\label{omega0123}
		\bigl(Q_j v R_0^\pm(\lambda^4)\bigr)(\cdot,z)& = \lambda^{-1-\delta_{j 0}}\Omega_{j,\pm}(\lambda,\cdot,y),\\
		\label{omega234}	\bigl(Q_3 v R_0^\pm(\lambda^4)\bigr)(x,y)&=\mathcal{J}_{3, \pm} (\lambda,x,y),
	\end{align}
	where $\delta_{j 0}=1$ if $j=0$ and $0$ otherwise. Moreover, by \eqref{J} and \eqref{lemma_projection_1}, we have for $\ell=0,1,2,$
	\begin{align}
	\bigl\|\partial_\lambda^\ell \bigl(e^{\mp i\lambda|z|}\Omega_{j,\pm}(\lambda,\cdot,z)\bigr)\bigr\|_{L^2} &\lesssim \lambda^{-\ell} \langle \lambda z \rangle^{-1/2}, \label{est:Omega-0-1}\\
	\bigl\|\partial_\lambda^\ell \bigl(e^{\mp i\lambda|z|}\mathcal{J}_{3,\pm}(\lambda,\cdot,z)\bigr)\bigr\|_{L^2} &\lesssim \lambda^{-\ell} |\log\lambda| \, \langle \lambda z \rangle^{-1/2}. \label{est:J-0-1}
\end{align}
We distinguish two cases: $(\alpha,\beta) \neq (3,3)$ and $(\alpha,\beta) = (3,3)$. 
For the first case we use the representation \eqref{omega0123}, and for the second case we use \eqref{omega234}.
	
	

\textbf{Case 1: $(\alpha, \beta) \neq (3,3)$.}
Using the relation $1 + \delta_{\alpha 0} = 2 - k_\alpha$ for $0 \le \alpha \le 3$ (see \eqref{tab:k_alpha}), we obtain
\[
\bigl\langle B^\pm(\lambda) Q_\beta v R_0^{\pm}(\lambda^4)(\cdot, y),\; Q_\alpha v R_0^{\mp}(\lambda^4)(\cdot, x) \bigr\rangle
= \lambda^{k_\alpha+k_\beta-4} \bigl\langle B^\pm(\lambda) \Omega_{\beta,\pm}(\lambda, \cdot, y),\; \Omega_{\alpha,\mp}(\lambda, \cdot, x) \bigr\rangle.
\]
Upon substituting this  
into \eqref{K_1} and \eqref{Lambda_1}, the factor $\lambda^{4-k_\alpha-k_\beta}$ cancels exactly, yielding the canonical oscillatory integrals
\begin{equation}\label{rf}
	\begin{aligned}
		\mathcal{K}_{B}^\pm(t,x,y) &= \frac{2}{\pi i}\int_0^{\infty} \cos(t \lambda^2)\,\lambda\, e^{\pm i\lambda r}\, \mathcal{E}^\pm(\lambda, x,y) \, d\lambda, \\[4pt]
		\mathcal{N}_{B}^\pm(t,x,y) &= \frac{2}{\pi i t}\int_0^{\infty} \sin(t\lambda^2)\,\lambda^{-1} e^{\pm i\lambda r}\, \mathcal{E}^\pm(\lambda, x,y) \, d\lambda,
	\end{aligned}
\end{equation}
where $r = |x| + |y|$, and the common amplitude factor is defined by
\[
\mathcal{E}^\pm(\lambda, x,y) = e^{\mp i\lambda r}\,\widetilde{\chi}_1(\lambda)\,
\bigl\langle B^\pm(\lambda) \Omega_{\beta,\pm}(\lambda,\cdot,y),\; \Omega_{\alpha,\mp}(\lambda,\cdot,x) \bigr\rangle.
\]

In this case the operator satisfies $\|\partial_\lambda^{\ell} B^\pm(\lambda)\|_{\mathbb{B}(L^2)} \lesssim \lambda^{-\ell}$ for $\ell = 0,1,2$. 
Applying Leibniz's rule together with H\"older's inequality gives, for $\ell = 0,1,2$,
\[
\begin{aligned}
	\bigl|\partial_\lambda^{\ell} \mathcal{E}^\pm(\lambda, x,y) \bigr|
	\lesssim \sum_{\ell_1+\ell_2+\ell_3 =\ell}
	&\bigl\|\partial_\lambda^{\ell_1}\bigl(\widetilde{\chi}_1(\lambda) B^\pm (\lambda)\bigr) \bigr\|_{\mathbb{B}(L^2)} \\
	&\times \bigl\|\partial_\lambda^{\ell_2} \bigl(e^{\mp i \lambda|y|}\Omega_{\beta,\pm}(\lambda,\cdot,y)\bigr) \bigr\|_{L^2}
	\bigl\|\partial_\lambda^{\ell_3} \bigl(e^{\pm i \lambda |x|}\Omega_{\alpha,\mp}(\lambda,\cdot,x)\bigr) \bigr\|_{L^2}.
\end{aligned}
\]
Combining this with the estimate \eqref{est:Omega-0-1} yields the amplitude factor bound
\[
\bigl|\partial_\lambda^{\ell} \mathcal{E}^\pm(\lambda, x,y) \bigr|
\lesssim \langle \lambda x\rangle^{-1/2} \langle \lambda y\rangle^{-1/2} \lambda^{-\ell} \widetilde{\chi}_1(\lambda/2)
\lesssim \langle \lambda r \rangle^{-1/2} \lambda^{-\ell} \widetilde{\chi}_1(\lambda/2).
\]
Finally, applying Lemma~\ref{oscillatory}(ii) with parameters $(\sigma, \nu) = (0,0)$ gives the uniform decay $O(\langle t\rangle^{-1})$.




\textbf{Case 2: $(\alpha, \beta) = (3,3)$.}
Here $\|\partial_\lambda^{\ell} B^\pm(\lambda)\|_{\mathbb{B}(L^2)} \lesssim \lambda^{-\ell}|\log\lambda|$. 
To neutralize this logarithmic divergence, we employ the representation \eqref{omega234}, 
which in contrast to \eqref{omega0123} provides an additional factor of $\lambda$ while preserving the  logarithmic singularity (see \eqref{est:J-0-1}).



	
	Substituting the representation  \eqref{omega234} into \eqref{K_1} and \eqref{Lambda_1}, and noting that $\lambda^{4 - k_3 - k_3} = \lambda^2$, we again obtain the canonical oscillatory integral forms \eqref{rf}, where $r = |x| + |y|$ and
\[
\mathcal{E}^\pm(\lambda, x,y) = \lambda^2 e^{\mp i\lambda r}\,\widetilde{\chi}_1(\lambda)\,
\bigl\langle B^\pm(\lambda) \mathcal{J}_{3,\pm}(\lambda,\cdot,y),\; \mathcal{J}_{3,\mp}(\lambda,\cdot,x) \bigr\rangle.
\]
	Applying Leibniz's rule, the $\lambda^2$ factor neutralizes the combined logarithmic singularities. Specifically,
	$$
	\begin{aligned}
			\bigl|\partial_\lambda^{\ell} \mathcal{E}^\pm(\lambda, x,y) \bigr|
		& \lesssim \sum_{\ell_1+\ell_2+\ell_3 =\ell}
		 \bigl\|\partial_\lambda^{\ell_1}\bigl(\widetilde{\chi}_1(\lambda) B^\pm (\lambda)\bigr) \bigr\|_{\mathbb{B}(L^2)} \\
		 & \ \ \ \ \ \ \ \ \ \ \  \times  \bigl\|\partial_\lambda^{\ell_2} \big(e^{\mp i \lambda|y|}\mathcal{J}_{3,\pm}(\lambda,\cdot,y)\big) \bigr\|_{L^2} \bigl\|\partial_\lambda^{\ell_3} \big(e^{\pm i \lambda |x|}\mathcal{J}_{3,\mp}(\lambda,\cdot,x)\big) \bigr\|_{L^2} \\
		&\lesssim \lambda^{2-\ell} |\log\lambda|^3 \langle \lambda x \rangle^{-1/2} \langle \lambda y \rangle^{-1/2} \widetilde{\chi}_1(\lambda/2)
		\lesssim \langle \lambda r \rangle^{-1/2}\lambda^{-\ell}\widetilde{\chi}_1(\lambda/2),
	\end{aligned}
	$$
	where in the last inequality we used $\lambda^2 |\log \lambda|^3 \lesssim 1$ for $0 < \lambda \ll 1$.
	Applying Lemma~\ref{oscillatory}(ii) once again yields the same $O(\langle t\rangle^{-1})$ decay, which completes the proof.
\end{proof}

\subsection{The $L^1_\omega \to L^\infty_{-\omega}$ estimates} \label{subsec:regular_first_2}

We now demonstrate that an improved decay rate of $(|t|\log|t|)^{-1}$ can be achieved by incorporating spatial weights.

Recall that $Q_0$ is the orthogonal projection onto the span of $v$ (i.e., $Q_0=P$); thus, it inherently lacks any cancellation against $v$. In contrast, for any  $(\alpha, \beta) \neq (0,0)$, at least one index is strictly positive. This allows us to exploit the exact cancellation property $Q_\alpha v = 0$ (or $Q_\beta v = 0$) on at least one side, thereby neutralizing the highly singular leading-order term of the free resolvent $R_0^\pm(\lambda^4)$.

Given this structural difference, we divide the analysis  into the following two cases:
\begin{itemize}
	\item[(I)] $B^\pm(\lambda)=\mathcal{M}_{\alpha,\beta}^\pm(\lambda)$ for $(\alpha,\beta)\neq(0,0)$ (see Proposition \ref{prop_K_weight_1});
	\item[(II)] $B^\pm(\lambda)=\mathcal{M}_{0,0}^\pm(\lambda)$ (see Proposition \ref{prop_K_weight_2}).
\end{itemize}

By  Theorem~\ref{thm:M_inverse} (I), we have for $(\alpha,\beta)\neq(0,0)$:
\begin{align}\label{Mneq00}
\bigl\|\partial_\lambda^\ell\mathcal{M}_{\alpha,\beta}^\pm(\lambda)\bigr\|_{\mathbb{B}(L^2)}
	\lesssim   \lambda^{-\ell}|\log \lambda|,  \quad  \ell = 0,1,2,
\end{align}
and
\begin{equation}\label{eq:M_00}
	\mathcal{M}_{0,0}^\pm(\lambda) =(a_0^\pm)^{-1} \|V\|_{L^1(\mathbb{R}^2)}^{-1}Q_0 + \lambda (-\log \lambda)^{3/2}\Lambda^{\pm}(\lambda),
\end{equation}		
where $\Lambda^{\pm}(\lambda) \in \mathbb{B}(L^2)$ satisfies
\(
\big\|\partial_\lambda^{\ell} \Lambda^{\pm}(\lambda) \big\|_{\mathbb{B}(L^2)} \lesssim \lambda^{-\ell}
\)
for $\ell = 0,1,2.$

To obtain the improved decay rates in the weighted setting, we use the following asymptotic expansion of the free resolvent from Lemma \ref{lem-reso}:
\begin{equation}\label{eq:expansion}
	R_0^{\pm}(\lambda^4)(x,y) = a_0^{\pm} \lambda^{-2}G_0(x,y) + E_1^{\pm}(\lambda,x,y)=a_0^{\pm} \lambda^{-2}+ E_1^{\pm}(\lambda,x,y),
\end{equation}
where  $
E_1^{\pm}(\lambda,x,y) = a_1^\pm |x-y|^2 + b_0 |x-y|^2 \log(\lambda |x-y|)
+ O_2\bigl(\lambda^{2}|x-y|^{4}\bigr).
$ For $0<\lambda\ll1$ and $\ell=0,1,2,$ one can verify that
\begin{equation}\label{derivative_1}
	\bigl| \partial_\lambda^\ell E_1^\pm(\lambda,x,y)\bigr| \lesssim \lambda^{-\ell-}\langle x \rangle^{4} \langle y \rangle^{4}.
\end{equation}

In this subsection, we set
\begin{equation}\label{eq:def_mathcal_E}
	\mathcal{E}^\pm(\lambda,x,y)
	= \widetilde{\chi}_1(\lambda) \lambda^{4-k_\alpha-k_\beta}
	\bigl\langle B^\pm(\lambda) \bigl( Q_\beta v R_0^{\pm}(\lambda^4) \bigr)(\cdot,y),
	\bigl( Q_\alpha v R_0^{\pm} (\lambda^4)\bigr)(\cdot,x) \bigr\rangle,
\end{equation}
then  the kernels  $\mathcal{K}_{B}^\pm$ and $\mathcal{N}_{B}^\pm$ defined in \eqref{K_1} and \eqref{Lambda_1} can be expressed as
\begin{align}
	\label{K-N-weig}
	& \mathcal{K}_{B}^\pm=\frac{2}{\pi i}\int_0^{\infty} \cos(t \lambda^2)\lambda  \mathcal{E}^\pm(\lambda,x,y)\, d \lambda, \  \ \mathcal{N}_{B}^\pm=\frac{2}{\pi i t}\int_0^{\infty} \sin(t\lambda^2)\lambda^{-1} \mathcal{E}^\pm(\lambda,x,y) \, d \lambda. \end{align}
Hence, we will focus on the estimate for  $\mathcal{E}^\pm(\lambda,x,y)$ to obtain the desired bound of $\mathcal{K}_{B}^\pm$ and $\mathcal{N}_{B}^\pm$.
\begin{proposition}\label{prop_K_weight_1}
	Let $B^\pm(\lambda)=\mathcal{M}_{\alpha,\beta}^\pm(\lambda)$ with $(\alpha,\beta) \neq (0,0)$. Then for $|t| \ge 2,$
	\[
	\bigl|\mathcal{K}_{B}^\pm(t,x,y)\bigr| +
	\bigl|\mathcal{N}_{B}^\pm(t,x,y)\bigr| \lesssim \frac{\omega(x)\omega(y)}{|t| \log|t|}.
	\]
\end{proposition}
\begin{proof}
	For $\alpha \ge 1$, the exact cancellation property $Q_\alpha v = 0$ (Remark~\ref{cancellationQ}) ensures that the highly singular term $a_0^{\pm} \lambda^{-2}$ vanishes in \eqref{eq:expansion}, yielding
	\[
	\bigl( Q_\alpha v R_0^{\pm}(\lambda^4)\bigr)(x,y) = \bigl( Q_\alpha v E_1^{\pm}(\lambda,\cdot,y) \bigr)(x).
	\]
	Thus, applying \eqref{derivative_1}, we obtain for $0 < \lambda \ll 1$ and  $\ell=0,1,2$:
	\[
	\|\partial_\lambda^\ell  \bigl(Q_\alpha v R_0^{\pm}(\lambda^4)\bigr)(\cdot,y)\|_{L^2}
	\lesssim\|v \partial_\lambda^\ell E_1^\pm(\lambda,\cdot,y)\|_{L^2}
	\lesssim \langle y\rangle^4\lambda^{-\ell-}.
	\]
	Conversely, for $\alpha = 0$ (where no such cancellation occurs), the $a_0^{\pm} \lambda^{-2}$ term remains, then
	\[
	\|\partial_\lambda^\ell \bigl(Q_0 v R_0^{\pm}(\lambda^4)\bigr)(\cdot,y)\|_{L^2}
	\lesssim \lambda^{-2-\ell}\|v\|_{L^2} + \langle y\rangle^4\lambda^{-\ell-}
	\lesssim \langle y\rangle^4\lambda^{-2-\ell}.
	\]
	Combine these bounds with the  estimates \eqref{Mneq00} for $\mathcal{M}_{\alpha,\beta}^\pm$ and note that the worst case
	occurs when exactly one of the indices $\alpha,\beta$ is zero, then via Leibniz's rule and H\"{o}lder's inequality,
	\begin{align}\label{eq:estimate_E}
		\bigl| \partial_\lambda^\ell \mathcal{E}^\pm(\lambda,x,y) \bigr|
		\lesssim \lambda^{1/2-\ell}\widetilde{\chi}_1(\lambda/2) \langle x\rangle^{4} \langle y\rangle^{4}, \ \  \ell = 0,1,2, \ \text{and } \ (\alpha,\beta) \neq (0,0).
	\end{align}
Applying Lemma~\ref{oscillatory} with  $r(x,y)=0$ and  $(\sigma,\nu)=(-1/2,0)$, we obtain that
	\begin{equation}\label{eq:bound_K_N}
		\bigl|\mathcal{K}_{B}^\pm(t,x,y)\bigr| + \bigl|\mathcal{N}_{B}^\pm(t,x,y)\bigr| \lesssim |t|^{-5/4}\langle x\rangle^4\langle y\rangle^4.
	\end{equation}
	
	On the other hand, Proposition \ref{prop_K} provides the uniform estimate
	\[
	\sup_{x,y \in \mathbb{R}^2} \left| \mathcal{K}_{B}^\pm(t,x,y) \right|
	+
	\sup_{x,y \in \mathbb{R}^2} \left| \mathcal{N}_{B}^\pm(t,x,y) \right| \lesssim
	|t|^{-1} .
	\]
	Combining this with \eqref{eq:bound_K_N}
	gives that
	\begin{equation}\label{eq:inequality_3}
		\min\Bigl(\frac{1}{|t|},
		\frac{\langle x\rangle^{4}\langle y\rangle^{4}}{|t|^{5/4}}\Bigr)
		=\frac{1}{|t|}\min\!\Bigl(1,
		\frac{\langle x\rangle^{4}\langle y\rangle^{4}}{|t|^{1/4}}\Bigr)
		\lesssim\frac{\omega(x)\omega(y)}{|t|\log|t|}, \qquad |t| \ge 2.
	\end{equation}
	This completes the proof.
\end{proof}


\begin{proposition}\label{prop_K_weight_2}
	Let $B^\pm(\lambda)=\mathcal{M}_{0,0}^\pm(\lambda)$. Then for $|t| \ge 2$,
	\begin{align}\label{K-N-cri}
		\bigl|\mathcal{K}_{B}^\pm(t,x,y)\bigr|  \lesssim \frac{\omega(x)\omega(y)}{|t|\log|t|},  \quad
		(\mathcal{N}_{B}^+ - \mathcal{N}_{B}^-)(t,x,y)= \frac{1}{8|t|} +O\left(\frac{\omega(x)\omega(y)}{|t|\log|t|}\right).
	\end{align}
\end{proposition}

\begin{proof}
	Substituting  $R_0^{\pm}(\lambda^4)(x,y) = a_0^{\pm} \lambda^{-2}G_0(x,y) + E_1^{\pm}(\lambda,x,y)$ into \eqref{eq:def_mathcal_E} and observing that $(Q_0vG_0)(\cdot,y)=(Q_0v)(\cdot)=v(\cdot)$, we separate an $x,y$-independent principal part and a remainder:
	\[
	\mathcal{E}^\pm(\lambda,x,y) = \widetilde{\chi}_1(\lambda) (a_0^\pm)^2 \bigl\langle \mathcal{M}_{0,0}^\pm(\lambda) v, v \bigr\rangle + \mathcal{E}_{\mathrm{rem}}^\pm(\lambda,x,y).
	\]
	Since $E_1^\pm$ satisfies the bounds \eqref{derivative_1}, $\mathcal{E}_{\mathrm{rem}}^\pm$ obeys the estimate \eqref{eq:estimate_E}, ensuring that its contribution to both kernels is bounded by $|t|^{-5/4}\langle x\rangle^4\langle y\rangle^4$.
	
	For the principal part, using the asymptotic expansion \eqref{eq:M_00} for $\mathcal{M}_{0,0}^\pm(\lambda)$, we obtain that
	\[
	\widetilde{\chi}_1(\lambda) (a_0^\pm)^2 \bigl\langle \mathcal{M}_{0,0}^\pm(\lambda) v, v \bigr\rangle = a_0^\pm \|V\|_{L^1}^{-1} \bigl\langle Q_0 v, v \bigr\rangle \widetilde{\chi}_1(\lambda) + \lambda (-\log \lambda)^{3/2} \widetilde{\chi}_1(\lambda) (a_0^\pm)^2 \bigl\langle \Lambda^{\pm}(\lambda) v, v \bigr\rangle.
	\]
	The second term on the right satisfies the bounds \eqref{eq:estimate_E}, yielding that its contribution to both kernels is bounded by $|t|^{-5/4}\langle x\rangle^4\langle y\rangle^4$.
	
	It remains to evaluate the integrals  \eqref{K-N-weig} with  $\mathcal{E}^\pm=a_0^\pm \|V\|_{L^1}^{-1} \bigl\langle Q_0 v, v \bigr\rangle\widetilde{\chi}_1(\lambda)=a_0^\pm\widetilde{\chi}_1(\lambda)$. For the cosine kernel,  by integrating by parts, we have
	\[
	a_0^\pm \int_0^\infty \cos(t\lambda^2)\lambda\widetilde{\chi}_1(\lambda)\,d\lambda  = O(|t|^{-N}),  \ \  \forall N>0.
	\]
	Combining this with the remainder bounds above gives $|\mathcal{K}_{B}^\pm(t,x,y)| \lesssim |t|^{-5/4}\langle x\rangle^4\langle y\rangle^4$. Interpolating this with the uniform $|t|^{-1}$ bound exactly as in \eqref{eq:inequality_3} yields the desired weighted estimate for $\mathcal{K}_{B}^\pm$.
	
	For the sine kernel difference, extracting the explicit leading-order integral gives
	$$
	\begin{aligned}
		(\mathcal{N}_{B}^+ - \mathcal{N}_{B}^-)(t,x,y) = (a_0^+ - a_0^-) \frac{2}{\pi i t} \int_0^\infty \sin(t\lambda^2) \lambda^{-1} \widetilde{\chi}_1(\lambda) \, d\lambda + O\Bigl(\frac{\langle x\rangle^4\langle y\rangle^4}{|t|^{5/4}}\Bigr).
	\end{aligned}
	$$
	Using the substitution $\rho = |t|\lambda^2$ and $\int_0^\infty \frac{\sin u}{u} du = \frac{\pi}{2},$  we obtain
	\begin{equation}\label{eq:integral_sin}
		\frac{1}{t}\int_0^{\infty} \sin(t\lambda^2) \lambda^{-1} \widetilde{\chi}_1(\lambda) \, d\lambda
		= \frac{\pi}{4|t|}
		- \frac{1}{2|t|} \int_0^{\infty} \frac{\sin\rho}{\rho}  \widetilde{\chi}_2\bigl(\sqrt{\frac{\rho}{|t|}}\bigr) \, d\rho=\frac{\pi}{4|t|}+O\bigl(\frac{1}{|t|^{2}}\bigr),
	\end{equation}
	where an integration by parts shows that the $\rho$-integral is   $O(|t|^{-1})$.
	Combining this with $a_0^{+} -a_0^{-}= \frac{i}{4}$ yields
	$$
	(\mathcal{N}_{B}^+ - \mathcal{N}_{B}^-)(t,x,y)=\frac{1}{8 |t|}+ O\Bigl(\frac{\langle x\rangle^4\langle y\rangle^4}{|t|^{5/4}}\Bigr).
	$$
	Combine this spatial remainder with the uniform bound $|t|^{-1}$, then by \eqref{eq:inequality_3}, we  obtain  the desired asymptotic expansion.
\end{proof}
\begin{remark}
{\rm
In the regular and first-kind resonance cases, all $\mathcal{N}_{B}^\pm$ satisfy the desired weighted bound $\omega(x)\omega(y)(|t|\log|t|)^{-1}$ for all pairs, save for the  contribution \eqref{K-N-cri} (i.e., the case $B^\pm(\lambda)=\mathcal{M}_{0,0}^\pm(\lambda)$).
Combining this with Proposition~\ref{prop_free_sin}, we see that the leading term ${1}/{(8|t|)}$ of the free sine propagator is exactly annihilated by the leading term in \eqref{K-N-cri}, yielding the weighted sine bound in \eqref{eq:thm_weighted}. The  cosine bound follows immediately, as both the free and perturbed cosine kernels
intrinsically exhibit the required decay.}
\end{remark}


\section{The second kind resonance case}\label{sec:second}
This section establishes the low-energy estimates in the
second-kind resonance case. We first translate the moment conditions
in Theorem \ref{main_theorem-2} into the projection framework used
in the resolvent expansion.

By Proposition \ref{characterizations-1} (also see Proposition \ref{characterizations}), when $V\neq0$, zero is a
second-kind resonance if and only if
$
S_3L^2\neq\{0\}$, and $
S_4L^2=\{0\}.
$
Moreover, the moment conditions on the resonance space
$\mathcal E$ correspond exactly to the following algebraic
conditions on the projections introduced in Definitions \ref{defS} and \ref{defQ}:
\begin{itemize}
\item 
$\ \ \tau\not\equiv0
\quad\Longleftrightarrow\quad
S_3L^2\neq S_4^0L^2
\quad\Longleftrightarrow\quad
Q_4^0=S_3-S_4^0\neq0;
$
\vskip0.2cm
\item $
\ \ \ker\mathfrak M_2\neq\{0\}
\quad\Longleftrightarrow\quad
S_4^1L^2\neq\{0\}
\quad\Longleftrightarrow\quad
Q_{42}^1=S_4^1-S_4\neq0,
$
\end{itemize}
where the last equivalence uses $S_4=0$ in the second-kind
resonance case.

Thus the mixed second-moment configuration in
Theorem \ref{main_theorem-2} is equivalent to
$
Q_4^0\neq0$ and
$Q_{42}^1\neq0,
$
whereas its complement is characterized by
$
Q_4^0=0$ or
$Q_{42}^1=0.
$


In view of this reduction, together with the free estimates in
Section \ref{sec:free_case} and the high-energy estimates in
Theorem \ref{main_theorem_high}, it remains to prove the following
Theorem \ref{main_theorem_low_2}.




\begin{theorem}\label{main_theorem_low_2}
Let $H = \Delta^2 + V$ with $|V(x)| \lesssim \langle x \rangle^{-14-}$.
Assume that $H$ has no positive embedded eigenvalues and that zero is a second-kind resonance of $H$. We distinguish two cases:
\begin{itemize}
    \item If $Q_4^0 \neq 0$ and $Q_{42}^1 \neq 0$, then
$$
\begin{aligned}
& \left\|\cos (t \sqrt{H}) P_{\mathrm{ac}}(H)\chi_1(H)\right\|_{L^1 \rightarrow L^{\infty}}
+ \bigg \|\frac{\sin (t \sqrt{H})}{t\sqrt{H}} P_{\mathrm{ac}}(H)\chi_1(H)\bigg \|_{L^1 \to L^{\infty}}
\lesssim \langle t \rangle^{-1}(\log (2+|t|))^2,
\end{aligned}
$$
and in the weighted setting $L^1_\omega\to L^\infty_{-\omega}$, for $|t| \ge 2,$
     \begin{align*}
		\left\|
			\cos(t\sqrt H)P_{\mathrm{ac}}(H)\chi_1(H)
			\right\|_{L^1_\omega\to L^\infty_{-\omega}}\lesssim \frac{\log|t|}{|t|},\quad \bigg\|
		\frac{\sin(t\sqrt H)}{t\sqrt H}
			P_{\mathrm{ac}}(H)\chi_1(H)
			\bigg\|_{L^1_\omega\to L^\infty_{-\omega}}
			\sim
			\frac{(\log|t|)^2}{|t|}. 
		\end{align*}

    \item If $Q_4^0 = 0$ or $Q_{42}^1 = 0$,  then  the unweighted  bound improves to
\begin{equation}\label{2-t}
    \left\|\cos (t \sqrt{H}) P_{\mathrm{ac}}(H)\chi_1(H)\right\|_{L^1 \to L^{\infty}} +
    \bigg\|\frac{\sin (t \sqrt{H})}{t\sqrt{H}} P_{\mathrm{ac}}(H) \chi_1(H) \bigg\|_{L^1 \to L^{\infty}} \lesssim \langle t \rangle^{-1},
\end{equation}
and in the weighted setting $L^1_\omega\to L^\infty_{-\omega}$, for $|t| \ge 2,$
\[
\left \|\cos (t \sqrt{H}) P_{\mathrm{ac}}(H)\chi_1(H)\right\|_{L^1_\omega\to L^\infty_{-\omega}} \lesssim\frac{1}{|t|\log |t|},  \ \ \bigg\|\frac{\sin (t \sqrt{H})}{t\sqrt{H}} P_{\mathrm{ac}}(H)\chi_1(H)\bigg\|_{L^1_\omega\to L^\infty_{-\omega}}\sim|t|^{-1}.
\]
      \end{itemize}
      \end{theorem}

\subsection{The estimates for the case $Q_{4}^0 \neq 0$ and $Q_{42}^1 \neq 0$} \label{subsec:second_1}  
Building upon the $L^1\to L^\infty$ dispersive bounds already established for the regular and first-kind resonance cases in Subsection \ref{subsec:regular_first_1}, we only need to bound the kernel differences $(\mathcal{K}_{B}^+-\mathcal{K}_{B}^-)(t,x,y)$ and $(\mathcal{N}_{B}^+-\mathcal{N}_{B}^-)(t,x,y)$ defined in \eqref{K_1} and \eqref{Lambda_1}, restricting our attention to the operators $B^\pm(\lambda)=\mathcal{M}_{\alpha, \beta}^{\pm}(\lambda)$ with the index pairs $(\alpha, \beta)$ satisfying either $\alpha=4$ with $0 \le \beta \le 4$, or $0 \le \alpha \le 4$ with $\beta=4$.  

Recall from Theorem \ref{thm:M_inverse} (II) that for $\alpha = 1,2,3$:  
\[
	\|\partial_\lambda^\ell \mathcal{M}_{4,\alpha}^\pm(\lambda)\|_{\mathbb{B}(L^2)} + \|\partial_\lambda^\ell \mathcal{M}_{\alpha,4}^\pm(\lambda)\|_{\mathbb{B}(L^2)} \lesssim \lambda^{1-\ell}|\log \lambda|^{4}, \quad \ell = 0,1,2.  
\]
Moreover, for pairs involving indices 0 and 4, we have:  
\begin{align}  
	\mathcal{M}_{0,4}^{\pm}(\lambda) &= (\log \lambda)^{-1} Q_0 \Lambda^\pm(\lambda) Q_{4}^0 + Q_0 \Lambda^\pm(\lambda) Q_{4}^1, \nonumber\\  
	\mathcal{M}_{4,0}^{\pm}(\lambda) &= (\log \lambda)^{-1} Q_{4}^0 \Lambda^\pm(\lambda) Q_0 + Q_{4}^1 \Lambda^\pm(\lambda) Q_0, \nonumber\\ 
	\label{M44-1}
    	\mathcal{M}_{4,4}^{\pm}(\lambda) &= Q_{4} \Lambda^\pm(\lambda) Q_{4} + (\log\lambda)\bigl(Q_{4} \Lambda^\pm(\lambda) Q_{4}^1+Q_{4}^1 \Lambda^\pm(\lambda) Q_{4}\bigr) + (\log\lambda)^2 Q_{4}^1 \Lambda^\pm(\lambda) Q_{4}^1.
\end{align} 
Throughout the paper, $\Lambda^\pm(\lambda)$ denotes  a generic operator in $\mathbb{B}(L^2)$, possibly different at each occurrence, satisfying
$
\|\partial_\lambda^\ell \Lambda^\pm(\lambda)\|_{\mathbb{B}(L^2)} \lesssim \lambda^{-\ell}, \  \ell=0,1,2.
$

In this subsection, we first establish the unweighted estimate in
Proposition~\ref{prop_second_2_3}. We then derive the weighted estimate and  the asymptotic
expansion associated with 
\(B^\pm(\lambda)=\mathcal M_{4,4}^\pm(\lambda)\), whose contribution exhibits 
slowest decay, in Proposition~\ref{prop_second_1}. Finally, in
Proposition~\ref{sine-sharp}, we prove the sharp decay rate for the
sine evolution.

\begin{proposition}\label{prop_second_2_3}
	Let $B^\pm(\lambda) \in \{\mathcal{M}_{4,\alpha}^{\pm}(\lambda), \mathcal{M}_{\alpha,4}^{\pm}(\lambda)\}$ with $\alpha \in \{0,1,2,3,4\}$. Then
	\[  
\sup _{x, y \in \mathbb{R}^2}\left| \mathcal{K}_{B}^\pm(t,x,y)\right| + \sup _{x, y \in \mathbb{R}^2}\left|\mathcal{N}_{B}^\pm (t,x,y)\right| \lesssim
	\begin{cases}
	\langle t\rangle^{-1}, & \alpha\in\{0,1,2,3\},\\
	\langle t\rangle^{-1}\bigl(\log(2+|t|)\bigr)^2, & {\alpha} = 4.
\end{cases}  
	\]  
\end{proposition}

\begin{proof}
	By Lemma \ref{lemma_projection}, the projected free resolvents relevant
	to the present proposition satisfy
	\begin{align}
		\bigl(Q_\alpha vR_0^\pm(\lambda^4)\bigr)(\cdot,z)
		&=
		\lambda^{-1-\delta_{\alpha0}}
		\Omega_{\alpha,\pm}(\lambda,\cdot,z),
		\label{Qalpha}\\
		\bigl(Q_4vR_0^\pm(\lambda^4)\bigr)(\cdot,z)
		&=
		\mathcal{J}_{4,\pm}(\lambda,\cdot,z),
		\label{De all Q}\\
		\bigl(Q_4^1vR_0^\pm(\lambda^4)\bigr)(\cdot,z)
		&=
		(b_0Q_4^1vG_2^0)(\cdot,z)
		+
		\mathcal{T}_4^1(\lambda,\cdot,z)
		+
		\mathcal{T}_{4,\pm}^1(\lambda,\cdot,z).
		\label{eq:resolvent_decomp-Q4}
	\end{align}
	Moreover, for $\ell=0,1,2$,
	\begin{equation}
		\begin{split}
			&	\big\|
			\partial_\lambda^\ell
			\bigl(
			e^{\mp i\lambda|z|}
			\mathcal{J}_{4,\pm}(\lambda,\cdot,z)
			\bigr) 
			\big\|_{L^2} \lesssim 
			\lambda^{-\ell} 	|\log\lambda|
			\langle\lambda z\rangle^{-1/2},\\
		&
		\big\|
		\partial_\lambda^\ell
		\bigl(
		e^{\mp i\lambda|z|}
		\Omega_{\alpha,\pm}(\lambda,\cdot,z)
		\bigr)
		\big\|_{L^2}
		+
		\big\|
		\partial_\lambda^\ell
		\bigl(
		e^{\mp i\lambda|z|}
		\mathcal{T}_{4,\pm}^1(\lambda,\cdot,z)
		\bigr)
		\big\|_{L^2}
		\lesssim
		\lambda^{-\ell}
		\langle\lambda z\rangle^{-1/2},\\
	&	\|
		\partial_\lambda^\ell
		\mathcal{T}_4^1(\lambda,\cdot,z)
		\|_{L^2}
		\lesssim
		\lambda^{-\ell},\quad 	\|(b_0Q_4^1vG_2^0)(\cdot,z)\|_{L^2}
		\lesssim1.
		\label{eq:E_bound-1}
		\end{split}
	\end{equation}
	We divide the proof of this proposition into three cases: $\alpha=0; 1\leq\alpha\leq3;$ and $\alpha=4.$
	\begin{comment}
	We shall also repeatedly use
	\begin{equation}\label{eq:weight-product}
		\langle\lambda x\rangle^{-1/2}
		\langle\lambda y\rangle^{-1/2}
		\lesssim
		\langle\lambda(|x|+|y|)\rangle^{-1/2}.
	\end{equation}
	Derivatives falling on $\widetilde{\chi}_1(\lambda)$ will always be
	absorbed into the slightly enlarged cutoff
	$\widetilde{\chi}_1(\lambda/2)$.
	\end{comment}
	\medskip
	
	
	\textbf{Case 1: $B^\pm=\mathcal{M}_{0,4}^\pm$ or
		$\mathcal{M}_{4,0}^\pm$.}
	We consider $\mathcal{M}_{0,4}^\pm$  (the symmetric case $\mathcal{M}_{4,0}^{\pm}$ follows identically). Write
	$$
	\mathcal{M}_{0,4}^\pm(\lambda)
	=
	B_0^\pm(\lambda)+B_1^\pm(\lambda),
	$$
	where
	$
	B_0^\pm(\lambda)
	=
	(\log\lambda)^{-1}
	Q_0\Lambda^\pm(\lambda)Q_4^0$
	and $
	B_1^\pm(\lambda)
	=
	Q_0\Lambda^\pm(\lambda)Q_4^1.
	$ 	By linearity, it suffices to estimate the contributions of
	$B_0^\pm$ and $B_1^\pm$ separately.
	
	We first treat $B_0^\pm$. Since
	$
	4-k_0-k_4=2,
	$
	substituting \eqref{Qalpha} and \eqref{De all Q} into both
	\eqref{K_1} and \eqref{Lambda_1} gives
	\begin{align}
		\mathcal{K}_{B_0}^\pm(t,x,y)
		={}&
		\frac{2}{\pi i}
		\int_0^\infty
		\cos(t\lambda^2)
		e^{\pm i\lambda(|x|+|y|)}
		\lambda
		\mathcal{E}_{0}^\pm(\lambda,x,y)
		\,d\lambda,
		\label{eq:K-B0-standard}\\
			\mathcal{N}_{B_0}^\pm(t,x,y)
		={}&
		\frac{2}{\pi it}
		\int_0^\infty
		\sin(t\lambda^2)
		e^{\pm i\lambda(|x|+|y|)}
		\lambda^{-1}
		\mathcal{E}_{0}^\pm(\lambda,x,y)
		\,d\lambda,
		\label{eq:N-B0-standard}
	\end{align}
	where
	\begin{align}
		\mathcal{E}_{0}^\pm(\lambda,x,y)
		={}&
		\widetilde{\chi}_1(\lambda)
		(\log\lambda)^{-1}
		\Bigl\langle
		Q_0\Lambda^\pm(\lambda)Q_4^0
		\Bigl(
		e^{\mp i\lambda|y|}
		\mathcal{J}_{4,\pm}(\lambda,\cdot,y)
		\Bigr),
		e^{\pm i\lambda|x|}
		\Omega_{0,\mp}(\lambda,\cdot,x)
		\Bigr\rangle.
		\label{eq:E-B0}
	\end{align}
	Indeed, the factor $\lambda^2$ arising from
	$\lambda^{4-k_0-k_4}$ exactly cancels the factor $\lambda^{-2}$
	in the $Q_0$-resolvent expansion.
	
	Observe that  the cosine and sine kernels have exactly the forms required in
	Lemma \ref{oscillatory}, with
	$
	r(x,y)=|x|+|y|
	$
	and the common amplitude factor $\mathcal{E}_{0}^\pm$ given by
	\eqref{eq:E-B0}.
Since $\|\partial_\lambda^\ell \Lambda^\pm(\lambda)\|_{\mathbb{B}(L^2)} \lesssim \lambda^{-\ell}$ for $\ell=0,1,2$ and the factor $(\log\lambda)^{-1}$ exactly compensates for the logarithmic growth of $\mathcal{J}_{4,\pm}$ in \eqref{eq:E-B0}, it follows from Leibniz's rule, H\"older's inequality, and \eqref{eq:E_bound-1} that
	\begin{equation*}
		%\label{eq:E-B0-bound}
		\left|
		\partial_\lambda^\ell
		\mathcal{E}_{0}^\pm(\lambda,x,y)
		\right|
		\lesssim
		\widetilde{\chi}_1(\lambda/2)
		\langle\lambda(|x|+|y|)\rangle^{-1/2}
		\lambda^{-\ell},
		\qquad
		\ell=0,1,2.
	\end{equation*}
 Lemma
	\ref{oscillatory}(ii) with $(\sigma,\nu)=(0,0)$ therefore yields
	the desired $\langle t\rangle^{-1}$ bound for both
	$\mathcal{K}_{B_0}^\pm$ and $\mathcal{N}_{B_0}^\pm$.
	
	We next treat $B_1^\pm$. By
	\eqref{eq:resolvent_decomp-Q4}, the corresponding kernels 	$\mathcal{K}_{B_1}^\pm$ and $\mathcal{N}_{B_1}^\pm$ split into
	three contributions, associated respectively with
	$Q_4^1vG_2^0$, $\mathcal{T}_4^1$, and
	$\mathcal{T}_{4,\pm}^1$. The same calculation leading to
	\eqref{eq:K-B0-standard} and \eqref{eq:N-B0-standard} shows that each
	contribution has the form required in Lemma \ref{oscillatory}. We
	therefore only specify the corresponding phases and amplitude factors.
	
	The contribution of $Q_4^1vG_2^0$ has the phase
	$
	r(x,y)=|x|
	$
	and the amplitude factor
	\begin{equation*}
		%\label{eq:E-B1-G}
		\mathcal{E}_{G}^\pm(\lambda,x,y)
		=
		\widetilde{\chi}_1(\lambda)
		\Bigl\langle
		Q_0\Lambda^\pm(\lambda)Q_4^1
	(Q_4^1vG_2^0)(\cdot,y),
		e^{\pm i\lambda|x|}
		\Omega_{0,\mp}(\lambda,\cdot,x)
		\Bigr\rangle.
	\end{equation*}
	The contribution of $\mathcal{T}_4^1$ has the same phase
	$r(x,y)=|x|$ and  the amplitude factor
	\begin{equation*}
		%\label{eq:E-B1-T}
		\mathcal{E}_{\mathcal{T}}^\pm(\lambda,x,y)
		=
		\widetilde{\chi}_1(\lambda)
		\Bigl\langle
		Q_0\Lambda^\pm(\lambda)Q_4^1
		\mathcal{T}_4^1(\lambda,\cdot,y),
		e^{\pm i\lambda|x|}
		\Omega_{0,\mp}(\lambda,\cdot,x)
		\Bigr\rangle.
	\end{equation*}
	Finally, the contribution of $\mathcal{T}_{4,\pm}^1$ has the  phase
	$
	r(x,y)=|x|+|y|
	$
	and the  amplitude factor
	\begin{align*}
		\mathcal{E}_{osc}^\pm(\lambda,x,y)
		={}&
		\widetilde{\chi}_1(\lambda)
		\Bigl\langle
		Q_0\Lambda^\pm(\lambda)Q_4^1
		\Bigl(
		e^{\mp i\lambda|y|}
		\mathcal{T}_{4,\pm}^1(\lambda,\cdot,y)
		\Bigr),
		e^{\pm i\lambda|x|}
		\Omega_{0,\mp}(\lambda,\cdot,x)
		\Bigr\rangle.
	%	\label{eq:E-B1-osc}
	\end{align*}
	
	By \eqref{eq:E_bound-1}, the derivative
	bounds for $\Lambda^\pm$, Leibniz's rule, and H\"older's inequality, each of these amplitude factors satisfies
	\begin{equation*}
		%\label{eq:E-B1-bound}
		\left|
		\partial_\lambda^\ell
		\mathcal{E}^\pm(\lambda,x,y)
		\right|
		\lesssim
		\widetilde{\chi}_1(\lambda/2)
		\langle\lambda r\rangle^{-1/2}
		\lambda^{-\ell},
		\qquad
		\ell=0,1,2,
	\end{equation*}
	where $\mathcal{E}^\pm$ denotes any of
	$\{
	\mathcal{E}_{G}^\pm,
	\mathcal{E}_{\mathcal{T}}^\pm,
	\mathcal{E}_{osc}^\pm\}
	$
	and $r$ is the corresponding phase specified above. Then applying Lemma
	\ref{oscillatory}(ii) with $(\sigma,\nu)=(0,0)$ yields  the desired $\langle t\rangle^{-1}$ bound for both
	$\mathcal{K}_{B_1}^\pm$ and $\mathcal{N}_{B_1}^\pm$.
	
	Combining the estimates for $B_0^\pm$ and $B_1^\pm$ gives
	$$
	\sup_{x,y\in \mathbb{R}^2}
	\bigl|
	\mathcal{K}_{\mathcal{M}_{0,4}}^\pm(t,x,y)
	\bigr|
	+
	\sup_{x,y\in \mathbb{R}^2}
	\bigl|
	\mathcal{N}_{\mathcal{M}_{0,4}}^\pm(t,x,y)
	\bigr|
	\lesssim
	\langle t\rangle^{-1}.
	$$

	

	\textbf{Case 2: $B^\pm=\mathcal{M}_{\alpha,4}^\pm$ or
		$\mathcal{M}_{4,\alpha}^\pm$, $1\leq\alpha\leq3$.}
	We consider $\mathcal{M}_{\alpha,4}^\pm$ (the symmetric case $\mathcal{M}_{4,\alpha}^{\pm}$ follows identically). Since
	$
	4-k_\alpha-k_4=1
	$
	and
	$
	\bigl(
	Q_\alpha vR_0^\mp(\lambda^4)
	\bigr)(\cdot,x)
	=
	\lambda^{-1}
	\Omega_{\alpha,\mp}(\lambda,\cdot,x),
	$
	the explicit factor $\lambda$ arising from
	$\lambda^{4-k_\alpha-k_4}$  in
	\eqref{K_1} and \eqref{Lambda_1} exactly cancels this
	$\lambda^{-1}$ singularity. By repeating the calculation in
	\eqref{eq:K-B0-standard}--\eqref{eq:N-B0-standard}, the two kernels $	\mathcal{K}_{\mathcal{M}_{\alpha,4}}^\pm$ and $	\mathcal{N}_{\mathcal{M}_{\alpha,4}}^\pm$ 
	take the forms required in Lemma \ref{oscillatory} with the  phase
	$
	r(x,y)=|x|+|y|
	$
	and the  amplitude factor
	\begin{align*}
		\mathcal{E}_{\alpha4}^\pm(\lambda,x,y)
		={}&
		\widetilde{\chi}_1(\lambda)
		\Bigl\langle
		\mathcal{M}_{\alpha,4}^\pm(\lambda)
		\Bigl(
		e^{\mp i\lambda|y|}
		\mathcal{J}_{4,\pm}(\lambda,\cdot,y)
		\Bigr),
		e^{\pm i\lambda|x|}
		\Omega_{\alpha,\mp}(\lambda,\cdot,x)
		\Bigr\rangle.
	%	\label{eq:E-alpha4}
	\end{align*}
	Recall that
$
\|\partial_\lambda^\ell \mathcal{M}_{\alpha, 4}^{\pm}(\lambda)\|_{\mathbb{B}(L^2)} \lesssim \lambda^{1-\ell}|\log\lambda|^{4}$ for   $\ell=0,1,2.$ 
Applying Leibniz's rule and H\"older's inequality alongside \eqref{eq:E_bound-1}, we see that the additional $\lambda$ factor from $ \mathcal{M}_{\alpha, 4}^{\pm}(\lambda)$ absorbs the logarithmic singularities:
\[  
\bigl| \partial_\lambda^{\ell} \mathcal{E}_{\alpha4}^\pm(\lambda,x,y) \bigr| \lesssim \widetilde{\chi}_1(\lambda/2) \langle \lambda r \rangle^{-1/2} \lambda^{1-\ell}|\log\lambda|^5 \lesssim \widetilde{\chi}_1(\lambda/2) \langle \lambda r \rangle^{-1/2} \lambda^{-\ell}, \quad \ell = 0,1,2. 
\]   
	Hence Lemma \ref{oscillatory}(ii) again gives
	$$
	\sup_{x,y}
	\bigl|
	\mathcal{K}_{\mathcal{M}_{\alpha,4}}^\pm(t,x,y)
	\bigr|
	+
	\sup_{x,y}
	\bigl|
	\mathcal{N}_{\mathcal{M}_{\alpha,4}}^\pm(t,x,y)
	\bigr|
	\lesssim
	\langle t\rangle^{-1}.
	$$
	
	
	\textbf{Case 3: $B^\pm=\mathcal{M}_{4,4}^\pm$.}
	This case  relies on the representations \eqref{De all Q} and \eqref{eq:resolvent_decomp-Q4}.
	Note that the expansion for $Q_4v R_0^\pm(\lambda^4)$ contains an $|\log\lambda|$ singularity, whereas the expansion for $Q_{4}^1v R_0^\pm(\lambda^4)$ does not. Moreover, none of the bounds in \eqref{eq:E_bound-1} carry the spatial weight $\omega(z)$.
	Note that  $4-k_4-k_4=0$, the intrinsic $\lambda$-factor in the definitions \eqref{K_1} and \eqref{Lambda_1} is $\lambda^0 = 1$.
Substituting  \eqref{De all Q}, \eqref{eq:resolvent_decomp-Q4} and  the  expansion 
		\begin{align}\label{eq:M44_rearranged-1}
		\mathcal{M}_{4,4}^{\pm}(\lambda) = Q_{4} \Lambda^\pm(\lambda) Q_{4} + (\log\lambda)\bigl(Q_{4} \Lambda^\pm(\lambda) Q_{4}^1+Q_{4}^1 \Lambda^\pm(\lambda) Q_{4}\bigr) + (\log\lambda)^2 Q_{4}^1 \Lambda^\pm(\lambda)Q_{4}^1,
	\end{align}
	 into \eqref{K_1} and \eqref{Lambda_1}, the kernels $\mathcal{K}_B^\pm$ and $\mathcal{N}_B^\pm$ reduce to finite linear combinations of canonical oscillatory integrals of the form:
	\begin{equation}\label{eq:canonical_oscillatory_prop1}
		\frac{2}{\pi i}	\int_0^\infty \lambda \cos(t\lambda^2) e^{\pm i\lambda r(x,y)} \mathcal{E}_{f,g}^\pm(\lambda, x, y) \, d\lambda \quad \text{and} \quad
		\frac{2}{\pi it} \int_0^\infty \lambda^{-1} \sin(t\lambda^2) e^{\pm i\lambda r(x,y)} \mathcal{E}_{f,g}^\pm(\lambda, x, y) \, d\lambda,
	\end{equation}
	where the amplitude factor $\mathcal{E}_{f,g}^\pm$ is defined via the inner product:
	\begin{equation}\label{eq:amplitude_def_prop1}
		\mathcal{E}_{f,g}^\pm(\lambda, x, y) = e^{\mp i\lambda r(x,y)} \widetilde{\chi}_1(\lambda) \big\langle M_{sub}^\pm(\lambda) f(\lambda, \cdot, y), \, g(\lambda, \cdot, x) \big\rangle.
	\end{equation}
	Here, $M_{sub}^\pm(\lambda)$ is a constituent operator from \eqref{eq:M44_rearranged-1}, while $f$ and $g$ are corresponding states selected from the expansions in \eqref{De all Q} and \eqref{eq:resolvent_decomp-Q4}. The phase function $r(x,y) \in \{0, |x|, |y|, |x|+|y|\}$ neutralizes the exponential phases embedded within $f$ and $g$.

	To uniformly bound \eqref{eq:canonical_oscillatory_prop1}, we apply H\"older's inequality to \eqref{eq:amplitude_def_prop1}, tracing the $\lambda$-logarithmic singularities inherited from each component:
	\begin{itemize}
		\item For $M_{sub}^\pm(\lambda) = Q_{4} \Lambda^\pm(\lambda) Q_{4}$, it acts between the states $f = \mathcal{J}_{4,\pm}$ and $g = \mathcal{J}_{4,\mp}$. The product of their individual $L^2$-bounds yields an amplitude factor constrained by $ |\log\lambda|^2$.
		\vskip0.2cm
		\item For $M_{sub}^\pm(\lambda) = (\log\lambda) Q_{4} \Lambda^\pm(\lambda) Q_{4}^1$ (the symmetric case $(\log\lambda) Q_{4}^1 \Lambda^\pm(\lambda)Q_{4}$ follows identically),
		the operator pairs $f\in\{b_0 Q_{4}^1 v G_2^0, \mathcal{T}_{4,\pm}^{1}, \mathcal{T}_{4}^{1}\}$ and $g = \mathcal{J}_{4,\mp}(\lambda, \cdot, x)$. The inherent logarithm in the operator multiplies the singularity of $g$, yielding a total $\lambda$-singularity of $ |\log\lambda|^2$.
		\vskip0.2cm
		\item For $M_{sub}^\pm(\lambda) = (\log\lambda)^2 Q_{4}^1 \Lambda^\pm(\lambda) Q_{4}^1$, this operator pairs with the bounded, $\log$-free $Q_{4}^1$-expansions on both sides, i.e., $f,g\in\{b_0 Q_{4}^1 v G_2^0, \mathcal{T}_{4,\pm}^{1}, \mathcal{T}_{4}^{1}\}.$  The $\lambda$-singularity arises entirely from the operator itself and is bounded by $|\log\lambda|^2$.
	\end{itemize}
	Thus, for all cross-terms, applying Leibniz's rule and H\"older's inequality to the amplitude factor yields
	\[
	\bigl| \partial_\lambda^{\ell} \mathcal{E}_{f,g}^\pm(\lambda,x,y) \bigr| \lesssim \widetilde{\chi}_1(\lambda/2) \langle \lambda r(x,y) \rangle^{-1/2} \lambda^{-\ell} |\log\lambda|^2, \quad \ell = 0,1,2.
	\]
	An application of Lemma \ref{oscillatory}(ii) with parameters $(\sigma,\nu)=(0,-2)$ then yields $$
	\sup_{x,y\in \mathbb{R}^2}
	\bigl|
	\mathcal{K}_{\mathcal{M}_{4,4}}^\pm(t,x,y)
	\bigr|
	+
	\sup_{x,y\in \mathbb{R}^2}
	\bigl|
	\mathcal{N}_{\mathcal{M}_{4,4}}^\pm(t,x,y)
	\bigr|
	\lesssim
	\langle t\rangle^{-1}\bigl(\log(2+|t|)\bigr)^2.
	$$
\end{proof}
Next, we turn to establishing the asymptotic expansion in the weighted setting \(L^1_\omega\to L^\infty_{-\omega}\), for which a more precise description of \(\mathcal M_{4,4}^\pm(\lambda)\) than that provided by \eqref{M44-1} is required. Returning to the proof of Theorem~\ref{thm:M_inverse} in Section~\ref{subsec:proof_inverse}, we combine \eqref{D40D04D44_2}, \eqref{Upsilon44}, and 
$
\mathcal M_{4,4}^\pm(\lambda)=\mathcal D_{4,4}^\pm(\lambda)+\Upsilon_{4,4}^\pm(\lambda)
$
to derive the refined expansion:
\begin{align}\label{huaM44}
\mathcal{M}_{4,4}^{\pm}(\lambda) &= Q_4^0\mathcal{C}^\pm Q_4^0 +(\log \lambda) \bigl(Q_4^0 \mathcal{C}_1^\pm Q_{42}^1 + Q_{42}^1 \mathcal{C}_2^\pm Q_4^0 \bigr)+ (\log \lambda)^2 (a_2^\pm)^{-1} Q_{42}^1 \mathcal{S}^{-1} Q_{42}^1 \nonumber \\
&  \quad + (\log \lambda)^{-1} Q_4^0\Lambda^\pm(\lambda) Q_4^0 + \bigl(Q_4^0 \Lambda^\pm(\lambda) Q_{4}^1 + Q_{4}^1 \Lambda^\pm(\lambda) Q_4^0\bigr)  +(\log \lambda) Q_{4}^1 \Lambda^\pm(\lambda) Q_{4}^1.  
\end{align}  
Here, with $\mathcal{W}_1 = (b_0 Q_0 v G_2 v Q_4^0)^{-1} Q_0 T_0 Q_{42}^1,$ the operators $\mathcal{C}^\pm, \mathcal{C}_1^\pm, \mathcal{C}_2^\pm$ are given by  
\begin{equation}\label{eq:mathcal_C}  
\mathcal{C}^\pm := (a_2^\pm)^{-1}\mathcal{W}_1 \mathcal{S}^{-1} \mathcal{W}_1^*, \quad  
\mathcal{C}_1^\pm := - (a_2^\pm)^{-1}\mathcal{W}_1 \mathcal{S}^{-1}, \quad  
\mathcal{C}_2^\pm := - (a_2^\pm)^{-1}\mathcal{S}^{-1}\mathcal{W}_1^*,
\end{equation}  
and $\mathcal{S}$ is the strictly positive operator on $Q_{42}^1 L^2$
defined by  $$
\mathcal{S} =\mathcal{W}_1^* \bigl[Q_4^0 v G_4 v Q_4^0 - Q_4^0 v G_4 v Q_{41}^1 (Q_{41}^1 v G_4 v Q_{41}^1)^{-1} Q_{41}^1 v G_4 v Q_4^0\bigr] \mathcal{W}_1,
$$
where  the inverse $(Q_{41}^1 v G_4 v Q_{41}^1)^{-1}$ is interpreted
as the zero operator  
if $Q_{41}^1L^2=\{0\}$.






Since \(\mathcal K_B^\pm\) and \(\mathcal N_B^\pm\) are uniformly
\(O(|t|^{-1})\) whenever
\(B^\pm(\lambda)=\mathcal M_{\alpha,\beta}^\pm(\lambda)\) with
\((\alpha,\beta)\neq(4,4)\), the desired weighted estimate reduces
to the contribution of
\(B^\pm(\lambda)=\mathcal M_{4,4}^\pm(\lambda)\), which is treated in
Proposition~\ref{prop_second_1}.

\begin{proposition}\label{prop_second_1}  
	Let $B^\pm(\lambda)= \mathcal{M}_{4,4}^{\pm}(\lambda)$. Then  
 for $|t|\geq 2$, 
\begin{align*}
	%\label{weight-44}	
\mathcal{K}_{B}^\pm(t,x,y)&=O\bigl(\omega(x)\omega(y)|t|^{-1} \log|t|\bigr),\nonumber\\
\mathcal{N}_{B}^\pm(t,x,y)&=	\mathcal{A}^\pm(x,y)\frac{(\log|t|)^2}{8i|t|}+{O}\bigl(\omega(x)\omega(y)|t|^{-1}\log|t|\bigr), 
	\end{align*} 
	where the  kernel $\mathcal{A}^\pm(x,y)$ is defined by:
	\begin{align}\label{eq:mathcal_A_pm}
		\mathcal{A}^\pm(x,y) :={}& b_0^2 \langle \mathcal{C}^\pm Q_4^0(v |\cdot|^2),  Q_4^0(v |\cdot|^2)\rangle 
		+ b_0^2 \langle  \mathcal{C}_1^\pm (Q_{42}^1 v G_2^0)(\cdot,y),  Q_4^0(v |\cdot|^2)\rangle  \nonumber\\
		&+ b_0^2 \langle  \mathcal{C}_2^\pm Q_4^0(v |\cdot|^2),  (Q_{42}^1 v G_2^0)(\cdot,x)\rangle + b_0^2 (a_2^\pm)^{-1} \bigl\langle  \mathcal{S}^{-1} (Q_{42}^1 v G_2^0)(\cdot,y), (Q_{42}^1 v G_2^0)(\cdot,x) \bigr\rangle.
	\end{align}
\end{proposition}  



\begin{proof} 
	We require the following detailed  resolvent expansions from  Lemma \ref{lemma_projection}:
	\begin{equation}\label{eq:resolvent_decomp}  \begin{split}
			\bigl(Q_4^0 v R_0^\pm(\lambda^4)\bigr)(\cdot,z) &=  b_0 (\log \lambda) Q_4^0(v|\cdot|^2) + (b_0 Q_4^0 v G_2^0)( \cdot,z) + \mathcal{T}_{4,\pm}^0(\lambda, \cdot,z)+\mathcal{T}_{4}^0(\lambda, \cdot,z),  \\
			\bigl(Q_{4}^1v R_0^\pm(\lambda^4)\bigr)(\cdot,z) & =(b_0 Q_{4}^1 v G_2^0)( \cdot,z) + \mathcal{T}_{4,\pm}^{1}(\lambda, \cdot,z)+\mathcal{T}_{4}^{1}(\lambda, \cdot,z), \\
			\bigl(Q_{42}^1v	R_0^\pm(\lambda^4)\bigr)(\cdot,z) &=  (b_0 Q_{42}^1v G_2^0)( \cdot,z) + \lambda\mathcal{T}_{4,\pm}^{1,2}(\lambda, \cdot,z)+\lambda\mathcal{T}_{4}^{1,2}(\lambda, \cdot,z).
		\end{split}
	\end{equation}  
	By \eqref{lemma_projection_G}--\eqref{lemma_projection_1}, we have \begin{equation}\label{eq:E_bound}  
		\begin{split}   
			&\|(b_0 Q_4^0 v G_2^0)(\cdot, z)\|_{L^2} \lesssim \omega(z) \   \text{and}\ \|(b_0 Qv G_2^0)(\cdot, z)\|_{L^2} \lesssim1\  \text{with}\ Q=Q_{4}^1, Q_{42}^1,
		\end{split}  
	\end{equation}      and for $\ell=0,1,2,$
    \begin{equation}\label{eq:E_bound-2} 
    \begin{split}
    &\big\| \partial_\lambda^\ell \bigl( e^{\mp i\lambda|z|}\mathcal{T}_{\pm}(\lambda, \cdot, z) \bigr) \big\|_{L^2}  \lesssim \lambda^{-\ell} \langle \lambda z \rangle^{-1/2} , \quad \mathcal{T}_\pm \in \{  \mathcal{T}_{4, \pm}^0, \, \mathcal{T}_{4, \pm}^{1},\ \mathcal{T}_{4, \pm}^{1,2}\}\\ 
   & \big\| \partial_\lambda^\ell \mathcal{T}_4^0(\lambda, \cdot, z) \big\|_{L^2} \lesssim \lambda^{-\ell}\omega(z),\ \ 
    \big\| \partial_\lambda^\ell \mathcal{T}(\lambda, \cdot, z) \big\|_{L^2} \lesssim \lambda^{-\ell}, \quad \mathcal{T} \in \{  \mathcal{T}_4^{1},\ \mathcal{T}_4^{1,2}\}.\end{split}
	\end{equation}
	We partition the operator \eqref{huaM44} into subdominant and dominant categories based on their singular behavior.
	\vskip0.2cm
	\textbf{Case 1: Subdominant terms.} 
		Consider the following terms in $\mathcal{M}_{4,4}^{\pm}(\lambda)$:   
	\begin{equation} \label{some-op}
		\begin{split}
			& (\log \lambda)^{-1} Q_4^0\Lambda^\pm(\lambda) Q_4^0, \quad Q_4^0 \Lambda^\pm(\lambda) Q_{4}^1, \quad Q_{4}^1 \Lambda^\pm(\lambda) Q_4^0, \quad
			(\log \lambda) Q_{4}^1 \Lambda^\pm(\lambda)Q_{4}^1.
		\end{split}
	\end{equation} 
	Similarly, upon substitution of \eqref{eq:resolvent_decomp} and the term $B^\pm(\lambda)$ from \eqref{some-op} into the integral kernels \eqref{K_1} and \eqref{Lambda_1}, the kernels $\mathcal{K}_B^\pm$ and $\mathcal{N}_B^\pm$ expand into finite linear combinations of canonical oscillatory integrals of the same form as \eqref{eq:canonical_oscillatory_prop1}.

	
	Considering that in \eqref{eq:resolvent_decomp} only the term $b_0 (\log \lambda) Q_4^0(v|\cdot|^2)$ carries the $(\log\lambda)$ singularity, and combining this with \eqref{eq:E_bound}, \eqref{eq:E_bound-2} and \eqref{some-op}, for all cross-terms, applying Leibniz's rule and H\"older's inequality to the corresponding amplitude factor yields
	\[
	\bigl| \partial_\lambda^{\ell} \mathcal{E}_{f,g}^\pm(\lambda,x,y) \bigr| \lesssim \widetilde{\chi}_1(\lambda/2) \langle \lambda r(x,y) \rangle^{-1/2} \lambda^{-\ell} |\log\lambda| \, \omega(x)\omega(y), \quad \ell = 0,1,2.
	\]
Applying Lemma \ref{oscillatory}(ii) with $(\sigma,\nu)=(0,-1)$ guarantees that all contributions from this class are  absorbed into the  bound ${O}\bigl(\omega(x)\omega(y)\langle t \rangle^{-1}\log(2+|t|)\bigr)$.
	\vskip0.2cm
	\textbf{Case 2: Dominant terms.} 
	The exact asymptotic expansions are driven exclusively by the four principal operators generating $(\log\lambda)^2$ singularities:
	\[
	Q_4^0\mathcal{C}^\pm Q_4^0, \quad (\log \lambda)Q_4^0 \mathcal{C}_1^\pm Q_{42}^1, \quad (\log \lambda)Q_{42}^1 \mathcal{C}_2^\pm Q_4^0, \quad (\log \lambda)^2 (a_2^\pm)^{-1} Q_{42}^1 \mathcal{S}^{-1} Q_{42}^1.
	\]
	To maximize the logarithmic singularity, these operators must rigidly pair with the spatial states that do not carry an extra $\lambda$ decay. Given the precise structure of the $Q_4^0v R_0^\pm(\lambda^4)$ and $Q_{42}^1v R_0^\pm(\lambda^4)$ expansions in \eqref{eq:resolvent_decomp}, the dominant amplitude  components strictly evaluate as follows:
	\begin{itemize}
		\item The operator $Q_4^0\mathcal{C}^\pm Q_4^0$ pairs with $b_0(\log\lambda)Q_4^0(v |\cdot|^2)$ from both expansions, yielding the amplitude factor:
		\[
		b_0^2 (\log\lambda)^2\widetilde{\chi}_1(\lambda) \langle Q_4^0\mathcal{C}^\pm Q_4^0(v |\cdot|^2), \, Q_4^0(v |\cdot|^2)\rangle.
		\]
		\item For $(\log \lambda)Q_4^0 \mathcal{C}_1^\pm Q_{42}^1$, since
		the remainders $\lambda\mathcal{T}_{4,\pm}^{1,2}$ and $ \lambda\mathcal{T}_{4}^{1,2}$ in the $Q_{42}^1v	R_0^\pm(\lambda^4)$ expansion  carry explicit $\lambda$ decay, the $y$-variable state is forced to pair with  $(b_0 Q_{42}^1 v G_2^0)(\cdot, y)$, while the $x$-variable state pairs with $b_0(\log\lambda)Q_4^0(v |\cdot|^2)$. The resulting amplitude factor is:
		\[
		b_0^2 (\log\lambda)^2\widetilde{\chi}_1(\lambda) \langle Q_4^0 \mathcal{C}_1^\pm (Q_{42}^1 v G_2^0)(\cdot,y), \, Q_4^0(v |\cdot|^2)\rangle.
		\]
		\item Symmetrically, the operator $(\log \lambda)Q_{42}^1 \mathcal{C}_2^\pm Q_4^0$ contributes the amplitude factor:
		\[
		b_0^2 (\log\lambda)^2\widetilde{\chi}_1(\lambda) \langle Q_{42}^1 \mathcal{C}_2^\pm Q_4^0(v |\cdot|^2), \, (Q_{42}^1 v G_2^0)(\cdot,x)\rangle.
		\]
		\item Finally, the intrinsic $(\log\lambda)^2$ factor in $(a_2^\pm)^{-1}(\log \lambda)^2  Q_{42}^1 \mathcal{S}^{-1} Q_{42}^1$ dictates that maximum growth is attained if and only if $Q_{42}^1$ pairs with $(b_0 Q_{42}^1 vG_2^0)$ on both sides. This yields:
		\[
		b_0^2 (\log\lambda)^2 (a_2^\pm)^{-1}\widetilde{\chi}_1(\lambda) \langle \mathcal{S}^{-1} (Q_{42}^1 v G_2^0)(\cdot,y), \, (Q_{42}^1 v G_2^0)(\cdot,x) \rangle.
		\]
	\end{itemize}
	Summing these dominant contributions (all possessing $r(x,y) = 0$), and recognizing that all other cross-terms exhibit at most $|\log\lambda|$ singularities falling into Case 1, we  finally derive that  for $B^\pm(\lambda)= \mathcal{M}_{4,4}^{\pm}(\lambda),$
	\begin{equation}  \label{K-N}	\begin{split}
    \mathcal{K}_{B}^{\pm}(t,x,y) &= \frac{2}{\pi i}\mathcal{A}^\pm(x,y) I(t) + {O}\bigl(\omega(x)\omega(y)\langle t \rangle^{-1}\log(2+|t|)\bigr), \\ \mathcal{N}_{B}^{\pm}(t,x,y) &= \frac{2}{\pi i}\mathcal{A}^\pm(x,y) J(t) + O\bigl(\omega(x)\omega(y)\langle t \rangle^{-1}\log(2+|t|)\bigr).  \end{split}
	\end{equation}    
	The temporal integrals are given by:
	\begin{align*}
		I(t) &= \int_0^\infty\cos(t\lambda^2)\lambda\widetilde{\chi}_1(\lambda)(\log\lambda)^2 \, d\lambda, \quad\text{and}\quad
		J(t) = t^{-1}\int_0^\infty \sin(t\lambda^2) \lambda^{-1} \widetilde{\chi}_1(\lambda) (\log \lambda)^2 \, d\lambda.
	\end{align*}
	By Lemma \ref{lemma:I_J}, these integrals admit the following asymptotic behaviors for $|t| \ge 2$:
	\begin{align*}    
		I(t) &= \frac{\pi}{8} \frac{\log |t|}{|t|} + {O}\big(|t|^{-1}\big), \quad \text{and} \quad J(t) = \frac{\pi}{16} \frac{(\log |t|)^2}{|t|} + {O}\Big(\frac{\log |t|}{|t|}\Big).
	\end{align*}
    Substituting the evaluations for $I(t)$ and $J(t)$ into \eqref{K-N} completes the proof.
\end{proof}


Next, we are devoted to the proof of the sharp bound for the sine evolution.

\begin{proposition}\label{sine-sharp}
	If $Q_4^0 \neq 0$ and $Q_{42}^1 \neq 0,$  then
\begin{equation}\label{optimality_0}
	\Big\|\frac{\sin (t \sqrt{H})}{t\sqrt{H}} P_{\mathrm{ac}}(H)\chi_1(H)\Big\|_{L^1_\omega \to L^{\infty}_{-\omega}} \sim \frac{(\log |t|)^2}{|t|},\ \ |t|\ge2.
\end{equation}  
\end{proposition}

\begin{proof}  
	Collating the results from Propositions \ref{prop_second_2_3}--\ref{prop_second_1} and the regular and first-kind resonance cases (Subsection~\ref{subsec:regular_first_1}),    we derive   the  following  asymptotic expansions  in the weighted space $L^1_\omega \to L^\infty_{-\omega}:$
	\begin{align}\label{expan-sin}
	\frac{\sin(t\sqrt{H})}{t\sqrt{H}} P_{\mathrm{ac}}(H)\chi_1(H)  
	= -\frac{(\log|t|)^2}{8i|t|}(\mathcal{A}^+ - \mathcal{A}^-) + \mathcal{O}\!\left(|t|^{-1}\log |t|\right),\ \ |t|\ge 2,
	\end{align}
	where $\mathcal{A}^\pm$ is the integral operator with the kernel $\mathcal{A}^\pm(x,y).$ 
	From the expression \eqref{eq:mathcal_A_pm} for the kernel $\mathcal{A}^\pm(x,y)$ and the fact that $G_2(x,y)=|x-y|^2$ satisfies $(Q_4^0 v G_2)(x,y) = Q_4^0(v|\cdot|^2)(x)$,  we derive 
	\[
	\begin{aligned}
		\mathcal{A}^+ - \mathcal{A}^- &= b_0^2 G_2 v Q_4^0(\mathcal{C}^+-\mathcal{C}^-) Q_4^0 v G_2
		+ b_0^2 G_2 v Q_4^0(\mathcal{C}_1^+-\mathcal{C}_1^-) Q_{42}^1 v G_2^0 \\  
		&\quad + b_0^2 G_2^0 v Q_{42}^1 (\mathcal{C}_2^+-\mathcal{C}_2^-) Q_4^0 v G_2  
		+ b_0^2\bigl[(a_2^+)^{-1}-(a_2^-)^{-1}\bigr] G_2^0 v Q_{42}^1 \mathcal{S}^{-1} Q_{42}^1 v G_2^0.  
	\end{aligned}
	\]  
	Consequently, by  \eqref{expan-sin}, the proof of \eqref{optimality_0} reduces to verifying that $\mathcal{A}^+ - \mathcal{A}^-$ is not identically zero.

Introduce the operators  
\[  
\mathcal{W}_1 = (b_0 Q_0 v G_2 v Q_4^0)^{-1} Q_0 T_0 Q_{42}^1, \quad  
A = G_2 v Q_4^0, \quad  
\mathcal{B} = G_2^0 v Q_{42}^1,  
\]  
and substitute the expressions defined in \eqref{eq:mathcal_C} for $\mathcal{C}^\pm, \mathcal{C}_1^\pm, \mathcal{C}_2^\pm$, then we derive
\[  
\mathcal{A}^+ - \mathcal{A}^- = b_0^2 \bigl[(a_2^+)^{-1} - (a_2^-)^{-1}\bigr] (A \mathcal{W}_1 - \mathcal{B}) \mathcal{S}^{-1} (A \mathcal{W}_1 -\mathcal{B})^*.  
\]  
Since $b_0^2[(a_2^+)^{-1} - (a_2^-)^{-1}] \neq 0$ and $\mathcal{S}^{-1}$ is positive definite on $Q_{42}^1 L^2$, the problem reduces to demonstrating that the operator $A\mathcal{W}_1 - \mathcal{B}$ is injective on the finite-dimensional space $Q_{42}^1 L^2$.  
	
	For any $\Psi \in Q_4^0 L^2$, orthogonality dictates that $\int_{\mathbb{R}^2} y^\alpha v(y) \Psi(y)  dy = 0$ for $|\alpha| \le 1$. Recalling $G_2(x,y) = |x|^2 + |y|^2 - 2x \cdot y$, we see that $A$ maps $\Psi$ to a constant:  
	$$  
	(A \Psi)(x) = \int_{\mathbb{R}^2} |y|^2 v(y) \Psi(y) \, dy =: C(\Psi).  
	$$  
	Thus, setting $\Psi = \mathcal{W}_1 \psi \in Q_4^0 L^2$ for any $\psi \in Q_{42}^1 L^2$, we obtain $(A \mathcal{W}_1 \psi)(x) \equiv C(\mathcal{W}_1 \psi)$.  
	
	Suppose that $(A\mathcal{W}_1 - \mathcal{B})\psi=0$ for some
	$\psi\in Q_{42}^1L^2$. Since
	$
	A \mathcal{W}_1 \psi
	$
	is constant,
	$
	G_2^0 v Q_{42}^1\psi
	$
	is constant. Applying $\Delta^2$ in the distributional sense and
	using
	$
	\Delta^2G_2^0
	=
	8\pi\delta_0,
	$
	we obtain
	$
	vQ_{42}^1\psi=v\psi=0.
	$ Thus, \begin{align}\label{psi1}
\psi=0,
\qquad\text{almost everywhere on }\{x:v(x)\neq0\}.
\end{align} Moreover,
	 by Proposition \ref{characterizations-1}(ii) and the inclusion $\psi \in Q_{42}^1 L^2 \subset S_3 L^2$, there exists $\phi \in W_{-1}(\mathbb{R}^2)$ such that $\psi = U v \phi$, implying that \begin{align}\label{psi2}
\psi=0,
\qquad\text{almost everywhere on }\{x:v(x)=0\}.
\end{align}
 Consequently, $\psi = 0$.
Therefore, $A\mathcal{W}_1 - \mathcal{B}$ is injective, then $\mathcal{A}^+ - \mathcal{A}^-$ represents a non-trivial operator, concluding the proof of optimality.  
\end{proof}  

\subsection{The estimates for the case $Q_{4}^0 = 0$ or $Q_{42}^1 = 0$}\label{subsec:second_2}
This condition $Q_{4}^0 = 0$ or $Q_{42}^1 = 0$ allows us to establish an unweighted $| t |^{-1}$ decay (i.e., \eqref{2-t}) and the following weighted estimates:  
\[ 
	\left \|\cos (t \sqrt{H}) P_{\mathrm{ac}}(H)\chi_1(H)\right\|_{L^1_\omega\to L^\infty_{-\omega}} \lesssim\frac{1}{|t|\log |t|},  \quad \bigg\|\frac{\sin (t \sqrt{H})}{t\sqrt{H}} P_{\mathrm{ac}}(H)\chi_1(H)\bigg\|_{L^1_\omega\to L^\infty_{-\omega}}\sim|t|^{-1}.  
\]  
We divide the analysis into two parts, starting with the unweighted $| t |^{-1}$ estimate.  

\vspace{2mm}  
\textbf{Part I. The $L^1\to L^\infty$ estimate.}  
Recall from Theorem \ref{thm:M_inverse} (II) that for $\alpha = 1,2,3$,
\[
\|\partial_\lambda^\ell \mathcal{M}_{4,\alpha}^\pm(\lambda)\|_{\mathbb{B}(L^2)} + \|\partial_\lambda^\ell \mathcal{M}_{\alpha,4}^\pm(\lambda)\|_{\mathbb{B}(L^2)} \lesssim \lambda^{1-\ell}|\log \lambda|^{4}, \quad \ell = 0,1,2;
\]
for the pairs $(\alpha, \beta) = (0,4), (4,0), (4,4)$,
        \begin{itemize}
        \item if $Q_4^0 = 0$, then $Q_4 = Q_4^1$ and
        \begin{align}\label{Q4=Q41}
         \mathcal{M}_{0, 4}^\pm(\lambda) =  Q_0 \Lambda^\pm(\lambda) Q_4^1, \ \  \mathcal{M}_{4, 0}^\pm(\lambda) =  Q_4^1 \Lambda^\pm(\lambda) Q_0, \ \
         \mathcal{M}_{4,4}^\pm(\lambda) = Q_4^1 \Lambda^\pm(\lambda) Q_4^1.
        \end{align}
       
    \item If $Q_4^0 \neq 0$ and $Q_{42}^1 = 0$, then $Q_4 = Q_4^0 \oplus Q_{41}^1$, and
        \begin{align}
		& \mathcal{M}_{4,0}^\pm(\lambda)=(\log \lambda)^{-1}  Q_4 \Lambda^\pm(\lambda) Q_0,  \quad
          \mathcal{M}_{0,4}^\pm(\lambda)=(\log \lambda)^{-1}  Q_0 \Lambda^\pm(\lambda) Q_4, \nonumber   \\ \label{2-M4-4}
		& \mathcal{M}_{4,4}^\pm(\lambda) = \sum_{h,l\in\{0,1\}}(\log \lambda)^{-(2-h-l)} Q_4^h \Lambda^\pm(\lambda) Q_4^l.
        \end{align}
        \end{itemize}
Given the $L^1\to L^\infty$ dispersive bounds established for the regular and first-kind resonance cases in Subsection \ref{subsec:regular_first_1}, it remains to prove Proposition \ref{prop_second_2_1}.


\begin{proposition}\label{prop_second_2_1}
	Let $B^\pm(\lambda) \in \{\mathcal{M}_{4,\alpha}^{\pm}(\lambda), \mathcal{M}_{\alpha, 4}^{\pm}(\lambda)\}$ with $\alpha \in \{0,1,2,3,4\}$. Then
	\[  
\sup _{x, y \in \mathbb{R}^2}\left| \mathcal{K}_{B}^\pm(t,x,y)\right| + \sup _{x, y \in \mathbb{R}^2}\left|\mathcal{N}_{B}^\pm (t,x,y)\right| \lesssim \langle t\rangle^{-1}.  
	\]  
\end{proposition}

\begin{proof}
	For $B^\pm(\lambda) \in \{\mathcal{M}_{4,\alpha}^{\pm}(\lambda), \mathcal{M}_{\alpha, 4}^{\pm}(\lambda)\}$ with $\alpha \in \{0,1,2,3\}$, the desired estimates are already covered by the arguments in Proposition \ref{prop_second_2_3}. Hence, it suffices to consider the  case $B^\pm(\lambda) = \mathcal{M}_{4,4}^{\pm}(\lambda)$.
	
	The key observation is that, among the components of
	$Q_4vR_0^\pm(\lambda^4)$, only the $Q_4^0$-component carries a
	logarithmic singularity. More precisely,
	$
	Q_4^0vR_0^\pm(\lambda^4)
	=
	Q_4^0\mathcal{J}_{4,\pm}(\lambda)
    	$ and 
	 $\| e^{\mp i\lambda|z|}\mathcal{J}_{4,\pm}(\lambda,\cdot,z)\|_{L^2}$ is bounded by $|\log\lambda|\langle\lambda z\rangle^{-1/2}$, whereas
	$Q_4^1vR_0^\pm(\lambda^4)$ has no logarithmic singularity.

    Under the condition that either 
	 $Q_4^0=0$ or $Q_4^0\neq0$ and $Q_{42}^1=0$,  the expansions
\eqref{Q4=Q41} and \eqref{2-M4-4} show that every occurrence of
$Q_4^0$ in $\mathcal{M}_{4,4}^\pm(\lambda)$ is accompanied by one
factor $(\log\lambda)^{-1}$. This factor exactly compensates for
the logarithmic singularity of the corresponding
$Q_4^0vR_0^\pm(\lambda^4)$ term, while the $Q_4^1$-components
require no such compensation.
    
    Using 
	$
	\|
	\partial_\lambda^\ell
	\Lambda^\pm(\lambda)
	\|_{\mathbb{B}(L^2)}
	\lesssim
	\lambda^{-\ell}
	$ for $
	\ell=0,1,2,$
	together with the projected resolvent bounds from Lemma
	\ref{lemma_projection}, Leibniz's rule, and H\"older's inequality, we obtain that every resulting amplitude factor
	$\mathcal{E}^\pm(\lambda,x,y)$ satisfies
	$$
	\left|
	\partial_\lambda^\ell
	\mathcal{E}^\pm(\lambda,x,y)
	\right|
	\lesssim
	\widetilde{\chi}_1(\lambda/2)
	\langle\lambda r(x,y)\rangle^{-1/2}
	\lambda^{-\ell},
	\qquad
	\ell=0,1,2.
	$$
	Therefore Lemma \ref{oscillatory}(ii), applied with
	$(\sigma,\nu)=(0,0)$, yields the uniform bound $\langle t\rangle^{-1}.$
\end{proof}












\begin{comment}

\subsection{The estimates for the subcase $Q_{4}^0 = Q_{41}^1=0$}\label{subsec:second_3}  
This condition $Q_4^0 = Q_{41}^1 = 0$ (i.e., $Q_4 = Q_{42}^1$) allows us to establish an unweighted $| t |^{-1}$ decay (i.e., \eqref{2-t}) and the following weighted estimates:  
\[ 
	\left \|\cos (t \sqrt{H}) P_{\mathrm{ac}}(H)\chi_1(H)\right\|_{L^1_\omega\to L^\infty_{-\omega}} \lesssim\frac{1}{|t|\log |t|},  \quad \bigg\|\frac{\sin (t \sqrt{H})}{t\sqrt{H}} P_{\mathrm{ac}}(H)\chi_1(H)\bigg\|_{L^1_\omega\to L^\infty_{-\omega}}\sim|t|^{-1}.  
\]  
We divide the analysis into two parts, starting with the unweighted $| t |^{-1}$ estimate.  

\vspace{2mm}  
\textbf{Part I. The $L^1\to L^\infty$ estimate.}   
Recall that Proposition \ref{prop_second_2_1} gives the weighted bound $\omega(x)\omega(y)\langle t \rangle^{-1}$. The spatial weight $\omega$  arose exclusively because the $L^2$ bounds \eqref{lemma_projection_2} for  $\mathcal{T}_{4}, \mathcal{T}_{4}^{0}, \mathcal{T}_{4}^{1},$ and $\mathcal{T}_{4}^{1,1}$ (associated with $Q_4, Q_4^0, Q_{4}^1,$ and $Q_{41}^1$) inherently carry the weight $\omega$.   

In the present subcase, since $Q_4 = Q_{42}^1$, Lemma~\ref{lemma_projection} (iii) simplifies the resolvent expansion to:  
\[
	\bigl(Q_{42}^1v R_0^\pm(\lambda^4)\bigr)(\cdot,z) =  (b_0 Q_{42}^1v G_2^0)( \cdot,z) + \lambda\mathcal{T}_{4,\pm}^{1,2}(\lambda, \cdot,z) + \lambda\mathcal{T}_{4}^{1,2}(\lambda, \cdot,z).  
\]   
Crucially, all terms in this expansion satisfy the  $L^2$ bounds free of $\omega(z)$. Specifically, 
\begin{equation*}  
\|\partial_\lambda^\ell\mathcal{T}_4^{1,2}(\lambda,\cdot,z)\|_{L^2} \lesssim \lambda^{-\ell}, \quad \ell=0,1,2,  
\end{equation*}  
while  $\big\|\partial_\lambda^\ell\bigl(e^{\mp i\lambda|z|}\mathcal{T}_{4,\pm}^{1,2}(\lambda,\cdot,z)\bigr)\big\|_{L^2} \lesssim \lambda^{-\ell}\langle\lambda z\rangle^{-1/2}$ and  $\bigl\|(b_0 Q_{42}^1 v G_2^0)(\cdot,z)\bigr\|_{L^2} \lesssim 1$. Because the entire expansion is free of $\omega(z)$, we obtain the following Proposition \ref{prop_second_3-un}.  
\begin{proposition}\label{prop_second_3-un}  
	Let $B^\pm(\lambda) \in \{\mathcal{M}_{4,\alpha}^{\pm}(\lambda), \mathcal{M}_{\alpha,4}^{\pm}(\lambda)\}$ with $\alpha \in \{0,1,2,3,4\}$. If $Q_4^0 = Q_{41}^1 = 0$, then  
	\[  
	\sup_{x,y\in\mathbb{R}^2} \bigl|\mathcal{K}_{B}^\pm(t,x,y)\bigr| + \sup_{x,y\in\mathbb{R}^2} \bigl|\mathcal{N}_{B}^\pm(t,x,y)\bigr| \lesssim \langle t \rangle^{-1}.  
	\]  
\end{proposition}  
 Combining Proposition \ref{prop_second_3-un} with the $L^1\to L^\infty$ dispersive bounds  established for the regular and first-kind resonance cases in Subsection \ref{subsec:regular_first_1}, we obtain the estimate \eqref{2-t}.
\end{comment}
\vspace{2mm}  
\textbf{Part II. Weighted $L^1_\omega \to L^\infty_{-\omega}$ estimate.} 
Building on the weighted estimates established for the regular and
first-kind resonance cases in
Subsection~\ref{subsec:regular_first_2}, we now consider the additional cases
$
B^\pm(\lambda)\in
\left\{
\mathcal M_{4,\alpha}^\pm(\lambda),
\mathcal M_{\alpha,4}^\pm(\lambda)
\right\}$ for
$\alpha\in\{0,1,2,3,4\}.$
The case
$
B^\pm(\lambda)=\mathcal M_{0,0}^\pm(\lambda)
$
must also be analyzed separately, since the leading-order expansion
of $\mathcal M_{0,0}^\pm(\lambda)$ in the second-kind resonance case may differ from that in the regular
and first-kind resonance cases.


 By arguments analogous to those in Proposition \ref{prop_K_weight_1}, we can obtain that for $|t|\ge 2,$
\begin{equation}\label{eq:fast_decay}  
\left| \mathcal{K}_{B}^\pm(t,x,y)\right| +
\left|\mathcal{N}_{B}^\pm(t,x,y)\right| \lesssim \frac{\omega(x)\omega(y)}{|t|\log |t|}, \  B^\pm(\lambda) \in \{\mathcal{M}_{4,\alpha}^{\pm}(\lambda), \mathcal{M}_{\alpha,4}^{\pm}(\lambda)\}\ \text{with}\ \alpha \in \{1,2,3\}.  
\end{equation}  
Consequently, the analysis is reduced to the case
$B^\pm(\lambda)= \mathcal{M}_{\alpha,\beta}^{\pm}(\lambda)$
with $(\alpha,\beta) \in \{0,4\} \times \{0,4\}$.
Within this setting, we distinguish two scenarios:
\textbf{Case 1:} $Q_4^0 = 0$;
\textbf{Case 2:} $Q_4^0 \neq 0$ and $Q_{42}^1 = 0$. 




\textbf{Case 1: \(Q_4^0=0\).}
In this case, \(Q_4=Q_4^1\neq0\). Tracing the argument in the proof of
Theorem~\ref{thm:M_inverse} in
Section~\ref{subsec:proof_inverse}, and combining
\eqref{D00}, \eqref{Upsilon44-2} with $
\mathcal M_{\alpha,\beta}^\pm(\lambda)=\mathcal D_{\alpha,\beta}^\pm(\lambda)+\Upsilon_{\alpha,\beta}^\pm(\lambda)
$ for $\alpha,\beta \in \{0,4\}$, we obtain 
the following refined expansions:
\[
	\mathcal{M}_{\alpha,\beta}^{\pm}(\lambda) = \mathcal{D}_{\alpha,\beta}^{\pm} + \lambda (-\log \lambda)^{3/2} \Lambda^\pm(\lambda), \qquad \alpha,\beta \in \{0,4\}, 
\] where
the operators $\mathcal{D}_{\alpha,\beta}^{\pm}$ are explicitly given by  
\begin{align}\label{Doo}
	\mathcal{D}_{0,0}^{\pm} &= (a_0^\pm)^{-1} \|V\|_{L^1}^{-1}Q_0 + (a_0^\pm)^{-2}\mathcal{W} (\boldsymbol{d}^\pm)^{-1} \mathcal{W}^* := (a_0^\pm)^{-1} \|V\|_{L^1}^{-1}Q_0 + \widetilde{\mathcal{D}}_{0,0}^{\pm}, \\
	\label{eq:expre_D_00} 
	\mathcal{D}_{4,4}^{\pm} &= (\boldsymbol{d}^\pm)^{-1}, \quad  
	\mathcal{D}_{0,4}^{\pm} = - (a_0^\pm)^{-1}  \mathcal{W} (\boldsymbol{d}^\pm)^{-1}, \quad  
	\mathcal{D}_{4,0}^{\pm} = - (a_0^\pm)^{-1} (\boldsymbol{d}^\pm)^{-1} \mathcal{W}^*,
\end{align}  
with $\mathcal{W} = \|V\|_{L^1}^{-1}Q_0 T_0 Q_4^1$ and  
$
(\boldsymbol{d}^\pm)^{-1} = [a_2^\pm Q_4^1 v G_4 v Q_4^1 - (a_0^\pm)^{-1} \|V\|_{L^1}^{-1} Q_4^1 T_0  Q_0 T_0 Q_4^1]^{-1}.  
$  In particular, in the present second-kind resonance case, the leading
term of \(\mathcal M_{0,0}^\pm(\lambda)\) is
$
(a_0^\pm)^{-1}\|V\|_{L^1}^{-1}Q_0
+
\widetilde{\mathcal D}_{0,0}^\pm,
$
rather than merely
$
(a_0^\pm)^{-1}\|V\|_{L^1}^{-1}Q_0,
$
as in the regular and first-kind resonance cases.

\begin{proposition}\label{prop_second_critical_block}  
	Let $B^\pm(\lambda) = \mathcal{M}_{\alpha,\beta}^{\pm}(\lambda)$ for $(\alpha,\beta) \in \{0,4\} \times \{0,4\}$. If $Q_4^0 = 0$ $(\text{i.e.,}\  Q_4=Q_4^1 )$, then  
	\[  
\bigl|\mathcal{K}_{B}^\pm(t,x,y)\bigr| \lesssim \frac{\omega(x)\omega(y)}{|t|\log|t|}, \quad  
\mathcal{N}_{B}^\pm(t,x,y) = \frac{1}{2 i |t|} C^\pm_{\alpha,\beta}(x,y) + {O}\!\left(\frac{\omega(x)\omega(y)}{|t|\log|t|}\right),  \ \ |t| \ge 2,
	\]  
    where the integral kernels $C^\pm_{\alpha,\beta}(x,y)$ 
	are uniformly bounded and explicitly given by:    
	\begin{equation*}  
		\begin{split}
		C^\pm_{0,0}(x,y) &= (a_0^\pm)^2\bigl\langle  \mathcal{D}_{0,0}^{\pm}(Q_0v G_0)(\cdot), (Q_0v G_0)(\cdot) \bigr\rangle,\\
        C^\pm_{0,4}(x,y) &= b_0a_0^\pm \bigl\langle  \mathcal{D}_{0,4}^{\pm} (Q_{4}^1 v G_2^0)(\cdot,y), (Q_0v G_0)(\cdot) \bigr\rangle, \\  
		C^\pm_{4,0}(x,y) &= b_0a_0^\pm \bigl\langle  \mathcal{D}_{4,0}^{\pm}(Q_0v G_0)(\cdot), (Q_{4}^1 v G_2^0)(\cdot,x) \bigr\rangle,\\
		C^\pm_{4,4}(x,y)&= b_0^2 \bigl\langle \mathcal{D}_{4,4}^{\pm} (Q_{4}^1 v G_2^0)(\cdot,y), (Q_{4}^1 v G_2^0)(\cdot,x) \bigr\rangle.  
		\end{split}
	\end{equation*}  
\end{proposition} 

\begin{proof}      
	For the case $(\alpha,\beta)=(0,0)$, the estimates are obtained by replacing 
    $(a_0^\pm)^{-1} \|V\|_{L^1(\mathbb{R}^2)}^{-1}Q_0$ with
    $\mathcal{D}_{0,0}^{\pm}$ in 
   Proposition~\ref{prop_K_weight_2}.   

	For the remaining cases involving the index $4$, we utilize the expansion \eqref{eq:expansion} for $R_0^\pm(\lambda^4)(x,y) $:    
	\[    
	R_0^\pm(\lambda^4)(x,y) - b_0G_2^0(x,y) = \frac{a_0^\pm}{\lambda^2}G_0(x,y) + g_0^\pm(\lambda)|x-y|^2 + \widetilde{E}_1^{\pm}(\lambda,x,y),   
	\]    
	where the remainder satisfies $\widetilde{E}_1^{\pm}(\lambda,x,y) = O_2(\lambda^{2}|x-y|^{4})$.     
	Crucially, the exact cancellation property $Q_{4}^1(x^\gamma v)=Q_{4}^1(|x|^2 v)= 0$ for $|\gamma| \le 1$ annihilates the first two terms, namely, applying $Q_{4}^1 v$ to the expansion yields an expression governed entirely by the $\lambda^2$ remainder $Q_{4}^1 v \widetilde{E}_1^\pm$: 
	\[    
	\bigl[Q_{4}^1 v \bigl(R_0^\pm(\lambda^4) - b_0 G_2^0\bigr)(\cdot, y)\bigr](x) = \bigl(Q_{4}^1 v \widetilde{E}_1^\pm\bigr)(\lambda,x,y),  
	\]    
    where $\|\partial_\lambda^\ell\bigl(Q_{4}^1 v \widetilde{E}_1^\pm\bigr)(\lambda,x,y)\|_{L_x^2}\lesssim \lambda^{2-\ell}\langle y\rangle^4$ for $\ell=0,1,2.$	Consequently, the dominant low-energy singularity arises from the product of three components: $b_0Q_{4}^1vG_2^0,$ $a_0^\pm\lambda^{-2}Q_0vG_0$ and the leading term $\mathcal{D}_{\alpha,\beta}^{\pm}$ in $\mathcal{M}_{\alpha,\beta}^{\pm}(\lambda).$  
	
	All other cross terms gain a relative energy factor of at least $\lambda^{1-}$. 
	By employing the same oscillatory integral strategy as in Proposition \ref{prop_K_weight_1}, these  remainder combinations lead to a rapid temporal decay of at least $|t|^{-5/4}\langle x\rangle^4\langle y\rangle^4$, which is readily bounded by $\omega(x)\omega(y)(|t|\log|t|)^{-1}$ via \eqref{eq:inequality_3}.    
	
	Extracting the leading singular constituents perfectly recovers the spatial kernels $C^\pm_{\alpha,\beta}(x,y)$, which are uniformly bounded since $\|(b_0 Q_{4}^1 v G_2^0)(\cdot,y)\|_{L^2} \lesssim 1$. For the sine evolution, this  kernel $C^\pm_{\alpha,\beta}(x,y)$ couples with the following integral (see \eqref{eq:integral_sin}):    
	\[
	\frac{2}{\pi i t}\int_0^\infty \sin(t\lambda^2)\lambda^{-1}\widetilde{\chi}_1(\lambda)\,d\lambda = \frac{1}{2 i |t|}  +{O}(|t|^{-2}).
	\]   
	Conversely, the corresponding integral for the cosine evolution  decays rapidly:    
	\[    
	\int_0^\infty \cos(t\lambda^2)\lambda\widetilde{\chi}_1(\lambda)\,d\lambda = O(|t|^{-N}), \quad \forall N>0.    
	\]    
	This directly establishes the desired  bound and asymptotic expansions, completing the proof.    
\end{proof}  

We are now ready to prove the sharpness of the sine estimate. 
\begin{proposition}\label{sine-sharp-2}
	If $Q_4^0= 0,$  then
\begin{equation}
\label{si-optimality-3}
	\Big\|\frac{\sin (t \sqrt{H})}{t\sqrt{H}} P_{\mathrm{ac}}(H)\chi_1(H)\Big\|_{L^1_\omega \to L^{\infty}_{-\omega}} \sim |t|^{-1},\ \ |t|\ge 2.
\end{equation}  
\end{proposition}
\begin{proof}
Following the framework established for the regular and first-kind resonance cases in Subsection \ref{subsec:regular_first_2}, and using the estimates from \eqref{eq:fast_decay} and in Proposition \ref{prop_second_critical_block}, we can obtain that 
\begin{align}\label {expan-si-2} 
\frac{\sin(t\sqrt{H})}{t\sqrt{H}} P_{\mathrm{ac}}(H)\chi_1(H)    
= -\frac{1}{2 i |t|}\widetilde{\mathcal{A}} + \mathcal{O}\!\left(\frac{1}{|t|\log |t|}\right),   
\end{align} 
in the sense of $\mathbb{B}(L^1_\omega, L^\infty_{-\omega})$, where the bounded  operator $\widetilde{\mathcal{A}}$ is defined as 
	\begin{equation}  \label{A-ex}  
    \begin{split}
		\widetilde{\mathcal{A}} = &G_0 vQ_0 \bigl((a_0^+)^2 \widetilde{\mathcal{D}}_{0,0}^{+}-(a_0^-)^2 \widetilde{\mathcal{D}}_{0,0}^{-} \bigr)Q_0v G_0  + b_0  G_0 v Q_0 \bigl(a_0^+\mathcal{D}_{0,4}^{+}-a_0^-\mathcal{D}_{0,4}^{-} \bigr) Q_{4}^1 v G_2^0 \\  
		& + b_0  G_2^0 v Q_{4}^1 \bigl(a_0^+\mathcal{D}_{4,0}^{+}-a_0^-\mathcal{D}_{4,0}^{-}\bigr) Q_{0} v G_0 + b_0^2 G_2^0 v Q_{4}^1 \bigl( \mathcal{D}_{4,4}^{+} - \mathcal{D}_{4,4}^{-} \bigr) Q_{4}^1 v G_2^0.
        \end{split}
	\end{equation}  
Here $\widetilde{\mathcal{D}}_{0,0}^{\pm}$ and $\mathcal{D}_{\alpha,\beta}^{\pm}$ are given in \eqref{Doo} and \eqref{eq:expre_D_00}. Note that we use $\widetilde{\mathcal{D}}_{0,0}^{\pm}$ rather than $\mathcal{D}_{0,0}^{\pm}$  in \eqref{A-ex}, because the $Q_0$-component $(a_0^\pm)^{-1} \|V\|_{L^1}^{-1} Q_0$ neutralizes the leading term ${1}/{(8|t|)}$ from $(\sin(t\Delta))/({t\Delta})\chi_1(\Delta)$, exactly as in the regular and first-kind resonance cases.   
By \eqref{expan-si-2}, to establish \eqref{si-optimality-3}, it suffices to prove   $\widetilde{\mathcal{A}} \neq 0$. 

	Substituting the expressions \eqref{Doo} and \eqref{eq:expre_D_00} of $\widetilde{\mathcal{D}}_{0,0}^{\pm}$ and $\mathcal{D}_{\alpha,\beta}^{\pm}$ into \eqref{A-ex}, we deduce that 
	\[  
	\widetilde{\mathcal{A}} = \bigl(G_0vQ_0\mathcal{W}-b_0G_2^0vQ_{4}^1\bigr)\bigl((\boldsymbol{d}^+)^{-1} - (\boldsymbol{d}^-)^{-1}\bigr)\bigl(G_0vQ_0\mathcal{W}-b_0G_2^0vQ_{4}^1\bigr)^*,
	\]  
	where $\mathcal{W} = \|V\|_{L^1}^{-1}Q_0 T_0 Q_4^1$ and  
$
(\boldsymbol{d}^\pm)^{-1} = [a_2^\pm Q_4^1 v G_4 v Q_4^1 - (a_0^\pm)^{-1} \|V\|_{L^1}^{-1} Q_4^1 T_0  Q_0 T_0 Q_4^1]^{-1}.  
$ 

	Note that 
	$
	i \bigl((\boldsymbol{d}^+)^{-1} - (\boldsymbol{d}^-)^{-1}\bigr)  
	$ is strictly positive definite on $Q_{4}^1 L^2$. Consequently, in order to show $\widetilde{\mathcal{A}} \neq 0$, it reduces to proving $G_0vQ_0\mathcal{W}-b_0G_2^0vQ_{4}^1$ is injective  on the finite-dimensional space $Q_{4}^1 L^2$.  
	
	Assume $\bigl(G_0vQ_0\mathcal{W}-b_0G_2^0vQ_{4}^1\bigr)\psi = 0$ for some $\psi \in Q_{4}^1 L^2$. Then
	\[
		b_0 \bigl(G_2^0 v Q_{4}^1 \psi\bigr)(x) = \langle \mathcal{W} \psi, v \rangle.  
	\]  
Applying $\Delta^2$ in the distributional sense and
	using
	$
	\Delta^2G_2^0
	=
	8\pi\delta_0,
	$
	we obtain
	$
	vQ_4^1\psi=v\psi=0.
	$ Then, by the same argument used to derive \eqref{psi1} and \eqref{psi2}, we get \(\psi=0\).
Thus, the operator $G_0vQ_0\mathcal{W}-b_0G_2^0vQ_{4}^1$ is injective on $Q_{4}^1 L^2$, confirming $\widetilde{\mathcal{A}} \neq 0$.  
\end{proof}  


\textbf{Case 2:}  $Q_4^0\neq0$ and $Q_{42}^1=0.$ In this case,
$
Q_4
=
Q_4^0\oplus Q_4^1,\ 
Q_4^1=Q_{41}^1.
$
By tracing the contributions of $\Lambda^\pm(\lambda)$ in \eqref{D40D04D44_1}, and using \eqref{Upsilon44-2} together with  $
\mathcal M_{\alpha,\beta}^\pm(\lambda)=\mathcal D_{\alpha,\beta}^\pm(\lambda)+\Upsilon_{\alpha,\beta}^\pm(\lambda)
$ for $\alpha,\beta \in \{0,4\},$  we obtain
\begin{align*}
	\mathcal{M}_{0,0}^\pm(\lambda)=&(\log \lambda)^{-2}  Q_0 \Lambda^\pm(\lambda) Q_0,\\
	\mathcal{M}^\pm_{0,4}(\lambda)=&(\log\lambda)^{-1}Q_0\mathfrak{D}Q_4^0+(\log\lambda)^{-1}Q_0\Lambda^\pm(\lambda)Q_4^1+(\log\lambda)^{-2}Q_0\Lambda^\pm(\lambda)Q_4,\\
	\mathcal{M}^\pm_{4,0}(\lambda)=&(\log\lambda)^{-1}Q_4^0\mathfrak{D}^*Q_0+(\log\lambda)^{-1}Q_4^1\Lambda^\pm(\lambda)Q_0+(\log\lambda)^{-2}Q_4\Lambda^\pm(\lambda)Q_0,\\
	\mathcal{M}^\pm_{4,4}(\lambda)=&(\log\lambda)^{-2}Q_4^0\mathfrak{C}_{00}^\pm Q_4^0+(\log\lambda)^{-1}\bigl(Q_4^0\mathfrak{C}_{01}^\pm Q_4^1+Q_4^1\mathfrak{C}_{10}^\pm Q_4^0\bigr)+Q_4^1\mathfrak{C}_{11}^\pm Q_4^1\\&+(\log\lambda)^{-3}Q_4^0\Lambda^\pm(\lambda)Q_4^0+(\log\lambda)^{-2}\bigl(Q_4^1\Lambda^\pm(\lambda)Q_4^0+Q_4^0\Lambda^\pm(\lambda)Q_4^1\bigr)+(\log\lambda)^{-1}Q_4^1\Lambda^\pm(\lambda)Q_4^1,
\end{align*}
where $\mathfrak{D}=(b_0Q_4^0vG_2vQ_0)^{-1},$ $\mathfrak{C}^\pm_{11}=(a_2^{\pm})^{-1}D_0$ with $D_0=\bigl(Q_{41}^1vG_4vQ_{41}^1\bigr)^{-1}$ and 
\begin{align*}
	\mathfrak{C}^\pm_{00}&=\mathfrak{D}^*	\Bigl((a_2^{\pm})^{-1}
	Q_0T_0Q_4^1D_0Q_4^1T_0Q_0
	-a_2^{\pm}Q_0vG_0vQ_0
	\Bigr)\mathfrak{D},\\
	\mathfrak{C}^\pm_{01}&=-(a_2^{\pm})^{-1}\mathfrak{D}^*Q_0T_0Q_4^1D_0,\quad 	\mathfrak{C}^\pm_{10}=-(a_2^{\pm})^{-1}D_0Q_4^1T_0Q_0\mathfrak{D}.
\end{align*}
Here,  if $Q_{41}^1L^2=\{0\}$, the inverse $(Q_{41}^1 v G_4 v Q_{41}^1)^{-1}$ is interpreted as the zero operator.

 Note that
$(Q_0vR_0^\pm(\lambda^4))(x,y)=a_0^{\pm}\lambda^{-2}(Q_0vG_0)(x,y)+(Q_0vE_1^\pm)(\lambda,x,y)$ (see \eqref{eq:expansion})  and
\begin{align*}
	\bigl(Q_{4}^0 v R_0^\pm(\lambda^4)\bigr)(x,y) &=b_0(\log\lambda) (Q_4^0vG_2)(x,y)+a_1^\pm (Q_4^0vG_2)(x,y)+(b_0 Q_4^0vG_2^0)(x, y)+\bigl(Q_{4}^0 v \widetilde{E}_1^\pm)(\lambda,x,y),\\
		\bigl(Q_{4}^1 v R_0^\pm(\lambda^4)\bigr)(x,y) &= (b_0 Q_4^1vG_2^0)(x, y) +\bigl(Q_{4}^1 v \widetilde{E}_1^\pm)(\lambda,x,y),
\end{align*}
where $\bigl\|(Q_0 v G_0)(\cdot,y)\bigr\|_{L^2}+\bigl\|(Q_4^0v G_2)(\cdot,y)\bigr\|_{L^2}+\bigl\|(Q_4^1v G_2^0)(\cdot,y)\bigr\|_{L^2}\lesssim1,$ $\bigl\|(b_0Q_4^0 v G_2^0)(\cdot,y)\bigr\|_{L^2}\lesssim\omega(y),$ and for $\ell=0,1,2,$
\begin{align*}	\bigl\|\partial_\lambda^\ell(Q_0v E_1^\pm)(\lambda,\cdot,y)\bigr\|_{L^2}\lesssim\lambda^{-\ell-}\langle y\rangle^4,\quad 	\bigl\|\partial_\lambda^\ell(Qv \widetilde{E}_1^\pm)(\lambda,\cdot,y)\bigr\|_{L^2}\lesssim\lambda^{2-\ell}\langle y\rangle^4\ \text{for}\ Q\in\{Q_4^0, Q_4^1\}.
\end{align*}
Combining the above estimates with the integral bound
\begin{align*}
	\ \ \ \ \ \left|
	\int_0^\infty e^{-it\lambda^2} \lambda \mathcal{E}(\lambda,x,y)\,d\lambda
	\right|
	\lesssim
	\begin{cases}	\dfrac{\omega(x)\omega(y)}{|t|\log|t|}, 
		& \text{if } 
		\begin{minipage}[t]{0.6\linewidth}
			$|\partial_\lambda^\ell \mathcal{E}(\lambda,x,y)|\lesssim\widetilde{\chi}_1(\lambda/2)|\log\lambda|^{-1} \lambda^{-\ell}\omega(x)\omega(y),$
		\end{minipage}\\[2ex]
		\dfrac{\langle x\rangle^4\langle y\rangle^4}{|t|^{5/4}}, 
		& \text{if } 
		\begin{minipage}[t]{0.6\linewidth}
			$|\partial_\lambda^\ell \mathcal{E}(\lambda,x,y)|\lesssim\widetilde{\chi}_1(\lambda/2)\lambda^{1/2-\ell} \langle x\rangle^4\langle y\rangle^4,$
		\end{minipage}
	\end{cases}
\end{align*}
for $\ell=0,1,2$, and proceeding analogously to Proposition \ref{prop_second_critical_block}, we obtain
$$
|\mathcal{K}_{\mathcal{M}_{0,0}}^\pm(t,x,y)|+|\mathcal{N}_{\mathcal{M}_{0,0}}^\pm(t,x,y)|\lesssim\omega(x)\omega(y)(|t|\log|t|)^{-1},
$$
and for $B^\pm(\lambda)\in\{\mathcal{M}^\pm_{0,4}(\lambda), \mathcal{M}^\pm_{4,0}(\lambda), \mathcal{M}^\pm_{4,4}(\lambda)\}$,
\[
\bigl|\mathcal{K}_{B}^\pm(t,x,y)\bigr| \lesssim \frac{\omega(x)\omega(y)}{|t|\log|t|}, \quad
\mathcal{N}_{B}^\pm(t,x,y) = \frac{1}{2 i |t|} \mathfrak{T}^\pm_{\alpha,\beta}(x,y) + O\!\left(\frac{\omega(x)\omega(y)}{|t|\log|t|}\right), \quad |t| \ge 2,
\]
where the integral kernels $\mathfrak{T}^\pm_{\alpha,\beta}(x,y)$ are uniformly bounded and explicitly given by:   
\begin{equation}  \label{sharp-operator}
	\begin{split}
	\mathfrak{T}^\pm_{0,4}(x,y) =& b_0a_0^\pm \bigl\langle  \mathfrak{D} (Q_{4}^0 v G_2)(\cdot,y), (Q_0v G_0)(\cdot) \bigr\rangle, \ \  
	\mathfrak{T}^\pm_{4,0}(x,y)= b_0a_0^\pm \bigl\langle  \mathfrak{D}^*(Q_0v G_0)(\cdot), (Q_{4}^0v G_2)(\cdot,x) \bigr\rangle,\\
		\mathfrak{T}^\pm_{4,4}(x,y)= &b_0^2 \Bigl(\bigl\langle \mathfrak{C}_{00}^\pm (Q_{4}^0 v G_2)(\cdot,y), (Q_{4}^0 v G_2)(\cdot,x) \bigr\rangle+\bigl\langle \mathfrak{C}_{01}^\pm (Q_{4}^1 v G_2^0)(\cdot,y), (Q_{4}^0 v G_2)(\cdot,x) \bigr\rangle\\
		&\ \ \ +\bigl\langle \mathfrak{C}_{10}^\pm (Q_{4}^0 v G_2)(\cdot,y), (Q_{4}^1 v G^0_2)(\cdot,x) \bigr\rangle+\bigl\langle \mathfrak{C}_{11}^\pm (Q_{4}^1 v G_2^0)(\cdot,y), (Q_{4}^1 v G^0_2)(\cdot,x) \bigr\rangle\Bigr).  
	\end{split}
\end{equation}  
Hence, 
it remains to prove the following sharp bound for the sine evolution.
\begin{proposition}\label{2-sine-sharp-2}
	If $Q_4^0\neq 0$  and $Q_{42}^1=0,$ then
	\begin{equation*}
		\Big\|\frac{\sin (t \sqrt{H})}{t\sqrt{H}} P_{\mathrm{ac}}(H)\chi_1(H)\Big\|_{L^1_\omega \to L^{\infty}_{-\omega}} \sim |t|^{-1},\ \ |t|\ge 2.
	\end{equation*}  
\end{proposition}
\begin{proof}
	By an argument similar to that in the proof of Proposition \ref{sine-sharp-2},  combining  \eqref{sharp-operator}  with  $$\frac{\sin(t\Delta)}{t\Delta}\chi_1(\Delta)
	=\frac{1}{8 |t|}G_0+\mathcal{O}\bigl(({|t|\log|t|})^{-1}\bigr),$$
we reduce the problem to proving that the operator $\mathfrak{A}\neq0,$ where
	\begin{align*}
		\mathfrak{A}=\frac{1}{8}G_0-\frac{1}{2i}\Bigl[ & b_0(a_0^+-a_0^-)\bigl(G_0vQ_0\mathfrak{D}Q_4^0vG_2+G_2vQ_4^0\mathfrak{D}^*Q_0vG_0\bigr)\\
		&+b_0^2\Bigl(G_2vQ_4^0\bigl(\mathfrak{C}_{00}^+-\mathfrak{C}_{00}^-\bigr)Q_4^0vG_2+G_2vQ_4^0\bigl(\mathfrak{C}_{01}^+-\mathfrak{C}_{01}^-\bigr)Q_4^1vG^0_2\\
		&+G^0_2vQ_4^1\bigl(\mathfrak{C}_{10}^+-\mathfrak{C}_{10}^-\bigr)Q_4^0vG_2+G^0_2vQ_4^1\bigl(\mathfrak{C}_{11}^+-\mathfrak{C}_{11}^-\bigr)Q_4^1vG^0_2\Bigr)\Bigr].
	\end{align*}
	By the cancellation of $Q_4^0$ and the  definition of 
	$\mathfrak{D}$, we obtain
	\begin{equation*}
		%\label{eq:D-inverse-identities}
		b_0G_0vQ_0\mathfrak{D}Q_4^0vG_2
		=
		G_0,
		\qquad
		b_0G_2vQ_4^0\mathfrak{D}^*Q_0vG_0
		=
		G_0.
	\end{equation*}
Hence, by calculating, we derive that 
		\begin{equation*}
			%\label{eq:A-factorization-final}
		\mathfrak{A}
		=
		-\frac{63}{512}G_0
		+
		\frac{8}{\pi^2}
		\mathcal{X}D_0\mathcal{X}^*,
	\end{equation*}
	where $	\mathcal{X}
	=
	G_2vQ_4^0\mathfrak{D}^*Q_0T_0Q_4^1
	-
	G_2^0vQ_4^1.$ If $Q_4^1L^2=\{0\},$ then $	\mathfrak{A}=	-\frac{63}{512}G_0\neq0.$ Hence, it reduces to showing $\mathfrak{A}\neq0$ when $Q_4^1L^2=Q_{41}^1L^2\neq\{0\}.$
	
	
	
We first  prove that
	$\mathcal{X}$ is
	injective on $Q_4^1L^2\neq\{0\}.$
	Suppose that $\mathcal{X}\psi=0$ for some
	$\psi\in Q_4^1L^2=Q_{41}^1L^2$. Since
	$
	G_2vQ_4^0\mathfrak{D}^*Q_0T_0Q_4^1\psi
	$
	is constant, 
	$
	G_2^0vQ_4^1\psi
	$
	is constant. Applying $\Delta^2$ in the distributional sense and
	using
	$
	\Delta^2G_2^0
	=
	8\pi\delta_0,
	$
	we obtain
	$
	vQ_4^1\psi=v\psi=0.
	$  Proceeding as in the derivation of \eqref{psi1} and \eqref{psi2}, we conclude that \(\psi=0\).
	 Hence, $\mathcal{X}$ is
	injective on $Q_4^1L^2$.
	


	Suppose, for contradiction, that
	$
	\mathfrak{A}=0.
	$
	Then $
		({8}/{\pi^2})
		\mathcal{X}D_0\mathcal{X}^*=({63}/{512})G_0.$
	Applying $\Delta^2$ to both sides, and observing that
	$
	\Delta^2G_0=0,
$ and $
	\Delta^2\mathcal{X}
	=
	-8\pi vQ_4^1,
	$
	we obtain
	$
	vQ_4^1D_0\mathcal{X}^*=0.
	$
	For any test function $f$, set
	$
	\psi
	=
	D_0\mathcal{X}^*f
	\in Q_4^1L^2.
	$
	Then
	$
	vQ_{4}^1\psi=0,
	$
	and hence
	$$
	D_0^{-1}\psi
	=
	Q_{41}^1vG_4vQ_{41}^1\psi=	Q_{4}^1vG_4vQ_{4}^1\psi
	=
	0.
	$$
	Since $D_0^{-1}$ is invertible on $Q_4^1L^2$, we conclude that
	$
	\psi=0.
	$
	Thus
	$
	D_0\mathcal{X}^*=0.
	$
	The invertibility of $D_0$ on $Q_{41}^1L^2=Q_{4}^1L^2\neq\{0\}$ implies
	$
	\mathcal{X}^*=0,
	$
	and hence
	$
	\mathcal{X}=0.
	$
	This contradicts the injectivity of $\mathcal{X}$ on the nonzero
	space $Q_4^1L^2$. Therefore
	$
	\mathfrak{A}\neq0.
	$
\end{proof}
\section{The third kind resonance and eigenvalue cases}\label{sec:third_fourth}
This section focuses on the low-energy part for the  third kind resonance and eigenvalue cases, that is, the following Theorems \ref{main_theorem_low_31} and \ref{main_theorem_low_32}.

\begin{theorem}\label{main_theorem_low_31}
Let $H = \Delta^2 + V$ with $|V(x)| \lesssim \langle x \rangle^{-18-}$.
Assume that $H$ has no positive embedded eigenvalues and that zero is either a third-kind resonance or an eigenvalue accompanied by a d-wave resonance $(\text{i.e.,}\  Q_5 \neq 0 )$. Then
    \begin{equation}\label{optimality_31}
        \left\|\cos (t \sqrt{H}) P_{\mathrm{ac}}(H)\chi_1(H)\right\|_{L^1\to L^{\infty}} +
        \bigg\|\frac{\sin (t \sqrt{H})}{t\sqrt{H}} P_{\mathrm{ac}}(H)\chi_1(H)\bigg\|_{L^1 \to L^{\infty}} \sim \frac{1}{\log |t|}, \quad |t| \gg 1.
    \end{equation}
\end{theorem}

\begin{theorem}\label{main_theorem_low_32}
	Let $H = \Delta^2 + V$ with $|V(x)| \lesssim \langle x \rangle^{-22-}$.  
	Assume that $H$ has no positive embedded eigenvalues and that zero is an eigenvalue of $H$ but exhibits no d-wave resonance $(\text{i.e.,}\  Q_5 = 0 )$.
	\begin{enumerate}  
		\item If $H$ admits a p-wave resonance $(\text{i.e.,}\  Q_4 \neq 0 )$, then 
		$$  
		\begin{aligned}  
		\left\|\cos (t \sqrt{H}) P_{\mathrm{ac}}(H)\chi_1(H)\right\|_{L^1 \to L^{\infty}} + 
		\bigg\|\frac{\sin (t \sqrt{H})}{t\sqrt{H}} P_{\mathrm{ac}}(H)\chi_1(H)\bigg\|_{L^1 \to L^{\infty}} &\lesssim\langle t \rangle^{-1} {(\log (2+|t|))^2}.  
		\end{aligned}  
		$$  
		\item If $H$ lacks a p-wave resonance $(\text{i.e.,}\  Q_4 = 0 )$, then  
		\begin{equation*}  
		\left\|\cos(t\sqrt{H}) P_{\mathrm{ac}}(H)\chi_1(H)\right\|_{L^1 \to L^\infty} +  \bigg\|\frac{\sin(t\sqrt{H})}{t\sqrt{H}} P_{\mathrm{ac}}(H)\chi_1(H)\bigg\|_{L^1 \to L^\infty} \lesssim \langle t\rangle^{-1}.  
		\end{equation*}  
		Furthermore, we distinguish two subcases for the weighted estimate when $|t| \ge 2$:  
		\begin{itemize}  
			\item If ${S}_5L^2\neq\widetilde{S}_5L^2 $, then   
			\[   \left \|\cos (t \sqrt{H}) P_{\mathrm{ac}}(H)\chi_1(H)\right\|_{L^1_\omega\to L^\infty_{-\omega}} \lesssim\frac{1}{|t|\log |t|},\  \ \bigg\|\frac{\sin (t \sqrt{H})}{t\sqrt{H}} P_{\mathrm{ac}}(H)\chi_1(H)\bigg\|_{L^1_\omega\to L^\infty_{-\omega}}\sim|t|^{-1}.  
			\]  
			\item If ${S}_5L^2=\widetilde{S}_5L^2 $, then
			\[ 
			\left\|\cos (t \sqrt{H}) P_{\mathrm{ac}}(H)\chi_1(H)\right\|_{L^1_\omega \to L^{\infty}_{-\omega}} +  
			\bigg\|\frac{\sin (t \sqrt{H})}{t\sqrt{H}} P_{\mathrm{ac}}(H)\chi_1(H)\bigg\|_{L^1_\omega \to L^{\infty}_{-\omega}} \lesssim \frac{1}{|t| \log|t|}.  
			\] 
            \end{itemize}
		Here, the subspace characterizing the isotropic property is defined as:  
		\begin{equation} \label{S5wide}   
		\widetilde{S}_5L^2 = \left\{ \psi \in S_5 L^2 \;\middle|\;     
		\begin{aligned}    
		&\langle x_1^3 x_2 v, \psi \rangle = 0, \quad \langle x_1 x_2^3 v, \psi \rangle = 0, \\  
		&\langle x_1^4 v, \psi \rangle = \langle x_2^4 v, \psi \rangle = 3 \langle x_1^2 x_2^2 v, \psi \rangle    
		\end{aligned}    
		\right\}.    
		\end{equation}  
	\end{enumerate}  
\end{theorem}  

\subsection{The estimates for the case $Q_{5} \neq 0$}     
 This subsection is devoted to the proof of Theorem \ref{main_theorem_low_31}. Assume that zero is either a third-kind resonance or an eigenvalue with a $d$-wave resonance. Following the strategy from the previous sections, we only need to estimate the integral kernels $\mathcal{K}_{B}^+-\mathcal{K}_{B}^-$ and $\mathcal{N}_{B}^+-\mathcal{N}_{B}^-$ defined in \eqref{K_1} and \eqref{Lambda_1}, restricted to the operators $B^\pm(\lambda)=\mathcal{M}_{\alpha, \beta}^{\pm}(\lambda)$ where either $\alpha \in \{5,6\}$ or $\beta \in \{5,6\}$.  

Recall from Theorem \ref{thm:M_inverse} (III) that
\[
	 \begin{aligned}
&\mathcal{M}_{\alpha,\beta}^\pm(\lambda)= \mathcal{A}_{\alpha,\beta}(\lambda)+(\log \lambda)^{-2}\Gamma_{\alpha,\beta}^{\pm}(\lambda), \ (\alpha,\beta)=(5,5),(5,6),(6,5), (6,6),
	\end{aligned}
    \]
with
$$
\left\|\partial_\lambda^{\ell} \mathcal{A}_{\alpha,\beta}(\lambda)\right\|_{\mathbb{B}\left(L^2\right)}
+\big\|\partial_\lambda^{\ell} \Gamma_{\alpha, \beta}^{ \pm}(\lambda) \big\|_{\mathbb{B}\left(L^2\right)} \lesssim \lambda^{-\ell},
\ \ell=0,1,2,\ \text{and}\ \alpha,\beta=5,6.
$$
Moreover, for $ \alpha=5,6, \ 0 \le \beta \le 4$ and $\ell=0,1,2,$
\begin{align*}
\big\|\partial_\lambda^{\ell} \mathcal{M}_{\alpha,\beta}^\pm(\lambda)\big\|_{\mathbb{B}\left(L^2\right)}
+\big\|\partial_\lambda^{\ell} \mathcal{M}_{\beta,\alpha}^\pm(\lambda)\big\|_{\mathbb{B}\left(L^2\right)}
\lesssim \lambda^{-\ell}.
\end{align*}

Therefore it suffices to analyze $\mathcal{K}_{B}^+-\mathcal{K}_{B}^-$ and $\mathcal{N}_{B}^+-\mathcal{N}_{B}^-$ for the following cases:
\begin{itemize}
	\item $B^\pm(\lambda)=\mathcal{M}_{\alpha,\beta}^\pm(\lambda) \ \text{with} \ (\alpha,\beta)=(5,5),(5,6),(6,5),(6,6);$
	\item $ B^\pm(\lambda)= \mathcal{M}_{\alpha, \beta}^\pm(\lambda) \ \text{with} \  (\alpha,\beta)  \in \{ \mathbb{Z}^2 \mid \alpha=5,6,\ 0 \le \beta \le 4 \ \text{or} \ \beta=5,6,\ 0 \le \alpha \le 4\}.$
\end{itemize}

   
\begin{proposition}\label{third_merged_1} 
	Let $B^\pm(\lambda)=\mathcal{M}_{\alpha, \beta}^{\pm}(\lambda)$ with $\alpha,\beta \in \{5,6\}$. Then      
	\[      
	\sup _{x, y \in \mathbb{R}^2}\left|\left(\mathcal{K}_{B}^+-\mathcal{K}_{B}^-\right)(t,x,y)\right| +      
	\sup _{x, y \in \mathbb{R}^2}\left|\left(\mathcal{N}_{B}^+-\mathcal{N}_{B}^-\right)(t,x,y)\right|  \lesssim \frac{1}{\log(2+|t|)}.      
	\]      
\end{proposition}      
\begin{proof}   
	For $\alpha, \beta \in \{5,6\}$, the operators $B^\pm(\lambda)$ are decomposed into a $\pm$-independent principal term $\mathcal{A}_{\alpha,\beta}(\lambda)$ and a $\pm$-dependent remainder: $(\log \lambda)^{-2}\Gamma_{\alpha,\beta}^{ \pm}(\lambda)$. We analyze the integral kernels generated by these two parts separately.  
	Recall from Lemma \ref{lemma_projection} (iii) that  
	\begin{align}\label{ex-Q56}  
		\bigl(Q_j v R_0^\pm(\lambda^4)\bigr)(\cdot,z) &= (b_0 Q_j v G_2^0)( \cdot,z) +\lambda^{m_j}\bigl(\mathcal{T}_{j}(\lambda, \cdot,z)+\mathcal{T}_{j,\pm}(\lambda, \cdot,z)\bigr), \quad j=5,6,
	\end{align}  
	where $m_5=1$ and $m_6=2$. 	By \eqref{lemma_projection_G}--\eqref{lemma_projection_1},  for $j=5,6,$  one has $\|Q_jvG_2^0(\cdot, z)\|_{L^2}\lesssim 1$ and
	\begin{equation*}
	\begin{split}       
	\big\| \partial_\lambda^\ell \bigl( e^{\mp i\lambda|z|}\mathcal{T}_{j,\pm}(\lambda, \cdot, z) \bigr) \big\|_{L^2}& \lesssim \lambda^{-\ell} \langle \lambda z \rangle^{-1/2}, \ \ 
	\big\| \partial_\lambda^\ell \mathcal{T}_j(\lambda, \cdot, z) \big\|_{L^2} \lesssim \lambda^{-\ell}, \ \   \ell=0,1,2.
	\end{split}    
	\end{equation*}        
	By \eqref{ex-Q56}, the kernels $\mathcal{K}_B^\pm$ and $\mathcal{N}_B^\pm$ defined in \eqref{K_1} and \eqref{Lambda_1} decompose into finite sums of canonical oscillatory integrals of the form:
	\begin{align}\label{3-4res-integral}
	\frac{2}{\pi i}\int_0^\infty \lambda \cos(t\lambda^2) e^{\pm i\lambda r(x,y)} \mathcal{E}_{fg}^\pm(\lambda, x, y) \, d\lambda \quad \text{and} \quad  
	\frac{2}{\pi i t} \int_0^\infty \lambda^{-1} \sin(t\lambda^2) e^{\pm i\lambda r(x,y)} \mathcal{E}_{fg}^\pm(\lambda, x, y) \, d\lambda,  
	\end{align}
	where the common amplitude factor $	\mathcal{E}_{fg}^\pm(\lambda, x, y) $ is defined by  
\begin{equation}\label{eq:amplitude_def}  
		\mathcal{E}_{fg}^\pm(\lambda, x, y) = e^{\mp i\lambda r(x,y)} \lambda^\gamma \widetilde{\chi}_1(\lambda) \big\langle B^\pm(\lambda) f(\lambda, \cdot, y), g(\lambda, \cdot, x) \big\rangle.    
	\end{equation}    
	Here, $f, g \in\{b_0 Q_j v G_2^0, \mathcal{T}_{j,\pm}, \mathcal{T}_{j}\}$ for $j\in\{5,6\}$. Moreover, $r(x,y)$ isolates the inherent spatial oscillations of $f$ and $g$: it contributes $|y|$ (resp. $|x|$) if the state is $\mathcal{T}_{j,\pm}(\lambda, \cdot, y)$ (resp. $\mathcal{T}_{j,\mp}(\lambda, \cdot, x)$), and contributes $0$ otherwise.   The exponent $\gamma$ is given by $\gamma = 4 - k_\alpha - k_\beta + p=-2+p$ for $\alpha,\beta \in \{5,6\}$, where $p \ge 0$ is the total $\lambda$-power contributed by $f$ and $g$. For example if $(f,g)=(b_0Q_\beta v G_2^0, \mathcal{T}_{\alpha,\mp}),$  then $r=|x|$, $p=m_\alpha.$
	
	\vspace{1.5mm}      
\textbf{Case 1: The $\pm$-independent principal terms $\mathcal{A}_{\alpha,\beta}(\lambda)$.} \\  
	For these terms, $B^+(\lambda)=B^-(\lambda) := B(\lambda)$. The kernel difference $\mathcal{K}_B^+-\mathcal{K}_B^-$ can be  
    rewritten as a sum of terms involving
    the resolvent difference:  
	\begin{align*} 
	&\bigl[R_0^+ (\lambda^4)v Q_\alpha B(\lambda) Q_\beta v R_0^+(\lambda^4)\bigr](x,y) - \bigl[R_0^- (\lambda^4)v Q_\alpha B(\lambda) Q_\beta v R_0^-(\lambda^4)\bigr](x,y)  \\  
	&= \bigl[(R_0^+(\lambda^4)-R_0^-(\lambda^4))v  Q_\alpha B(\lambda) Q_\beta v R_0^+(\lambda^4) \bigr](x,y)+\big[ R_0^-(\lambda^4) v  Q_\alpha B(\lambda) Q_\beta v (R_0^+(\lambda^4)-R_0^-(\lambda^4))\bigr](x,y).  
	\end{align*}  
	By \eqref{ex-Q56}, we obtain that
		\begin{align}\label{difference5,6}
			\bigl[Q_\alpha v\bigl(R_0^+(\lambda^4)-R_0^-(\lambda^4)\bigr)\bigr](\cdot, z)&=\lambda^{m_\alpha}\bigl(\mathcal{T}_{\alpha,+}(\lambda, \cdot,z)-\mathcal{T}_{\alpha,-}(\lambda,\cdot,z)\bigr).
	\end{align}  
 Consequently, in the amplitude factor \eqref{eq:amplitude_def}, at least one state in the pairing $(f,g)$ must be a $\mathcal{T}_{j,\pm}$ remainder with $j\in\{5,6\}$, which guarantees a  factor $\lambda^{m_j}$ ($m_j \ge 1$).   
	Thus, we deduce that $p \ge 1$, yielding $\gamma \ge -1$.  
	
	The most singular amplitude factor \eqref{eq:amplitude_def} occurs  when $p=1$. Since $m_5=1$ and $m_6=2$, this requires $\mathcal{T}_{5,+}-\mathcal{T}_{5,-}$ to be paired with $b_0 Q_j v G_2^0$  ($j\in\{5,6\}$)  (e.g., $f = \mathcal{T}_{5,\pm}$, $g = b_0 Q_j v G_2^0$, yielding $r(x,y) = |y|$ and $B(\lambda)=\mathcal{A}_{j,5}(\lambda)$). Such amplitude factor satisfies:  
	\begin{align*}   
	\bigl|\partial_\lambda^{\ell} \mathcal{E}_{fg}^\pm(\lambda,x,y) \bigr| \lesssim \langle \lambda r(x,y)\rangle^{-1/2} \lambda^{-1-\ell} \widetilde{\chi}_1 (\lambda/2), \quad \ell=0,1,2.      
	\end{align*}    
	Applying Lemma \ref{oscillatory} (ii) with parameters $(\sigma,\nu)=(1,0)$  yields the bound $\langle t\rangle^{-1/2}$.   
	
	In particular, when we take $B(\lambda)=\mathcal{A}_{6,6}(\lambda)$, the expansion \eqref{difference5,6} guarantees a factor $\lambda^2$ (note that $m_6=2$), yielding $\gamma \ge 0.$ Hence, we can apply Lemma \ref{oscillatory}(ii) with  $(\sigma,\nu)=(0,0)$  to obtain that 
	\begin{equation}\label{mathcalA66}
	\sup _{x, y \in \mathbb{R}^2}\left|\left(\mathcal{K}_{B}^+-\mathcal{K}_{B}^-\right)(t,x,y)\right| +      
	\sup _{x, y \in \mathbb{R}^2}\left|\left(\mathcal{N}_{B}^+-\mathcal{N}_{B}^-\right)(t,x,y)\right|  \lesssim \langle t\rangle^{-1}, \ \text{with} \ B^\pm(\lambda)=\mathcal{A}_{6,6}(\lambda).
	\end{equation}
	
	\vspace{1.5mm}      
	\textbf{Case 2: The $\pm$-dependent remainder terms $(\log\lambda)^{-2}\Gamma_{\alpha,\beta}^\pm(\lambda)$.} \\  
	Since these operators depend on $\pm$, the algebraic identity \eqref{difference5,6} is inapplicable.  Thus we evaluate the kernels directly using the expansion \eqref{ex-Q56}. The dominant contribution to the decay rate arises from $f = b_0 Q_\beta v G_2^0$ and $g = b_0 Q_\alpha v G_2^0$.  
	
	In this scenario, no $\lambda$-powers are gained ($p=0$, hence $\gamma = -2$), and  $r(x,y)=0$. Incorporating the operator norm $\|\partial_\lambda^\ell \Gamma_{\alpha,\beta}^\pm\| \lesssim \lambda^{-\ell}$, the corresponding  amplitude factor obeys:  
	\[    
	|\partial_\lambda^\ell \mathcal{E}_{fg}^\pm(\lambda, x, y)| \lesssim |\log\lambda|^{-2}\lambda^{-2-\ell} \widetilde{\chi}_1(\lambda/2), \quad \ell=0,1,2.    
	\]    
	Applying Lemma \ref{oscillatory}(ii) with  $(\sigma,\nu)=(2,2)$  dictates the slowest  decay rate $(\log(2+|t|))^{-1}$. All other cross-pairings contain remainder powers ($p \ge 1$), which yield the bound $\langle t \rangle^{-1/2} (\log(2+|t|))^{-2}$ by  Lemma \ref{oscillatory}(ii) with parameters $(\sigma,\nu)=(1,2).$   
	
	 Summing the bounds from Case 1 and Case 2 concludes the proof. In particular, we derive
for	$B^\pm(\lambda)=\mathcal{M}_{\alpha, \beta}^{\pm}(\lambda)$ with $\alpha,\beta \in \{5,6\}$, 
	 \[
	 	\begin{aligned}
	 \left(\mathcal{K}_{B}^+-\mathcal{K}_{B}^-\right)(t,x,y)&=\frac{2}{\pi i} K_{\alpha,\beta}(t,x,y)+O\bigl(\langle t \rangle^{-\frac{1}{2}} \bigr),\\
  \left(\mathcal{N}_{B}^+-\mathcal{N}_{B}^-\right)(t,x,y)&=\frac{2}{\pi i} N_{\alpha,\beta}(t,x,y)+O\bigl(\langle t \rangle^{-\frac{1}{2}} \bigr),
  \end{aligned}
     \]
	 where 
	 \begin{equation}\label{kn-sharp}
	 \begin{split}
	 	K_{\alpha,\beta}(t,x,y)&=b_0^2\int_0^\infty  \cos(t\lambda^2) \lambda^{-1}(\log\lambda)^{-2}\widetilde{\chi}_1(\lambda) \big\langle \bigl(\Gamma_{\alpha,\beta}^+
        -\Gamma_{\alpha,\beta}^-\bigr)(\lambda)(Q_\beta v G_2^0)( \cdot, y), (Q_\alpha v G_2^0)( \cdot, x) \big\rangle d\lambda,\\
		N_{\alpha,\beta}(t,x,y)&=\frac{b_0^2}{t}\int_0^\infty  \sin(t\lambda^2) \lambda^{-3}(\log\lambda)^{-2}\widetilde{\chi}_1(\lambda) \big\langle \bigl(\Gamma_{\alpha,\beta}^+
        -\Gamma_{\alpha,\beta}^-\bigr)(\lambda)(Q_\beta v G_2^0)( \cdot, y), (Q_\alpha v G_2^0)( \cdot, x) \big\rangle d\lambda,
	 	\end{split}
	 \end{equation}
	 satisfying $|K_{\alpha,\beta}(t,x,y)|+|N_{\alpha,\beta}(t,x,y)|\lesssim(\log(2+|t|))^{-1},$ uniformly in $x,y$.
\end{proof}   


\begin{proposition}\label{third_merged_2} 
	Let $B^\pm(\lambda)=\mathcal{M}_{\alpha,\beta}^{\pm}(\lambda)$ with $\alpha=5,6$, $0\le\beta\le4$, or $0\le\alpha\le4$, $\beta=5,6$. Then
	\[
\sup_{x,y\in\mathbb{R}^2}\left|\mathcal{K}_{B}^\pm(t,x,y)\right|
	+ \sup_{x,y\in\mathbb{R}^2}\left|\mathcal{N}_{B}^\pm(t,x,y)\right|
	\lesssim{\langle t\rangle^{-1/2}} \log(2+|t|).
	\]
\end{proposition}
\begin{comment}
\begin{proof}
	Without loss of generality, we consider $\alpha \in \{5,6\}$ and $0 \le \beta \le 4$ (the symmetric case follows identically).  For the resolvent adjacent to $Q_5$ or $Q_6$, we employ the expansion \eqref{ex-Q56}.
	For the resolvent adjacent to $Q_4$  and $Q_\beta$ with $\beta\in\{0,1,2,3\}$, we use the following  expansions, respectively:
	\[	\bigl(Q_4 v R_0^\pm(\lambda^4)\bigr)(\cdot,z) = \mathcal{J}_{4,\pm}(\lambda, \cdot,z),\quad
	\bigl(Q_\beta v R_0^\pm(\lambda^4)\bigr)(\cdot,z) = \lambda^{-1-\delta_{\beta0}} \Omega_{\beta,\pm}(\lambda,\cdot,z),
	\]
	where $\delta_{\beta0}=1$ if $\beta=0$ and $\delta_{\beta0}=0$ otherwise.
	
	Following the same method as before, we obtain the canonical oscillatory integral forms \eqref{3-4res-integral}. 
	Moreover,
	recall  the operator estimates  $\|\partial_\lambda^\ell \mathcal{M}_{\alpha,\beta}^\pm\|_{\mathbb{B}(L^2)} \lesssim \lambda^{-\ell}$ for $\ell=0,1,2,$ and the crucial exponents are as follows:
	\[
	4 - k_\alpha - k_0 = 1\;\;(\alpha=5,6),\quad
	4 - k_\alpha - k_\beta = 0\;\;(\alpha=5,6,\; \beta=1,2,3),\quad
	4 - k_\alpha - k_4 = -1\;\;(\alpha=5,6).
	\] 	These dictate the behavior of the amplitude factor, leading to
	the following estimates. 
	
	(i) For $0\le\alpha\le3$, $\beta=5,6$ or $\alpha=5,6$, $0\le\beta\le3$,  Lemma~\ref{oscillatory}(ii) with $(\sigma,\nu)=(1,0)$ yields
	\[
\sup_{x,y\in\mathbb{R}^2}\left|\mathcal{K}_{B}^\pm(t,x,y)\right|
	+ \sup_{x,y\in\mathbb{R}^2}\left|\mathcal{N}_{B}^\pm(t,x,y)\right|
	\lesssim \langle t\rangle^{-1/2}.
	\]
	
	(ii) For $\alpha=4$, $\beta=5,6$ or $\alpha=5,6$, $\beta=4$, Lemma~\ref{oscillatory} (ii) with $(\sigma,\nu)=(1,-1)$ gives
	\[
\sup_{x,y\in\mathbb{R}^2}\left|\mathcal{K}_{B}^\pm(t,x,y)\right|
	+ \sup_{x,y\in\mathbb{R}^2}\left|\mathcal{N}_{B}^\pm(t,x,y)\right|
	\lesssim \langle t\rangle^{-1/2} \log(2+|t|) .\]
\end{proof}
\end{comment}
\begin{proof}
The proof follows the same argument as in Proposition
\ref{third_merged_1}. Indeed, using \eqref{ex-Q56} for the
$Q_5$- and $Q_6$-projected free resolvents, together with the
corresponding estimates for the $Q_\beta$-projected resolvent
($0\leq\beta\leq4$) (see Lemma \ref{lemma_projection} (i)-(ii)), one reduces the kernels to the canonical
oscillatory integrals in \eqref{3-4res-integral}. The resulting
amplitude factors satisfy the bounds required by Lemma
\ref{oscillatory}(ii) with $(\sigma,\nu)=(1,0)$ when
$0\leq\beta\leq3$, and with $(\sigma,\nu)=(1,-1)$ when
$\beta=4$, the latter logarithmic loss coming from
$\mathcal J_{4,\pm}$. Hence the desired
$\langle t\rangle^{-1/2}\log(2+|t|)$ bound follows. The symmetric case is treated in the same way.
\end{proof}







	It remains to establish the sharpness of the $(\log|t|)^{-1}$ bound.
	Synthesizing the preceding analysis, we conclude that for $|t|\gg 1$, 
the decay rate is determined by the kernels $\mathcal{K}_{\alpha,\beta}(t,x,y)$ and $\mathcal{N}_{\alpha,\beta}(t,x,y)$ defined in \eqref{kn-sharp}.
	To identify their leading-order behavior, we need
	a more precise decomposition of
	\(\Gamma_{\alpha,\beta}^\pm(\lambda)\).    Returning to the proof of Theorem~\ref{thm:M_inverse} in Section~\ref{subsec:proof_inverse}, we combine \eqref{Gamma55}, \eqref{Gamma56} \eqref{Gamma65}, and \eqref{Gamma66}
	to derive the refined expansion:  
	\[    
	\Gamma_{\alpha,\beta}^{\pm}(\lambda) = \mathcal{L}_{\alpha,\beta}^\pm + (-\log \lambda)^{-1/2} \Lambda^\pm(\lambda),
	\]    
where the leading terms are given by  
	\begin{equation}\label{mathcal_L} 
		\begin{aligned}  
			\mathcal{L}_{5,5}^{\pm} &= -a_3^\pm b_1^{-2}(Q_5 v G_6 v Q_5)^{-1}, \ \  
			&&\mathcal{L}_{5,6}^{\pm} = -a_3^\pm \mathcal{D}_{5,5} (Q_5 v G_6^0 v Q_6) \mathcal{D}_{6,6}, \\  
			\mathcal{L}_{6,5}^{\pm} &= -a_3^\pm \mathcal{D}_{6,6} (Q_6 v G_6^0 v Q_5) \mathcal{D}_{5,5}, \ \  
			&&\mathcal{L}_{6,6}^{\pm} = a_3^\pm b_1 \mathcal{D}_{6,6} (Q_6 v G_6^0 v Q_5) \mathcal{D}_{5,5} (Q_5 v G_6^0 v Q_6) \mathcal{D}_{6,6},  
		\end{aligned}  
	\end{equation}  
	with $\mathcal{D}_{5,5} = -(b_1 Q_5 v G_6 v Q_5)^{-1}$ and $\mathcal{D}_{6,6} = (b_1 Q_6 v G_6^0 v Q_6)^{-1}$.
	
	This decomposition isolates the dominant singularity:    
	\begin{itemize}
	\item Replacing $\Gamma_{\alpha,\beta}^{\pm}(\lambda)$ with $(-\log\lambda)^{-1/2}\Lambda^{\pm}(\lambda)$ in \eqref{kn-sharp} yields an improved decay rate of $(\log(2+|t|))^{-3/2}$, obtained via Lemma~\ref{oscillatory} (ii) with $(\sigma, \nu) = (2, 5/2)$.
    \vskip0.2cm
	\item The term $\mathcal{L}_{\alpha,\beta}^{\pm}$ in $\Gamma_{\alpha,\beta}^{\pm}(\lambda)$  produces the slowest decay $(\log(2+|t|))^{-1}$, given by Lemma~\ref{oscillatory} (ii) with $(\sigma, \nu) = (2, 2)$. And by Remark~\ref{remark:osci}, this rate comes from low-frequency region $|t|\lambda^2 \ll 1$, while high-frequency region $|t|\lambda^2 \gtrsim 1$ yields the faster $(\log(2+|t|))^{-2}$ decay.
\end{itemize} 



By isolating this primary low-frequency contribution (specifically, the regime $|t|\lambda^2 \ll 1$), we deduce the asymptotic expansion in $\mathbb{B}(L^1, L^\infty)$:      
\[      
	\begin{aligned}      
		\cos (t \sqrt{H}) P_{\mathrm{ac}}(H)\chi_1(H) &= -\frac{2}{\pi i}\sum_{\alpha,\beta = 5}^{6} T_{\mathcal{K}_{\alpha,\beta}}(t) + {O}\Big((\log|t|)^{-3/2}\Big), \\      
		\frac{\sin (t \sqrt{H})}{t\sqrt{H}} P_{\mathrm{ac}}(H)\chi_1(H) &= -\frac{2}{\pi i}\sum_{\alpha,\beta = 5}^{6} T_{\mathcal{N}_{\alpha,\beta}}(t) + {O}\Big((\log|t|)^{-3/2}\Big),   
	\end{aligned}      
\]    
where $T_{\mathcal{K}_{\alpha,\beta}}(t)$ and $T_{\mathcal{N}_{\alpha,\beta}}(t)$ are integral operators defined by the respective kernels:    
\[
\begin{aligned}      
	\mathcal{K}_{\alpha,\beta}(t,x,y) &= c_{\alpha,\beta}(x,y)\int_0^{\infty} \cos(t\lambda^2) \lambda^{-1} \widetilde{\chi}_1(\lambda)\chi_1(t \lambda^2) (\log \lambda)^{-2} \, d\lambda, \\      
	\mathcal{N}_{\alpha,\beta}(t,x,y) &= {c_{\alpha,\beta}(x,y)}{t}^{-1} \int_0^{\infty} \sin(t\lambda^2) \lambda^{-3} \widetilde{\chi}_1(\lambda) \chi_1(t \lambda^2) (\log \lambda)^{-2} \, d\lambda,   
\end{aligned}
\]
with  $c_{\alpha,\beta}(x,y) = b_0^2 \big\langle (\mathcal{L}_{\alpha,\beta}^+ -\mathcal{L}_{\alpha,\beta}^-) (Q_\beta v G_2^0)(\cdot, y), (Q_\alpha v G_2^0)(\cdot, x) \big\rangle$.   
Here the summation over $\alpha, \beta \in \{5,6\}$ accommodates the third-kind resonance scenario, where $Q_6 = 0$ implies $c_{\alpha,\beta} \equiv 0$ for all pairs $(\alpha,\beta) \neq (5,5)$.  
  
Thus, establishing the optimal bounds in \eqref{optimality_31} reduces to proving the following Proposition \ref{sharp-3-4}.


\begin{proposition}\label{sharp-3-4}  
	For $|t| \gg 1$, the leading operators satisfy the lower bounds:  
\begin{equation}\label{low_bound_weighted_1}  
		\Big\|\sum_{\alpha,\beta = 5}^{6}T_{ \mathcal{K}_{\alpha,\beta}}(t)\Big\|_{L^1 \to L^\infty} \gtrsim \frac{1}{\log |t|}, \quad   
		\Big\|\sum_{\alpha,\beta = 5}^{6} T_{\mathcal{N}_{\alpha,\beta}}(t)\Big\|_{L^1 \to L^\infty} \gtrsim \frac{1}{\log |t|}.  
	\end{equation}  
\end{proposition}  

\begin{proof}
	 By continuous embeddings $L^{2,\sigma}(\mathbb{R}^2) \hookrightarrow L^1(\mathbb{R}^2)$ and $L^\infty(\mathbb{R}^2) \hookrightarrow L^{2,-\sigma}(\mathbb{R}^2)$ for $\sigma > 1$, it is  sufficient to construct a specific test function $\varphi \in L^{2,\sigma}(\mathbb{R}^2)$ such that:  
\begin{equation}\label{eq:lower_bounds}  
		\Big|\Big\langle \sum_{\alpha,\beta = 5}^{\mathbf{k}+2} T_{\mathcal{K}_{\alpha,\beta}}(t) \varphi, \varphi \Big\rangle \Big| \gtrsim \frac{1}{\log |t|}, \quad \text{and} \quad \Big|\Big\langle \sum_{\alpha,\beta = 5}^{\mathbf{k}+2} T_{\mathcal{N}_{\alpha,\beta}}(t) \varphi, \varphi \Big\rangle \Big| \gtrsim \frac{1}{\log |t|}. 
	\end{equation}  
	
	Under the hypothesis that zero is a third-kind resonance or an eigenvalue accompanied by a $d$-wave resonance, the subspace $Q_5 L^2\neq\{0\}$. We fix  $\psi \in Q_5 L^2$ and explicitly construct our test function as $\varphi := v U \psi $. By $T_0Q_5=0,$ one has $U\psi=-b_0vG_2^0v\psi.$ Thus $\varphi= v U \psi = -b_0 |V| G_2^0 v\psi.$ By Lemma \ref{estimate_G_2^0},
	 the estimate $|\varphi(x)| \lesssim \langle x \rangle^{-\mu-1}(1+\log\langle x\rangle)$ guarantees that $\varphi \in L^{2,\sigma}(\mathbb{R}^2)$.  
	 
	 Since $Q_5 T_0 = 0$, we have $Q_5(U + b_0 v G_2^0 v)U \psi = 0$, which dictates that  
	\[  
	b_0 Q_5 v G_2^0 \varphi = b_0 Q_5 v G_2^0 v U \psi = -Q_5 \psi = -\psi.  
	\]  
Furthermore, by identical logic, using $Q_6T_0=0$ and the orthogonality of $\psi \in Q_5 L^2$ to $Q_6 L^2$, we have
	$$
	b_0 Q_6 v G_2^0 \varphi=b_0 Q_6 v G_2^0 v U \psi =  -Q_6 \psi=0.
	$$
 Consequently, $\langle c_{\alpha,\beta} \varphi, \varphi \rangle=b_0^2\big\langle (\mathcal{L}_{\alpha,\beta}^+ -\mathcal{L}_{\alpha,\beta}^-) (Q_\beta v G_2^0)\varphi, (Q_\alpha v G_2^0)\varphi \big\rangle=0$
  for all $(\alpha,\beta) \neq (5,5)$.  
	The only surviving term corresponds to $(\alpha,\beta)=(5,5)$. Substituting \(\mathcal{L}_{5,5}^{\pm}=-a_3^\pm b_1^{-2}(Q_5vG_6vQ_5)^{-1}\) (defined in \eqref{mathcal_L}), and utilizing the strict negative definiteness of \((Q_5 v G_6 v Q_5)^{-1}\) on $Q_5L^2$, together with $\text{Im}(a_3^\pm) \neq 0$, we obtain
	\begin{align*}  
		\langle c_{5,5} \varphi, \varphi \rangle &= b_0^2 \big\langle (\mathcal{L}_{5,5}^+ -\mathcal{L}_{5,5}^-) Q_5 v G_2^0 \varphi, Q_5 v G_2^0 \varphi \big\rangle \\  
		&= -(a_3^+ - a_3^-) b_1^{-2}\big\langle (Q_5 v G_6 v Q_5)^{-1} \psi, \psi \big\rangle := c_0 \neq 0.  
	\end{align*}  
	
	Finally, evaluating the temporal integrals for $|t| \gg 1$, we explicitly extract the logarithmic decay:  
	\begin{align*}  
		\left|\int_0^{\infty} \cos(t\lambda^2) \lambda^{-1} \widetilde{\chi}_1(\lambda)  \chi_1(t \lambda^2) (\log \lambda)^{-2} \, d\lambda \right| &\ge \cos 1 \int_0^{|t|^{-1/2} \lambda_0^{1/2}} \lambda^{-1} (\log \lambda)^{-2} \, d\lambda \gtrsim \frac{1}{\log |t|}.  
	\end{align*}  
		An analogous estimate using ${\sin x}/{x} \ge \sin 1$ on $[0,1]$ holds for the sine integral. Summing the non-vanishing $(5,5)$ contribution secures the lower bounds \eqref{eq:lower_bounds}, which, via the embeddings,  yields \eqref{low_bound_weighted_1}. 
\end{proof}  


\subsection{The estimates for the case $Q_{5} = 0$}
We next turn to the case where zero is an eigenvalue of $H$ (i.e., $\mathbf{k} = 4$) and all solutions $\phi \in W_{0}(\mathbb{R}^2)$ to $H\phi = 0$ are in $L^2(\mathbb{R}^2)$, i.e., $Q_5 L^2 = \{0\}$.

The condition $Q_5 L^2 = \{0\}$ causes certain terms in $(M^\pm(\lambda))^{-1}$ to vanish, leading to the improved decay rate in Theorem \ref{main_theorem_low_32}.
Since the asymptotic expansion of $(M^\pm(\lambda))^{-1}$ under the assumption $Q_5 L^2 = \{0\}$ follows by an argument analogous to that in the proof of Theorem~\ref{thm:M_inverse}, we omit the details and directly state the result in Lemma~\ref{M_inverse_optimal} below.


\begin{lemma}\label{M_inverse_optimal}  
	Assume that zero is an eigenvalue of $H$ but exhibits no $d$-wave resonance $(\text{i.e.,}\  Q_5 = 0)$. Suppose that $|V(x)| \lesssim \langle x \rangle^{-22-}$ and $0<\lambda\ll1$. Then 
	\begin{itemize}  
		\item[(i)] For $0 \le \alpha, \beta \le 4$, the estimates for $\mathcal{M}_{\alpha,\beta}^\pm(\lambda)$ established in Theorem~\ref{thm:M_inverse} remain valid.
		
		\item[(ii)] For $1 \le \alpha \le 3$ and $\ell = 0,1,2$,  
		\begin{align}  \label{M6123}
		\bigl\|\partial_\lambda^{\ell} \mathcal{M}_{6,\alpha}^\pm(\lambda)\bigr\|_{\mathbb{B}(L^2)}  
		+ \bigl\|\partial_\lambda^{\ell} \mathcal{M}_{\alpha,6}^\pm(\lambda)\bigr\|_{\mathbb{B}(L^2)}  
		\lesssim \lambda^{1-\ell}.  
		\end{align}
		
		\item[(iii)] For $(\alpha,\beta) =(6,0),\,(0,6),\,(6,4),\,(4,6),\,(6,6)$, the operators admit the following expansions:  
		\[  
		\begin{aligned}  
			\mathcal{M}_{6,0}^{\pm}(\lambda) &= \lambda (\log \lambda)Q_6\Lambda^\pm(\lambda)Q_0, \quad 
			\mathcal{M}_{0,6}^{\pm}(\lambda) = \lambda (\log \lambda)Q_0\Lambda^\pm(\lambda)Q_6,  \\  
			\mathcal{M}_{6,4}^{\pm}(\lambda) &=  \lambda (\log \lambda)^2 Q_6 \Lambda^\pm(\lambda)Q_{42}^1+\lambda (\log \lambda) Q_6 \Lambda^\pm(\lambda)Q_4,\\
		\mathcal{M}_{4,6}^{\pm}(\lambda) &= \lambda (\log \lambda)^2 Q_{42}^1 \Lambda^\pm(\lambda)Q_{6} +\lambda (\log \lambda) Q_4 \Lambda^\pm(\lambda)Q_6, \\ 
		\mathcal{M}_{6,6}^{\pm}(\lambda) &= \mathcal{D}_{6,6}
			+  \lambda^2 (\log \lambda)^2Q_6 \Lambda^\pm(\lambda)Q_6.
		\end{aligned}  
		\]  
	\end{itemize}
	
	Moreover, if we additionally assume $Q_4 = 0$, then the operators $\mathcal{M}_{\alpha,4}^{\pm}(\lambda)$ and $\mathcal{M}_{4,\alpha}^{\pm}(\lambda)$ vanish identically for all $0\leq \alpha\leq 6$, and the following refined expansions hold:
	\[  
	\begin{aligned}  
		\mathcal{M}_{6,0}^{\pm}(\lambda) &= -a_2^\pm (a_0^\pm)^{-1} \|V\|_{L^1}^{-1} \lambda \mathcal{D}_{6,6} Q_6 vG_4 vQ_0   
		+ \lambda^2 (-\log \lambda)^{3/2}Q_6 \Lambda^\pm(\lambda)Q_0, \\  
		\mathcal{M}_{0,6}^{\pm}(\lambda) &= -a_2^\pm (a_0^\pm)^{-1} \|V\|_{L^1}^{-1} \lambda Q_0 vG_4 vQ_6 \mathcal{D}_{6,6}  
		+ \lambda^2 (-\log \lambda)^{3/2} Q_0\Lambda^\pm(\lambda)Q_6,  \\  
		\mathcal{M}_{6,6}^{\pm}(\lambda) &= \mathcal{D}_{6,6} -a_4^\pm \lambda^2 \mathcal{D}_{6,6} Q_6 vG_8 vQ_6 \mathcal{D}_{6,6} \\  
		&\quad +(a_2^\pm)^2 (a_0^\pm)^{-1} \|V\|_{L^1}^{-1} \lambda^2 \mathcal{D}_{6,6} Q_6 vG_4 v  Q_0 vG_4 vQ_6 \mathcal{D}_{6,6} +\lambda^3 (-\log \lambda)^{3/2} Q_6\Lambda^\pm(\lambda)Q_6.  
	\end{aligned}  
	\]  
In the above, $\mathcal{D}_{6,6} = (b_1 Q_6 v G_6^0 v Q_6)^{-1}$ and  $\Lambda^\pm(\lambda)$ denotes a generic $\lambda$-dependent operator in $\mathbb{B}(L^2)$ (possibly differing at each occurrence) satisfying
$\|\partial_\lambda^{\ell} \Lambda^\pm(\lambda)\|_{\mathbb{B}(L^2)} \lesssim \lambda^{-\ell}$ for $ \ell=0,1,2.
$
\end{lemma}


Based on the expansions provided by Lemma~\ref{M_inverse_optimal}, we now proceed to prove Theorem~\ref{main_theorem_low_32}. In view of the extensive analysis carried out in the previous sections, the proof reduces to the study of the difference kernels
$
\mathcal{K}_{B}^+ - \mathcal{K}_{B}^-$ and $\mathcal{N}_{B}^+ - \mathcal{N}_{B}^-$,
with
\[
B^\pm(\lambda) \in \bigl\{ \mathcal{M}_{\alpha, \beta}^{\pm}(\lambda) \mid \alpha=6,\ 0 \le \beta \le 4 \ \text{or} \ 0 \le \alpha \le 4,\ \beta=6 \ \text{or} \ \alpha=\beta=6 \bigr\},
\]
where the kernels $\mathcal{K}_{B}^\pm$, $\mathcal{N}_{B}^\pm$ are as defined in \eqref{K_1} and \eqref{Lambda_1}, respectively. 
 We divide the proof into two parts according to
whether a $p$-wave resonance is present or absent.

\vspace{0.3cm}    
\textbf{Part I: The Case $Q_4 \neq 0$ (Presence of $p$-wave resonance).}    
 We establish the  kernel estimates for all components involving the index 6 in the following Proposition \ref{prop_part1}.  

\begin{proposition}\label{prop_part1}    
	Assume $Q_5 = 0$ and $Q_4 \neq 0$. For any operator $B^\pm(\lambda)$ containing the index $6$, the associated integral kernels satisfy the following estimates:  
	\begin{itemize}
		\item[(i)] If $B^\pm(\lambda) \in \{ \mathcal{M}_{\alpha,6}^{\pm}(\lambda), \mathcal{M}_{6,\alpha}^{\pm}(\lambda) \}$ for $1 \le \alpha \le 3$, then  
		\[
\sup_{x,y\in\mathbb{R}^2}\left|\mathcal{K}_{B}^\pm(t,x,y)\right|
+\sup_{x,y\in\mathbb{R}^2}\left|\mathcal{N}_{B}^\pm(t,x,y)\right|
	\lesssim \langle t\rangle^{-1}.
	\] 
	Moreover,  both $\mathcal{K}_{B}^\pm$ and $\mathcal{N}_{B}^\pm$  satisfy the  improved weighted bound $\omega(x)\omega(y)(|t|\log|t|)^{-1}$ for $|t| \ge 2$.
	In particular, the same unweighted and weighted bounds remain valid under the alternative assumption $Q_5 = Q_4 = 0$, since in this case $\mathcal{M}_{\alpha,6}^{\pm}(\lambda)$ and $\mathcal{M}_{6,\alpha}^{\pm}(\lambda)$ also satisfy the estimate \eqref{M6123}.
		\item[(ii)] If $B^\pm(\lambda) \in \{ \mathcal{M}_{0,6}^{\pm}(\lambda), \mathcal{M}_{6,0}^{\pm}(\lambda) \}$, then
		\[
\sup_{x,y\in\mathbb{R}^2}\left|\mathcal{K}_{B}^\pm(t,x,y)\right|
	+ \sup_{x,y\in\mathbb{R}^2}\left|\mathcal{N}_{B}^\pm(t,x,y)\right|
	\lesssim \langle t\rangle^{-1} \log(2+|t|).
	\]
		\item[(iii)] If $B^\pm(\lambda) \in \{ \mathcal{M}_{4,6}^{\pm}(\lambda), \mathcal{M}_{6,4}^{\pm}(\lambda), \mathcal{M}_{6,6}^{\pm}(\lambda) \}$, then  
		\[
\sup _{x, y \in \mathbb{R}^2}\left|\left(\mathcal{K}_{B}^+-\mathcal{K}_{B}^-\right)(t,x,y)\right| +
	\sup _{x, y \in \mathbb{R}^2}\left|\left(\mathcal{N}_{B}^+-\mathcal{N}_{B}^-\right)(t,x,y)\right|
	\lesssim \langle t\rangle^{-1} (\log(2+|t|))^{2}.
	\]
	\end{itemize}  
\end{proposition}   
\begin{comment}
\begin{proof}   
 (i)  Let $B^\pm(\lambda) \in \{\mathcal{M}_{\alpha,6}^{\pm}(\lambda), \mathcal{M}_{6,\alpha}^{\pm}(\lambda)\}$ for $1 \le \alpha \le 3$.
	We utilize the following expansions for the projected free resolvents:
	\begin{align}
		\bigl(Q_6 v R_0^\pm(\lambda^4)\bigr)(\cdot,z) &= (b_0 Q_6 v G_2^0)( \cdot,z) + \lambda^{2}\bigl(\mathcal{T}_{6}(\lambda, \cdot,z)+\mathcal{T}_{6,\pm}(\lambda, \cdot,z)\bigr), \label{Q666} \\
		\bigl(Q_\alpha v R_0^\pm(\lambda^4)\bigr)(\cdot,z) &= \lambda^{-1}\Omega_{\alpha,\pm}(\lambda,\cdot,z). \label{Q123}
	\end{align}
	Upon substitution into \eqref{K_1} and \eqref{Lambda_1}, the $\lambda^{-1}$ singularity originating from the $Q_\alpha$-expansion \eqref{Q123} is canceled by $4-k_6-k_\alpha = 0$ and $\|\partial_\lambda^{\ell} \mathcal{M}_{6,\alpha}^\pm\|_{\mathbb{B}(L^2)}+\|\partial_\lambda^{\ell} \mathcal{M}_{\alpha,6}^\pm\|_{\mathbb{B}(L^2)} \lesssim \lambda^{1-\ell}$ from Lemma \ref{M_inverse_optimal}(ii). Applying Lemma~\ref{oscillatory}(ii) with  $(\sigma,\nu) = (0,0)$ yields the uniform unweighted bound $O\bigl(\langle t \rangle^{-1}\bigr)$.
	
	To establish the improved weighted estimate, we replace \eqref{Q123} with the  expansion $\bigl(Q_\alpha vR_0^\pm(\lambda^4)\bigr) = \bigl(Q_\alpha v {E}_1^{\pm}\bigr)$ (see \eqref{eq:expansion}) to extract an extra $\lambda^{1-}$ factor at the cost of  spatial weights, mirroring the procedure in Proposition~\ref{prop_K_weight_1}. This  ensures the accelerated decay rate $O\bigl(\omega(x)\omega(y) |t|^{-1} (\log|t|)^{-1}\bigr)$.
	
	\vspace{0.15cm}  
	 (ii) Let $B^\pm(\lambda) \in \{\mathcal{M}_{6,0}^{\pm}(\lambda), \mathcal{M}_{0,6}^{\pm}(\lambda)\}$.
	We pair the $Q_6$ expansion \eqref{Q666} with the singular $Q_0$ expansion $\bigl(Q_0 v R_0^\pm(\lambda^4)\bigr)(\cdot,z) = \lambda^{-2}\Omega_{0,\pm}(\lambda,\cdot,z)$.
The inherent  factor for these indices is $\lambda^{4-k_6-k_0} = \lambda$ in  \eqref{K_1} and \eqref{Lambda_1}. According to Lemma \ref{M_inverse_optimal}(iii), the operator introduces an extra $\lambda$ and a logarithmic singularity: $\mathcal{M}_{6,0}^{\pm}(\lambda) = \lambda (\log \lambda)Q_6\Lambda^\pm(\lambda)Q_0$ (and symmetrically for $\mathcal{M}_{0,6}^\pm$).
	Multiplying these factors together, and applying Lemma~\ref{oscillatory}(ii) with parameters $(\sigma,\nu) = (0,-1)$ yields the uniform bound ${O}\big(\langle t\rangle^{-1} \log(2+|t|)\big)$.
	
	\vspace{0.15cm}
    (iii) Let $B^\pm(\lambda) \in \{\mathcal{M}_{6,4}^{\pm}(\lambda), \mathcal{M}_{4,6}^{\pm}(\lambda), \mathcal{M}_{6,6}^{\pm}(\lambda)\}$.
We employ \eqref{Q666} for $Q_6$, the identity $\bigl(Q_4 v R_0^\pm(\lambda^4)\bigr)(\cdot,z) = \mathcal{J}_{4,\pm}(\lambda, \cdot,z)$ for $Q_4$, and the expansion
$
\bigl(Q_{42}^1 v R_0^\pm(\lambda^4)\bigr)(\cdot,z) = (b_0 Q_{42}^1 v G_2^0)(\cdot,z) + \lambda(\mathcal{T}_{4,\pm}^{1,2} + \mathcal{T}_4^{1,2})(\lambda,\cdot,z)
$
for $Q_{42}^1$. Recall from Lemma \ref{lemma_projection} that $\mathcal{J}_{4,\pm}$ exhibits a $|\log \lambda|$ growth, whereas $\mathcal{T}_{4,\pm}^{1,2}$ and $\mathcal{T}_4^{1,2}$ carry no logarithmic singularity.

For $B^\pm(\lambda) =\mathcal{M}_{6,4}^{\pm}(\lambda) = \lambda (\log \lambda)^2 Q_6 \Lambda^\pm(\lambda) Q_{42}^1 + \lambda (\log \lambda) Q_6 \Lambda^\pm(\lambda) Q_4$ (and symmetrically for $\mathcal{M}_{4,6}^{\pm}$), the operator $R_0^\pm(\lambda^4) vQ_6\mathcal{M}_{6,4}^{\pm}(\lambda)Q_4vR_0^\pm(\lambda^4) $ yields the bound $\lambda (\log \lambda)^2$ from both terms: the first term $\lambda (\log \lambda)^2 Q_6 \Lambda^\pm(\lambda) Q_{42}^1$ explicitly, and the second term $\lambda (\log \lambda) Q_6 \Lambda^\pm(\lambda) Q_4$ also yields $\lambda (\log \lambda)^2$ because the $|\log \lambda|$ growth of $\mathcal{J}_{4,\pm}$ multiplies the existing $\lambda (\log \lambda)$ factor. Combining this with $\lambda^{4-k_6-k_4} = \lambda^{-1}$, we apply Lemma~\ref{oscillatory}(ii) with parameters $(\sigma,\nu)=(0,-2)$. This yields the desired bound $O\bigl(\langle t\rangle^{-1} (\log(2+|t|))^{2}\bigr)$.

Finally, for $B^\pm(\lambda) =\mathcal{M}_{6,6}^{\pm}(\lambda) = \mathcal{D}_{6,6} + \lambda^2 (\log \lambda)^2 Q_6 \Lambda^\pm(\lambda) Q_6$, we have $\lambda^{4-k_6-k_6} = \lambda^{-2}$. The leading term $\mathcal{D}_{6,6}$ is independent of the sign $\pm$; consequently, the difference $\mathcal{K}_{\mathcal{D}_{6,6}}^+ - \mathcal{K}_{\mathcal{D}_{6,6}}^-$ is bounded by $\langle t \rangle^{-1}$ (see \eqref{mathcalA66} in Proposition~\ref{third_merged_1} for details). For the remaining term involving $\lambda^2(\log\lambda)^2 Q_6 \Lambda^\pm(\lambda) Q_6$, applying Lemma~\ref{oscillatory}(ii) with $(\sigma,\nu)=(0,-2)$ gives the uniform bound $\langle t\rangle^{-1} (\log(2+|t|))^{2}$. Summing these contributions completes the proof of (iii).
\end{proof}
\end{comment}
\begin{proof}
	The proof follows the same canonical oscillatory-integral reduction
	used in the preceding propositions.  We use
	$
	Q_6vR_0^\pm(\lambda^4)
	=
	b_0Q_6vG_2^0
	+
	\lambda^2(\mathcal T_6+\mathcal T_{6,\pm}),
	$
	together with the projected resolvent expansions adjacent to
	$Q_0$, $Q_\alpha$ $(1\leq\alpha\leq3)$, and $Q_4$. Inserting
	these expansions and the bounds from Lemma
	\ref{M_inverse_optimal} into \eqref{K_1} and \eqref{Lambda_1}, we see that the
	resulting amplitudes satisfy the hypotheses of Lemma
	\ref{oscillatory}(ii) with $(\sigma,\nu)=(0,0)$ for the blocks
	$\mathcal M_{6,\alpha}^\pm,\mathcal M_{\alpha,6}^\pm$
	$(1\leq\alpha\leq3)$, with $(\sigma,\nu)=(0,-1)$ for the blocks
	$\mathcal M_{6,0}^\pm,\mathcal M_{0,6}^\pm$, and with
	$(\sigma,\nu)=(0,-2)$ for the blocks
	$\mathcal M_{6,4}^\pm,\mathcal M_{4,6}^\pm,\lambda^2 (\log \lambda)^2Q_6 \Lambda^\pm(\lambda)Q_6$.
	These three cases give respectively
	$\langle t\rangle^{-1}$,
	$\langle t\rangle^{-1}\log(2+|t|)$, and
	$\langle t\rangle^{-1}(\log(2+|t|))^2$. The sign-independent
	principal part $\mathcal D_{6,6}$ is handled as in
	\eqref{mathcalA66} and satisfies the stronger
	$\langle t\rangle^{-1}$ bound. 
    The weighted estimate in part (i)
	is obtained exactly as in Proposition \ref{prop_K_weight_1}.
\end{proof}


\begin{proof}[\textbf{Proof of Theorem \ref{main_theorem_low_32} (1)}]  
	Combining  Proposition \ref{prop_part1} with the  estimates established  in the regular, first-kind, and second-kind resonance cases,
	we conclude the proof of Theorem \ref{main_theorem_low_32}(1).  
\end{proof}  

\vspace{0.3cm}  
\textbf{Part II: The Case $Q_4 = 0$ (Absence of $p$-wave resonance).}  
In this part, based on the expansions detailed in Lemma \ref{M_inverse_optimal}, the absence of $Q_4$ guarantees that  $\mathcal{M}_{6,0}^{\pm}(\lambda)$ and $\mathcal{M}_{0,6}^{\pm}(\lambda)$ generate an explicit  $\lambda$ factor rather than $\lambda\log\lambda$. Concurrently, the remainder in $\mathcal{M}_{6,6}^{\pm}(\lambda)$ scales as ${O}(\lambda^2)$ instead of ${O}(\lambda^2(\log\lambda)^2)$:
\begin{align}\label{Q4=Q5}
\mathcal{M}_{6,0}^{\pm}(\lambda) = \lambda Q_6 \Lambda^\pm(\lambda)Q_0, \quad 
\mathcal{M}_{0,6}^{\pm}(\lambda) = \lambda Q_0 \Lambda^\pm(\lambda)Q_6, \quad 
\mathcal{M}_{6,6}^{\pm}(\lambda) = \mathcal{D}_{6,6} + \lambda^2 Q_6 \Lambda^\pm(\lambda)Q_6.
\end{align}
Following the same method as in Proposition \ref{prop_part1}, we obtain  for all $B^\pm(\lambda) \in \{ \mathcal{M}_{\alpha,6}^{\pm}(\lambda), \mathcal{M}_{6,\alpha}^{\pm}(\lambda) \}$ with $\alpha \in \{0,1,2,3,6\}$, the associated kernels exhibit a uniform ${O}(\langle t \rangle^{-1})$ decay. 

Synthesizing this with the regular, first-kind, and second-kind resonance cases, immediately establishes the unweighted estimate ${O}(\langle t \rangle^{-1})$  stated in Theorem \ref{main_theorem_low_32}(2).

To obtain the refined weighted asymptotic expansion, we isolate the precise leading-order constants generated by these operators. We formalize this in the following proposition.

\begin{proposition}\label{prop_part2}  
	Assume $Q_5 = 0$ and $Q_4 = 0$. For $|t|\ge 2$, the low-energy part  of the time evolution admits the following asymptotic expansions in $\mathbb{B}(L^1_\omega, L^\infty_{-\omega})$:
	\[
		\begin{aligned}
			\cos(t\sqrt{H}) P_{\mathrm{ac}}(H) \chi_1(H) &= -\frac{2}{\pi i}\widetilde{{\mathcal{B}}} I_c(t) + \mathcal{O}\!\bigl( |t|^{-1} (\log |t|)^{-1} \bigr), \\[4pt]  
			\frac{\sin(t\sqrt{H})}{t\sqrt{H}} P_{\mathrm{ac}}(H) \chi_1(H) &= -\frac{2}{\pi i}\widetilde{\mathcal{B}} I_s(t) + \mathcal{O}\!\bigl( |t|^{-1} (\log |t|)^{-1} \bigr),  
		\end{aligned}
	\]
	where  $I_c(t) = {O}(|t|^{-N})$ for   $\forall N>0$,  $I_s(t) = \frac{\pi}{4|t|} + {O}(|t|^{-2})$ and the time-independent  operator $\widetilde{\mathcal{B}}$ is given by:  
	\begin{equation}  \label{mathcal B}
		\begin{split}
			\widetilde{\mathcal{B}} &= - b_0^2(a_4^+ - a_4^-) G_2^0 v \mathcal{D}_{6,6} Q_6 v G_8 v Q_6 \mathcal{D}_{6,6} v G_2^0 \\  
			&\quad + b_0^2\bigl((a_2^+)^2 (a_0^+)^{-1} - (a_2^-)^2 (a_0^-)^{-1}\bigr)\|V\|_{L^1}^{-1} G_2^0 v \mathcal{D}_{6,6} Q_6 v G_4 v Q_0 v G_4 v Q_6 \mathcal{D}_{6,6} v G_2^0.  
		\end{split}
	\end{equation}
\end{proposition}

\begin{proof}  
	To establish the weighted estimates, it suffices to focus  on  $B^\pm \in \{\mathcal{M}_{6,0}^{\pm}, \mathcal{M}_{0,6}^{\pm}, \mathcal{M}_{6,6}^{\pm}\}$. By  Proposition \ref{prop_part1}(i) and the results for the regular and first-kind resonance cases, all the remaining terms have already been shown to obey the  bound $\mathcal{O}\!\bigl( |t|^{-1} (\log |t|)^{-1} \bigr)$. 

We start to evaluate  $T_{\mathcal{K}_{B}^{+}}(t)-T_{\mathcal{K}_{B}^{-}}(t)$ for  $B^\pm \in \{\mathcal{M}_{6,0}^{\pm}, \mathcal{M}_{0,6}^{\pm}, \mathcal{M}_{6,6}^{\pm}\}$ when $|t|\ge 2$. Proceeding analogously to the proof of Proposition \ref{prop_second_critical_block}, we utilize  the free resolvent expansion in Lemma \ref{lem-reso}:
\begin{equation*}
	\begin{split}
		\bigl(Q_0v R^\pm_0(\lambda^4)\bigr)(\cdot, y) &= \frac{a_0^\pm}{\lambda^2}(Q_0vG_0)(\cdot, y) + (Q_0vE^\pm_1)(\lambda, \cdot, y), \\
		\bigl(Q_6v R^\pm_0(\lambda^4)\bigr)(\cdot, y) &= b_0(Q_6vG_2^0)(\cdot, y) + a_2^\pm\lambda^2(Q_6vG_4)(\cdot, y) + (Q_6vE^\pm_2)(\lambda, \cdot, y),
	\end{split}
\end{equation*}
where  $|\partial_\lambda^\ell E_1^\pm(\lambda,x,y)| \lesssim \lambda^{-\ell-}\langle x \rangle^{4} \langle y \rangle^{4}$ and $|\partial_\lambda^\ell E_2^\pm(\lambda,x,y)| \lesssim \lambda^{-\ell+4-}\langle x \rangle^{6} \langle y \rangle^{6}$ for $\ell=0,1,2$. 
Moreover, for the leading sign-independent term $\mathcal{D}_{6,6}$ within $\mathcal{M}_{6,6}^{\pm}(\lambda)$, we  also need to exploit  the difference:  
\begin{align*} 
	&\bigl[R_0^+ (\lambda^4)v Q_6 \mathcal{D}_{6,6} Q_6 v R_0^+(\lambda^4)\bigr](x,y) - \bigl[R_0^- (\lambda^4)v Q_6 \mathcal{D}_{6,6} Q_6 v R_0^-(\lambda^4)\bigr](x,y)  \\  
	&= \bigl[(R_0^+(\lambda^4)-R_0^-(\lambda^4))v  Q_6 \mathcal{D}_{6,6} Q_6 v R_0^+(\lambda^4) \bigr](x,y) + \bigl[ R_0^-(\lambda^4) v  Q_6 \mathcal{D}_{6,6} Q_6 v (R_0^+(\lambda^4)-R_0^-(\lambda^4))\bigr](x,y).  
\end{align*}  

Synthesizing this analysis and extracting the dominant constants, we derive the leading-order expansions for the integral operators:
\begin{itemize}  
	\item For $\mathcal{M}_{6,0}^{\pm}(\lambda)$:   
	$T_{\mathcal{K}_{B}^{+}} - T_{\mathcal{K}_{B}^{-}} = -\frac{2 b_0}{\pi i}  (a_2^+ - a_2^-)\|V\|_{L^1}^{-1} G_2^0 v \mathcal{D}_{6,6} Q_6 vG_4vQ_0vG_0 \, I_c(t) + \mathcal{O}\!\bigl( |t|^{-1} (\log |t|)^{-1} \bigr)$.  
	
	\item For $\mathcal{M}_{0,6}^{\pm}(\lambda)$:  
	$T_{\mathcal{K}_{B}^{+}} - T_{\mathcal{K}_{B}^{-}} = -\frac{2 b_0}{\pi i}  (a_2^+ - a_2^-)\|V\|_{L^1}^{-1} G_0vQ_0v G_4 vQ_6 \mathcal{D}_{6,6} v G_2^0 \, I_c(t) + \mathcal{O}\!\bigl( |t|^{-1} (\log |t|)^{-1} \bigr)$.  
	
	\item For $\mathcal{M}_{6,6}^{\pm}(\lambda)$:  
	$T_{\mathcal{K}_{B}^{+}} - T_{\mathcal{K}_{B}^{-}} = \frac{2}{\pi i} \Theta I_c(t) + \mathcal{O}\!\bigl( |t|^{-1} (\log |t|)^{-1} \bigr)$, where
	\begin{align*}  
		\Theta &= b_0 (a_2^+ - a_2^-) G_2^0 v \mathcal{D}_{6,6} Q_6 vG_4 + b_0 (a_2^+ - a_2^-) G_4 vQ_6 \mathcal{D}_{6,6} v G_2^0 \\  
		&\quad - b_0^2(a_4^+ - a_4^-) G_2^0 v \mathcal{D}_{6,6} Q_6 v G_8 v Q_6 \mathcal{D}_{6,6} v G_2^0 \\  
		&\quad + b_0^2 \bigl((a_2^+)^2 (a_0^+)^{-1} - (a_2^-)^2 (a_0^-)^{-1}\bigr) \|V\|_{L^1}^{-1}  G_2^0 v \mathcal{D}_{6,6} Q_6 v G_4 v Q_0 v G_4 v Q_6 \mathcal{D}_{6,6} v G_2^0.
	\end{align*}  
\end{itemize} 
	Here  $T_{\mathcal{K}_{B}^{\pm}}(t)$ denotes the integral operator with the kernel ${\mathcal{K}_{B}^{\pm}}(t,x,y)$ and
	$$I_c(t) = \int_0^\infty \cos(t\lambda^2) \lambda \widetilde{\chi}_1(\lambda) \, d\lambda = O(|t|^{-N}), \  \quad\forall N>0.$$
	By summing these contributions, the total coefficient of $I_c(t)$ is $\frac{2}{\pi i}\widetilde{\mathcal{B}}$ where
	$$  
	\widetilde{\mathcal{B}} = -b_0 (a_2^+ - a_2^-)\|V\|_{L^1}^{-1} G_2^0 v \mathcal{D}_{6,6} Q_6 vG_4vQ_0vG_0 - b_0 (a_2^+ - a_2^-) \|V\|_{L^1}^{-1} G_0vQ_0v G_4 vQ_6 \mathcal{D}_{6,6} v G_2^0 + \Theta.  
	$$  
	Crucially, since $Q_0$ is the projection onto $v$, by the cancellation property $Q_6(x^\gamma v)= 0$ for $|\gamma| \le 3$, we have the identities: 
	$$G_2^0 v \mathcal{D}_{6,6} Q_6 vG_4vQ_0vG_0 = \|V\|_{L^1}G_2^0 v \mathcal{D}_{6,6} Q_6 vG_4,\quad G_0vQ_0v G_4 vQ_6 \mathcal{D}_{6,6} v G_2^0 = \|V\|_{L^1}G_4 vQ_6 \mathcal{D}_{6,6} v G_2^0.$$  
	Thus we obtain the expression \eqref{mathcal B} for $\widetilde{\mathcal{B}}.$ The argument for the sine evolution is identical by substituting $I_c(t)$ with $I_s(t)$, where $$I_s(t) = t^{-1} \int_0^\infty \sin(t\lambda^2) \lambda^{-1} \widetilde{\chi}_1(\lambda) \, d\lambda = \frac{\pi}{4|t|} + O(|t|^{-2}),$$ (see \eqref{eq:integral_sin}). This completes the proof.  
\end{proof}  


\begin{proof}[\textbf{Proof of Theorem \ref{main_theorem_low_32} (2)}]  We have established  the unweighted estimate ${O}(\langle t \rangle^{-1})$ by utilizing \eqref{Q4=Q5}.
	Regarding the weighted estimates, Proposition \ref{prop_part2} establishes that the time decay is  dictated by whether $\widetilde{\mathcal{B}}$ vanishes. If $\widetilde{\mathcal{B}} \neq 0$, the $I_s(t)$ term dominates, yielding $\sim |t|^{-1}$ decay for the sine evolution. If $\widetilde{\mathcal{B}} = 0$, the remaining terms yield the improved bound $\mathcal{O}(|t|^{-1} (\log |t|)^{-1})$. The theorem is thus established if we can prove the algebraic characterization of $\ker \widetilde{\mathcal{B}}$ detailed in Lemma \ref{lemma:non-zero_1}.  
\end{proof}  

\begin{lemma}\label{lemma:non-zero_1}  
	The operator $\widetilde{\mathcal{B}} \colon L^1_\omega \to L^\infty_{-\omega}$ is non-zero if and only if $S_5L^2 \neq \widetilde{S}_5L^2$.   
\end{lemma}  

\begin{proof}  
	Plugging  $a_0^\pm = \pm \frac{i}{8}$, $a_2^\pm = \pm \frac{i}{512}$, and $a_4^\pm = \pm \frac{i}{1179648}$ into the expression for $\widetilde{\mathcal{B}}$ yields:  
	\[  
	\widetilde{\mathcal{B}} = -b_0^2 \cdot \frac{i}{16384} \, G_2^0 v Q_6\mathcal{D}_{6,6} \mathbf{K} \mathcal{D}_{6,6} Q_6v G_2^0,  
	\]  
	where the operator $\mathbf{K} \colon Q_6 L^2 \to Q_6 L^2$ is defined by  
	\[  
	\mathbf{K} = \frac{1}{36}Q_6 v G_8 v Q_6 -  \|V\|_{L^1}^{-1} Q_6 v G_4 v Q_0 v G_4 v Q_6.  
	\]  
	We first prove $G_2^0 v Q_6 \colon Q_6 L^2 \to L^\infty_{-\omega}$ is injective. Suppose $G_2^0 v \psi = 0$ for some $\psi \in Q_6 L^2$. Applying $\Delta^2$ in the distributional sense yields $v \psi = 0$. Thus \(\psi=0\) a.e. on $\{x:v(x)\neq0\}.$
    By Proposition~\ref{characterizations-1}(iv) and the fact that $\psi \in Q_6 L^2 = S_5 L^2$, there exists $\phi \in L^2(\mathbb{R}^2)$ such that $\psi = U v \phi$, implying  \(\psi=0\) a.e.  on \(\{x: v(x)=0\}\). Hence $\psi  = 0$, confirming injectivity. By duality and $\dim(Q_6 L^2)< \infty$, $Q_6 v G_2^0$ is surjective onto $Q_6 L^2$. 

	Since $\mathcal{D}_{6,6}$ is invertible on $Q_6 L^2$, the operator $\widetilde{\mathcal{B}}$ is the zero operator if and only if $\mathbf{K} = 0$ on $Q_6 L^2$, i.e., $\ker\mathbf{K} = Q_6L^2 = S_5L^2$. Thus, to prove this lemma, it suffices to show $\ker\mathbf{K}=\widetilde{S}_5L^2$.  
	
Indeed,	for any $f \in Q_6 L^2$, a direct computation of the associated quadratic form yields:  
	\[  
	\langle \mathbf{K} f, f \rangle =  \frac{4}{9}  \sum_{i,j,k,l=1}^2 \bigl| \langle x_i x_j x_k x_l v, f \rangle \bigr|^2 + \frac{4}{3}  \sum_{i,j=1}^2 \bigl| \langle x_i x_j |x|^2 v, f \rangle \bigr|^2 - \frac{5}{6}  \bigl| \langle |x|^4 v, f \rangle \bigr|^2.  
	\]  
	By parameterizing the moments as $A = \langle x_1^4 v, f \rangle$, $B = \langle x_2^4 v, f \rangle$, $C = \langle x_1^2 x_2^2 v, f \rangle$, $D = \langle x_1^3 x_2 v, f \rangle$, and $E = \langle x_1 x_2^3 v, f \rangle$, we decouple the form into non-negative squares:  
	\begin{equation}\label{eq:negative_squares}  
		\langle \mathbf{K} f, f \rangle =  \frac{8}{9} |A-B|^2 + \frac{1}{18} |A+B-6C|^2 + \frac{32}{9} |D+E|^2 + \frac{8}{9} |D-E|^2\geq0.  
	\end{equation}  
	Therefore, $\mathbf{K}$ is positive semi-definite yielding $f \in \ker \mathbf{K} \iff \langle \mathbf{K}f, f \rangle = 0$. According to \eqref{eq:negative_squares}, this happens if and only if $A = B = 3C$ and $D = E = 0$. These equations coincide exactly with the definition of $\widetilde{S}_5L^2$ (see \eqref{S5wide}), yielding that $\ker\mathbf{K} = \widetilde{S}_5L^2$. The proof is complete.  
\end{proof}  


\section{A characterization of zero resonances}\label{sec:proof_characterization}
In this section, we give a precise characterization of the zero-energy resonances of $H = \Delta^2 + V$ on $\mathbb{R}^2$, from which Proposition~\ref{characterizations} follows.
For $\sigma \in \mathbb{R}$, recall the intersection space
\[
W_{-\sigma}(\mathbb{R}^2) = \bigcap_{s>\sigma} L^{2,-s}(\mathbb{R}^2),
\]
and the projection subspaces $S_jL^2$ associated with the orthogonal projections $S_j$ for $0 \le j \le 5$ given in Definition~\ref{defS}:
\begin{itemize}
	\item $S_0L^2 = \{ f \in L^2 \mid \langle v, f \rangle = 0\}$;
	\item $S_1L^2 = \{ f \in S_0L^2 \mid \langle x_j v, f \rangle = 0,\ j = 1,2 \}$;
	\item $S_2L^2 = \{ f \in S_1L^2 \mid S_1T_0f = 0 \}$;
	\item $S_3L^2 = \{ f \in S_2L^2 \mid S_0T_0f = 0 \}$;
	\item $S_4L^2 = \{ f \in S_3L^2 \mid \langle x_ix_j v, f \rangle = 0,\ i,j = 1,2,\ \text{and}\ PT_0f = 0 \}$;
	\item $S_5L^2 = \{ f \in S_4L^2 \mid \langle x_ix_jx_k v, f \rangle = 0,\ i,j,k = 1,2 \}$.
\end{itemize}
Here $v = \sqrt{|V|}$, $P = \|V\|_{L^1(\mathbb{R}^2)}^{-1} \langle \cdot, v\rangle v$, and $T_0 = U + b_0 vG_2^0v$.

By the definition of the projection subspaces $S_jL^2$ and the zero-energy types introduced in Definition~\ref{definition1}, Proposition~\ref{characterizations} follows directly from Proposition~\ref{characterizations-1} below. Therefore, establishing Proposition~\ref{characterizations-1} is the primary focus of this section.

\begin{proposition}\label{characterizations-1}
	Assume that $H = \Delta^2 + V$ and $|V(x)| \lesssim \langle x \rangle^{-\mu}$ with $\mu > 0$.
	\begin{itemize}
		\item[(i)] If $\mu > 11$, then
		$\psi \in S_{2} L^2$ if and only if $\psi = Uv\phi$ with $\phi \in W_{-2}(\mathbb{R}^2)$ satisfying $H\phi = 0$ and
		\begin{equation}\label{phi-phi0}
			\phi(x) = \phi_0(x) + \sum_{|\alpha|=1} C_\alpha x^\alpha \quad \text{for some } \phi_0 \in W_{-1}(\mathbb{R}^2).
		\end{equation}
		\item[(ii)] If $\mu > 14$, then
		$\psi \in S_{3} L^2$ if and only if $\psi = Uv\phi$ with $\phi \in W_{-1}(\mathbb{R}^2)$ satisfying $H\phi = 0$.
		\item[(iii)] If $\mu > 18$, then
		$\psi \in S_{4} L^2$ if and only if $\psi = Uv\phi$ with $\phi \in W_{0}(\mathbb{R}^2)$ satisfying $H\phi = 0$.
		\item[(iv)] If $\mu > 18$, then
		$\psi \in S_{5} L^2$ if and only if $\psi = Uv\phi$ with $\phi \in L^2(\mathbb{R}^2)$ satisfying $H\phi = 0$.
	\end{itemize}
\end{proposition}
From Proposition \ref{characterizations-1}, we also obtain the following  spectral characterizations:
\begin{itemize}
\item \textbf{$p$-wave resonance:} $S_3 \neq S_4$ iff there exists a nonzero solution $\phi \in W_{-1}(\mathbb{R}^2) \setminus W_{0}(\mathbb{R}^2)$ to $H\phi = 0$. In other words, $Q_4 \neq 0$ is  equivalent to  $H$ admitting a $p$-wave resonance.
	
	\item \textbf{$d$-wave resonance:} $S_4 \neq S_5$ iff there exists a nonzero solution $\phi \in W_{0}(\mathbb{R}^2) \setminus L^2(\mathbb{R}^2)$ to $H\phi = 0$. In other words, $Q_5 \neq 0$ is  equivalent to  $H$ admitting a $d$-wave resonance.
\end{itemize}


\subsection{Preliminary estimates for $G_2^0$}
Prior to the proof of Proposition \ref{characterizations-1}, we establish two technical lemmas concerning  the fundamental solution $G_2^0(x) = |x|^2 \log |x|$. These are crucial for deriving the decay rates and cancellation properties.

\begin{lemma} \label{estimate_G_2^0}
	Let $j \in \mathbb{N}^+$ and $f \in L^{2,\sigma}(\mathbb{R}^2)$ with $\sigma > \max\{j, 2\} + 1$. Suppose $f$ satisfies 
\begin{equation}\label{eq:cancellation}
		\int_{\mathbb{R}^2} y^\alpha f(y) \, dy = 0 \quad \text{for all multi-indices } \alpha \text{ with } |\alpha| \le j-1.
	\end{equation}
	Then 
	\[
	\left| (G_2^0f)(x) \right| \lesssim \|f\|_{L^{2,\sigma}} \, \langle x \rangle^{2-j} (1 + \log \langle x \rangle), \quad \forall x \in \mathbb{R}^2.
	\]
\end{lemma}

\begin{proof}
	For $|x| \le 2$, fix $0<\epsilon< \min(\sigma-3, 2)$. Then $|G_2^0(x-y)|\lesssim \langle y \rangle^{2+\epsilon}$. By the Cauchy--Schwarz inequality, $|G_2^0 f(x)|\lesssim \|f\|_{L^{2,\sigma}}$, which verifies the bound locally.
	
	For $|x| > 2$, note that $\langle x \rangle \sim |x|$. We decompose the integration domain into $A_1(x) = \{ y : |y| > |x|/2 \}$ and $A_2(x) = \{ y : |y| \le |x|/2 \}$, writing $G_2^0f(x) =\mathcal {I}_1(x) + \mathcal {I}_2(x)$.
	
	\underline{Estimate of $\mathcal{I}_1(x)$:} For $y \in A_1(x)$, $|G_2^0(x-y)| \lesssim |y|^2 (1 + \log |y|)$. By Cauchy--Schwarz and switching to polar coordinates $r = |x|\rho$:
	\begin{align*}
		|\mathcal{I}_1(x)| &\lesssim \|f\|_{L^{2,\sigma}} \Bigl( \int_{|y| > |x|/2} |y|^4 (1 + \log |y|)^2 \langle y \rangle^{-2\sigma} \, dy \Bigr)^{1/2} \\
		&\lesssim \|f\|_{L^{2,\sigma}} \, |x|^{3-\sigma} \Bigl( \int_{1/2}^\infty \rho^{5-2\sigma} (1 + \log(|x|\rho))^2 \, d\rho \Bigr)^{1/2}.
	\end{align*}
	Using $1 + \log(ab) \lesssim (1 + \log a)(1 + \log b)$ for $a,b>1$ and $\sigma > \max\{j, 2\} + 1$,  we obtain
 	\[
	|\mathcal{I}_1(x)| \lesssim \|f\|_{L^{2,\sigma}} \langle x \rangle^{3-\sigma} \Big(1 + \log \frac{|x|}{2}\Big) \Big( \int_{1/2}^\infty \rho^{5-2\sigma} (1 + \log (2\rho))^2 d\rho \Big)^{1/2} \lesssim \|f\|_{L^{2,\sigma}}  \langle x\rangle^{2-j} (1 + \log \langle x\rangle).
	\]
 

 
\underline{Estimate of $\mathcal{I}_2(x)$:} Taylor expanding $G_2^0(x-y)$ about $y=0$ to order $j-1$, and utilizing the cancellation \eqref{eq:cancellation} and the relation $A_2(x) = \mathbb{R}^2 \setminus A_1(x)$, we obtain
	\begin{align*}
	\mathcal{I}_2(x) 
    &=-\sum_{|\alpha|\le j-1} \frac{(-1)^{|\alpha|}}{\alpha!}\partial_x^\alpha G_2^0(x)\int_{A_1(x)} y^\alpha f(y)dy + \mathcal{R}(x),
	\end{align*}
	where the remainder $\mathcal{R}(x)$ is given by
	\[
	\mathcal{R}(x) = \sum_{|\alpha|=j} \frac{(-1)^{|\alpha|}}{\alpha!} \int_{A_2(x)}  y^\alpha f(y) \Bigl( \int_0^1 j(1-\theta)^{j-1} (\partial_x^\alpha G_2^0)(x-\theta y) d\theta \Bigr) dy.
	\]
	Since $| \int_{A_1(x)} y^\alpha f(y) dy| \lesssim |x|^{|\alpha|+1-\sigma} \|f\|_{L^{2,\sigma}}$ and $|\partial_x^\alpha G_2^0(x)| \lesssim |x|^{2-|\alpha|} \log|x|$,
then the first part of $\mathcal{I}_2(x)$ is bounded by
$ \|f\|_{L^{2,\sigma}} |x|^{3-\sigma} \log|x|$.
	For the remainder $\mathcal{R}(x)$, since $|x - \theta y| \sim |x|$ for $y\in A_2(x)$ and $\theta\in[0,1]$, we have $|\partial_x^\alpha G_2^0(x-\theta y)| \lesssim |x|^{2-j}(1+\log|x|)$ for $|\alpha|=j$. The condition $\sigma > j+1$ ensures integrability, yielding $|\mathcal{R}(x)| \lesssim \|f\|_{L^{2,\sigma}} \langle x \rangle^{2-j}(1+\log\langle x \rangle)$.
\end{proof}
\begin{lemma} \label{S_5canc}
	Let $f \in L^{2,\sigma}(\mathbb{R}^2)$ with $\sigma > 5$. If $G_2^0 f \in L^2(\mathbb{R}^2)$, then $f$ must satisfy 
 \[
\int_{\mathbb{R}^2} y^\alpha f(y) \, dy = 0, \quad \text{for every multi-index}  \   \alpha \  \text{with}\  |\alpha| \leq 3.
\]
\end{lemma}

\begin{proof}
	For $|x|>2$, following the argument in Lemma \ref{estimate_G_2^0} with $j=4$ and $\sigma > 5$, we can obtain
	\begin{align*}
		(G_2^0 f)(x) = \sum_{|\alpha| \leq 3} \frac{(-1)^{|\alpha|}}{\alpha!} M_\alpha \partial_x^\alpha G_2^0(x) +  O(|x|^{-2} \log |x|) =: \Psi_0(x) + O(|x|^{-2} \log |x|),
	\end{align*}
where $M_\alpha= \int_{\mathbb{R}^2} y^\alpha f(y) \, dy$.
	Since both $G_2^0 f$ and the remainder are in $L^2(\{ |x| > 2 \})$, we must have $\Psi_0(x) \in L^2(\{ |x| > 2 \})$. Computing $\partial_x^\alpha G_2^0(x)$ explicitly gives:
	\begin{align*}
		\Psi_0(x) &= |x|^2\log|x|M_0+ (-2\log|x|-1)\sum_{|\alpha|=1}x^{\alpha}M_\alpha \\
		&\quad + \Big[(\frac{1}{2}+\log|x|)\int_{\mathbb{R}^2}|y|^2 f(y)dy + \sum_{|\alpha|=2}\frac{x^{\alpha}}{|x|^2}M_\alpha\Big] \\
		&\quad + \Big[-\sum_{i=1}^{2}\frac{x_i}{|x|^2}\int_{\mathbb{R}^2}y_i|y|^2 f(y)dy + \frac{2}{3}\sum_{|\alpha|=3}\frac{x^{\alpha}}{|x|^4}M_\alpha\Big].
	\end{align*}
Note that for $|x|>2,$ we have $|x|^2 \log |x| \in W_{-3} \setminus W_{-2}$ and
\begin{align*}
	(2 \log |x| + 1)x^\alpha &\in W_{-2}\setminus W_{-1} \ \ \text{for}\ |\alpha|=1, \\
	\frac{1}{2} + \log |x|,\  \frac{x^\alpha}{|x|^2}  &\in W_{-1} \setminus W_0 \ \ \text{for}\ |\alpha|=2, \\
	\frac{x_i}{|x|^2}, \  \frac{x^\alpha}{|x|^4}  &\in W_0 \setminus L^2 \ \ \text{for}\ |\alpha|=3.
\end{align*}
	Hence, the only possibility for $\Psi_0(x) \in L^2(\{ |x| > 2 \})$ is that all the constant coefficients multiplying these profiles must be exactly zero, i.e., $M_\alpha = 0$ for all  $|\alpha| \leq 3$. This yields the desired cancellations.  
\end{proof}

\subsection{Proof of Proposition \ref{characterizations-1}}
We present a unified proof of the equivalence. Specifically, we will establish that $\psi \in S_{k+1} L^2$ if and only if $H\phi = 0$, where  $\psi = Uv\phi \in W_{k-3}(\mathbb{R}^2)$ for $k \in \{1, 2, 3\}$, and  $\psi = Uv\phi \in L^2(\mathbb{R}^2)$ for $k = 4$. Furthermore, for the  case $k=1$, we will show that this equivalence additionally requires $\phi$ to satisfy the structural form \eqref{phi-phi0}.

\textbf{Proof of ``$\bm{\Longrightarrow}$'' (From Subspace to Resonance):} 
Assume $\psi \in S_{{k}+1}L^2$. 
For ${k}=1,2$, Definition \ref{defS} dictates $S_{2-{k}}T_0\psi = 0$. Thus $T_0\psi = (I - S_{2-{k}})T_0\psi = v \sum_{|\alpha| \le 2-{k}} C_{\alpha} x^\alpha$. Then
\begin{align}\label{k=1,2}
	U\psi = -b_0 v G_2^0 v\psi + T_0\psi = -b_0 v G_2^0 v\psi + v\sum_{|\alpha| \le 2-{k}} C_{\alpha }x^\alpha.
\end{align}
For ${k}=3,4$, the condition $PT_0\psi = 0$ yields $T_0\psi =(P+S_0) T_0\psi=0$, implying
\begin{align}\label{k=3,4}
	U\psi = -b_0 v G_2^0 v\psi.
\end{align}
By the cancellation properties of $S_{{k}+1}$ (Definition \ref{defS}), $f = v\psi$ satisfies \eqref{eq:cancellation} with $j=2$ (for ${k}=1$) and $j={k}$ (for ${k}=2,3,4$). Applying Lemma \ref{estimate_G_2^0}, we obtain:
\begin{equation}\label{estimate_phi_G}
	\left| G_2^0 v\psi(x) \right| \lesssim
	\begin{cases}
		1 + \log \langle x \rangle, & {k} = 1,2, \\
		\langle x \rangle^{2 - {k}} \big( 1 + \log \langle x \rangle \big), & {k} = 3,4.
	\end{cases}
\end{equation}
This confirms $G_2^0 v\psi \in W_{-1}$ (for ${k}=1$) and $G_2^0 v\psi \in W_{{k}-3}$ (for ${k}=2,3,4$).

We now define the function $\phi(x)$ 
\begin{equation}\label{expression_psi}
	\phi(x) :=
	\begin{cases}
		-b_0 G_2^0 v\psi(x) + \sum_{|\alpha| \le 2 - {k}} C_\alpha x^\alpha, & {k} = 1,2, \\
		-b_0 G_2^0 v\psi(x), & {k} = 3,4.
	\end{cases}
\end{equation}
Combine this with  \eqref{k=1,2} and \eqref{k=3,4},  then $\psi = Uv\phi$. Moreover, $\phi$ falls exactly into the required spaces $W_{{k}-3}$ (or $L^2$ for ${k}=4$), with the special structure for ${k}=1$ automatically satisfied. Finally, since $\Delta^2 G_2^0 = b_0^{-1} \delta_0$, we verify the equation:
\[
(\Delta^2 + V)\phi = -b_0 \Delta^2 G_2^0 v\psi + V\phi = -v( Uv\phi) + V\phi = -V\phi + V\phi= 0,
\]
which completes the forward implication.


\vspace{0.2cm}
\textbf{Proof of ``$\bm{\Longleftarrow}$'' (From Resonance to Subspace):} 
Conversely, suppose there exists $\phi$ satisfying $(\Delta^2 + V)\phi = 0$ in the respective space for index ${k}$. Define $\psi = U v \phi$. We proceed in three steps to show $\psi \in S_{{k}+1}L^2$.


\textbf{Step 1: Establish the cancellation conditions.} 
We first prove $\langle x^{\alpha}v, \psi\rangle=0$ for $|\alpha| \leq 1$ when ${k}=1$, for $|\alpha| \leq {k}-1$ when ${k}=2,3$, and for $|\alpha| \leq 2$ when ${k}=4$.
Let $\eta \in C_c^\infty(\mathbb{R}^2)$ be a cutoff function satisfying $\eta(x) = 1$ for $|x| \le 1$ and $\eta(x) = 0$ for $|x| \ge 2$.  Using $(\Delta^2 + V)\phi = 0$ and integrating by parts against $\eta(\rho x)$ with $\rho \in (0,1)$, we obtain
\begin{align}
	\int_{\mathbb{R}^2} x^\alpha v(x) \psi(x) \eta(\rho x)  d x
	&= -\int_{\mathbb{R}^2} x^\alpha \eta(\rho x) \Delta^2 \phi(x)  d x \nonumber\\\label{psi-phi}
	&= \sum_{\substack{\alpha_i \geq \beta_i \text{ for } i=1,2 \\ |\beta|+|\gamma|=4, \beta, \gamma \in \mathbb{N}_0^2}} C_{\beta, \gamma}  \rho^{|\gamma|} \int_{\mathbb{R}^2} x^{\alpha-\beta} \phi(x) (\partial^\gamma \eta)(\rho x)  d x.
\end{align}
Then applying the Cauchy-Schwarz inequality and noting that the support of $(\partial^\gamma \eta)(y)$ is contained in $\{ |y| \le 2\}$ for all $\gamma\in \mathbb{N}_0^2$, we deduce that for $|\alpha|-|\beta|+s \ge 0$ (ensured by taking $s \ge 0$),
\begin{align*}
	\left| \int_{\mathbb{R}^2} x^{\alpha-\beta} \phi(x) (\partial^\gamma \eta)(\rho x) \, dx \right| 
	&\le \left\| \langle x\rangle^{|\alpha|-|\beta|+s} (\partial^\gamma \eta)(\rho x) \right\|_{L^2} \left\| \langle x\rangle^{-s} \phi \right\|_{L^2} \lesssim \rho^{-|\alpha|+|\beta|-s-1} \left\| \langle x\rangle^{-s} \phi \right\|_{L^2}.
\end{align*}
Summing up, we obtain the integral \eqref{psi-phi} is  bounded by $\rho^{3-|\alpha|-s} \| \langle x\rangle^{-s} \phi \|_{L^2}$. 
Since $\phi \in W_{{k}-3}$, we can select $s = 3 - {k} + \varepsilon$ with $0 < \varepsilon \ll 1$ such that $\phi \in L^{2,-s}$ and $3 - |\alpha| - s > 0$ for all strictly required $|\alpha|$. Taking the limit $\rho \to 0$, the integral vanishes, giving the cancellations for ${k}=2,3$. (For ${k}=1$, substituting $\phi$ with $\phi_0 \in W_{-1}$ and taking $s=1+\varepsilon$ yields the same).


For ${k}=4$, since $\phi \in L^2$, then $s=0$. The exponent $3 - |\alpha| > 0$ dictates that the integral vanishes for $|\alpha| \le 2$. The missing cancellation for $|\alpha|=3$ will be deduced via a bootstrap argument below.


\textbf{Step 2:  Deduce the structural forms.}
Using the cancellation properties established in Step~1 and arguing as in the derivation of \eqref{estimate_phi_G} for ${k}=1,2,3$, we conclude that $G_2^0 v\psi(x) \in W_{-1}(\mathbb{R}^2)$ for ${k}=1$, and $G_2^0 v\psi(x) \in W_{{k}-3}(\mathbb{R}^2)$ for ${k}=2,3$. Similarly, for ${k}=4$, the vanishing of the moments $\langle x^{\alpha}v, \psi\rangle$ for $|\alpha| \le 2$, already obtained in Step~1, at least yields $G_2^0 v\psi(x) \in W_{0}(\mathbb{R}^2)$.
Thus by  $\phi \in W_{{k}-3}(\mathbb{R}^2)$ for ${k}=1,2,3$  and $\phi\in L^2(\mathbb{R}^2),$
we have
\begin{equation}\label{estimate}
	\begin{split}
		\phi+b_0 G_2^0 v \psi \in
		W_{{k}-3}(\mathbb{R}^2)\ \text{for} \ {k}=1,2,3, \ \text{and at least } 	\phi+b_0 G_2^0 v \psi \in
		W_{0}(\mathbb{R}^2) \ \text{for} \ {k}=4.
	\end{split}
\end{equation}
Note that
$
	\Delta^2\left(\phi+b_0 G_2^0 v \psi\right) =\Delta^2 \phi+ v \psi=\left(\Delta^2+V\right) \phi=0.
$
It follows that Fourier transform of the function $\phi+b_0 G_2^0 v \psi $ is supported
at the origin, which is the sum of finite derivatives of Dirac distribution $\delta$ (see, e.g., \cite{GG21b}).
This immediately implies that $\phi+b_0 G_2^0 v \psi $ is a polynomial on $\mathbb{R}^2$.
Combining this with  \eqref{estimate}, we conclude that $\phi$ indeed admits the representation stated in \eqref{expression_psi} for ${k}=1,2,3,4$.


\textbf{Step 3: Finalize ${k}=4$ cancellations and projection conditions.}
Note that $v\psi=V\phi$ and
\begin{align*}
	|(V\phi)(x)|=|b_0V(x) G_2^0 v\psi(x)|\lesssim\langle x \rangle^{-\mu}\langle x \rangle^{-1} \big( 1 + \log \langle x \rangle \big)\lesssim\langle x \rangle^{-\mu}.
\end{align*}
Together with $\phi=-b_0G_2^0 v \psi \in L^2$ for ${k} = 4$, Lemma~\ref{S_5canc} with $f = v \psi$  grants the remaining cancellation condition for $|\alpha|=3$.
Finally, we verify $S_{2-{k}}T_0\psi$ for ${k}=1,2$ and $PT_0\psi=0$ for ${k}=3,4.$ Indeed,
for ${k}=1,2$, applying $S_{2-{k}}$ to $T_0\psi$ yields:
\begin{align*}
	S_{2-{k}}T_0\psi &= S_{2-k}U\psi + b_0 S_{2-{k}}v G_2^0 v\psi = S_{2-{k}}v\phi + b_0 S_{2-{k}}v G_2^0 v\psi \\
	&= S_{2-{k}}\Big( -b_0 vG_2^0 v\psi + v\sum_{|\alpha| \le 2 - {k}} C_\alpha x^\alpha \Big) + b_0 S_{2-{k}}v G_2^0 v\psi = 0,
\end{align*}
due to the intrinsic orthogonality of $S_{2-{k}}$. For ${k}=3,4$, since  $\phi = -b_0 G_2^0 v \psi$, then
\begin{align*}
	PT_0\psi(x)&=PU\psi(x)+b_0PvG_2^0 v\psi(x)\\
	&=Pv\phi+b_0PvG_2^0 v\psi(x)=-b_0PvG_2^0 v\psi(x)+b_0PvG_2^0 v\psi(x)=0. 
\end{align*}

This verifies all required conditions for $S_{{k}+1}L^2$, concluding the proof.
\hfill $\square$

\section{Asymptotic expansions of $(M^{\pm}(\lambda))^{-1}$ at zero energy}
\label{subsec:proof_inverse}

We begin by presenting two abstract algebraic lemmas that are essential for inverting the operator matrices in the proof of Theorem \ref{thm:M_inverse}.

\begin{lemma}[{\cite[Lemma 3.12]{JK}}]\label{lemma_1}
	Let $\mathfrak{X}, \mathfrak{Y}, X, Y$ be vector spaces. Let $\mathbf{A} : \mathfrak{X} \rightarrow \mathfrak{Y}, \mathbf{B} : X \rightarrow \mathfrak{X}, \mathbf{C} : \mathfrak{Y} \rightarrow Y$ be linear operators. Define $\mathbf{D}=\mathbf{C}\mathbf{A}\mathbf{B}$. If $\mathbf{D}^{-1}$ exists, then
	$$
\mathbf{A}^{-1}=\mathbf{B}\mathbf{D}^{-1}\mathbf{C}
	$$
	provided $\mathbf{B}$ is surjective and $\mathbf{C}$ is injective.
\end{lemma}

\begin{lemma}[{\cite[Lemma 2.3]{JN}}]\label{lemma_2}
	Let $A$ be an operator matrix on $\mathcal{H}=\mathcal{H}_1 \medoplus \mathcal{H}_2$:
	$$
	A=\left(\begin{array}{ll}
		a_{11} & a_{12} \\
		a_{21} & a_{22}
	\end{array}\right), \quad a_{i j}: \mathcal{H}_j \rightarrow \mathcal{H}_i, 1 \leq i, j \leq 2,
	$$
	where $a_{11}$ and $a_{22}$ are closed, and $a_{12}$, $a_{21}$ are bounded. Suppose $a_{11}$ has a bounded inverse. Define 
	$$
	\boldsymbol{d} := a_{22}-a_{21} a_{11}^{-1} a_{12}.
	$$
	Then the following statements hold:
	\begin{enumerate}
		\item[(i)] $A$ has a bounded inverse if and only if $\boldsymbol{d}$
		has a bounded inverse. Furthermore, if $\boldsymbol{d}$ has a bounded inverse, we have
		\begin{equation}\label{A-m}
		A^{-1}=\left(\begin{array}{cc}
			a_{11}^{-1} a_{12} \boldsymbol{d}^{-1} a_{21} a_{11}^{-1}+a_{11}^{-1} & -a_{11}^{-1} a_{12} \boldsymbol{d}^{-1}  \\
			-\boldsymbol{d}^{-1} a_{21} a_{11}^{-1}& \boldsymbol{d}^{-1}
		\end{array}\right) .
		\end{equation}
		\item[(ii)] If $A$ is self-adjoint, then $A > 0$ (strictly positive) if and only if $a_{11} > 0$ and $\boldsymbol{d} > 0$.
	\end{enumerate}
\end{lemma}

Equipped with these algebraic tools, we now proceed with the proof of Theorem \ref{thm:M_inverse}.

\begin{proof}[Proof of Theorem \ref{thm:M_inverse}]
We first recall from Proposition~\ref{characterizations Q} that if zero is a $\mathbf{k}$-th kind resonance of $H$, the space $L^2$ admits the orthogonal decomposition $L^2 \cong \bigoplus_{\alpha=0}^{\mathbf{k}+2} Q_\alpha L^2$.
	 For  $0<\lambda\ll1$, we define an isomorphism $\mathbf{B}_\lambda: \bigoplus_{\alpha=0}^{\mathbf{k}+2} Q_\alpha L^2 \to L^2$ by
	\[
\mathbf{B}_\lambda(f_\alpha)_{\alpha=0}^{\mathbf{k}+2} := \sum_{\alpha=0}^{\mathbf{k}+2} \lambda^{-k_\alpha} \sigma_\alpha(\lambda) Q_\alpha f_\alpha,
	\]
	where the scaling factors $\sigma_\alpha(\lambda)$ are given by	\begin{equation}\label{tab:sigma_alpha}
		\begin{array}{c|ccccccc}
			\alpha & 0 & 1 & 2 & 3 & 4 & 5 & 6 \\
			\hline
			\sigma_\alpha(\lambda) & 1 & (-\log \lambda)^{-1/2} & 1 & (-\log \lambda)^{1/2} & 1 & (-\log \lambda)^{-1/2} & 1
		\end{array}
	\end{equation}
	and the exponents $k_\alpha$ are defined by
    \begin{equation*}
\begin{tabular}{c|ccccccc}
    \( \alpha \)    & 0 & 1 & 2 & 3 & 4 & 5 & 6 \\
    \hline
    \( k_\alpha \)  & 0 & 1 & 1 & 1 & 2 & 3 & 3
\end{tabular}
\end{equation*}
	Since $\mathbf{B}_\lambda$ is bijective, Lemma~\ref{lemma_1} dictates that if $\mathbf{B}_\lambda^*M^{\pm}(\lambda) \mathbf{B}_\lambda$ is invertible on $\bigoplus_{\alpha=0}^{\mathbf{k}+2} Q_\alpha L^2$ for sufficiently small $\lambda > 0$, then
	\begin{align}\label{M-nin}
	(M^{\pm}(\lambda))^{-1} = \mathbf{B}_\lambda (\mathbf{B}_\lambda^* M^{\pm}(\lambda) \mathbf{B}_\lambda)^{-1} \mathbf{B}_\lambda^*,
	\end{align}
	where $\mathbf{B}_\lambda^*$ denotes the adjoint of
\(\mathbf B_\lambda\). 
	Consequently, inverting $M^\pm(\lambda)$ on $L^2$  reduces to studying the invertibility of $\mathbf{B}_\lambda^* M^\pm(\lambda) \mathbf{B}_\lambda$.
	
	When zero is a $\mathbf{k}$-th kind resonance for some integer $0 \le \mathbf{k} \le 4$, combining the expansions \eqref{Mpm-1}--\eqref{Mpm-3} of $M^\pm(\lambda)$ with the cancellation properties of $Q_\alpha$ (Remark~\ref{cancellationQ}) yields the decomposition:
	\[
	\lambda^{-k_\alpha-k_\beta}\sigma_\alpha(\lambda)\sigma_\beta(\lambda)Q_\alpha M^\pm(\lambda)Q_\beta = \lambda^{-2}\left(E_{\alpha, \beta}^{\pm} + r_{\alpha, \beta}^{\pm}(\lambda)\right),
	\]
	where $E_{\alpha, \beta}^{\pm}$ and $r_{\alpha, \beta}^{\pm}(\lambda)$ are given by following \eqref{D_k} and \eqref{partial}-\eqref{partial2}, respectively.
	Consequently, on the direct sum space $\bigoplus_{\alpha=0}^{\mathbf{k}+2} Q_\alpha L^2$, the operator matrix $\mathbf{B}_\lambda^* M^{\pm}(\lambda) \mathbf{B}_\lambda$ admits the expansion
	\[
	\mathbf{B}_\lambda^* M^{\pm}(\lambda) \mathbf{B}_\lambda= \lambda^{-2} \left( \mathbf{D}_{\mathbf{k}}^{\pm} + \mathbf{S}_{\mathbf{k}}^{\pm}(\lambda) \right),
	\]
	where $\mathbf{D}_{\mathbf{k}}^{\pm} = \bigl( E_{\alpha, \beta}^{\pm} \bigr)_{\alpha, \beta=0}^{\mathbf{k}+2}$ is the leading matrix and $\mathbf{S}_{\mathbf{k}}^{\pm}(\lambda) = \bigl( r_{\alpha, \beta}^{\pm}(\lambda) \bigr)_{\alpha, \beta=0}^{\mathbf{k}+2}$ is the remainder. 
	
	We specify the exact structures of these matrices below:
	\begin{itemize}
		\item \textbf{Asymptotic expansion of the remainder.} The entries $r_{\alpha,\beta}^{\pm}(\lambda)$ exhibit the following asymptotic behaviors:
		\begin{equation}\label{partial}
			\begin{split}
				r_{\alpha,\beta}^{\pm}(\lambda) &= (-\log \lambda)^{-1/2} Q_\alpha T_0 Q_\beta + \lambda^2 (-\log \lambda)^{-1/2} \Lambda^\pm(\lambda),
				\quad \text{for } (\alpha,\beta) \in \{(1,2), (2,1)\}; \\
                r_{\alpha,\beta}^{\pm}(\lambda) &= b_1 (-\log \lambda)^{-1/2} Q_\alpha v G_6^0 v Q_\beta + \lambda^2 (-\log \lambda)^{-1/2} \Lambda^\pm(\lambda),
				\quad \text{for } (\alpha,\beta) \in \{(5,6), (6,5)\}; \\
                r_{1,1}^{\pm}(\lambda) &= (-\log \lambda)^{-1} \left( a_1^{\pm} Q_1 v G_2 v Q_1 + Q_1 T_0 Q_1 \right) + \lambda^2 (-\log \lambda)^{-1} \Lambda^\pm(\lambda); \\
				r_{\alpha,\beta}^{\pm}(\lambda) &= (-\log \lambda)^{-1} a_2^\pm Q_\alpha v G_4 v Q_\beta + \lambda^2 \Lambda^\pm(\lambda),
				\quad \text{for } (\alpha,\beta) \in \{(1,5), (5,1)\};\\
                r_{5,5}^{\pm}(\lambda) &= (-\log \lambda)^{-1} \left( a_3^\pm Q_5 v G_6 v Q_5 + b_1 Q_5 v G_6^0 v Q_5 \right) + \lambda^2 (-\log \lambda)^{-1} \Lambda^\pm(\lambda); \\
        r_{\alpha,\beta}^\pm(\lambda) &= b_1 \lambda \log \lambda Q_\alpha v G_6 v Q_\beta + \lambda \Lambda^\pm(\lambda),
				\quad \text{for } (\alpha,\beta) \in \{(4,6), (6,4)\};\\
				r_{\alpha,\beta}^\pm(\lambda) &=  \lambda a_2^\pm Q_\alpha v G_4 v Q_\beta + \lambda^3 (-\log \lambda) \Lambda^\pm(\lambda),
				\quad \text{for } (\alpha,\beta) \in \{(0,6), (6,0)\};  
			\end{split}
		\end{equation}
		and for all other index pairs $(\alpha,\beta)$,
		\begin{align}\label{partial2}
			r_{\alpha,\beta}^{\pm}(\lambda) = \lambda (-\log \lambda)^{3/2} \Lambda^\pm(\lambda).
		\end{align}
			Here and throughout the subsequent proofs, $\Lambda^\pm(\lambda)$ denotes a generic operator in $\mathbb{B}(L^2)$ (possibly varying at each occurrence) that satisfies the estimates: 
	\[  
	\left\|\partial_\lambda^\ell \Lambda^\pm(\lambda)\right\|_{\mathbb{B}(L^2)} \lesssim \lambda^{-\ell}, \quad \ell = 0, 1, 2.  
	\]  
	 
		\item \textbf{Explicit form for $\mathbf{k}=4$.} The leading matrix $\mathbf{D}_{4}^{\pm} = (E_{\alpha, \beta}^{\pm})_{\alpha, \beta=0}^{6}$ admits the structure:
		\begin{equation}\label{D_k}
			\mathbf{D}_{4}^{\pm} =
			\begin{pmatrix}
				a_0^\pm Q_0vG_0vQ_0 & 0 & E_{0,4}^{\pm} & 0 \\
				0 & \mathbf{D}_{1,1} & 0 & 0 \\
				E_{4,0}^{\pm} & 0 & a_2^\pm Q_4vG_4vQ_4 & 0 \\
				0 & 0 & 0 & \mathbf{D}_{2,2}
			\end{pmatrix},
		\end{equation}
		where $E_{\alpha, \beta}^{\pm} = g_0^{\pm}(\lambda) Q_\alpha vG_2 vQ_\beta + Q_\alpha T_0 Q_\beta$ for $(\alpha, \beta) \in \{(0,4), (4,0)\}$, and 
		\begin{equation}\label{D}
			\mathbf{D}_{1,1} = \begin{pmatrix}
				-b_0 Q_1vG_2vQ_1 & 0 & Q_1T_0Q_3 \\
				0 & Q_2T_0Q_2 & 0 \\
				Q_3T_0Q_1 & 0 & 0
			\end{pmatrix}, \
			\mathbf{D}_{2,2} = \begin{pmatrix}
				-b_1 Q_5vG_6vQ_5 & 0 \\
				0 & b_1 Q_6vG_6^0vQ_6
			\end{pmatrix}.
		\end{equation}
		We abbreviate the entries of $\mathbf{D}_{1,1}$ and $\mathbf{D}_{2,2}$ as $E_{\alpha,\beta}$ (dropping the $\pm$), since they are independent of the sign $\pm$.
		
		\item \textbf{General case for $0 \leq \mathbf{k} \leq 3$.}
		For general $\mathbf{k}$, the operator matrix $\mathbf{D}_{\mathbf{k}}^{\pm}$
		can be formally understood as  the restriction of $\mathbf{D}_{4}^{\pm}$ onto $\bigoplus_{\alpha \le \mathbf{k}+2} Q_\alpha L^2$.
	\end{itemize}
	To streamline the presentation, we will detail the subsequent inversion explicitly for the maximal resonance case $\mathbf{k}=4$. By the algebraic nature of our construction, the explicit formula for $(M^{\pm}(\lambda))^{-1}$ when $\mathbf{k}=4$  recovers the general $\mathbf{k}$ case  upon setting $Q_\alpha = 0$ for all $\alpha > \mathbf{k}+2$.
    
	We now establish the theorem for the  case $\mathbf{k}=4$, i.e., zero is an eigenvalue of $H.$  We organize the proof into three steps.

{\textbf{Step 1. Exact Invertibility of the Leading Matrix $\mathbf{D}_{4}^{\pm}$.}}

	By the structure of $\mathbf{D}_{4}^{\pm}$, its invertibility reduces to establishing the invertibility of the following three independent components:
	\begin{enumerate}
		\item[(i)] $\mathbf{D}_{1,1}$ on $Q_1 L^2 \medoplus Q_2 L^2 \medoplus Q_3 L^2$;
		\item[(ii)] $\mathbf{D}_{2,2}$ on $Q_5 L^2 \medoplus Q_6 L^2$;
		\item[(iii)] $\mathbf{A}^{\pm} = \begin{pmatrix} E_{0,0}^{\pm} & E_{0,4}^{\pm} \\ E_{4,0}^{\pm} & E_{4,4}^{\pm} \end{pmatrix}$ on $Q_0 L^2 \medoplus Q_4 L^2$.
	\end{enumerate}
Note that because the relevant subspaces $Q_j L^2$ are finite-dimensional for $j \neq  2$, injectivity is equivalent to invertibility, a fact we utilize implicitly below.
	
		\textbf{(i) Invertibility of $\mathbf{D}_{1,1}$:}
	First, $E_{2,2} = Q_2 T_0 Q_2$ is invertible on $Q_2 L^2$ by the definition of $Q_2$ and Riesz-Fredholm theory. By Lemma \ref{lemma_2}, it reduces to showing that the operator $Q_{3} T_0 Q_{1} (Q_{1}v G_{2} v Q_{1})^{-1} Q_{1}T_0 Q_{3}$ is invertible on $Q_{3} L^2$. For any nonzero $f \in Q_1 L^2$, we have
	$$
	\langle Q_1 v G_2 v Q_1 f, f \rangle = -2 \sum_{j=1}^2 |\langle x_j v, f \rangle|^2 < 0,
	$$
	implying that $(Q_1 v G_2 v Q_1)^{-1}$ is strictly negative definite. If $f \in Q_3 L^2$ satisfies $$\langle Q_3 T_0 Q_1 (Q_1 v G_2 v Q_1)^{-1} Q_1 T_0 Q_3 f, f \rangle = 0,$$ this strict negativity enforces $Q_1 T_0 Q_3 f = 0$. Since $f \in Q_3 L^2 \subseteq S_2 L^2$, we already have $S_1 T_0 f = 0$. Hence, $0 = Q_1 T_0 f = (S_0 - S_1) T_0 f = S_0 T_0 f$. This implies $f \in S_3 L^2$, forcing $f \in Q_3 L^2 \cap S_3 L^2 = \{0\}$. Thus, $\mathbf{D}_{1,1}$ is invertible.
	
\textbf{(ii) Invertibility of $\mathbf{D}_{2,2}$.} Since $\mathbf{D}_{2,2}$ is a diagonal operator matrix, it suffices to show  $E_{5,5}$ and $E_{6,6}$ are invertible on $Q_5 L^2$ and $Q_6 L^2,$ respectively.
	 For any nonzero $f \in Q_5 L^2$, we have
	\[
	\langle E_{5,5}f, f \rangle = -b_1 \langle Q_5 v G_6 v Q_5 f, f \rangle
	= b_1 \Big( 12 \sum_{j=1}^{2} |\langle |x|^2 x_j v, f \rangle|^2 + 8 \sum_{i,j,\ell=1}^{2} |\langle x_i x_j x_\ell v, f \rangle|^2 \Big)>0,
	\]
	with $ b_1 = {1}/({4608\pi}).$ Hence $ E_{5,5}$ is strictly positive definite on $Q_5L^2$.
	We next verify the invertibility of $E_{6,6} = b_1 Q_6 v G_6^0 v Q_6$. For any $f \in Q_6 L^2 \subset S_5 L^2$, the associated quadratic form evaluates to
	\[
	\langle E_{6,6} f, f \rangle = \int_{\mathbb{R}^2} |(\Delta^2)^{-1} v f(x)|^2 \, dx = \|b_0 G_2^0 v f\|_{L^2}^2.
	\]
	Proposition~\ref{characterizations-1}(iv) and \eqref{expression_psi} guarantee that $b_0 G_2^0 v f \in L^2$ and that $f =-b_0 U v G_2^0 v f$. Thus, if $\langle E_{6,6} f, f \rangle = 0$, then $G_2^0 v f = 0$, which yields $f = 0$. This establishes the strict positivity of $E_{6,6}$ on $Q_6 L^2$, ensuring its invertibility (see also \cite[Lemma 4.9]{LSY21} for proof).
		
	\textbf{(iii) Invertibility of $\mathbf{A}^{\pm}$.} In this part, we not only establish the invertibility of $\mathbf{A}^{\pm}$, but also derive the formula for  $(\mathbf{A}^{\pm})^{-1}$, which is essential for obtaining the refined expansion of $(M^{\pm}(\lambda))^{-1}$ in the subsequent analysis.  
Our discussion proceeds according to whether $Q_4^0 = 0$ or not.
	
	\textbf{Case 1: $Q_4^0 = 0$.} Here $Q_4 = Q_4^1$ (see Definition \ref{defQ} that  $Q_4 =Q_4^0 +Q_4^1$ ), so
	\begin{equation}\label{eq:cancel_Q_4}
		Q_4 v = Q_4(x_j v) = Q_4(|x|^2 v) = 0.
	\end{equation}
	Consequently, $Q_4 v G_2 v Q_0 = Q_0 v G_2 v Q_4 = 0$, reducing $\mathbf{A}^{\pm}$ to the  form  
	\[
	\mathbf{A}^{\pm} = \begin{pmatrix}
		a_0^\pm Q_0 v G_0 v Q_0 & Q_0 T_0 Q_4^1 \\
		Q_4^1 T_0 Q_0 & a_2^\pm Q_4^1 v G_4 v Q_4^1
	\end{pmatrix}.
	\]
Note that $a_0^\pm Q_0 v G_0 v Q_0= a_0^\pm\|V\|_{L^1}Q_0$ with $a_0^\pm=\pm i/8$, then $(a_0^\pm Q_0 v G_0 v Q_0)^{-1}=(a_0^\pm)^{-1}\|V\|_{L^1}^{-1}Q_0.$
	By Lemma~\ref{lemma_2}, $\mathbf{A}^{\pm}$ is invertible if and only if the Schur complement 
	\[
	\boldsymbol{d}^\pm := a_2^\pm Q_4^1 v G_4 v Q_4^1 - Q_4^1 T_0 Q_0 (a_0^\pm Q_0 v G_0 v Q_0)^{-1} Q_0 T_0 Q_4^1
	\]
	is invertible on $Q_4L^2$ (here $Q_4=Q_4^1$). For any $f \in Q_4 L^2$, using \eqref{eq:cancel_Q_4}, we 
	derive that
	\[
		\langle \boldsymbol{d}^\pm f, f \rangle = \pm i \Bigl[ \sum_{j,k=1}^2 \tfrac{1}{128} |\langle x_j x_k v, f \rangle|^2 + \langle 8 (Q_0 v G_0 v Q_0)^{-1} Q_0 T_0 Q_4^1 f, Q_0 T_0 Q_4^1 f \rangle \Bigr].
	\]
	Because $(Q_0 v G_0 v Q_0)^{-1}$ is strictly positive on $Q_0 L^2$, the condition $\langle \boldsymbol{d}^\pm f, f \rangle = 0$ forces $\langle x_j x_k v, f \rangle = 0$ for all $j,k$, as well as $Q_0 T_0 Q_4^1 f = 0$. This implies $f \in S_4 L^2$. Since $Q_4 L^2 \cap S_4 L^2 = \{0\}$, it follows that $f = 0$, thereby confirming the invertibility of $\boldsymbol{d}^\pm$. The inverse $(\mathbf{A}^{\pm})^{-1} := \bigl(\mathcal{D}_{i,j}^{\pm}\bigr)_{i,j\in\{0,4\}}$ can then be explicitly constructed via \eqref{A-m}, yielding  
\begin{equation}\label{D00}  
	\begin{aligned}  
		\mathcal{D}_{0,0}^{\pm} &= (a_0^\pm)^{-1}\|V\|_{L^1}^{-1}Q_0 + (a_0^\pm)^{-2} \mathcal{W} (\boldsymbol{d}^\pm)^{-1} \mathcal{W}^*,  \quad  
		\mathcal{D}_{0,4}^{\pm} = - (a_0^\pm)^{-1} \mathcal{W} (\boldsymbol{d}^\pm)^{-1},  \\  
		\mathcal{D}_{4,0}^{\pm} &= - (a_0^\pm)^{-1} (\boldsymbol{d}^\pm)^{-1} \mathcal{W}^*,\quad  
		\mathcal{D}_{4,4}^{\pm} = (\boldsymbol{d}^\pm)^{-1},  
	\end{aligned}  
\end{equation}  
where we have denoted $\mathcal{W} := \|V\|_{L^1}^{-1} Q_0 T_0 Q_4^1$.  
	
\textbf{Case 2: $Q_4^0 \neq 0$.} Utilizing the projection splitting $Q_4 = Q_4^0 + Q_4^1$, we represent $\mathbf{A}^{\pm}$ as a $2 \times 2$ block matrix $\mathbf{A}^{\pm} = \bigl( \mathbf{A}_{ij}^{\pm}\bigr)_{i,j=1}^2$ with respect to the orthogonal decomposition $(Q_0 L^2 \oplus Q_4^0 L^2) \oplus Q_{4}^1 L^2$. The block components are explicitly given by:
	\begin{align}\label{Aij}
		\mathbf{A}_{11}^{\pm} &= \begin{pmatrix}
			a_0^\pm Q_0 v G_0 v Q_0 & g_0^{\pm}(\lambda) Q_0 v G_2 v Q_4^0 + Q_0 T_0 Q_4^0 \\[6pt]
			g_0^{\pm}(\lambda) Q_4^0 v G_2 v Q_0 + Q_4^0 T_0 Q_0 & a_2^\pm Q_4^0 v G_4 v Q_4^0
		\end{pmatrix}, \nonumber\\[8pt]
		\mathbf{A}_{12}^{\pm} &= \begin{pmatrix} Q_0 T_0 Q_{4}^1  \\ a_2^\pm Q_4^0 v G_4 v Q_{4}^1 \end{pmatrix}, \quad
		\mathbf{A}_{21}^{\pm} = \begin{pmatrix} Q_{4}^1 T_0 Q_0 & a_2^\pm Q_4^1 v G_4 v Q_4^0 \end{pmatrix}, \quad
		\mathbf{A}_{22}^{\pm} = a_2^\pm Q_4^1 v G_4 v Q_4^1.
	\end{align}
	
	First, we show that $\mathbf{A}_{11}^{\pm}$ is invertible on $Q_0 L^2 \oplus Q_4^0 L^2$. Write $\mathbf{A}_{11}^{\pm} = (\log \lambda) \mathbf{A}_0 + \mathbf{A}_1^\pm$ with
	\[
	\mathbf{A}_0 = \begin{pmatrix}
		0 & b_0 Q_0 v G_2 v Q_4^0 \\
		b_0 Q_4^0 v G_2 v Q_0 & 0
	\end{pmatrix},\quad
	\mathbf{A}_1^\pm = \begin{pmatrix}
		a_0^\pm Q_0 v G_0 v Q_0 & Q_0 T_0 Q_4^0 + a_1^{\pm} Q_0 v G_2 v Q_4^0 \\
		Q_4^0 T_0 Q_0 + a_1^{\pm} Q_4^0 v G_2 v Q_0 & a_2^\pm Q_4^0 v G_4 v Q_4^0
	\end{pmatrix}.
	\]
	Observe that the operators $Q_4^0 v G_2 v Q_0 \colon Q_0 L^2 \to Q_4^0 L^2$ and $Q_0 v G_2 v Q_4^0 \colon Q_4^0 L^2 \to Q_0 L^2$ are both invertible. Indeed, for any $f\in Q_0 L^2$,
	$$
	Q_4^0 v G_2 v Q_0 f = \langle f, v \rangle Q_4^0 (|x|^2 v)=0 \Longleftrightarrow
	\langle f, v \rangle = 0 \Longleftrightarrow f = 0.
	$$
	Similarly, for any $f \in Q_4^0 L^2$,
		$$
	Q_0 v G_2 v Q_4^0 f = \langle f, |x|^2 v \rangle v = 0 \Longleftrightarrow \langle f, |x|^2 v \rangle = 0 \Longleftrightarrow f \in S_4^0 L^2 \Longleftrightarrow f=0.
	$$
    Because $\dim(Q_0 L^2) = \dim(Q_4^0 L^2) = 1$, this injectivity guarantees that both operators are invertible. Thus $\mathbf{A}_0$ is invertible with
	\[
	\mathbf{A}_0^{-1} = \begin{pmatrix}
		0 & (b_0 Q_4^0 v G_2 v Q_0)^{-1} \\
		(b_0 Q_0 v G_2 v Q_4^0)^{-1} & 0
	\end{pmatrix}.
	\]
	For sufficiently small $\lambda>0$, the Neumann series yields
	\begin{equation}\label{eq:inverse_A_11}
		(\mathbf{A}_{11}^{\pm})^{-1} = (\log \lambda)^{-1} \mathbf{A}_0^{-1} - (\log \lambda)^{-2} \mathbf{A}_0^{-1} \mathbf{A}_1^\pm \mathbf{A}_0^{-1} + (\log \lambda)^{-3}\Lambda^\pm(\lambda).
	\end{equation}
	By Lemma~\ref{lemma_2}, $\mathbf{A}^{\pm} = \bigl( \mathbf{A}_{ij}^{\pm}\bigr)_{i,j=1}^2$ is invertible if and only if
	$$
	\begin{aligned}
		\boldsymbol{d}^\pm 
        %& := \mathbf{A}_{22}^{\pm} - \mathbf{A}_{21}^{\pm} (\mathbf{A}_{11}^{\pm})^{-1} \mathbf{A}_{12}^{\pm} \\
		&= a_2^\pm Q_4^1 v G_4 v Q_4^1
		- \begin{pmatrix} Q_4^1 T_0 Q_0 & a_2^\pm Q_4^1 v G_4 v Q_4^0 \end{pmatrix}
		(\mathbf{A}_{11}^{\pm})^{-1}
		\begin{pmatrix} Q_0 T_0 Q_4^1 \\ a_2^\pm Q_4^0 v G_4 v Q_4^1 \end{pmatrix}
	\end{aligned}
	$$
	is invertible on $Q_4^1 L^2$.   We analyze this  invertibility  by splitting $Q_4^1 L^2 = Q_{41}^1 L^2 \oplus Q_{42}^1 L^2$  (see Definition  \ref{defQ} and  Remark \ref{cancellationQ} for more details):

	$\bullet$ \textbf{Subcase 2a ($Q_{42}^1 = 0$):}  
		Here $Q_4^1 = Q_{41}^1$. We first consider the nontrivial case: $
		Q_{41}^1L^2\neq\{0\}.
		$  Since
\begin{equation}\label{eq:positive}
			a_2^\pm \langle Q_{4}^1 v G_4 v Q_{4}^1 f, f \rangle	=a_2^\pm \langle Q_{41}^1 v G_4 v Q_{41}^1 f, f \rangle = a_2^\pm \sum_{j,\ell=1}^2 4 |\langle x_j x_\ell v, f \rangle|^2,\quad \forall f\in Q_4^1 L^2=Q_{41}^1 L^2,
		\end{equation}
		with $a_2^\pm = \pm i/512$, by definition of $Q_{41}^1,$  this quadratic form is nondegenerate. Thus, $a_2^\pm Q_{4}^1 v G_4 v Q_{4}^1$ is invertible on $Q_4^1 L^2$. Combining  this with $\|(\mathbf{A}_{11}^{\pm})^{-1}\|_{\mathbb{B}(L^2)} = O(|\log \lambda|^{-1})$
		yields
		 $\boldsymbol{d}^\pm$ is invertible and 
		$$
		(\boldsymbol{d}^\pm)^{-1} = (a_2^\pm Q_{41}^1 v G_4 v Q_{41}^1)^{-1} + (\log \lambda)^{-1}\Lambda^\pm(\lambda):= (a_2^\pm)^{-1} D_0+ (\log \lambda)^{-1}\Lambda^\pm(\lambda).
		$$
		Substituting the expressions \eqref{eq:inverse_A_11} for $(\mathbf{A}_{11}^{\pm})^{-1}$ and \eqref{Aij} for $\mathbf{A}_{ij}^{\pm}$ into \eqref{A-m}, we deduce the block representation
		$
		(\mathbf{A}^\pm)^{-1} = \big( \mathcal{D}_{\alpha,\beta}^{\pm}(\lambda) \big)_{\alpha,\beta \in \{0,4\}},
		$
		where
	\begin{equation}\label{D40D04D44_1}
			\begin{split}
				\mathcal{D}_{0,0}^{\pm}(\lambda) & = (\log \lambda)^{-2} Q_0 \Lambda^\pm(\lambda) Q_0, \\
				\mathcal{D}_{0,4}^{\pm}(\lambda) &= (\log \lambda)^{-1} \left[ Q_0 \Lambda^\pm(\lambda) Q_4^0 + Q_0 \Lambda^\pm(\lambda) Q_{4}^1 \right] , \\
				\mathcal{D}_{4,0}^{\pm}(\lambda) &= (\log \lambda)^{-1} \left[ Q_4^0 \Lambda^\pm(\lambda) Q_0 + Q_{4}^1 \Lambda^\pm(\lambda) Q_0 \right], \\
				\mathcal{D}_{4,4}^{\pm}(\lambda) &= (\log \lambda)^{-2} Q_4^0 \Lambda^\pm(\lambda) Q_4^0
				+ (\log \lambda)^{-1} \left[ Q_4^0 \Lambda^\pm(\lambda) Q_{4}^1 + Q_{4}^1 \Lambda^\pm(\lambda) Q_4^0 \right] \\
				& \quad + (\log \lambda)^{-1}Q_{4}^1 \Lambda^\pm(\lambda) Q_{4}^1+(a_2^\pm)^{-1} D_0.
			\end{split}
		\end{equation}
        
	It remains to observe that the same representation \eqref{D40D04D44_1} also applies to the degenerate case
		$Q_{41}^1L^2=\{0\}$.
		Indeed, since $Q_{42}^1=Q_4^1=0$, the $Q_4^1$-block is absent. Then
		$(\mathbf A^\pm)^{-1}=(\mathbf A_{11}^\pm)^{-1}$.
		Thus, by \eqref{eq:inverse_A_11}, the block representation \eqref{D40D04D44_1} remains valid, where we interpret the term $D_0$ as
		$
		D_0
		=
		\left(
		Q_{41}^1vG_4vQ_{41}^1
		\right)^{-1}=0.
		$
		
		$\bullet$ \textbf{Subcase 2b ($Q_{42}^1 \neq 0$):} Observe that $Q_4 v G_4 v Q_{42}^1 = 0$. With respect to the decomposition $Q_4^1 L^2 = Q_{41}^1 L^2 \oplus Q_{42}^1 L^2$, we can express $\boldsymbol{d}^\pm=\bigl(d^\pm_{ij}\bigr)_{i,j=1}^2$, where  
		\begin{equation}\label{eq:d_ij}  
			\begin{split}  
				d_{11}^\pm &= a_2^\pm Q_{41}^1 v G_4 v Q_{41}^1 + (\log \lambda)^{-1}\Lambda^\pm(\lambda), \ \  
				d_{12}^\pm = -(\log \lambda)^{-1} a_2^\pm B + (\log \lambda)^{-2}\Lambda^\pm(\lambda), \\  
				d_{21}^\pm &= -(\log \lambda)^{-1} a_2^\pm B^* + (\log \lambda)^{-2}\Lambda^\pm(\lambda), \ \  
				d_{22}^\pm = (\log \lambda)^{-2} a_2^\pm \mathcal{D} + (\log \lambda)^{-3}\Lambda^\pm(\lambda),  
			\end{split}  
		\end{equation}  
		with  
		\[  
		\begin{aligned}    
			B &= Q_{41}^1 v G_4 v Q_4^0 (b_0 Q_0 v G_2 v Q_4^0)^{-1} Q_0 T_0 Q_{42}^1, \\  
			\mathcal{D} &= Q_{42}^1 T_0 Q_0 (b_0 Q_4^0 v G_2 v Q_0)^{-1} Q_4^0 v G_4 v Q_4^0 (b_0 Q_0 v G_2 v Q_4^0)^{-1} Q_0 T_0 Q_{42}^1.  
		\end{aligned}  
		\] 
        
	We first assume that
		$
		Q_{41}^1L^2\neq\{0\}.
		$
		By \eqref{eq:positive}, the leading operator $Q_{41}^1 v G_4 v Q_{41}^1$ is invertible on $Q_{41}^1 L^2$. Define $D_0 := (Q_{41}^1 v G_4 v Q_{41}^1)^{-1}$. Then for sufficiently small $\lambda>0$, $d_{11}^\pm$ is invertible on $Q_{41}^1 L^2$ with its inverse given by  
		\[  
		(d_{11}^\pm)^{-1} = (a_2^\pm)^{-1} D_0 + (\log \lambda)^{-1}\Lambda^\pm(\lambda).
		\]  
		 By Lemma~\ref{lemma_2}, $\boldsymbol{d}^\pm=\bigl(d^\pm_{ij}\bigr)_{i,j=1}^2$ is invertible if and only if
		\begin{equation}\label{d_1}
			\boldsymbol{d}^\pm_1 := d_{22}^\pm - d_{21}^\pm (d_{11}^\pm)^{-1} d_{12}^\pm = (\log \lambda)^{-2} a_2^\pm \mathcal{S} + (\log \lambda)^{-3}\Lambda^\pm(\lambda)
		\end{equation}
		is invertible on $Q_{42}^1 L^2$, where
\begin{equation*}
			\mathcal{S} = \mathcal{D} - B^* D_0 B=\mathcal{W}_1^* \boldsymbol{d}_2 \mathcal{W}_1,
		\end{equation*}
		with $\mathcal{W}_1 = (b_0 Q_0 v G_2 v Q_4^0)^{-1} Q_0 T_0 Q_{42}^1$ and
		$$
		\boldsymbol{d}_2 = Q_4^0 v G_4 v Q_4^0 - Q_4^0 v G_4 v Q_{41}^1 (Q_{41}^1 v G_4 v Q_{41}^1)^{-1} Q_{41}^1 v G_4 v Q_4^0.
		$$
		Observe that $Q_{41}^1 v G_4 v Q_{41}^1$ is strictly positive on $Q_{41}^1 L^2$ and the block matrix
		\begin{equation}\label{QvGvQ}
			\begin{pmatrix}
				Q_{4}^0 v G_4 v Q_{4}^0 & Q_{4}^0 v G_4 v Q_{41}^1 \\
				Q_{41}^1 v G_4 v Q_{4}^0 & Q_{41}^1 v G_4 v Q_{41}^1
			\end{pmatrix}
		\end{equation}
		is strictly positive on $Q_{4}^0 L^2 \oplus Q_{41}^1 L^2$. Here, \eqref{QvGvQ} is exactly the block matrix representation of the strictly positive operator $(Q_4^0+Q_{41}^1)v G_4 v (Q_4^0+Q_{41}^1)$ with respect to $Q_{4}^0 L^2 \oplus Q_{41}^1 L^2$. By Lemma~\ref{lemma_2} (ii), it follows that $\boldsymbol{d}_2$ is strictly positive on $Q_{4}^0 L^2$.
		
		Therefore, for any $f \in Q_{42}^1 L^2$, the condition $\langle \mathcal{S} f, f \rangle = \langle \boldsymbol{d}_2 \mathcal{W}_1 f, \mathcal{W}_1 f \rangle = 0$ implies that $\mathcal{W}_1 f = 0$. This yields $Q_0 T_0 Q_{42}^1 f = P T_0 f = 0$, which means $f \in S_4 L^2$. Then  $f \in S_4 L^2 \cap Q_{42}^1 L^2 = \{0\}$. 
		This proves that $\mathcal{S}$ is positive definite. Consequently, by \eqref{d_1}, $\boldsymbol{d}_1^\pm$ is invertible on $Q_{42}^1 L^2$  and 
		\[
		(\boldsymbol{d}_1^\pm)^{-1} = (\log \lambda)^{2} (a_2^\pm)^{-1} \mathcal{S}^{-1} + (\log \lambda)\Lambda^\pm(\lambda).
		\]
		Thus $\boldsymbol{d}_1^\pm$ is invertible on $Q_{42}^1 L^2$, which in turn guarantees the invertibility of $\boldsymbol{d}^\pm$ and $\mathbf{A}^\pm$.
        
		Next, we turn to deriving the formula for $	(\mathbf{A}^{\pm})^{-1}.$
		Under the decomposition $Q_4^1 L^2 = Q_{41}^1 L^2 \oplus Q_{42}^1 L^2$, using the expression \eqref{eq:d_ij} for $d_{ij}^\pm$ together with \eqref{A-m}, we obtain
		\begin{align*}
			(\boldsymbol{d}^\pm)^{-1} 
			= \begin{pmatrix}
				(a_2^\pm)^{-1}\bigl(D_0 + D_0 B \mathcal{S}^{-1} B^* D_0\bigr) + (\log \lambda)^{-1}\Lambda^\pm(\lambda) & (\log \lambda) (a_2^\pm)^{-1}D_0 B \mathcal{S}^{-1} + \Lambda^\pm(\lambda) \\[4pt]
				(\log \lambda) (a_2^\pm)^{-1}\mathcal{S}^{-1} B^* D_0 + \Lambda^\pm(\lambda) & (\log \lambda)^2 (a_2^\pm)^{-1}\mathcal{S}^{-1} + (\log \lambda)\Lambda^\pm(\lambda)
			\end{pmatrix}.
		\end{align*}
		Moreover, recall that $\boldsymbol{d}^\pm = \mathbf{A}_{22}^{\pm} - \mathbf{A}_{21}^{\pm} (\mathbf{A}_{11}^{\pm})^{-1} \mathbf{A}_{12}^{\pm}$ and $\mathbf{A}^{\pm} = \bigl( \mathbf{A}_{ij}^{\pm}\bigr)_{i,j=1}^2$, where  $\mathbf{A}_{ij}^{\pm}$ are defined in \eqref{Aij}. Then, by \eqref{A-m}, with respect to the decomposition \( (Q_0 L^2 \oplus Q_4^0 L^2)\oplus (Q_{41}^1 L^2 \oplus Q_{42}^1 L^2)\), we obtain
		\begin{align*}
		(\mathbf{A}^{\pm})^{-1} = \begin{pmatrix}
			(\mathbf{A}_{11}^{\pm})^{-1} + (\mathbf{A}_{11}^{\pm})^{-1} \mathbf{A}_{12}^{\pm} (\boldsymbol{d}^\pm)^{-1} \mathbf{A}_{21}^{\pm} (\mathbf{A}_{11}^{\pm})^{-1} & -(\mathbf{A}_{11}^{\pm})^{-1} \mathbf{A}_{12}^{\pm} (\boldsymbol{d}^\pm)^{-1} \\ 
			-(\boldsymbol{d}^\pm)^{-1} \mathbf{A}_{21}^{\pm} (\mathbf{A}_{11}^{\pm})^{-1} & (\boldsymbol{d}^\pm)^{-1}
		\end{pmatrix}.
		\end{align*}
	Direct computations based on the previous inverse bounds reveal that  
		\[
		\begin{aligned}
			(\mathbf{A}_{11}^{\pm})^{-1} + (\mathbf{A}_{11}^{\pm})^{-1} \mathbf{A}_{12}^{\pm} (\boldsymbol{d}^\pm)^{-1} \mathbf{A}_{21}^{\pm} (\mathbf{A}_{11}^{\pm})^{-1}
			&= \begin{pmatrix}
				(\log \lambda)^{-2}\Lambda^\pm(\lambda) & (\log \lambda)^{-1}\Lambda^\pm(\lambda) \\[4pt]
				(\log \lambda)^{-1}\Lambda^\pm(\lambda) & P_{22}^\pm(\lambda)
			\end{pmatrix}, \\[8pt]
			-(\mathbf{A}_{11}^{\pm})^{-1} \mathbf{A}_{12}^{\pm} (\boldsymbol{d}^\pm)^{-1}
			&= \begin{pmatrix}
				(\log \lambda)^{-1}\Lambda^\pm(\lambda) & \Lambda^\pm(\lambda) \\[4pt]
				\Lambda^\pm(\lambda) & K_{22}^\pm(\lambda)
			\end{pmatrix}, \\[8pt]
			-(\boldsymbol{d}^\pm)^{-1} \mathbf{A}_{21}^{\pm} (\mathbf{A}_{11}^{\pm})^{-1}
			&= \begin{pmatrix}
				(\log \lambda)^{-1}\Lambda^\pm(\lambda) & \Lambda^\pm(\lambda) \\[4pt]
				\Lambda^\pm(\lambda) & N_{22}^\pm(\lambda)
			\end{pmatrix},
		\end{aligned}
		\]
		with  $\mathcal{W}_1 = (b_0 Q_0 v G_2 v Q_4^0)^{-1} Q_0 T_0 Q_{42}^1,$ and
		\[
		\begin{aligned}
			K_{22}^\pm(\lambda) &= -(\log \lambda) (a_2^\pm)^{-1}\mathcal{W}_1 \mathcal{S}^{-1}+\Lambda^\pm(\lambda)
			:=(\log \lambda)\mathcal{C}_1^\pm+ \Lambda^\pm(\lambda), \\
			N_{22}^\pm(\lambda) &= -(\log \lambda) (a_2^\pm)^{-1}\mathcal{S}^{-1} \mathcal{W}_1^* +\Lambda^\pm(\lambda)
			:=(\log \lambda)\mathcal{C}_2^\pm+ \Lambda^\pm(\lambda), \\
			P_{22}^\pm(\lambda) & = (a_2^\pm)^{-1}\mathcal{W}_1 \mathcal{S}^{-1} \mathcal{W}_1^* +(\log \lambda)^{-1}\Lambda^\pm(\lambda)
			:=\mathcal{C}^\pm+(\log \lambda)^{-1}\Lambda^\pm(\lambda).
		\end{aligned}
		\]
	Consequently, condense back to the representation $(\mathbf{A}^\pm)^{-1} = \big( \mathcal{D}_{\alpha,\beta}^{\pm}(\lambda) \big)_{\alpha,\beta \in \{0,4\}}$ on \( Q_0 L^2 \oplus Q_4 L^2 \), then
	\begin{equation}\label{D40D04D44_2}
        \begin{split}
				\mathcal{D}_{0,0}^{\pm}(\lambda) &= (\log \lambda)^{-2} Q_0 \Lambda^\pm(\lambda) Q_0,   \\
				\mathcal{D}_{0,4}^{\pm}(\lambda) &= (\log \lambda)^{-1} Q_0 \Lambda^\pm(\lambda) Q_{4}^0+  Q_0 \Lambda^\pm(\lambda) Q_{4}^1,   \\
				\mathcal{D}_{4,0}^{\pm}(\lambda) &= (\log \lambda)^{-1} Q_{4}^0 \Lambda^\pm(\lambda) Q_0 +Q_{4}^1 \Lambda^\pm(\lambda) Q_0,  \\
				\mathcal{D}_{4,4}^{\pm}(\lambda)  &=Q_4^0\mathcal{C}^\pm Q_4^0+(\log \lambda)^{-1} Q_4^0\Lambda^\pm(\lambda) Q_4^0
				+\bigl(Q_4^0 \Lambda^\pm(\lambda) Q_{41}^1 + Q_{41}^1 \Lambda^\pm(\lambda) Q_4^0\bigr) \\
				& + (\log \lambda)\Bigl[(Q_4^0 \mathcal{C}_1^\pm Q_{42}^1 +Q_{42}^1 \mathcal{C}_2^\pm Q_4^0) +(Q_{41}^1 \Lambda^\pm(\lambda) Q_{42}^1 +Q_{42}^1 \Lambda^\pm(\lambda) Q_{41}^1) \Bigr]     \\
				&+\bigl(Q_4^0 \Lambda^\pm(\lambda) Q_{42}^1 +Q_{42}^1 \Lambda^\pm(\lambda)  Q_4^0\bigr)+ Q_{41}^1 \Lambda^\pm(\lambda) Q_{41}^1 +(\log \lambda)^2 (a_2^\pm)^{-1} Q_{42}^1 \mathcal{S}^{-1} Q_{42}^1 \\
				& + (\log \lambda) Q_{42}^1 \Lambda^\pm(\lambda)Q_{42}^1.
        \end{split}
        \end{equation}

	It remains only to mention the  case $ Q_{41}^1L^2=\{0\}. $ In this case the $Q_{41}^1$-block is absent. Equivalently, the preceding formula remains valid if one interprets $D_0 = \left( Q_{41}^1vG_4vQ_{41}^1 \right)^{-1} $ as the zero operator.







        
	Thus, we have established the invertibility of $\mathbf{A}^{\pm}$ across all three possible cases: $Q^0_4=0$; $Q^0_4\neq 0$ with $Q_{42}^1=0$; and $Q^0_4\neq 0$ with $Q_{42}^1\neq 0$. Based on the invertibility of the respective components established in parts (i)--(iii), we conclude that $\mathbf{D}_{4}^{\pm}$ is invertible on $\medoplus_{\alpha=0}^{6} Q_\alpha L^2$, and its inverse is given by    
\begin{equation}\label{inverse_D4}    
\left(\mathbf{D}_{4}^{\pm}\right)^{-1} =  
	\begin{pmatrix}    
		\mathcal{D}_{0,0}^{\pm}(\lambda) & 0 & \mathcal{D}_{0,4}^{\pm}(\lambda) & 0 \\    
		0 & \left(\mathbf{D}_{1,1}\right)^{-1} & 0 & 0 \\    
		\mathcal{D}_{4,0}^{\pm}(\lambda) & 0 & \mathcal{D}_{4,4}^{\pm}(\lambda) & 0 \\    
		0 & 0 & 0 & \left(\mathbf{D}_{2,2}\right)^{-1}
	\end{pmatrix}    
	:= \left( \mathcal{D}_{\alpha,\beta}^\pm(\lambda) \right)_{\alpha,\beta=0}^6,    
\end{equation}    
where $\mathcal{D}_{\alpha,\beta}^{\pm}(\lambda) \in \mathbb{B}(L^2)$ for all $0\leq\alpha,\beta\leq6$. The entries $\mathcal{D}_{\alpha,\beta}^{\pm}(\lambda)$ with $\alpha,\beta \in \{0,4\}$ are defined in \eqref{D00}, \eqref{D40D04D44_1}, and \eqref{D40D04D44_2}, respectively corresponding to the three distinct cases.    
Given the structures of $\mathbf{D}_{1,1}$ and $\mathbf{D}_{2,2}$ in \eqref{D}, their inverses take the forms  
		\begin{align}\label{D-invert}  
	\left(\mathbf{D}_{1,1}\right)^{-1}=  
			\begin{pmatrix}  
				\mathcal{D}_{1,1} & 0 & \mathcal{D}_{1,3} \\  
				0 & \mathcal{D}_{2,2} & 0 \\  
				\mathcal{D}_{3,1} & 0 & \mathcal{D}_{3,3}  
			\end{pmatrix}, \quad   
			\left(\mathbf{D}_{2,2} \right)^{-1}=   
			\begin{pmatrix}  
				\mathcal{D}_{5,5} & 0 \\  
				0 & \mathcal{D}_{6,6}  
			\end{pmatrix}.  
		\end{align}  
		Here, we write $\mathcal{D}_{\alpha,\beta}$ instead of $\mathcal{D}_{\alpha,\beta}^\pm(\lambda)$ for brevity, as the entries of $\left(\mathbf{D}_{1,1}\right)^{-1}$ and $\left(\mathbf{D}_{2,2}\right)^{-1}$ depend on neither $\lambda$ nor the choice of sign $\pm$. Moreover, 
		\begin{align*}  
			\mathcal{D}_{2,2} = (Q_2 T_0 Q_2)^{-1}, \quad  
			\mathcal{D}_{5,5} = (-b_1 Q_5 v G_6 v Q_5)^{-1}, \quad  
			\mathcal{D}_{6,6} = (b_1 Q_6 v G_6^0 v Q_6)^{-1}.  
		\end{align*}  
		
	
    {\textbf{Step 2. Asymptotic expansion of $(M^{\pm}(\lambda))^{-1}$ via Neumann series.}}
    
We begin by estimating the matrix operators $\mathbf{S}_{4}^{\pm}(\lambda)\bigl(\mathbf{D}_{4}^{\pm}\bigr)^{-1}=\left(\sum_{j=0}^6 r_{\alpha,j}^\pm(\lambda)\mathcal{D}_{j,\beta}^{\pm}(\lambda)\right)_{\alpha,\beta=0}^{6}$. From \eqref{D00}, \eqref{D40D04D44_1}, and \eqref{D40D04D44_2}, we observe that all leading entries satisfy 
\[
\Big\| \partial_\lambda^\ell \mathcal{D}_{\alpha,\beta}^{\pm}(\lambda) \Big\|_{\mathbb{B}(L^2)} \lesssim
\begin{cases}
\lambda^{-\ell}, & \text{if } (\alpha, \beta) \neq (4,4), \\[4pt]
|\log \lambda|^2 \lambda^{-\ell}, & \text{if } (\alpha, \beta) = (4,4),
\end{cases}
\]
for $\ell = 0,1,2$. Combining this with the remainder estimates \eqref{partial}--\eqref{partial2} for $r_{\alpha,j}^\pm(\lambda)$, and noting that each factor $\mathcal{D}_{4,4}^{\pm}(\lambda)$ is paired with a remainder entry $r_{\alpha,4}^\pm(\lambda)$ satisfying $r_{\alpha,4}^\pm(\lambda) = O\bigl(\lambda (-\log \lambda)^{3/2}\bigr)$, we obtain
\[
\mathbf{S}_{4}^{\pm}(\lambda)\bigl(\mathbf{D}_{4}^{\pm}\bigr)^{-1} = O\bigl((-\log \lambda)^{-1/2}\bigr).
\]
We then employ the Neumann series expansion in the small parameter $(-\log \lambda)^{-1/2}$:
\begin{align}
\left(\mathbf{B}_\lambda^* M^{\pm}(\lambda) \mathbf{B}_\lambda\right)^{-1}
&=\lambda^2 \left(\mathbf{D}_{4}^{\pm}\right)^{-1}\left(\Id +\mathbf{S}_{4}^{\pm}(\lambda)\left(\mathbf{D}_{4}^{\pm}\right)^{-1}\right)^{-1} \nonumber\\
&= \lambda^2 \Bigl[ \left(\mathbf{D}_{4}^{\pm}\right)^{-1} + \sum_{n=1}^{\infty}(-1)^n \left(\mathbf{D}_{4}^{\pm}\right)^{-1} \left( \mathbf{S}_{4}^{\pm}(\lambda) \left(\mathbf{D}_{4}^{\pm}\right)^{-1} \right)^n \Bigr] \nonumber\\
&:= \lambda^2 \left( \mathcal{D}_{\alpha,\beta}^\pm(\lambda) + \Upsilon_{\alpha,\beta}^\pm(\lambda) \right)_{\alpha,\beta=0}^{6}. \label{MM-inver}
\end{align}

To bound the remainder operators $\Upsilon_{\alpha,\beta}^\pm(\lambda)$, we examine the structure of matrix multiplication in the Neumann series.
Any term contributing to $\Upsilon_{\alpha,\beta}^\pm(\lambda)$ that involves the factor $\mathcal{D}_{4,4}^{\pm}(\lambda)$ must be multiplied by an adjacent remainder entry of the form $r_{4,j}^\pm(\lambda)$ or $r_{j,4}^\pm(\lambda)$. According to \eqref{partial}--\eqref{partial2}, all such entries carrying an index $4$ satisfy $r_{4,j}^{\pm}(\lambda), r_{i,4}^{\pm}(\lambda) = O\bigl(\lambda (-\log \lambda)^{3/2}\bigr)$. This polynomial decay in $\lambda$ compensates the $|\log \lambda|^2$ growth originating from $\mathcal{D}_{4,4}^{\pm}(\lambda)$. Consequently, from \eqref{partial}--\eqref{partial2} it follows that every remainder component $\Upsilon_{\alpha,\beta}^\pm(\lambda)$ belongs to $\mathbb{B}(L^2)$ and obeys the decay estimate
\begin{equation}\label{Upsilon}
\left\|\partial_\lambda^\ell \Upsilon_{\alpha,\beta}^\pm(\lambda)\right\|_{\mathbb{B}(L^2)} \lesssim (-\log \lambda)^{-1/2} \lambda^{-\ell}, \qquad \ell = 0,1,2. 
\end{equation}

Finally, reverting the preconditioning step via \eqref{M-nin}, we obtain the full inverse operator on $L^2$:
\[
\left(M^{\pm}(\lambda)\right)^{-1} = \sum_{0 \le \alpha, \beta \le 6} \lambda^{2 - k_\alpha - k_\beta} Q_\alpha \mathcal{M}_{\alpha,\beta}^{\pm}(\lambda) Q_\beta,
\]
where the unified block operators are  defined by
\[
\mathcal{M}_{\alpha,\beta}^{\pm}(\lambda) := \sigma_\alpha(\lambda) \sigma_\beta(\lambda) \left(\mathcal{D}_{\alpha,\beta}^\pm(\lambda) + \Upsilon_{\alpha,\beta}^\pm(\lambda) \right).
\]
	

	{\textbf{Step 3. Extraction of Decay Rates for $\mathcal{M}_{\alpha,\beta}^{\pm}(\lambda)$.}}
    
		Our final task is to track the powers of $\lambda$ and $\log \lambda$ appearing in $\mathcal{D}_{\alpha,\beta}^\pm$, $\Upsilon_{\alpha,\beta}^\pm$ and establish the asymptotic expansion for $\mathcal{M}_{\alpha,\beta}^{\pm}(\lambda)$.
	 The analysis reduces to examining the structure of these components.  
	
	\vspace{1.5mm}  
	\textbf{(1) Analysis of $\mathcal{M}_{\alpha,\beta}^{\pm}(\lambda)$ for $0 \leq \alpha, \beta \leq 4.$}
From \eqref{inverse_D4}--\eqref{D-invert}, we have
\[
\mathcal{D}_{0,3}^{\pm} = \mathcal{D}_{3,0}^{\pm} = \mathcal{D}_{2,3}^{\pm} = \mathcal{D}_{3,2}^{\pm} = 0,
\]
which implies that $\mathcal{M}_{\alpha,\beta}^{\pm}(\lambda) = (-\log \lambda)^{1/2} \Upsilon_{\alpha,\beta}^\pm(\lambda)$ for the pairs $(\alpha,\beta) \in \{(0,3), (3,0), (2,3), (3,2)\}$.
Combining this with the definition of $\sigma_\alpha(\lambda)$ in \eqref{tab:sigma_alpha} and the bound \eqref{Upsilon} for $\Upsilon_{\alpha,\beta}^\pm(\lambda)$, we deduce that for $0 \le \alpha, \beta \le 3$ and $\ell = 0,1,2$,
\[
\Big\| \partial_\lambda^\ell \mathcal{M}_{\alpha,\beta}^\pm(\lambda) \Big\|_{\mathbb{B}(L^2)} \lesssim
\begin{cases}
	\lambda^{-\ell}, & \text{for } (\alpha, \beta) \neq (3,3), \\
	|\log \lambda| \lambda^{-\ell}, &  \text{for } (\alpha, \beta) = (3,3).
	\end{cases}\]

For $\mathcal{M}_{0,0}^{\pm}(\lambda)=\mathcal{D}_{0,0}^\pm(\lambda) + \Upsilon_{0,0}^\pm(\lambda)$, we now examine the refined structure.
Since $\mathcal{D}_{0,\beta}^\pm(\lambda)$ and $\mathcal{D}_{\beta,0}^\pm(\lambda)$ are nonzero only for $\beta \in \{0,4\}$, every term in $\Upsilon_{0,0}^\pm(\lambda)$ contains at least a factor $r_{i,j}^\pm$ with $i$ or $j$ in $\{0,4\}$. The estimates \eqref{partial}--\eqref{partial2} therefore yield
\begin{equation}\label{Upsilon-00-rough}
	\left\|
	\partial_\lambda^\ell
	\Upsilon_{0,0}^\pm(\lambda)
	\right\|_{\mathbb B(L^2)}
	\lesssim
	\lambda^{-\ell+1-},
	\qquad
	\ell=0,1,2.
\end{equation}
If $Q_4^0\neq0,$ by \eqref{D40D04D44_1} and \eqref{D40D04D44_2} that  $\mathcal{D}_{0,0}^{\pm}(\lambda) = (\log \lambda)^{-2} Q_0 \Lambda^\pm(\lambda) Q_0,$ combining with \eqref{Upsilon-00-rough}, one has $$\mathcal{M}_{0,0}^{\pm}(\lambda) = (\log \lambda)^{-2} Q_0 \Lambda^\pm(\lambda) Q_0.$$
If  $Q_4^0 = 0$, then the operator norms of all entries $\mathcal{D}_{\alpha,\beta}^{\pm}(\lambda)$ are bounded without any logarithmic factor (see \eqref{D00}). Since every term in $\Upsilon_{0,0}^\pm(\lambda)$ contains at least a factor $r_{i,j}^\pm$ with $i$ or $j$ in $\{0,4\}$ satisfying the form $r_{i,j}^\pm=\lambda(-\log\lambda)^{3/2}\Lambda^\pm(\lambda)$,  we obtain $\Upsilon_{0,0}^\pm(\lambda) = \lambda (-\log \lambda)^{3/2} \Lambda^\pm(\lambda)$. Consequently, by \eqref{D00}, we arrive at the refined expansion of \begin{align*}
%\label{M-00} 
\mathcal M_{0,0}^\pm(\lambda)
=
\mathcal D_{0,0}^\pm
+
\lambda(-\log\lambda)^{3/2}
Q_0\Lambda^\pm(\lambda)Q_0.\end{align*}

Next, we analyze the operators $\mathcal{M}_{\alpha,4}^{\pm}(\lambda)$ and $\mathcal{M}_{4,\alpha}^{\pm}(\lambda)$ for $0 \le \alpha \le 4$. By symmetry, we focus only on $\mathcal{M}_{\alpha,4}^{\pm}(\lambda) = \sigma_\alpha(\lambda) \mathcal{D}_{\alpha,4}^\pm(\lambda) + \sigma_\alpha(\lambda) \Upsilon_{\alpha,4}^\pm(\lambda)$. Since $\mathcal{D}_{j,4}^\pm(\lambda)$ is nonzero only for $j \in \{0,4\}$, every term in $\Upsilon_{\alpha,4}^\pm(\lambda)$ contains at least one factor $r_{i,j}^\pm(\lambda)$ with  $j \in\{0,4\}$. Combining this fact with the estimates in \eqref{partial}--\eqref{partial2}, and noting that all leading entries satisfy $\|\partial_\lambda^\ell \mathcal{D}_{\alpha,\beta}^{\pm}(\lambda)\|_{\mathbb{B}(L^2)} \lesssim \lambda^{-\ell}$, with the sole exception
\[
\|\partial_\lambda^\ell \mathcal{D}_{4,4}^{\pm}(\lambda)\|_{\mathbb{B}(L^2)} \lesssim 
\begin{cases}
|\log \lambda|^2 \lambda^{-\ell}, & \text{if } Q_4^0 \neq 0 \text{ and } Q_{42}^1 \neq 0, \\
\lambda^{-\ell}, & \text{otherwise},
\end{cases}
\]
for $\ell = 0,1,2$, we obtain the following refined behavior.

If $Q_4^0 = 0$ or $Q_{42}^1 = 0$, then
\begin{align}\label{Upsilon44-2}
\Upsilon_{\alpha,4}^\pm(\lambda) = \lambda (-\log \lambda)^{3/2} \Lambda^\pm(\lambda) \qquad \text{for all } \alpha \in \{0,1,2,3,4\}.
\end{align}

If $Q_4^0 \neq 0$ and $Q_{42}^1 \neq 0$, the slowest-decaying term in $\Upsilon_{\alpha,4}^\pm(\lambda)$ is
\begin{align*}
\mathcal{D}_{\alpha,j}^{\pm}(\lambda)r_{j,4}^\pm(\lambda) \mathcal{D}_{4,4}^{\pm}(\lambda) &= \lambda (-\log \lambda)^{7/2} \Lambda^\pm(\lambda) \quad \text{for } \alpha \in \{0,1,2,3\},\\
\mathcal{D}_{4,4}^{\pm}(\lambda) r_{4,4}^\pm(\lambda) \mathcal{D}_{4,4}^{\pm}(\lambda) &= \lambda (-\log \lambda)^{11/2} \Lambda^\pm(\lambda)\quad \text{for } \alpha = 4,
\end{align*}
which implies that
\begin{align}\label{Upsilon44}
\Upsilon_{\alpha,4}^\pm(\lambda)= \lambda (-\log \lambda)^{7/2} \Lambda^\pm(\lambda) \quad \text{for } \alpha \in \{0,1,2,3\},\ \
	\Upsilon_{4,4}^\pm(\lambda)= \lambda (-\log \lambda)^{11/2} \Lambda^\pm(\lambda)\quad \text{for } \alpha = 4,
\end{align}
Thus, combining these estimates with the definition of $\sigma_\alpha(\lambda)$ and the expansions of $\mathcal{D}_{i,j}^{\pm}(\lambda)$ ($i,j \in \{0,4\}$) from \eqref{D00}, \eqref{D40D04D44_1}, and \eqref{D40D04D44_2}, we obtain the desired results for all $\mathcal{M}_{\alpha,4}^{\pm}(\lambda) = \sigma_\alpha(\lambda) \mathcal{D}_{\alpha,4}^\pm(\lambda) + \sigma_\alpha(\lambda) \Upsilon_{\alpha,4}^\pm(\lambda)$ with $0 \le \alpha \le 4$.

\textbf{(2) Analysis of $\mathcal{M}_{\alpha,\beta}^{\pm}(\lambda)$ and $\mathcal{M}_{\beta,\alpha}^{\pm}(\lambda)$ for $\beta \in \{5,6\}$ and $0 \leq \alpha \leq 6.$}    
	
	Firstly, consider the cases $0 \leq \alpha \leq 4.$    
	Here we restrict our attention to $\mathcal{M}_{\beta,\alpha}^{\pm}(\lambda)$, as the corresponding results for $\mathcal{M}_{\alpha,\beta}^{\pm}(\lambda)$ follow from identical arguments by symmetry.   
	From \eqref{inverse_D4}--\eqref{D-invert}, when $\beta \in \{5,6\}$, the operators $\mathcal{D}_{\alpha,\beta}^\pm(\lambda)$ and $\mathcal{D}_{\beta,\alpha}^\pm(\lambda)$ are non-vanishing only if $\alpha = \beta$. Thus, for $0 \le \alpha \le 4$ and $\beta \in \{5,6\}$,    
	\begin{align*}    
		\mathcal{M}_{\beta,\alpha}^{\pm}(\lambda) &= \sigma_\beta(\lambda) \sigma_\alpha(\lambda) \Upsilon_{\beta,\alpha}^\pm(\lambda).
	\end{align*}    
By the definition of $\sigma_\alpha(\lambda)$ in  \eqref{tab:sigma_alpha} and the estimate $\eqref{Upsilon}$,  we deduce that for $0 \le \alpha \le 4$ and $\beta \in \{5,6\}$,    
	\[  
\left\|\partial_\lambda^\ell \mathcal{M}_{\beta,\alpha}^\pm(\lambda)\right\|_{\mathbb{B}(L^2)} \lesssim \lambda^{-\ell}, \quad \ell = 0, 1, 2.  
\]  
	
	Secondly, consider the case $(5,5).$  
 Note that $\mathcal{D}_{\alpha,5}^\pm(\lambda)$ and $\mathcal{D}_{5,\alpha}^\pm(\lambda)$ are nonzero only for $\alpha = 5$. Combining this with \eqref{MM-inver} and the structure of matrix multiplication, we derive
	\begin{align*}
		\mathcal{M}_{5,5}^{\pm}(\lambda) &= (-\log \lambda)^{-1} \Big[  \mathcal{D}_{5,5}-\mathcal{D}_{5,5}r_{5,5}^{\pm}(\lambda)\mathcal{D}_{5,5}+\sum_{j,l=0}^{6}\mathcal{D}_{5,5}r_{5,j}^{\pm}(\lambda)\mathcal{D}_{j,l}^{\pm}(\lambda)r_{l,5}^{\pm}(\lambda)\mathcal{D}_{5,5}+ (-\log \lambda)^
		{-\frac{3}{2}} \Lambda^\pm(\lambda)
		\Big].
	\end{align*}
	By the expression  \eqref{partial}--\eqref{partial2} of $r_{i,j}^{\pm}(\lambda)$ and most $\mathcal{D}_{j,l}^{\pm}(\lambda)=0$ from \eqref{inverse_D4}--\eqref{D-invert}, it follows that
	\begin{align*}
		&\mathcal{D}_{5,5}r_{5,5}^{\pm}(\lambda)\mathcal{D}_{5,5}= (-\log \lambda)^{-1}\Big[a_3^\pm b_1^{-2}(Q_5 v G_6 v Q_5)^{-1}+b_1  \mathcal{D}_{5,5}(Q_5 v G_6^0 v Q_5) \mathcal{D}_{5,5}\Big]+ \lambda^2 (-\log \lambda)^{-1} \Lambda^\pm(\lambda),\\
		&\sum_{j,l=0}^{6} \mathcal{D}_{5,5} r_{5,j}^{\pm}(\lambda) \mathcal{D}_{j,l}^{\pm}(\lambda) r_{l,5}^{\pm}(\lambda) \mathcal{D}_{5,5} = b_1^2 (-\log \lambda)^{-1} \mathcal{D}_{5,5} (Q_5 v G_6^0 v Q_6) \mathcal{D}_{6,6} (Q_6 v G_6^0 v Q_5) \mathcal{D}_{5,5} \\
		& \ \ \ \ \ \ \ \ \ \ \ \ \ \ \ \ \ \ \ \ \ \ \  \ \ \ \ \ \ \ \ \ \ \ \ \ \  \ \  \ \ \ \ \ \ \  + (-\log \lambda)^{-2} \Lambda^\pm(\lambda).
	\end{align*}
	Therefore, $\mathcal{M}_{5,5}^{\pm}(\lambda)$ admits the representation
	\[
	\mathcal{M}_{5,5}^\pm(\lambda) =  \mathcal{A}_{5,5}(\lambda) + (\log \lambda)^{-2} \Gamma_{5,5}^{\pm}(\lambda),
	\]
	where
	\begin{align}
		\mathcal{A}_{5,5}(\lambda) &= -(\log \lambda)^{-1}\mathcal{D}_{5,5} + (\log \lambda)^{-2} \left[ -b_1 \mathcal{D}_{5,5} (Q_5 v G_6^0 v Q_5) \mathcal{D}_{5,5}  + b_1^2 \mathcal{D}_{5,5} (Q_5 v G_6^0 v Q_6) \mathcal{D}_{6,6} (Q_6 v G_6^0 v Q_5) \mathcal{D}_{5,5} \right], \nonumber\\
		\Gamma_{5,5}^{\pm}(\lambda) &= -a_3^\pm b_1^{-2} (Q_5 v G_6 v Q_5)^{-1} + (-\log \lambda)^{-1/2} \Lambda^\pm(\lambda).\label{Gamma55}
	\end{align}
	This establishes the desired bound for $\mathcal{M}_{5,5}^{\pm}(\lambda)$.
	
	Thirdly, consider the cases $(5,6)$ and $(6,5).$ 
	The analysis follows a similar pattern to the previous case, so we omit some details.
	For $\mathcal{M}_{5,6}^{\pm}(\lambda)$, we have
	\begin{align*}
		\mathcal{M}_{5,6}^{\pm}(\lambda) &= (-\log \lambda)^{-1/2} \Upsilon_{5,6}^\pm(\lambda) \\
		&= (-\log \lambda)^{-1/2} \Big[ -\mathcal{D}_{5,5} r_{5,6}^{\pm}(\lambda) \mathcal{D}_{6,6} + \mathcal{D}_{5,5} r_{5,5}^{\pm}(\lambda) \mathcal{D}_{5,5} r_{5,6}^{\pm}(\lambda) \mathcal{D}_{6,6} \\ 
		&\quad - \mathcal{D}_{5,5} r_{5,6}^{\pm}(\lambda) \mathcal{D}_{6,6} r_{6,5}^{\pm}(\lambda) \mathcal{D}_{5,5} r_{5,6}^{\pm}(\lambda) \mathcal{D}_{6,6} + (-\log \lambda)^{-2} \Lambda^\pm(\lambda) \Big] \\
		&=  \mathcal{A}_{5,6}(\lambda) + (\log \lambda)^{-2} \Gamma_{5,6}^{\pm}(\lambda),
	\end{align*}
	where
	\begin{align}
		\mathcal{A}_{5,6}(\lambda) &=(\log \lambda)^{-1} b_1 \mathcal{D}_{5,5} (Q_5 v G_6^0 v Q_6) \mathcal{D}_{6,6} + (\log \lambda)^{-2} \Big[ b_1^2 \mathcal{D}_{5,5} (Q_5 v G_6^0 v Q_5) \mathcal{D}_{5,5} (Q_5 v G_6^0 v Q_6) \mathcal{D}_{6,6} \nonumber\\
		&\quad -b_1^3 \mathcal{D}_{5,5} (Q_5 v G_6^0 v Q_6) \mathcal{D}_{6,6} (Q_6 v G_6^0 v Q_5) \mathcal{D}_{5,5} (Q_5 v G_6^0 v Q_6) \mathcal{D}_{6,6} \Big], \nonumber\\
		\Gamma_{5,6}^{\pm}(\lambda) &= -a_3^\pm \mathcal{D}_{5,5} (Q_5 v G_6^0 v Q_6) \mathcal{D}_{6,6} + (-\log \lambda)^{-1/2} \Lambda^\pm(\lambda).\label{Gamma56}
	\end{align}
	This establishes the desired result for $\mathcal{M}_{5,6}^{\pm}(\lambda)$.
	Similarly, for $\mathcal{M}_{6,5}^{\pm}(\lambda)$, we obtain
	\begin{align*}
		\mathcal{M}_{6,5}^{\pm}(\lambda) &= \mathcal{A}_{6,5}(\lambda) + (\log \lambda)^{-2} \Gamma_{6,5}^{\pm}(\lambda),
	\end{align*}
	where
	\begin{align}
		\mathcal{A}_{6,5}(\lambda) &= (\log \lambda)^{-1}b_1 \mathcal{D}_{6,6} (Q_6 v G_6^0 v Q_5) \mathcal{D}_{5,5} + (\log \lambda)^{-2} \Big[ b_1^2 \mathcal{D}_{6,6} (Q_6 v G_6^0 v Q_5) \mathcal{D}_{5,5} (Q_5 v G_6^0 v Q_5) \mathcal{D}_{5,5} \nonumber\\
		&\quad -b_1^3 \mathcal{D}_{6,6} (Q_6 v G_6^0 v Q_5) \mathcal{D}_{5,5} (Q_5 v G_6^0 v Q_6) \mathcal{D}_{6,6} (Q_6 v G_6^0 v Q_5) \mathcal{D}_{5,5} \Big],\nonumber \\
		\Gamma_{6,5}^{\pm}(\lambda) &= -a_3^\pm \mathcal{D}_{6,6} (Q_6 v G_6^0 v Q_5) \mathcal{D}_{5,5} + (-\log \lambda)^{-1/2} \Lambda^\pm(\lambda).\label{Gamma65}
	\end{align}
	This completes the analysis for $\mathcal{M}_{6,5}^{\pm}(\lambda)$.
	
Finally,
	it remains to analyze the case $\mathcal{M}_{6,6}^{\pm}(\lambda)$. We have
	\begin{align*}
		\mathcal{M}_{6,6}^{\pm}(\lambda) &=\mathcal{D}_{6,6}+ \Upsilon_{6,6}^\pm(\lambda) \\
	&=\mathcal{D}_{6,6}+\mathcal{D}_{6,6} r_{6,5}^{\pm}(\lambda) \mathcal{D}_{5,5}r_{5,6}^{\pm}(\lambda) \mathcal{D}_{6,6} - \mathcal{D}_{6,6} r_{6,5}^{\pm}(\lambda) \mathcal{D}_{5,5} r_{5,5}^{\pm}(\lambda) \mathcal{D}_{5,5}r_{5,6}^{\pm}(\lambda) \mathcal{D}_{6,6} \\
		&\quad +\mathcal{D}_{6,6} r_{6,5}^{\pm}(\lambda) \mathcal{D}_{5,5} r_{5,6}^{\pm}(\lambda) \mathcal{D}_{6,6} r_{6,5}^{\pm}(\lambda) \mathcal{D}_{5,5} r_{5,6}^{\pm}(\lambda) \mathcal{D}_{6,6}+ (-\log \lambda)^{-5/2} \Lambda^\pm(\lambda)  \\
		&=  \mathcal{A}_{6,6}(\lambda) + (\log \lambda)^{-2} \Gamma_{6,6}^{\pm}(\lambda),
	\end{align*}
	where
	\begin{align}
		\mathcal{A}_{6,6}(\lambda) &=\mathcal{D}_{6,6}+(-\log \lambda)^{-1}
		b_1^2 \mathcal{D}_{6,6} (Q_6 v G_6^0 v Q_5) \mathcal{D}_{5,5} (Q_5 v G_6^0 v Q_6) \mathcal{D}_{6,6}\nonumber \\
		&+(-\log \lambda)^{-2} \Big[-b_1^3\mathcal{D}_{6,6} (Q_6 v G_6^0 v Q_5) \mathcal{D}_{5,5} (Q_5 v G_6^0 v Q_5) \mathcal{D}_{5,5}(Q_5 v G_6^0 v Q_6) \mathcal{D}_{6,6}\nonumber\\
	&+b_1^4\mathcal{D}_{6,6}\left((Q_6 v G_6^0 v Q_5) \mathcal{D}_{5,5} (Q_5 v G_6^0 v Q_6) \mathcal{D}_{6,6}\right)^2 \Big],\nonumber\\
		\Gamma_{6,6}^{\pm}(\lambda) &= a_3^{\pm }b_1 \mathcal{D}_{6,6} (Q_6 v G_6^0 v Q_5) \mathcal{D}_{5,5}(Q_5v G_6^0 v Q_6) \mathcal{D}_{6,6}  + (-\log \lambda)^{-1/2} \Lambda^\pm(\lambda).\label{Gamma66}
	\end{align}
	This establishes the desired result for $\mathcal{M}_{6,6}^{\pm}(\lambda)$.
\end{proof}


\section{The proof of Lemmas \ref{lemma_projection} and \ref{oscillatory}}\label{section 8}
This section is primarily devoted to completing the proofs of Lemma~\ref{lemma_projection} and Lemma~\ref{oscillatory}. Moreover, we also provide the asymptotic expansion of two special oscillatory integrals.

\subsection{Proof of Lemma \ref{lemma_projection}}
The proof of Lemma \ref{lemma_projection} relies heavily on the directional Taylor expansion provided by Lemma \ref{taylor} (see \cite[Lemma 3.5]{LSY21}), which allows us to systematically exploit the  cancellation properties of the projections $Q_{\alpha}$.     

\begin{lemma}[{\cite[Lemma 3.5]{LSY21}}]\label{taylor}
Assume that $\lambda>0$,  $\tilde{x}=(-x_2,x_1) \in \mathbb{R}^2$, $ \rho \equiv \rho(x)=\frac{x}{|x|}$ for $x\neq 0$ and $\rho(x)=0$ for $x=0$.
Let $\theta\in [0,1]$ and $ |x| \cos\xi=|x| \cos\xi(x, y, \theta )= \langle x, \rho(y-\theta x)\rangle$,
 where $\xi\equiv \xi(x, y, \theta )$ is the angle between the vectors $x$ and $y-\theta x$.
Then
\begin{enumerate}
\item[(i)] If $F \in C^{1}(\mathbb{R})$, then
$$
F(\lambda|x-y|) = F(\lambda|y|) - \lambda |x| \int_{0}^{1} \mathcal{L}_{1}[F](\lambda|y-\theta x|, \xi) \, d\theta.
$$

\item[(ii)] If $F \in C^{2}(\mathbb{R})$ with $F'(0) = 0$, then
$$
F(\lambda|x-y|)
= F(\lambda|y|) - \lambda F'(\lambda|y|) \langle \rho(y), x \rangle
 + \lambda^{2}|x|^{2} \int_{0}^{1} (1-\theta) \mathcal{L}_{2}[F](\lambda|y-\theta x|, \xi) \, d\theta.
$$

\item[(iii)] If $F \in C^{3}(\mathbb{R})$ with $F'(0)=F''(0)=0$, then
$$
\begin{aligned}
F(\lambda|x-y|)
&= F(\lambda|y|) - \lambda F'(\lambda|y|) \langle \rho(y), x \rangle
 + \frac{\lambda^{2}}{2} \Bigl( \frac{F'(\lambda|y|)}{\lambda|y|} \langle \rho(y), \tilde{x}\rangle^{2}
+ F''(\lambda|y|) \langle \rho(y), x\rangle^{2} \Bigr) \\
&\quad - \frac{\lambda^{3}|x|^{3}}{2} \int_{0}^{1} (1-\theta)^{2} \mathcal{L}_{3}[F](\lambda|y-\theta x|, \xi) \, d\theta,
\end{aligned}
$$
where $\langle \rho(y), \tilde{x}\rangle^{2} = |x|^{2} - \langle \rho(y), x\rangle^{2}$.

\item[(iv)] If $F \in C^{4}(\mathbb{R})$ with $F^{(k)}(0)=0$ for $k=1,2,3$, then
$$
\begin{aligned}
F(\lambda|x-y|)
&= F(\lambda|y|) - \lambda F'(\lambda|y|) \langle \rho(y), x \rangle  + \frac{\lambda^{2}}{2} \Bigl( \frac{F'(\lambda|y|)}{\lambda|y|} \langle \rho(y), \tilde{x}\rangle^{2}
+ F''(\lambda|y|) \langle \rho(y), x\rangle^{2} \Bigr) \\
&\quad - \frac{\lambda^{3}}{3!} \Bigl[ \Bigl( -\frac{F'(\lambda|y|)}{\lambda^{2}|y|^{2}}
+ \frac{F''(\lambda|y|)}{\lambda|y|} \Bigr) 3\langle \rho(y), x\rangle \langle \rho(y), \tilde{x}\rangle^{2}
 + F^{(3)}(\lambda|y|) \langle \rho(y), x\rangle^{3} \Bigr] \\
&\quad + \frac{\lambda^{4}|x|^{4}}{3!} \int_{0}^{1} (1-\theta)^{3} \mathcal{L}_{4}[F](\lambda|y-\theta x|, \xi) \, d\theta.
\end{aligned}
$$
\end{enumerate}
Here $F^{(\ell)}$ denotes the $\ell$-th derivative of $F$, and define
\begin{align*}
\mathcal{L}_{1}[F](r,\xi) &:= F'(r)\cos\xi, \ \ \
\mathcal{L}_{2}[F](r,\xi) := F''(r)\cos^{2}\xi + \frac{F'(r)}{r}\sin^{2}\xi, \\
\mathcal{L}_{3}[F](r,\xi) &:= \Bigl( -\frac{F'(r)}{r^{2}} + \frac{F''(r)}{r} \Bigr) 3\cos\xi\sin^{2}\xi + F^{(3)}(r)\cos^{3}\xi, \\
\mathcal{L}_{4}[F](r,\xi) &:= \Bigl( \frac{F'(r)}{r^{3}} - \frac{F''(r)}{r^{2}} \Bigr)(15\cos^{2}\xi\sin^{2}\xi - 3\sin^{4}\xi)
 + \frac{F^{(3)}(r)}{r} 6\cos^{2}\xi\sin^{2}\xi + F^{(4)}(r)\cos^{4}\xi.
\end{align*}
\end{lemma}
However, the proof of Lemma \ref{lemma_projection} below  involves the logarithmic profile
\[
	\widetilde{G}_1(p):=b_0p^2\log p\,\chi_1(p),
\]
which is  \(C^1\bigl([0, \infty)\bigr)\) but not \(C^2\bigl([0, \infty)\bigr)\) at the origin and hence does not satisfy the
hypotheses of Lemma~\ref{taylor}\textup{(ii)}. The following remark
shows that the conclusion of that part nevertheless remains valid for
\(\widetilde{G}_1(p)\).
\begin{remark}\label{remark:taylor-log}
{\rm	Although
$	\widetilde{G}_1(p)$
	does not belong to \(C^2([0,\infty))\), the conclusion of
	Lemma~\ref{taylor}\textup{(ii)} remains valid for
	\(F(\lambda|x-y|)=\widetilde{G}_1(\lambda|x-y|)\) whenever \(x\neq0\), with the integrand
	understood almost everywhere in \(\theta\).

    
	Indeed, 
	for \(\varepsilon>0\), set
	$
	r_\varepsilon(\theta)
	:=
	\sqrt{\varepsilon^2+|y-\theta x|^2}$ and 
	$h_\varepsilon(\theta)
	:=
	\widetilde{G}_1\bigl(\lambda r_\varepsilon(\theta)\bigr).
	$
 Since \(r_\varepsilon(\theta)>0\),
	\(h_\varepsilon\in C^2([0,1])\) and  \(\chi_1\in C_c^\infty(\mathbb{R})\), a direct calculation gives
	\[
	|h_\varepsilon''(\theta)|
	\lesssim
	\lambda^2|x|^2
	\Big(
	\bigl|\widetilde{G}_1''(\lambda r_\varepsilon(\theta))\bigr|
	+
	\frac{
		\bigl|\widetilde{G}_1'(\lambda r_\varepsilon(\theta))\bigr|
	}{
		\lambda r_\varepsilon(\theta)
	}
	\Big)\lesssim \lambda^2|x|^2\Bigl(1+\big|\log\bigl(\lambda r_\varepsilon(\theta)\bigr)\big|\chi_1\bigl(\lambda r_\varepsilon(\theta)\bigr)\Bigr).
	\]
	If the segment
	\(\{y-\theta x:0\le\theta\le1\}\) meets the origin, there is a
	unique \(\theta_0\in[0,1]\) such that \(y-\theta_0x=0\), and
	\[
	|y-\theta x|=|x|\,|\theta-\theta_0|.
	\]
	Consequently, for \(0<\varepsilon\le1\) and almost every
	\(\theta\in[0,1]\),
	\[
	|h_\varepsilon''(\theta)|
	\lesssim C_{\lambda,x}
	\bigl(1+|\log|\theta-\theta_0||\bigr)\in L_\theta^1\bigl([0,1]\bigr).
	\]
	If the segment does not meet
	the origin, then \(|y-\theta x|\) is bounded away from zero, and
	\(h_\varepsilon''\) is uniformly bounded for small \(\varepsilon\).
	Moreover, for almost every \(\theta\),
	\[
	h_\varepsilon''(\theta)
	\longrightarrow
	\lambda^2|x|^2
	\mathcal L_2[\widetilde{G}_1]
	\bigl(\lambda|y-\theta x|,\xi\bigr)
	\qquad\text{as }\varepsilon\to0.
	\]
	Although this pointwise limit  fails at the single point
	\(\theta=\theta_0\), this is immaterial for the integral remainder.
	Lebesgue's dominated convergence theorem  permits passing
	to the limit \(\varepsilon\to0\), proving the claimed identity.}
\end{remark}





Now, we start to use Lemma \ref{taylor} to prove Lemma \ref{lemma_projection}.






\begin{proof}
	We organize the proof according to the order of the cancellation
	properties of the projections.
	
	To streamline the  analysis, we first formulate a universal bound for the derivatives of the kernel. For any smooth profile $\mathcal{G}$ on $\mathbb{R}$ satisfying $|\partial_p^\ell \big(e^{\mp ip}\mathcal{G}(p)\big)| \lesssim \langle p \rangle^{-1/2}|p|^{-\ell}$,  applying Leibniz's rule together with the   inequality $\langle \lambda(\theta x-y) \rangle^{-1/2} \lesssim \langle \lambda y \rangle^{-1/2} \langle x \rangle^{1/2}$ ($\theta \in [0,1],  0<\lambda\ll1$) yields:
\begin{align}
	\left|\partial_\lambda^\ell \left( e^{\mp i\lambda|y|} \mathcal{G}(\lambda|\theta x-y|) \right)\right|
	&= \left|\partial_\lambda^\ell \left[ e^{\pm i\lambda(|\theta x-y|-|y|)} \left( e^{\mp i\lambda|\theta x-y|} \mathcal{G}(\lambda|\theta x-y|) \right) \right]\right| \nonumber \\
	&\lesssim \sum_{\ell_1+\ell_2=\ell} |x|^{\ell_1} \lambda^{-\ell_2} \langle \lambda(\theta x-y) \rangle^{-1/2} \lesssim \lambda^{-\ell} \langle \lambda y \rangle^{-1/2} \langle x \rangle^{\ell+1/2}. \label{eq:general_oscillatory_bound}
\end{align}
	The decay assumption \eqref{condition} ensures that all additional
	polynomial weights in $x$ below
	are integrable against $v$ in $L^2_x$. We shall use
	\eqref{eq:general_oscillatory_bound} repeatedly without further
	comment.
	
	\medskip
	
	%\noindent
	\textit{Proof of part~{\rm (i)}.}
	Recall from \eqref{def:F_pm} that $R^\pm_0(\lambda^4)(x,y) = \lambda^{-2} F_\pm(\lambda|x-y|)$, where 
\begin{align}\label{Fpm}
	{F}_\pm(p)=&\Bigl(a_0^\pm+b_0p^2\log p+a_1^\pm p^2
	+a_2^\pm p^4+O\bigl(p^{5}\bigr)\Bigr)\chi_1(p)+\frac{i}{8}\Bigl(\pm e^{\pm ip} w_\pm(p)-e^{-p}w_+(ip)\Bigr)\chi_2(p).
\end{align}
Observe that $F_\pm(p)$ comprising both low-frequency polynomial expansions and high-frequency oscillations satisfies:
\[
\bigl|\partial_p^\ell \bigl(e^{\mp ip} F_\pm(p)\bigr)\bigr| + \bigl|\partial_p^\ell \bigl(e^{\mp ip} F_\pm'(p)\bigr)\bigr| \lesssim \langle p \rangle^{-1/2} |p|^{-\ell}, \quad \ell = 0,1,2.
\]
For $\alpha = 0$, defining $\Omega_{0,\pm}(\lambda,x,y) := \left(Q_0 v F_\pm(\lambda|\cdot-y|)\right)(x)$ and directly applying \eqref{eq:general_oscillatory_bound} gives:
\begin{align*}
	\left\| \partial_\lambda^\ell \left(e^{\mp i \lambda|y|} \Omega_{0,\pm}(\lambda, x, y)\right) \right\|_{L^2_x}
	\lesssim \lambda^{-\ell} \langle \lambda y \rangle^{-1/2} \left\| v(x) \langle x \rangle^{\ell+1/2} \right\|_{L^2_x} \lesssim \lambda^{-\ell} \langle \lambda y \rangle^{-1/2}.
\end{align*}
For $\alpha \ge 1$, we exploit Lemma~\ref{taylor}(i) to write $$F_\pm(\lambda|x-y|) = F_\pm(\lambda|y|) - \lambda |x| \int_0^1 F'_\pm(\lambda|\theta x-y|) \cos\xi \, d\theta.$$
Since $Q_\alpha v = 0$ for $\alpha \ge 1$ (see Remark \ref{cancellationQ}), the leading term $F_\pm(\lambda|y|)$ is annihilated, yielding:
\begin{align*}
	\left(Q_\alpha v R_0^{\pm}(\lambda^4)(\cdot, y)\right)(x)
	&= -\lambda^{-1} Q_\alpha \Big(|\cdot| v \int_0^1 F'_\pm(\lambda| \theta \cdot-y|)\cos\xi \, d\theta \Big)(x) := \lambda^{-1} \Omega_{\alpha,\pm}(\lambda, x, y).
\end{align*}
Applying \eqref{eq:general_oscillatory_bound} to $F'_\pm$ guarantees $\|\partial_\lambda^\ell (e^{\mp i \lambda|y|} \Omega_{\alpha,\pm})\|_{L^2_x} \lesssim \lambda^{-\ell} \langle \lambda y \rangle^{-1/2}$, concluding Part (i).

	\medskip
	
	%\noindent
	\textit{Proof of part~{\rm (ii)}.}
	By Remark \ref{cancellationQ},
	$
	Q_\alpha v
	=
	Q_\alpha(z_jv)
	=
	0
	$ for $ j=1,2$ and $ 2\leq\alpha\leq6.$
 Lemma \ref{taylor}(ii) and Remark \ref{remark:taylor-log} with $z\neq0$ give
	\begin{align*}
		&
		Q_\alpha
		\bigl(
		vF_\pm(\lambda|z-y|)
		\bigr)(x)=
		\lambda^2
		Q_\alpha
		\Bigl(
		|z|^2v
		\int_0^1
		(1-\theta)
		\mathcal{L}_2[F_\pm]
		\bigl(
		\lambda|y-\theta z|,
		\xi
		\bigr)
		\,d\theta
		\Bigr)(x),
		%\label{J-representation}
	\end{align*}
	where \(Q_\alpha\) acts on the \(z\)-variable and outputs a function of \(x\). Modifying the input function on the
	null set \(\{z=0\}\) does not alter its image  in \(L_x^2\).
   Thus, define
	$
	\bigl(Q_\alpha vR_0^\pm(\lambda^4)\bigr)(x,y)
	:=
	\mathcal{J}_{\alpha,\pm}(\lambda,x,y),
	$
	where
\begin{equation}\label{def-J}
		\mathcal{J}_{\alpha,\pm}(\lambda,x,y)
		=
		Q_\alpha
		\Bigl(
		|z|^2v
		\int_0^1
		(1-\theta)
		\mathcal{L}_2[F_\pm]
		\bigl(
		\lambda|y-\theta z|,
		\xi
		\bigr)
		\,d\theta
		\Bigr)(x).
	\end{equation}
	The expansion \eqref{Fpm} of $F_\pm$ implies
	\begin{equation}\label{L2-decomposition}
		\mathcal{L}_2[F_\pm](p,\xi)
		=
		2b_0(\log p)\chi_1(p)
		+
		\Theta_\pm(p,\xi),
	\end{equation}
		where $\Theta^\pm$ consolidates the bounded low-frequency derivatives and the high-frequency oscillatory components, which  satisfies the bound \eqref{eq:general_oscillatory_bound}. 
	It therefore remains only to estimate the logarithmic term. For
 $\ell=0,1,2$,
	\begin{align}
		&
		\left|
		\partial_\lambda^\ell
		\Bigl[
		e^{\mp i\lambda|y|}
		\log(\lambda |y-\theta x|)
		\chi_1(\lambda |y-\theta x|)
		\Bigr]
		\right|
	\lesssim
		\lambda^{-\ell}
		\langle\lambda y\rangle^{-1/2}
		\langle x\rangle^{\ell+1/2}
		\Bigl(
		|\log\lambda|
		+
		|y-\theta x|^{-1/2}
		\Bigr).
		\label{log-profile-bound}
	\end{align}
	Moreover, the decay assumption \eqref{condition} yields
	\begin{equation}\label{weighted-singular-bound}
		\left\|
		v(x)|x|^2
		\langle x\rangle^{\ell+1/2}
		|y-\theta x|^{-1/2}
		\right\|_{L^2_x}
		\lesssim
		\theta^{-1/2}
		\left\langle\theta^{-1}y\right\rangle^{-1/2},
		\qquad \theta\in (0,1].
	\end{equation}
	Combining with  \eqref{def-J}--\eqref{weighted-singular-bound}, we obtain
	\begin{align*}
		&
		\left\|
		\partial_\lambda^\ell
		\Bigl(
		e^{\mp i\lambda|y|}
		\mathcal{J}_{\alpha,\pm}(\lambda,\cdot,y)
		\Bigr)
		\right\|_{L^2}
	\lesssim
		\lambda^{-\ell}
		|\log\lambda|
		\langle\lambda y\rangle^{-1/2}
		\int_0^1
		(1-\theta)\theta^{-1/2}
		\,d\theta
	\lesssim
		\lambda^{-\ell}
		|\log\lambda|
		\langle\lambda y\rangle^{-1/2}.
	\end{align*}
	This proves part~{\rm (ii)}.
	
	\medskip
	
	%\noindent
	\textit{Proof of part~{\rm (iii)}.}
	Define
	$
	\widetilde{G}_j(r)
	:=
	b_0\chi_j(r)r^2\log r
$ for $	 j=1,2.$
	Note that
	$$
	\lambda^{-2}
	\Bigl(
	\widetilde{G}_1(\lambda r)
	+
	\widetilde{G}_2(\lambda r)
	\Bigr)
	=
	b_0r^2\log r
	+
	b_0(\log\lambda)r^2.
	$$
	Therefore,
	\begin{align}
		R_0^\pm(\lambda^4)(x,y)
		={}&
		b_0G_2^0(x,y)
		+
		b_0(\log\lambda)G_2(x,y)
	-
		\lambda^{-2}
		\widetilde{G}_2(\lambda|x-y|)
		+
		\mathcal{R}_\pm(\lambda,x,y),
		\label{basic-four-term-decomposition}
	\end{align}
	where
	$
		\mathcal{R}_\pm(\lambda,x,y)
		:=
		R_0^\pm(\lambda^4)(x,y)
		-
		\lambda^{-2}
		\widetilde{G}_1(\lambda|x-y|).
$
	We focus on  the four terms in
	\eqref{basic-four-term-decomposition} separately.
	
	\medskip
	
	%\noindent
	\textit{Step 1: Deal with $(b_0QvG_2^0)(x,y)$.}
We decompose $b_0G_2^0(x,y)$ as 
	$
	b_0G_2^0(x,y)
	=
	\widetilde{G}_1(|x-y|)
	+
	\widetilde{G}_2(|x-y|).
	$
Note that  the compactly supported piece $\widetilde{G}_1$ is uniformly bounded,  trivially contributing ${O}(1)$ to the $L^2_x$ norm. Then  only the
	term
	$
	Q
	\bigl(
	v\widetilde{G}_2(|\cdot-y|)
	\bigr)
	$
	requires consideration.
	
First, for
	$
	Q
	\in
	\{Q_\alpha:2\leq\alpha\leq4\}
	\cup
	\{Q_4^0\},
	$
	the cancellations
	$
	Qv=Q(x_jv)=0
	$ ($j=1,2$)
	and Lemma \ref{taylor}(ii) give
	\begin{align}
		&
		Q
		\bigl(
		v\widetilde{G}_2(|\cdot-y|)
		\bigr)(x)
	=
		Q
		\Bigl(
		|\cdot|^2v
		\int_0^1
		(1-\theta)
		\mathcal{L}_2[\widetilde{G}_2]
		\bigl(
		|y-\theta\cdot|,
		\xi
		\bigr)
		\,d\theta
		\Bigr)(x).
		\label{QvG20-1}
	\end{align}
	 Moreover,
	the second-order derivative produces a logarithmic growth:
	$$
	\left|\mathcal{L}_2[\widetilde{G}_2](|\theta x-y|,\xi)\right| \lesssim 1 + |\log(|\theta x-y|)| \chi_2(|\theta x-y|) \lesssim \log(2+|x|)\log(2+|y|).
	$$
	Upon evaluating the $L^2_x$-norm against $v(x)$, the $\log(2+|x|)$ factor is absorbed, resulting in \begin{align}\label{static-log-bound}
		%\label{Q_4G}
		\|(Q v \widetilde{G}_2)(\cdot,y)\|_{L^2} \lesssim \log(2+|y|), \quad \text{for } 	Q
		\in
		\{Q_\alpha:2\leq\alpha\leq4\}
		\cup
		\{Q_4^0\}.\end{align}
	
	
	
	
	
	
	Next, for
	$
	Q
	\in
	\{Q_4^1,Q_{41}^1,Q_{42}^1,Q_5,Q_6\},
	$
	in addition to the zeroth- and first-order cancellations, these
	projections satisfy
	$
	Q\bigl(|x|^2v\bigr)=0.
	$ Consequently, 
	applying Lemma \ref{taylor}(iii), 
	we obtain
	\begin{align}
		Q
		\bigl(
		v\widetilde{G}_2(|\cdot-y|)
		\bigr)(x)
=&
		\frac{1}{2}
		\Bigl(
		\widetilde{G}_2''(|y|)
		-
		\frac{\widetilde{G}_2'(|y|)}{|y|}
		\Bigr)
		Q
		\Bigl(
		v\langle\rho(y),\cdot\rangle^2
		\Bigr)(x)
	\nonumber\\
	&	-
		\frac{1}{2}
		Q
		\Bigl(
		|\cdot|^3v
		\int_0^1
		(1-\theta)^2
		\mathcal{L}_3[\widetilde{G}_2]
		\bigl(
		|y-\theta\cdot|,
		\xi
		\bigr)
		\,d\theta
		\Bigr)(x).
		\label{QvG20-2}
	\end{align}
	Since
	$
	\big|
	\widetilde{G}_2''(r)
	-
	r^{-1}\widetilde{G}_2'(r)
	\big|
	\lesssim1,
	$
	and
	$
	\big|
	\mathcal{L}_3[\widetilde{G}_2](r,\xi)
	\big|
	\lesssim
	\langle r\rangle^{-1},
	$
	it follows that
	\begin{equation}\label{static-uniform-bound}
		\Big\|
		Q
		\bigl(
		v\widetilde{G}_2(|\cdot-y|)	\bigr)(x)
		\Big\|_{L_x^2}
		\lesssim1,\quad \text{for}	\ Q
		\in
		\{Q_4^1,Q_{41}^1,Q_{42}^1,Q_5,Q_6\}.
	\end{equation}
	Together with the bounded contribution of $\widetilde{G}_1$,
	\eqref{static-log-bound} and \eqref{static-uniform-bound} prove
	\eqref{lemma_projection_G}.
	
	\medskip
	
	%\noindent
	\textit{Step 2: Deal with $-\lambda^{-2}Q(v	\widetilde{G}_2(\lambda|\cdot-y|))(x)$.}
	Applying the same Taylor expansions as in
	\eqref{QvG20-1} and \eqref{QvG20-2}, now to 
	$\widetilde{G}_2(\lambda|\cdot-y|)$, gives
	\begin{equation*}
    %\label{scaled-low-order-bound}
	\left\|
		\partial_\lambda^\ell\mathcal{T}(\lambda,x,y)
		\right\|_{L_x^2}=	\left\|
		\partial_\lambda^\ell\left(
		-\lambda^{-2}
	Q
	\bigl(
	v\widetilde{G}_2(\lambda|\cdot-y|)
	\bigr)(x)\right)
		\right\|_{L_x^2}
		\lesssim
		\begin{cases}
			\lambda^{-\ell}\log(2+|y|),
			&
			Q\in
			\{Q_\alpha:2\leq\alpha\leq4\}
			\cup
			\{Q_4^0\},
			\\
			\lambda^{-\ell},
			&
			Q\in
			\{Q_4^1,Q_{41}^1\}.
		\end{cases}
	\end{equation*}
	Here $\mathcal{T}(\lambda, x,y ):=-\lambda^{-2}\bigl(Q \bigl(v\widetilde{G}_2(\lambda|\cdot-y|)\bigr)\bigr)(x)$ with $(\mathcal{T}, Q)\in\{ (\mathcal{T}_\alpha, Q_\alpha)|_{2\le\alpha\le4}, (\mathcal{T}_4^0, Q_4^0), (\mathcal{T}_4^1, Q_4^1), (\mathcal{T}_4^{1,1}, Q_{41}^1) \}.$
	
	
	
	
	
	For
	$
	Q\in\{Q_{42}^1,Q_5,Q_6\},
	$
	define
	$$
	m(Q)
	=
	\begin{cases}
		1,
		&
		Q=Q_{42}^1\ \text{or}\ Q=Q_5,
		\\
		2,
		&
		Q=Q_6.
	\end{cases}
	$$
	The projections $Q_{42}^1$ and $Q_5$ annihilate all quadratic
	moments, whereas $Q_6$ also annihilates all cubic moments.
	Accordingly, Lemma \ref{taylor}(iii)--(iv) gives
	\begin{align*}
	-\frac{1}{\lambda^2}Q\bigl(v	\widetilde{G}_2(\lambda|\cdot-y|)\bigr)(x)
		&=	\frac{(-1)^{m(Q)+1}\lambda^{m(Q)}}{(m(Q)+1)!}
		Q
		\Bigl(
		|\cdot|^{m(Q)+2}v
		\int_0^1
		(1-\theta)^{m(Q)+1}
		\mathcal{L}_{m(Q)+2}[\widetilde{G}_2]
		\bigl(
		\lambda|y-\theta\cdot|,
		\xi
		\bigr)
		\,d\theta
		\Bigr)(x)\\
		&:=\lambda^{m(Q)}\mathcal{T}(\lambda,x,y).
	%	\label{scaled-G2-high-cancellation}
	\end{align*}
Here $(\mathcal{T},Q)\in\{(\mathcal{T}_4^{1,2}, Q_{42}^1), (\mathcal{T}_5, Q_5), (\mathcal{T}_6, Q_6)\}.$
	The bound $|\partial_\lambda^\ell \mathcal{L}_{m(Q)+2}[\widetilde{G}_2](\lambda| \theta \cdot-y|, \xi)| \lesssim  \lambda^{-\ell}$  implies
	$$
	\left\|
	\partial_\lambda^\ell
	\mathcal{T}(\lambda,\cdot,y)
	\right\|_{L^2}
	\lesssim
	\lambda^{-\ell},
	\qquad
	\mathcal{T}
	\in
	\{\mathcal{T}_4^{1,2},\mathcal{T}_5,\mathcal{T}_6\}.
	$$
	Hence we obtain
	\eqref{lemma_projection_T}.
	
	\medskip
	
	%\noindent
	\textit{Step 3: Deal with $	b_0(\log\lambda)(QvG_2)(x,y)$.}
	The cancellations
	$
	Qv=Q(x_jv)=0
	$ ($ j=1,2$)
	yield
	\[
		(QvG_2)(x,y)
		=
		Q
		\Bigl(
		v(\cdot)|\cdot-y|^2
		\Bigr)(x)=
		Q
		\bigl(
		|\cdot|^2v
		\bigr)(x).
	\]
	Hence, for
	$
	Q
	\in
	\{Q_\alpha:2\leq\alpha\leq4\}
	\cup
	\{Q_4^0\},
	$
	the logarithmic contribution is precisely
	$
	b_0(\log\lambda)
	Q\bigl(|\cdot|^2v\bigr)(x).
	$
	
	For
	$
	Q
	\in
	\{Q_4^1,Q_{41}^1,Q_{42}^1,Q_5,Q_6\},
	$
	the additional cancellation
	$
	Q\bigl(|x|^2v\bigr)=0
	$
	shows that this term vanishes identically, i.e., $	b_0(\log\lambda)(QvG_2)(x,y)=0.$
	
	\medskip
	
	%\noindent
	\textit{Step 4: Deal with $Qv	\mathcal{R}_\pm(\lambda,x,y)$.}
	Define
	$
	F_1^\pm(p)
	:=
	F_\pm(p)-\widetilde{G}_1(p).
	$
	Then
	$
	\mathcal{R}_\pm(\lambda,x,y)
	=
	\lambda^{-2}
	F_1^\pm(\lambda|x-y|)
	$ and
		\begin{align*}
		F_1^\pm(p) &= \left(a_0^\pm  + a_1^\pm p^2 + a_2^\pm p^4 + O(p^5)\right)\chi_1(p) + \Big(\pm \frac{i}{8} e^{\pm ip} w_\pm(p) - \frac{i}{8} e^{-p} w_+(ip)\Big)\chi_2(p).
	\end{align*}
	
	For
	$
	Q
	\in
	\{Q_\alpha:2\leq\alpha\leq4\}
	\cup
	\{Q_4^0,Q_4^1,Q_{41}^1\},
	$
	Lemma \ref{taylor}(ii) gives
	\[
		Q
		\bigl(
		v\mathcal{R}_\pm(\lambda,\cdot,y)
		\bigr)(x)
	=
		Q
		\Bigl(
		|\cdot|^2v
		\int_0^1
		(1-\theta)
		\mathcal{L}_2[F_1^\pm]
		\bigl(
		\lambda|y-\theta\cdot|,
		\xi
		\bigr)
		\,d\theta
		\Bigr)(x).
	\]
	These kernels $Q
	\bigl(
	v\mathcal{R}_\pm(\lambda,\cdot,y)
	\bigr)(x)$  are denoted by
	$
	\mathcal{T}_{\alpha,\pm},
	\mathcal{T}_{4,\pm}^0,
	\mathcal{T}_{4,\pm}^1,
	\mathcal{T}_{4,\pm}^{1,1},
	$
	respectively. The estimate \eqref{lemma_projection_1} follows
	 from \eqref{eq:general_oscillatory_bound}.
	
	It remains to consider
	$
	Q\in\{Q_{42}^1,Q_5,Q_6\}.
	$
	Set
	$
	F_2^\pm(p)
	:=
	F_1^\pm(p)-a_1^\pm p^2,
	$ then $$
	(F_2^\pm)'(0)
	=
	(F_2^\pm)''(0)
	=
	(F_2^\pm)^{(3)}(0)
	=
	0.
	$$
	Moreover,  these projections annihilate all quadratic moments yielding 
	$
	Q
	\bigl(
	vF_1^\pm(\lambda|\cdot-y|)
	\bigr)
	=
	Q
	\bigl(
	vF_2^\pm(\lambda|\cdot-y|)
	\bigr).
	$
	For $Q=Q_{42}^1,Q_5$, set $m(Q)=1$, while for $Q=Q_6$, set
	$m(Q)=2$. Lemma \ref{taylor}(iii)--(iv) then yields
	\begin{align*}
		&
		Q
		\bigl(
		v\mathcal{R}_\pm(\lambda,\cdot,y)
		\bigr)(x)
	=
		\frac{(-1)^{m(Q)}\lambda^{m(Q)}}{(m(Q)+1)!}
		Q
		\Bigl(
		|\cdot|^{m(Q)+2}v
		\int_0^1
		(1-\theta)^{m(Q)+1}
		\mathcal{L}_{m(Q)+2}[F_2^\pm]
		\bigl(
		\lambda|y-\theta\cdot|,
		\xi
		\bigr)
		\,d\theta
		\Bigr)(x).
	%	\label{osc-rem-high-cancellation}
	\end{align*}
	We denote the corresponding kernels by
	$
	\lambda\mathcal{T}_{4,\pm}^{1,2},
	\lambda\mathcal{T}_{5,\pm},
	\lambda^2\mathcal{T}_{6,\pm},
	$
	respectively. Moreover since   $\mathcal{L}_{m(Q)+2}[F_2^\pm](\lambda|y-\theta\cdot|,\xi)$ obey the  bound \eqref{eq:general_oscillatory_bound}, then it ensures that 
	all $\mathcal{T}_{\pm} \in \{ \mathcal{T}_{4,\pm}^{1,2},\mathcal{T}_{5,\pm}, \mathcal{T}_{6,\pm} \}$
	satisfy the  estimate \eqref{lemma_projection_1}.
	
	
	
	
	
	
	
	Combining Steps 1--4 with the decomposition
	\eqref{basic-four-term-decomposition} proves all the expansions and
	estimates in part~{\rm (iii)}, and completes the proof.
\end{proof}









\subsection{Oscillatory integral estimates}
In this subsection we establish the estimates for the oscillatory integrals that arise in this paper. We begin by proving Lemma~\ref{oscillatory}.
\begin{comment}
	\begin{proof}    
	Assume $t > 0$ and $r \ge 0$ without loss of generality and abbreviate $\mathcal{E}(\lambda) := \mathcal{E}^\pm(\lambda, x,y)$ and $h := h(x,y)$. 

\textbf{Step 1: Reduction to a Model Oscillatory Integral.} 
	Using Euler's identities, both $\mathcal{K}^\pm$ and $\mathcal{N}^\pm$ can be expressed as linear combinations of integrals of the form    
	\begin{align}  \label{form-cos}  
		\int_{0}^\infty e^{i\Phi(\lambda)} \lambda \mathcal{A}(\lambda) \, d\lambda,  
	\end{align}  
	where the generalized phase  $\Phi(\lambda) := \alpha t\lambda^2 + \beta \lambda r$ for $\alpha, \beta \in \{+1, -1\}$. For $\mathcal{K}^\pm$, the effective amplitude factor is $\mathcal{A}(\lambda) = \mathcal{E}(\lambda)$. For $\mathcal{N}^\pm$, the corresponding amplitude factor is $\mathcal{A}(\lambda) = t^{-1} \lambda^{-2} \mathcal{E}(\lambda)$. 
	We claim that estimating $\mathcal{N}^\pm$ reduces entirely to the bounds for $\mathcal{K}^\pm$. To see this, we partition the frequency domain using $1 = \chi_1(t\lambda^2) + \chi_2(t\lambda^2)$. 
    
    In the low-frequency regime ($t\lambda^2 \le 2\lambda_0\ll1$), the inequality $|\sin(t\lambda^2)| \le t\lambda^2$ ensures the integrand of $\mathcal{N}^\pm$ is bounded by $\lambda |\mathcal{E}(\lambda)|$, precisely mirroring $\mathcal{K}^\pm$. In the high-frequency regime ($t\lambda^2 \gtrsim 1$), the factor $t^{-1}\lambda^{-2} \lesssim 1$. Thus, the modified amplitude factor  $t^{-1} \lambda^{-2} \mathcal{E}(\lambda)$ satisfies the exact same decay assumptions \eqref{condit 1} and \eqref{condit 2} as $\mathcal{E}(\lambda)$.    
    Consequently, we restrict our subsequent analysis solely to the model integral \eqref{form-cos} with $\mathcal{A}(\lambda) = \mathcal{E}(\lambda)$.
	
	The phase derivatives are $\Phi'(\lambda) = 2\alpha t\lambda + \beta r$ and $\Phi''(\lambda) = 2\alpha t$. If $\alpha\beta = -1$, a stationary point exists at $\rho = r/(2t)$. We decompose the high-frequency regime ($\lambda \geq t^{-1/2}\lambda_0^{1/2}$) into a stationary region $\mathcal{I}_{stat} = \{ \lambda : |\lambda - \rho| \le \lambda/2 \}$ and a non-stationary region $\mathcal{I}_{non} = \{ \lambda : |\lambda - \rho| > \lambda/2 \text{ or } \alpha\beta = 1 \}$. Note that $\frac{2}{3}\rho \le \lambda \le 2\rho$ (i.e., $\lambda \sim \rho$) in $\mathcal{I}_{stat}$, whereas $|\Phi'(\lambda)| \gtrsim t\lambda$ in $\mathcal{I}_{non}$.
	
	\vspace{0.5em}

\textbf{Step 2: Proof of Part (i).} 
	For the low-frequency part, assumption \eqref{condit 1} yields directly
	\begin{align}\label{c-1}
		\Big| \int_0^{\infty} e^{i\Phi} \chi_1(t\lambda^2)\lambda \mathcal{E}(\lambda) \, d\lambda \Big| \lesssim \int_0^{t^{-1/2}} \lambda |\mathcal{E}(\lambda)| \, d\lambda \lesssim h \int_0^{t^{-1/2}} \lambda \, d\lambda \lesssim h t^{-1}.
	\end{align}
	For the high-frequency part over $\mathcal{I}_{\mathrm{stat}}$, we apply the Van der Corput lemma (see, e.g., \cite[P. 334]{Stein}). Since $|\Phi''(\lambda)| \sim t$ and $\lambda \sim \rho = r/(2t)$,  setting $\psi(\lambda) :=\chi_2(t\lambda^2) \lambda \mathcal{E}(\lambda),$ it follows that
	\begin{equation}\label{c-2}
		\begin{split}
		\left| \int_{\mathcal{I}_{\mathrm{stat}}} e^{i\Phi}\chi_2(t\lambda^2) \lambda \mathcal{E}(\lambda) \, d\lambda \right|  
		&\lesssim t^{-1/2} \left( \|\psi\|_{L^\infty(\mathcal{I}_{\mathrm{stat}})} + \int_{\mathcal{I}_{\mathrm{stat}}} |\partial_\lambda \psi(\lambda)| \, d\lambda \right) \\
		&\lesssim t^{-1/2} \sup_{\lambda \sim \rho} \Big(\lambda |\mathcal{E}(\lambda)|+\rho |\mathcal{E}(\lambda)|+\rho\lambda|\partial_\lambda\mathcal{E}(\lambda)|\Big) \\
		&\lesssim t^{-1/2} (\rho h \langle\rho r \rangle^{-1/2}) \lesssim h t^{-1/2} \frac{r}{2t} \left(\frac{r^2}{2t}\right)^{-1/2} \lesssim h t^{-1}.
		\end{split}
	\end{equation}
	In $\mathcal{I}_{\mathrm{non}}$, performing two successive integrations by parts yields
	\begin{align} \label{integration part}
		\int_{\mathcal{I}_{\mathrm{non}}} e^{i\Phi} \chi_2(t\lambda^2)\lambda \mathcal{E}(\lambda) \, d\lambda 
		&= \left[ e^{i\Phi} \frac{\chi_2(t\lambda^2)\lambda \mathcal{E}(\lambda)}{i\Phi'} \right]_{\partial \mathcal{I}_{\mathrm{non}}} 
		- \left[ e^{i\Phi} \frac{1}{i\Phi'} \partial_\lambda \left( \frac{\chi_2(t\lambda^2)\lambda \mathcal{E}(\lambda)}{i\Phi'} \right) \right]_{\partial \mathcal{I}_{\mathrm{non}}} \nonumber\\
		&\quad + \int_{\mathcal{I}_{\mathrm{non}}} e^{i\Phi} \partial_\lambda \left( \frac{1}{i\Phi'} \partial_\lambda \left( \frac{\chi_2(t\lambda^2)\lambda \mathcal{E}(\lambda)}{i\Phi'} \right) \right) d\lambda.
	\end{align}
	Using assumption \eqref{condit 1} together with the bounds $|\Phi'| \gtrsim t\lambda$ and $|\Phi''| \sim t$, both boundary terms are bounded by $O(h t^{-1})$, while the integrand is bounded by $O(h t^{-2} \lambda^{-3})$. Consequently,
	\begin{align}\label{c-3}
		\left| \int_{\mathcal{I}_{\mathrm{non}}} e^{i\Phi} \chi_2(t\lambda^2) \lambda \mathcal{E}(\lambda) \, d\lambda \right| \lesssim h t^{-1} + h t^{-2} \int_{\lambda_0^{1/2}t^{-1/2}}^\infty \lambda^{-3} \, d\lambda \lesssim h t^{-1}.
	\end{align}
	Combining \eqref{c-1}, \eqref{c-2}, and \eqref{c-3} establishes Part (i).
	
	\vspace{0.5em}  
	 \textbf{Step 3: Proof of Part (ii).} 	For $t \le 2$, the oscillatory nature is not required. The support of $\widetilde{\chi}_1$ and assumption \eqref{condit 2} guarantee that the integral is bounded by $h \int_0^{1/4} \lambda^{1-\sigma} |\log \lambda|^{-\nu} d\lambda$. This converges to $O(h)$ under the  conditions $\sigma<2, \nu\in\mathbb{R}$ or $\sigma=2, \nu>1$.
	
	For $t > 2$, the low-frequency component extracts the precise logarithmic decay via the upper incomplete Gamma function $\Gamma(a, x) = \int_x^\infty u^{a-1} e^{-u} du$:
	\begin{align*}
		\int_0^{t^{-1/2}} \frac{\lambda}{\lambda^{\sigma} |\log \lambda|^{\nu}} d\lambda 
		= \int_{\frac{1}{2}\log t}^{\infty} \frac{u^{-\nu}}{e^{u(2-\sigma)}} du	
		= \begin{cases}
			\frac{2^{\nu-1}}{(\nu-1)(\log t)^{\nu-1}}, & \sigma=2, \nu>1, \\[4pt]
			(2-\sigma)^{\nu-1} \Gamma( 1-\nu, \frac{2-\sigma}{2} \log t ), & \sigma<2, \nu \in \mathbb{R}.
		\end{cases}
	\end{align*}
	Using the standard asymptotic expansion $\Gamma(a,x) \sim x^{a-1}e^{-x}$ as $x \to \infty$, we obtain:
	\begin{align}\label{low-fre}
		\Big| \int_0^{\infty} e^{i\Phi} \chi_1(t\lambda^2)\lambda \mathcal{E}(\lambda) d\lambda \Big| \lesssim 
		\begin{cases}
			h(\log t)^{-\nu+1}, & \sigma=2, \nu>1, \\[4pt]
			ht^{-1+\sigma/2} (\log t)^{-\nu}, & \sigma<2, \nu \in \mathbb{R}.
		\end{cases}
	\end{align}
	For the high-frequency part, note that under the given parameter conditions ($0 < \sigma < 2$, $\nu \in \mathbb{R}$; or $\sigma = 0$, $\nu \le 0$; or $\sigma = 2$, $\nu > 1$),
	 the map $\lambda \mapsto \lambda^{-\sigma}|\log \lambda|^{-\nu}$ is decreasing for $\lambda \ll 1$. Thus, 
	 \begin{align}\label{log}
	 	\lambda^{-\sigma}|\log \lambda|^{-\nu}\lesssim t^{\sigma/2}(\log t)^{-\nu}, \quad \text{on the support interval}\  \lambda\in[\lambda_0^{\frac{1}{2}}t^{-\frac{1}{2}}, 2(2\lambda_0)^{\frac{1}{4}}].\end{align}
	Applying the Van der Corput lemma in $\mathcal{I}_{stat}$ exactly as in \eqref{c-2}, by \eqref{log}, the bound becomes:
	\begin{align*}
		\left| \int_{\mathcal{I}_{stat}} e^{i\Phi}\chi_2(t\lambda^2) \lambda \mathcal{E}(\lambda) d\lambda \right| &\lesssim h t^{-1/2} \langle \rho r \rangle^{-1/2} \rho \cdot t^{\sigma/2}(\log t)^{-\nu} \lesssim h t^{-1+\sigma/2}(\log t)^{-\nu}.
	\end{align*}
	
	Similarly, integrating by parts twice over $\mathcal{I}_{\mathrm{non}}$ yields the representation \eqref{integration part}. Using assumption \eqref{condit 2} together with the bounds $|\Phi'| \gtrsim t\lambda$, $|\Phi''| \sim t$, and \eqref{log}, we deduce that the boundary terms are bounded by $O(t^{-1+\sigma/2} (\log t)^{-\nu} h)$, and  the integrand of the remainder integral is bounded by $O(t^{-2} \lambda^{-\sigma-3} |\log \lambda|^{-\nu} h)$. Consequently,
	\begin{align*}
		h \int_{\lambda_0^{1/2}t^{-1/2}}^{2(2\lambda_0)^{1/4}} t^{-2} \lambda^{-\sigma-3} |\log \lambda|^{-\nu} \, d\lambda 
		\lesssim h t^{\frac{\sigma}{2}-2} (\log t)^{-\nu} \int_{t^{-1/2}}^{\infty} \lambda^{-3} \, d\lambda 
		\lesssim \frac{h}{t^{1-\sigma/2}(\log t)^\nu}.
	\end{align*}
	Hence we obtain the desired estimates over $\mathcal{I}_{\mathrm{non}}$. Combining these with the estimates over $\mathcal{I}_{\mathrm{stat}}$, we deduce the following bound in the high-frequency regime:
	\begin{align}\label{high-fre}
		\Big| \int_0^{\infty} e^{i\Phi} \chi_2(t\lambda^2)\lambda \mathcal{E}(\lambda) \, d\lambda \Big| \lesssim \frac{h}{t^{1-\sigma/2}(\log t)^\nu}.
	\end{align}
	Summing the low- and high-frequency contributions yields the desired bound. 
	
	Finally, we emphasize that in the case $0 < \sigma < 2$, $\nu \in \mathbb{R}$, or $\sigma = 2$, $\nu > 1$, a single integration by parts (requiring only derivative indices $\ell=0,1$) is strictly sufficient for the estimate over $\mathcal{I}_{\mathrm{non}}$. Indeed, after one integration by parts, the remainder integral directly satisfies
	\begin{align*}  
		\left| \int_{\mathcal{I}_{\mathrm{non}}} e^{i\Phi} \partial_\lambda \left( \frac{\chi_2(t\lambda^2)\lambda \mathcal{E}(\lambda)}{i\Phi'} \right) d\lambda \right|  
		&\lesssim h \int_{\lambda_0^{1/2}t^{-1/2}}^{2(2\lambda_0)^{1/4}} t^{-1} \lambda^{-\sigma-1} |\log \lambda|^{-\nu} \, d\lambda \\  
		&\lesssim h t^{-1} t^{\sigma/4} (\log t)^{-\nu} \int_{t^{-1/2}}^\infty \lambda^{-1-\sigma/2} \, d\lambda 
		\lesssim \frac{h}{t^{1-\sigma/2}(\log t)^\nu},  
	\end{align*}  
	where we used the fact that the map $\lambda \mapsto \lambda^{-\sigma/2}|\log \lambda|^{-\nu}$ is decreasing for $\lambda \ll 1$.	
\end{proof}    


\begin{remark}\label{remark:osci} {\rm
	Based on the proof of Lemma~\ref{oscillatory}, we make the following two observations.
	\begin{itemize}
        \item[(i)] In the proof of Lemma~\ref{oscillatory}, we essentially establish that under assumption \eqref{condit 1},
\[
\int_{0}^{\infty} e^{i\Phi(\lambda)} \lambda \mathcal{E}(\lambda,x,y) \, d\lambda\lesssim h(x,y)|t|^{-1},
\]
while under assumption \eqref{condit 2}, we have
  \[
    \int_{0}^{\infty} e^{i\Phi(\lambda)} \lambda \mathcal{E}(\lambda,x,y) \, d\lambda \lesssim 
    \begin{cases}
        \dfrac{h(x,y)}{\langle t \rangle^{1-\frac{\sigma}{2}} \bigl(\log(2+|t|)\bigr)^{\nu}}, & 0 < \sigma < 2,\ \nu \in \mathbb{R}\ \text{or}\ \sigma = 0,\ \nu \le 0, \\[8pt]
h(x,y)\bigl(\log(2+|t|)\bigr)^{-\nu+1}, & \sigma = 2,\ \nu > 1,
    \end{cases}
    \]
   where $\Phi(\lambda)= \alpha t\lambda^2 + \beta \lambda r$ with  $\alpha, \beta \in \{+1, -1\}.$ 
        
		\item[(ii)] By estimates \eqref{low-fre} and \eqref{high-fre}, in the case where $\sigma = 2$ and $\nu > 1$, we obtain
		\[
		\Big| \int_0^{\infty} e^{i\Phi} \chi_1(t\lambda^2) \lambda \mathcal{E}(\lambda) \, d\lambda \Big| \lesssim \frac{h(x,y)}{(\log t)^{\nu-1}}, \qquad
		\Big| \int_0^{\infty} e^{i\Phi} \chi_2(t\lambda^2) \lambda \mathcal{E}(\lambda) \, d\lambda \Big| \lesssim \frac{h(x,y)}{(\log t)^{\nu}}.
		\]
		The high-frequency decay $O(h(x,y)(\log t)^{-\nu})$ is strictly faster (by one full logarithmic power) than the low-frequency integral $O(h(x,y)(\log t)^{-(\nu-1)})$.
	\end{itemize}}
\end{remark}
\end{comment}
\begin{proof}
Assume $t > 0$ without loss of generality and abbreviate $\mathcal{E}(\lambda) := \mathcal{E}^\pm(\lambda, x,y)$ and $h := h(x,y)$. 
	\medskip
	
	%\noindent
	\textbf{Step 1: Reduction to a model oscillatory integral.}
	By Euler's identities, the estimate for $\mathcal K^\pm$ reduces to
	estimating integrals of the form
	$$
	K(t,x,y)
	:=
	\int_0^\infty
	e^{i\Phi(\lambda)}
	\lambda\mathcal E(\lambda)d\lambda,
	\qquad
	\Phi(\lambda)
	:=
	\alpha t\lambda^2+\beta r\lambda,
	$$
	where $\alpha,\beta\in\{-1,1\}$.
	The same reduction applies to $\mathcal N^\pm$. Indeed, decompose
	$
	1
	=
	\chi_1(t\lambda^2)
	+
	\chi_2(t\lambda^2).
	$
	On the support of $\chi_1(t\lambda^2)$, it follows that
	$ 
	|\sin(t\lambda^2)|
	\lesssim
	t\lambda^2,
	$
	and hence
	$
	t^{-1}
	|\sin(t\lambda^2)|
	\lambda^{-1}
	|\mathcal E(\lambda)|
	\lesssim
	\lambda|\mathcal E(\lambda)|.
	$
	Thus the low-frequency contribution to $\mathcal N^\pm$ is estimated
	in the same way as the corresponding contribution to $K$.
	For the high-frequency contribution, after expanding the sine into
	exponentials, we obtain the relevant amplitude
	$
	\mathcal E_t(\lambda)
	:=
	t^{-1}\lambda^{-2}
	\chi_2(t\lambda^2)
	\mathcal E(\lambda).
	$
		On the support of $\chi_2(t\lambda^2)$ that
$	t\lambda^2\gtrsim1,
	$
it follows that
	$$
	\big|
	\partial_\lambda^\ell
	\left(
	t^{-1}\lambda^{-2}\chi_2(t\lambda^2)
	\right)
	\big|
	\lesssim
	\lambda^{-\ell}.
	$$
	Hence $\mathcal E_t$ satisfies the same
	amplitude estimates as $\mathcal E$.
	It therefore suffices to estimate
	$K$.
	
	
	We decompose
	$
	K(t,x,y)
	=
	K_{\mathrm{lo}}(t,x,y)
	+
	K_{\mathrm{hi}}(t,x,y)
	$
by
	$
	1
	=
	\chi_1(t\lambda^2)
	+
	\chi_2(t\lambda^2).
	$
Note that	$
	\Phi'(\lambda)
	=
	2\alpha t\lambda+\beta r,
$
then	a positive stationary point can occur only
 when \begin{align}\label{critial}
 \Phi'(\rho)
 =0
 \  \text{ with}
	\ 
	\rho
	=
	\frac{|r|}{2t}>0.
	\end{align}
	Choose $\eta\in C_c^\infty((0,\infty))$ such that $	\operatorname{supp}\eta
	\subset
	[\frac14,4]$ and
	$
	\eta(s)=1$ for
	$\frac12\leq s\leq2.
	$
	If there exists $\rho$ satisfying the condition \eqref{critial}, define
	$
	\eta_\rho(\lambda)
	=
	\eta\left({\lambda}/{\rho}\right);
	$
	otherwise, set $\eta_\rho=0$.
	 We then write
	$
	K_{\mathrm{hi}}
	=
	K_{\mathrm{stat}}
	+
	K_{\mathrm{non}}
	$
	where
	$$
	K_{\mathrm{stat}}
	=
	\int_0^\infty
	e^{i\Phi(\lambda)}
	\eta_\rho(\lambda)
	\chi_2(t\lambda^2)
	\lambda\mathcal E(\lambda)d\lambda
	,\quad 
	K_{\mathrm{non}}
	=
	\int_0^\infty
	e^{i\Phi(\lambda)}
	\bigl(1-\eta_\rho(\lambda)\bigr)
	\chi_2(t\lambda^2)
	\lambda\mathcal E(\lambda)d\lambda.
	$$
	On the support of $\eta_\rho$, one has
	$
	\lambda\sim\rho,
	$
	whereas on the support of the non-stationary amplitude, 
	$
	|\Phi'(\lambda)|
	\gtrsim
	t\lambda.
	$
	
	\medskip
	
	%\noindent
	\textbf{Step 2: Proof of Part \emph{(i)}.}
	For the low-frequency contribution, assumption \eqref{condit 1}
	implies
	$$
	\begin{aligned}
		|K_{\mathrm{lo}}|
		\lesssim
		\int_0^{Ct^{-1/2}}
		\lambda|\mathcal E(\lambda)| d\lambda 
		\lesssim
		h\int_0^{Ct^{-1/2}}\lambda d\lambda 
		\lesssim
		ht^{-1}.
	\end{aligned}
	$$
	
	For the stationary contribution in the high-frequency part, set
	$
	a_{\mathrm{stat}}(\lambda)
	:=
	\eta_\rho(\lambda)
	\chi_2(t\lambda^2)
	\lambda\mathcal E(\lambda).
	$
	Since $|\Phi''(\lambda)|=2t$, the van der Corput lemma yields
	\begin{align}\label{van-1}
	|K_{\mathrm{stat}}|
	\lesssim
	t^{-1/2}
	\left(
	\|a_{\mathrm{stat}}\|_{L^\infty}
	+
	\|\partial_\lambda a_{\mathrm{stat}}\|_{L^1}
	\right).
	\end{align}
	By the
	assumption \eqref{condit 1} and $\lambda\sim\rho$ on the support of $a_{\mathrm{stat}}$, it follows that
	$$
	\|a_{\mathrm{stat}}\|_{L^\infty}
	+
	\|\partial_\lambda a_{\mathrm{stat}}\|_{L^1}
	\lesssim
	h\rho
	\langle\rho r\rangle^{-1/2}\lesssim h  \frac{|r|}{2t} \left(\frac{r^2}{2t}\right)^{-1/2} \lesssim h t^{-1/2}.
	$$
	Therefore,
	$
	|K_{\mathrm{stat}}|
	\lesssim
	ht^{-1}.
	$
    
	It remains to estimate $K_{\mathrm{non}}$. Choose
	$\psi\in C_c^\infty((1/2,2))$ such that
	$$
	\sum_{j\in\mathbb Z}
	\psi(2^{-j}s)
	=
	1,
	\qquad
	s>0.
	$$
	Set
	$
	\Lambda_j
	:=
	2^j t^{-{1}/{2}}
	$ and $
	\psi_j(\lambda)
	:=
	\psi\left(\frac{\lambda}{\Lambda_j}\right),
	$
	together with $t^{-{1}/{2}}\lesssim\lambda,$ then
	\begin{align}\label{K-non}
	K_{\mathrm{non}}
	=
		\sum_{t^{-1/2}\lesssim\Lambda_j}
	K_j,
	\end{align}
	where
	$$
	K_j
	:=
	\int_0^\infty
	e^{i\Phi(\lambda)}
	a_j(\lambda) d\lambda,\quad
	a_j(\lambda)
	:=
	\psi_j(\lambda)
	\bigl(1-\eta_\rho(\lambda)\bigr)
	\chi_2(t\lambda^2)
	\lambda\mathcal E(\lambda).
	$$
	Each $a_j$ is compactly supported in $(0,\infty)$, so integration by
	parts produces no boundary terms.
	Hence two integrations by parts give
	\begin{align}\label{Ij}
	K_j
	=
	\int_0^\infty
	e^{i\Phi(\lambda)}
	L^2a_j(\lambda)d\lambda,
	\end{align}
where the operator is defined by $
Lf
=
\partial_\lambda
\Big(
\frac{f}{i\Phi'(\lambda)}
\Big).
$	
	 	Using assumption \eqref{condit 1}, together with 
	$
	|\Phi'(\lambda)|
	\gtrsim
	t\Lambda_j$
and 
$	|\Phi''(\lambda)|
	\lesssim
	t,
	$
we obtain the estimate
	$
	\big|
	\left(L\right)^2a_j(\lambda)
	\big|
	\lesssim
	h t^{-2}
	\Lambda_j^{-3}
	\langle\Lambda_j r\rangle^{-1/2}.
	$
	Since the support of $a_j$ has length $O(\Lambda_j)$, it follows that
	$$
	\begin{aligned}
		|K_j|
		&\leq
		\int_{\operatorname{supp}a_j}
		\left|
	(L^2a_j)(\lambda)
		\right|
		\,d\lambda\lesssim
		h t^{-2}
		\Lambda_j^{-2}
		\langle\Lambda_j r\rangle^{-1/2}\lesssim
		h t^{-2}
		\Lambda_j^{-2}.
	\end{aligned}
	$$
	Consequently,
	$$
	\begin{aligned}
		|K_{\mathrm{non}}|
		&\lesssim
		h t^{-2}
		\sum_{\Lambda_j\gtrsim t^{-1/2}}
		\Lambda_j^{-2}
	\lesssim
		h t^{-1}
		\sum_{j=0}^{\infty}
	2^{-2j} 
		\lesssim
		ht^{-1}.
	\end{aligned}
	$$
	
	
	
	%\noindent
	\textbf{Step 3: Proof of Part \emph{(ii)}.}
	For $0<t\lesssim1$,  by the compact
	support of $\widetilde{\chi}_1(\lambda/2)$ and assumption
	\eqref{condit 2},
	\begin{align}\label{t bound}
	|K|
	\lesssim
	h
	\int_0^{1/2}
	\lambda^{1-\sigma}
	|\log\lambda|^{-\nu} d\lambda.
	\end{align}
	This integral is finite when $\sigma<2$ or when $\sigma=2$ and
	$\nu>1$. Hence the asserted estimates hold.
	
	Suppose henceforth that $t\gg1$. 
The low-frequency component extracts the precise logarithmic decay via the upper incomplete Gamma function $\Gamma(a, x) = \int_x^\infty u^{a-1} e^{-u} du$:
	\begin{align*}
		|K_{\mathrm{lo}}|
	\lesssim h	\int_0^{t^{-1/2}} \frac{\lambda}{\lambda^{\sigma} |\log \lambda|^{\nu}} d\lambda 
		= h\int_{\frac{1}{2}\log t}^{\infty} \frac{u^{-\nu}}{e^{u(2-\sigma)}} du	
		= h\begin{cases}
			\frac{2^{\nu-1}}{(\nu-1)(\log t)^{\nu-1}}, & \sigma=2, \nu>1, \\[4pt]
			(2-\sigma)^{\nu-1} \Gamma( 1-\nu, \frac{2-\sigma}{2} \log t ), & \sigma<2, \nu \in \mathbb{R}.
		\end{cases}
	\end{align*}
	Using the standard asymptotic expansion $\Gamma(a,x) \sim x^{a-1}e^{-x}$ as $x \to \infty$, we obtain:
	\begin{align}\label{low-fre}
		|K_{\mathrm{lo}}| \lesssim 
		\begin{cases}
			h(\log t)^{-\nu+1},               & \sigma=2, \nu>1,              \\[4pt]
			ht^{-1+\sigma/2} (\log t)^{-\nu}, & \sigma<2, \nu \in \mathbb{R}.
		\end{cases}
	\end{align}
	
	

	
	For the high-frequency part, note that under the given parameter conditions ($0 < \sigma < 2$, $\nu \in \mathbb{R}$; or $\sigma = 0$, $\nu \le 0$; or $\sigma = 2$, $\nu > 1$),
	the map $\lambda \mapsto \lambda^{-\sigma}|\log \lambda|^{-\nu}$ is decreasing  on the sufficiently small support of
	$\widetilde{\chi}_1(\lambda/2)$.  Hence, whenever
	$\lambda\gtrsim t^{-1/2}$, it follows that
	$
	\lambda^{-\sigma}
	|\log\lambda|^{-\nu}
	\lesssim
	t^{\sigma/2}
	(\log t)^{-\nu}.
	$

	
	Applying the Van der Corput estimate  \eqref{van-1} and
	assumption \eqref{condit 2}, we obtain
	\begin{align}\label{Ksat}
		|K_{\mathrm{stat}}|
		&\lesssim
		h t^{-1/2}
		\rho^{1-\sigma}
		|\log\rho|^{-\nu}
		\langle\rho r\rangle^{-1/2}\lesssim
		h t^{-1+\sigma/2}
		(\log t)^{-\nu}.
	\end{align}
	



	
	
	We next estimate the non-stationary contribution. By the compact
	support of $\widetilde{\chi}_1(\lambda/2)$, the same dyadic
	decomposition as in \eqref{K-non} applies, except that only the scales
	satisfying
	$
	t^{-1/2}
	\lesssim
	\Lambda_j
	\leq
	\lambda_*
	$
	contribute, where $\lambda_*\in(0,1)$ is a fixed sufficiently small
	constant. 	After performing two integrations by parts with no boundary terms, we arrive at the same identity as in \eqref{Ij}. Applying assumption \eqref{condit 2} for \(\ell=0,1,2\), together with the estimates
	$
	|\Phi'(\lambda)| \gtrsim t\Lambda_j
	$ and
	$|\Phi''(\lambda)| \lesssim t,
	$
	we obtain
	$$
	|K_j|
	\lesssim
	h t^{-2}
	\Lambda_j^{-\sigma-2}
	|\log\Lambda_j|^{-\nu}
	\langle\Lambda_j r\rangle^{-1/2}\lesssim
	h t^{-2}
	\Lambda_j^{-\sigma-2}
	|\log\Lambda_j|^{-\nu}.
	$$
	Hence,  it follows that
	$$
	|K_{\mathrm{non}}|
	\lesssim
	h t^{-2}
	\sum_{t^{-1/2}\lesssim\Lambda_j\le \lambda_*}
	\Lambda_j^{-\sigma-2}
	|\log\Lambda_j|^{-\nu}.
	$$
It suffices to prove 
	\begin{align}\label{est-b}
	\sum_{t^{-1/2}\lesssim\Lambda_j\le \lambda_*}
	\Lambda_j^{-\sigma-2}
	|\log\Lambda_j|^{-\nu}
	\lesssim
	t^{1+\sigma/2}
	(\log t)^{-\nu}.
	\end{align}
	To prove this, we introduce the positive integer $J$ satisfying
	 $
	2^Jt^{-1/2}\sim\lambda_*.
	$
	Thus
	$
	J\sim\log t
	$ and 

	$$
	\sum_{t^{-1/2}\lesssim\Lambda_j\le \lambda_*}
\Lambda_j^{-\sigma-2}
|\log\Lambda_j|^{-\nu}\lesssim	t^{1+\frac{\sigma}{2}}
	\sum_{j=0}^{J}
	2^{-(\sigma+2)j}
	\left|
	\log\Lambda_j
	\right|^{-\nu}.
	$$
	We split the sum into the ranges $0\leq j\leq J/2$ and
	$J/2<j\leq J$. If $0\leq j\leq J/2$, then
	$
	\left|
	\log\Lambda_j
	\right|
	\sim
	\log t,
	$
	and hence for $\sigma>-2$,
	$$
	\sum_{0\leq j\leq J/2}
	2^{-(\sigma+2)j}
	\left|
	\log\Lambda_j
	\right|^{-\nu}
	\lesssim
	(\log t)^{-\nu}
	\sum_{j=0}^\infty2^{-(\sigma+2)j}
	\lesssim
	(\log t)^{-\nu}.
	$$
	For $J/2<j\leq J$, the fixed restriction
	$
t^{-\frac{1}{2}}	\lesssim\Lambda_j\leq\lambda_*<1
	$
	implies
	$
0<	|\log\lambda_*|\leq	|\log\Lambda_j|
\lesssim\log t
.
	$
Then, 
	$$
	\sum_{J/2<j\leq J}
	2^{-(\sigma+2)j}
	\left|\log
\Lambda_j
	\right|^{-\nu}\lesssim
	(\log t)^{\max\{-\nu,0\}}\sum_{J/2<j\leq J}
	2^{-(\sigma+2)j}
	\lesssim
	t^{-(\sigma+2)/4}
	(\log t)^{\max\{-\nu,0\}}	\lesssim
	(\log t)^{-\nu},
	$$
	for $\sigma>-2$. Therefore,
	we derive \eqref{est-b}.
	Consequently,
	\begin{align}\label{Knon}
	|K_{\mathrm{non}}|
	\lesssim
	h t^{-2}
	t^{1+\sigma/2}
	(\log t)^{-\nu}
	=
	h t^{-1+\sigma/2}
	(\log t)^{-\nu},\quad \text{for all}\ \sigma>-2 \ \text{and}\ \nu\in\mathbb{R}.
	\end{align}
	Combining the low-frequency, stationary, and non-stationary estimates (i.e., \eqref{low-fre}, \eqref{Ksat} and \eqref{Knon})
	proves the desired  estimate  of Part \emph{(ii)}. In particular, 
	if \(r=0\), only the
	low-frequency and non-stationary estimates  are needed. Thus, the bound  $h t^{-1+\sigma/2}
		(\log t)^{-\nu}$ is valid whenever
	$
	-2<\sigma<2$ and $
	\nu\in\mathbb R.
	$
\end{proof}
\begin{remark}\label{remark:osci} {\rm
	Based on the proof of Lemma~\ref{oscillatory}, we make the following two observations.
	\begin{itemize}
        \item[(i)] In the proof of Lemma~\ref{oscillatory}, we essentially establish that under assumption \eqref{condit 1},
\[
\int_{0}^{\infty} e^{i\Phi(\lambda)} \lambda \mathcal{E}(\lambda,x,y) \, d\lambda\lesssim h(x,y)|t|^{-1},
\]
while under assumption \eqref{condit 2}, we have
  \[
    \int_{0}^{\infty} e^{i\Phi(\lambda)} \lambda \mathcal{E}(\lambda,x,y) \, d\lambda \lesssim 
    \begin{cases}
        \dfrac{h(x,y)}{\langle t \rangle^{1-\frac{\sigma}{2}} \bigl(\log(2+|t|)\bigr)^{\nu}}, & 0 < \sigma < 2,\ \nu \in \mathbb{R}\ \text{or}\ \sigma = 0,\ \nu \le 0, \\[8pt]
h(x,y)\bigl(\log(2+|t|)\bigr)^{-\nu+1}, & \sigma = 2,\ \nu > 1,
    \end{cases}
    \]
   where $\Phi(\lambda)= \alpha t\lambda^2 + \beta \lambda r$ with  $\alpha, \beta \in \{+1, -1\}.$ 
   
        
    
        
        
        
		\item[(ii)] By estimates \eqref{t bound}, \eqref{low-fre}, \eqref{Ksat} and \eqref{Knon}, in the case where $\sigma = 2$ and $\nu > 1$, we obtain 
		\[
		\Big| \int_0^{\infty} e^{i\Phi} \chi_1(t\lambda^2) \lambda \mathcal{E}(\lambda) \, d\lambda \Big| \lesssim \frac{h(x,y)}{(\log |t|)^{\nu-1}}, \quad
		\Big| \int_0^{\infty} e^{i\Phi} \chi_2(t\lambda^2) \lambda \mathcal{E}(\lambda) \, d\lambda \Big| \lesssim \frac{h(x,y)}{(\log |t|)^{\nu}},\quad |t|\geq2.
		\]
		The high-frequency decay $O(h(x,y)(\log |t|)^{-\nu})$ is strictly faster (by one full logarithmic power) than the low-frequency integral $O(h(x,y)(\log |t|)^{-(\nu-1)})$.
	\end{itemize}}
\end{remark}
Next, we evaluate the following oscillatory integrals, which are important in the resonance analysis, especially in deriving sharp bounds when zero is a resonance of the second kind of $H$.
\begin{lemma}\label{lemma:I_J}    
	For $|t| \ge 2$, the following asymptotic expansions hold:    
	\[
    \begin{aligned}    
		I(t) &:= \int_0^\infty \cos(t\lambda^2) \lambda \widetilde{\chi}_1(\lambda) (\log \lambda)^2 \, d\lambda = \frac{\pi}{8} \frac{\log |t|}{|t|} + O\big(|t|^{-1}\big),  \\    
		J(t) &:= t^{-1} \int_0^\infty \sin(t\lambda^2) \lambda^{-1} \widetilde{\chi}_1(\lambda) (\log \lambda)^2 \, d\lambda = \frac{\pi}{16} \frac{(\log |t|)^2}{|t|} + O\Big(\frac{\log |t|}{|t|}\Big).    
	\end{aligned}
    \]
\end{lemma}    
\begin{comment}
\begin{proof}    
	We first consider $I(t)$. The change of variables $u = |t|\lambda^2$ yields     
	\begin{equation}\label{eq:I1_decomp}  
		I(t) = \frac{(\log |t|)^2}{8|t|} I_0(t) - \frac{\log |t|}{4|t|} I_1(t) + \frac{1}{8|t|} I_2(t),   
	\end{equation}    
	where $I_j(t) := \int_0^\infty \cos u \, \widetilde{\chi}_1(\sqrt{u/|t|}) (\log u)^j \, du$. Integration by parts via $\cos u = \frac{d}{du}\sin u$ gives    
	\[    
	I_j(t) = - \int_0^\infty \sin u \, \frac{d}{du}\Bigl[ \widetilde{\chi}_1(\sqrt{u/|t|}) (\log u)^j \Bigr] du.    
	\]    
	By the Leibniz rule, the contribution where the derivative acts on the cutoff function $\widetilde{\chi}_1$ is supported entirely in the regime $u \sim |t|$. Since $\widetilde{\chi}_1$ is smooth, repeated integration by parts shows that this boundary contribution decays rapidly, bounded by $O(|t|^{-N})$ for any integer $N \ge 1$.     
	Consequently, for $j=0$, we immediately obtain $I_0(t) = O(|t|^{-N})$.    
	
	For $j=1$, the derivative acting on $\log u$ yields the principal integral:    
	\[    
	I_1(t) = - \int_0^\infty \frac{\sin u}{u} \widetilde{\chi}_1(\sqrt{u/|t|}) \, du + O(|t|^{-N}) = -\frac{\pi}{2} + R_1(t) + O(|t|^{-N}),    
	\]    
	where the remainder is defined as $R_1(t) := \int_0^\infty \frac{\sin u}{u} \bigl(1 - \widetilde{\chi}_1(\sqrt{u/|t|})\bigr) du$. A further integration by parts on $R_1(t)$ yields:    
	\[    
	R_1(t) = \int_0^\infty \cos u \, \frac{d}{du}\Big[ \frac{1 - \widetilde{\chi}_1(\sqrt{u/|t|})}{u} \Big] du.    
	\]    
	The derivative inside the integral is bounded by $O(u^{-2})$, and the integrand is supported on $u \gtrsim |t|$. This directly evaluates to $R_1(t) = O(|t|^{-1})$. Hence, $I_1(t) = -\frac{\pi}{2} + O(|t|^{-1})$.    
	
	For $j=2$, observing that the untruncated integral $\int_0^\infty \frac{\sin u}{u} \log u \, du$ converges to a constant, an analogous integration by parts for the remainder confirms that $I_2(t) = O(1)$.     
	
	Substituting the respective bounds $I_0(t) = O(|t|^{-N})$, $I_1(t) = -\frac{\pi}{2} + O(|t|^{-1})$, and $I_2(t) = O(1)$ into \eqref{eq:I1_decomp} establishes \eqref{eq:I_asymp}.    
	
	For $J(t)$, noting that $t^{-1}\sin(t\lambda^2) \equiv |t|^{-1}\sin(|t|\lambda^2)$, the identical change of variables leads to:    
	\begin{equation}\label{eq:J_decomp}   
		J(t) = \frac{(\log |t|)^2}{8|t|} J_0(t) - \frac{\log |t|}{4|t|} J_1(t) + \frac{1}{8|t|} J_2(t),   
	\end{equation}    
	where $J_j(t) := \int_0^\infty \frac{\sin u}{u} \widetilde{\chi}_1(\sqrt{u/|t|}) (\log u)^j \, du$.    
	By analogous reasoning, $J_0(t)$ is the sole term providing the leading-order behavior, yielding $J_0(t) = \int_0^\infty \frac{\sin u}{u} \, du - R_1(t) = \frac{\pi}{2} + O(|t|^{-1})$. For $j=1, 2$, the untruncated integrals evaluate to convergent constants, ensuring $J_j(t) = O(1)$. Collecting these results into \eqref{eq:J_decomp} immediately yields \eqref{eq:J_asymp}, which completes the proof.    
\end{proof}  
\end{comment}
\begin{proof}
The asserted expansions follow from the change of variables
$
u=|t|\lambda^2,
$
the identity
$
(\log\lambda)^2
=
\frac14(\log u-\log|t|)^2,
$
and repeated integrations by parts. We omit the
routine details.
\end{proof}




%\cqd
%\bigskip

%%%%%%%%%%%%%%%%%%%%%%%%%%%%%%%%%%%%%%%%%%%%%%%%%%%%%%%%%%%%%%%%%%%%%%%%%%%%%%%%%%%%%%%%%%%%
%%%%%%%%%%%%%%%%%%%%%%%%%%%%%%%%%%%%%%%%%%%%%%%%%%%%%%%%%%%%%%%%%%%%%%%%%%%%%%%%%%%%%%%%%%%


{\bf Acknowledgements:} The authors are partially supported by NSFC (grants No. 12171182 and 12531005). The authors would like to express their thanks to Professor Avy Soffer for his interests and insightful discussions.


%%%%%%%%%%%%%%%%%%%%%%%%%%%%%%%%%%%%%%%%%%%%%%%%%%%%%%%%%%%%%%%%%%%%%%%	
\begin{thebibliography}{99}
\bibitem{AS64}
M. Abramowitz, I. A. Stegun,
\emph{Handbook of Mathematical Functions with Formulas, Graphs, and Mathematical Tables,}
National Bureau of Standards Applied Mathematics Series, vol. 55, U.S. Government
Printing Office, Washington, D.C., 1964. For sale by Superintendent of Documents.

\bibitem{Agmon}
S. Agmon, \emph{Spectral properties of Schr\"{o}dinger operators and
  scattering theory}, Ann. Scuola Norm. Sup. Pisa Cl. Sci., (4) \textbf{2} (1975), no. 2, 151--218.

  \bibitem{Beals}
M. Beals, \emph{Optimal $L^\infty$ decay for solutions to the wave equation with a potential},
 Comm. Partial Differential Equations, \textbf{19} (1994), no. 7--8, 1319--1369.

\bibitem{B-S}
M. Beals, W. Strauss, \emph{$L^p$ estimates for the wave equation with a potential},
 Comm. Partial Differential Equations, \textbf{18} (1993), no. 7--8, 1365--1397.

\bibitem{Bec-Gold}
M. Beceanu, M. Goldberg, \emph{Strichartz estimates and maximal operators  for the wave equation in $\mathbb{R}^3$},
J. Funct. Anal., \textbf{266} (2014), no. 3, 1476--1510.
\bibitem{Bec16}
M. Beceanu,
\emph{Decay estimates for the wave equation in two dimensions},
J. Differential Equations, \textbf{260} (2016), no.~6, 5378--5420.
\bibitem{BPSZ03}
N. Burq, F. Planchon, J. Stalker, A. S. Tahvildar-Zadeh,
\emph{Strichartz estimates for the wave and Schr\"{o}dinger equations with the inverse-square potential},
J. Funct. Anal., \textbf{203} (2003), no. 2, 519--549.

\bibitem{BPSZ04}
N. Burq, F. Planchon, J. Stalker, A. S. Tahvildar-Zadeh,
\emph{Strichartz estimates for the wave and Schr\"{o}dinger equations with potentials of critical decay},
Indiana Univ. Math. J., \textbf{53} (2004), no. 6, 1665--1680.

\bibitem{CV}
F. Cardoso, G. Vodev,
\emph{Optimal dispersive estimates for the wave equation with $C^{\frac{n-3}{2}}$ potentials in dimensions $4\leq n \leq 7$},
 Comm. Partial Differential Equations, \textbf{37} (2012), no. 1, 88--124.

\bibitem{CLSY_ArXiv} M. Chen, P. Li, A. Soffer, X. Yao, \emph{Decay estimates for Beam equations with potential in dimension three}, J. Funct. Anal., \textbf{288} (2025), no. 1, 110671.

\bibitem{CWY}
S. Chen, Z. Wan, X. Yao, \emph{Decay estimates for beam equations with potentials on the line}, J. Differential Equations, \textbf{439} (2025), 113418.

\bibitem{CSWY25}
H. Cheng, A. Soffer, Z. Wu, X. Yao,
\emph{The $L^p$-boundedness of wave operators for higher order Schr\"{o}dinger operator with zero singularities in low odd dimensions},
arXiv e-prints, (2025), arXiv:2505.07009.

\bibitem{DP}
P. D'Ancona, V. Pierfelice,
\emph{On the wave equation with a large rough potential}, J. Funct. Anal., \textbf{227} (2005),
no. 1, 30--77.

\bibitem{EGG14}
M. B. Erdo\u{g}an, M. Goldberg, W. R. Green,
\emph{Dispersive estimates for four dimensional Schr\"{o}dinger and wave equations with obstructions at zero energy}, Comm. Partial Differential Equations, \textbf{39} (2014), no. 10, 1936--1964.

\bibitem{EGG23} M. B. Erdo\u{g}an, M. Goldberg, W. R. Green, \emph{Counterexamples to $L^p$ boundedness of wave operators for classical and higher order Schr\"odinger operators}, J. Funct. Anal., \textbf{285} (2023), no. 5, 110008.

\bibitem{EGG_ArXiv} M. B. Erdo\u{g}an, M. Goldberg, W. R. Green, \emph{Dispersive estimates for higher order Schr\"odinger operators with scaling-critical potentials}, Math. Ann., \textbf{392} (2025), no. 2, 2225--2252.

\bibitem{EG-2013}
M. B. Erdo\u{g}an, W. R. Green, \emph{A weighted dispersive estimate for Schr\"{o}dinger
operators in dimension two}, Comm. Math. Phys., \textbf{319} (2013), 791--811.

\bibitem{Erdogan-Green}
M. B. Erdo\u{g}an, W. R. Green, \emph{Dispersive estimates for Schr\"{o}dinger operators
 in dimension two with obstructions at zero energy}, Trans. Amer. Math. Soc., \textbf{365} (2013), no. 12, 6403--6440.

\bibitem{EG21}
M. B. Erdo\u{g}an, W. R. Green, \emph{The $L^p$-continuity of wave operators for higher order
Schr\"{o}dinger operators}, Adv. Math., \textbf{404} (2022), 108450.

\bibitem{Erdogan-Green23} M. B. Erdo\u{g}an, W. R. Green, \emph{A note on endpoint $L^p$-continuity of wave operators for classical and higher order Schr\"odinger operators}, J. Differential Equations, \textbf{355} (2023), 144--161.

\bibitem{EGL}
 M. B. Erdo\u{g}an, W. R. Green, K. LaMaster, \emph
{$L^p$-continuity of wave operators for higher order Schr\"odinger operators with threshold eigenvalues in high dimensions}, Discrete Contin. Dyn. Syst.,
 \textbf{45} (2025), no. 9,  3258--3274.

\bibitem{EGT19}
M. B. Erdo\u{g}an, W. R. Green, E. Toprak, \emph{On the fourth order 
Schr\"{o}dinger equation in three dimensions: dispersive estimates and
 zero energy resonances}, J. Differential Equations, \textbf{271} (2021), 152--185.

\bibitem{FSWY20}
H. Feng, A. Soffer, Z. Wu, X. Yao, \emph{Decay estimates for higher-order
elliptic operators}, Trans. Amer. Math. Soc., \textbf{373} (2020), no. 4, 2805--2859.

\bibitem{FSY}
H. Feng, A. Soffer, X. Yao, \emph{Decay estimates and
 Strichartz estimates of fourth-order Schr\"{o}dinger operator}, J. Funct. Anal., \textbf{274} (2018), no. 2, 605--658.

\bibitem{FHHH82}
R. Froese, I. Herbst, M. Hoffmann-Ostenhof,
T. Hoffmann-Ostenhof, \emph{On the absence of positive eigenvalues for one-body
  Schr\"{o}dinger operators}, J. Anal. Math., \textbf{41} (1982), 272--284.

\bibitem{Yajima_2024JST}
  A. Galtbayar, K. Yajima, \emph{The $L^p$-boundedness of wave operators for fourth order Schr\"odinger operators
  	on $\R^4$}, J. Spectr. Theory, \textbf{14} (2024), no. 1, 271--354.

\bibitem{GV03}
V. Georgiev, N. Visciglia, \emph{Decay estimates for wave equation with potential},
 Comm. Partial Differential Equations, \textbf{28} (2003), no. 7--8, 1325--1369.

\bibitem{GV95}
J. Ginibre, G. Velo, \emph{Generalized Strichartz inequalities for the wave equation},
J. Funct. Anal., \textbf{133} (1995), no. 1, 50--68.

\bibitem{GG21b}
M. Goldberg,  W. R. Green, \emph{On the $L^p$ boundedness of the wave operators for fourth
order Schr\"{o}dinger operators}, Trans. Amer. Math. Soc., \textbf{374} (2021), no. 6, 4075--4092.

\bibitem{GV06}
M. Goldberg, M. Visan, \emph{A counterexample to dispersive estimates for Schr\"{o}dinger operators in higher
 dimensions}, Comm. Math. Phys., \textbf{266} (2006), 211--238.

\bibitem{Green14}
W. R. Green, \emph{Time decay estimates for the wave equation with potential in dimension two},
 J. Differential Equations, \textbf{257} (2014), no. 3, 868--919.

\bibitem{GT19}
W. R. Green, E. Toprak, \emph{On the fourth order Schr\"{o}dinger equation in four dimensions: Dispersive estimates and zero energy resonances},
J. Differential Equations, \textbf{267} (2019), no. 3, 1899--1954.

\bibitem{IJ03}
A. D. Ionescu, D. Jerison, \emph{On the absence of positive eigenvalues of
  {S}chr\"{o}dinger operators with rough potentials}, Geom. Funct. Anal.,
  \textbf{13} (2003), no. 5, 1029--1081.

\bibitem{JK}
A. Jensen, T. Kato, \emph{Spectral properties of Schr\"{o}dinger
  operators and time-decay of the wave functions}, Duke Math. J., \textbf{46} (1979), no. 3, 583--611.

\bibitem{JN}
A. Jensen, G. Nenciu, \emph{A unified approach to resolvent expansions at thresholds}, Rev. Math. Phys., \textbf{13} (2001), no. 6, 717--754.

\bibitem{Kat59}
T. Kato, \emph{Growth properties of solutions of the reduced wave equation
 with a variable coefficient}, Comm. Pure Appl. Math., \textbf{12} (1959), 403--425.

\bibitem{Kirchhoff1850}
G. Kirchhoff, \emph{\"Uber das Gleichgewicht und die Bewegung einer elastischen Scheibe},
J. Reine Angew. Math. \textbf{40} (1850), 51--88.

\bibitem{KT06}
H. Koch,  D. Tataru, \emph{Carleman estimates and absence of embedded eigenvalues}, Comm. Math. Phys., \textbf{267} (2006), no. 2, 419--449.

\bibitem{Kop}
E. Kopylova, \emph{Dispersive estimates for the 2D wave equation}, Russ. J. Math. Phys.,
\textbf{17} (2010), no. 2, 226--239.

\bibitem{LSY21}
P. Li, A. Soffer,  X. Yao,
\emph{Decay estimates for fourth-order Schr\"odinger operators in dimension two},
J. Funct. Anal., \textbf{284} (2023), no. 6, 109816.

\bibitem{Love1927}
A. E. H. Love, \emph{A Treatise on the Mathematical Theory of Elasticity},
4th ed., Cambridge University Press, Cambridge, 1927.




\bibitem{MSW}
B. Marshall, W. Strauss, S. Wainger,
 \emph{$L^p-L^q$ estimates for the Klein-Gordon equation}, J. Math. Pures Appl., (9) \textbf{59} (1980), no. 4, 417--440.



\bibitem{MWY22}
H. Mizutani, Z. Wan, X. Yao,
 \emph{$L^p$-boundedness of wave operators for bi-Schr\"{o}dinger operators on the line},
Adv. Math., \textbf{451} (2024), 109806.

\bibitem{MWY23}
H. Mizutani, Z. Wan, X. Yao,
\emph{Counterexamples and weak $(1,1)$ estimates of wave operators for fourth-order Schr\"odinger operators in dimension three}, J. Spectr. Theory, \textbf{14} (2024), no. 4, 1409--1450.


\bibitem{MWY_ArXiv23_2} H. Mizutani, Z. Wan, X. Yao, \emph{$L^p$-boundedness of wave operators for fourth order Schr\"odinger operators with zero resonances on $\R^3$}, J.  Funct. Anal., \textbf{289} (2025), no. 8, 111013.

\bibitem{Mou}
S. Moulin,
 \emph{High frequency dispersive estimates in dimension two}, Ann. Henri Poincar\'{e},
 \textbf{10} (2009), no. 2, 415--428.

\bibitem{Schlag-2005}
W. Schlag, \emph{Dispersive estimates for Schr\"odinger operators in dimension two}, Comm. Math.
Phys., \textbf{257} (2005), no. 1, 87--117.

\bibitem{Simon69}
B. Simon, \emph{On positive eigenvalues of one-body Schr\"{o}dinger operators}, Comm. Pure Appl. Math., \textbf{22}
 (1969), 531--538.

\bibitem{SWY22} A. Soffer, Z. Wu, X. Yao, \emph{Decay estimates for bi-Schr\"odinger operators in dimension one}, Ann. Henri Poincar\'{e}, \textbf{23} (2022), 2683--2744.

\bibitem{Stein}
E.~M. Stein, \emph{Harmonic analysis: real-variable methods, orthogonality, and oscillatory integrals}, Princeton Mathematical Series, vol. 43, Princeton University Press, Princeton, NJ, 1993. With the assistance of Timothy S. Murphy, Monographs in Harmonic Analysis, III.

\bibitem{Str05}
W. Strauss, \emph{Nonlinear Wave Equations}, CBMS Reg. Conf. Ser. Math., vol. 73, Amer. Math. Soc., Providence, RI, 1989.

\bibitem{TimoshenkoWoinowskyKrieger1959}
S. P. Timoshenko, S. Woinowsky-Krieger,
\emph{Theory of Plates and Shells},
2nd ed., McGraw--Hill, New York, 1959.



\end{thebibliography}

\end{document}






